% GENERATED FILE -- DO NOT EDIT.
%
% Produced by harness/assemble_contract.py from module.yaml and claims/*.tex.
% Edit the claim files and regenerate; `--check` fails if this file is stale.
\documentclass[11pt,a4paper]{article}
\pdftrailerid{d9ee3a7955a9b2738d73d78ca0914305}
%% --- 1. Mathematics (order matters: physics after braket wins \ket, \tr) ---
\usepackage{amsmath}
\usepackage{amssymb}
\usepackage{amsfonts}
\usepackage{amsthm}
\usepackage{mathtools}
\usepackage{dsfont}
\usepackage{braket}
\usepackage{physics}
%% --- 2. Page furniture ------------------------------------------------------
\usepackage{xcolor}
\usepackage{array}    % >{...} column prefixes in the wide tables of some claims
\usepackage{enumitem}
\usepackage{graphicx}
\usepackage{booktabs}
\usepackage{longtable}
\usepackage{xltabular}              % front- and back-matter tables across pages
\usepackage{needspace}
\usepackage{xurl}
\usepackage{tcolorbox}
\tcbuselibrary{skins,breakable}
%% --- 3. Links and navigation (hyperref late; bookmark, refcount after it) ---
\usepackage{xr-hyper}   % before hyperref: the contract's labels
\usepackage{hyperref}
\usepackage{bookmark}
\usepackage{refcount}
%% --- 4. Colours: every distinction also has a shape (rule, caps, size, weight)
\definecolor{nnGrey}{gray}{0.40}      % non-normative text
\definecolor{metaGrey}{gray}{0.30}    % the generated meta line
\definecolor{techGrey}{gray}{0.45}    % heading of a technical claim
\definecolor{mainRule}{RGB}{0,110,140}
\definecolor{reqAccent}{RGB}{165,70,0}
\definecolor{linkBlue}{RGB}{0,62,150}
\definecolor{extLink}{RGB}{120,40,140}  % links into another PDF
\hypersetup{colorlinks=true, linkcolor=linkBlue, citecolor=linkBlue,
  urlcolor=linkBlue, filecolor=extLink, bookmarksnumbered=true,
  bookmarksopen=true, bookmarksopenlevel=1, pdfdisplaydoctitle=true}
% A link into another PDF (\href to a .pdf, \hyperref{file}..., xr-hyper) is a
% relative URI, <file>#nameddest=<destination>, in the file-link colour:
% browser viewers (Edge, Chrome) ignore hyperref's default GoToR action.
\makeatletter
\renewcommand*{\hyper@linkfile}[3]{% anchor text, file name, destination
  \begingroup
    \ifx\\#3\\\def\CQ@target{#2}\else\def\CQ@target{#2\hyper@hash nameddest=#3}\fi
    \let\@urlcolor\@filecolor
    \hyper@linkurl{#1}{\CQ@target}%
  \endgroup}
\makeatother
%% --- 5. The mapped Unicode set (harness/contract_navigation.py, UNICODE_TEX)
\DeclareUnicodeCharacter{03B1}{\ensuremath{\alpha}}
\DeclareUnicodeCharacter{03B2}{\ensuremath{\beta}}
\DeclareUnicodeCharacter{03B3}{\ensuremath{\gamma}}
\DeclareUnicodeCharacter{03B4}{\ensuremath{\delta}}
\DeclareUnicodeCharacter{03B5}{\ensuremath{\varepsilon}}
\DeclareUnicodeCharacter{03F5}{\ensuremath{\epsilon}}
\DeclareUnicodeCharacter{03B6}{\ensuremath{\zeta}}
\DeclareUnicodeCharacter{03B7}{\ensuremath{\eta}}
\DeclareUnicodeCharacter{03B8}{\ensuremath{\theta}}
\DeclareUnicodeCharacter{03D1}{\ensuremath{\vartheta}}
\DeclareUnicodeCharacter{03B9}{\ensuremath{\iota}}
\DeclareUnicodeCharacter{03BA}{\ensuremath{\kappa}}
\DeclareUnicodeCharacter{03BB}{\ensuremath{\lambda}}
\DeclareUnicodeCharacter{03BC}{\ensuremath{\mu}}
\DeclareUnicodeCharacter{03BD}{\ensuremath{\nu}}
\DeclareUnicodeCharacter{03BE}{\ensuremath{\xi}}
\DeclareUnicodeCharacter{03C0}{\ensuremath{\pi}}
\DeclareUnicodeCharacter{03C1}{\ensuremath{\rho}}
\DeclareUnicodeCharacter{03F1}{\ensuremath{\varrho}}
\DeclareUnicodeCharacter{03C3}{\ensuremath{\sigma}}
\DeclareUnicodeCharacter{03C2}{\ensuremath{\varsigma}}
\DeclareUnicodeCharacter{03C4}{\ensuremath{\tau}}
\DeclareUnicodeCharacter{03C5}{\ensuremath{\upsilon}}
\DeclareUnicodeCharacter{03C6}{\ensuremath{\varphi}}
\DeclareUnicodeCharacter{03D5}{\ensuremath{\phi}}
\DeclareUnicodeCharacter{03C7}{\ensuremath{\chi}}
\DeclareUnicodeCharacter{03C8}{\ensuremath{\psi}}
\DeclareUnicodeCharacter{03C9}{\ensuremath{\omega}}
\DeclareUnicodeCharacter{0393}{\ensuremath{\Gamma}}
\DeclareUnicodeCharacter{0394}{\ensuremath{\Delta}}
\DeclareUnicodeCharacter{0398}{\ensuremath{\Theta}}
\DeclareUnicodeCharacter{039B}{\ensuremath{\Lambda}}
\DeclareUnicodeCharacter{039E}{\ensuremath{\Xi}}
\DeclareUnicodeCharacter{03A0}{\ensuremath{\Pi}}
\DeclareUnicodeCharacter{03A3}{\ensuremath{\Sigma}}
\DeclareUnicodeCharacter{03A5}{\ensuremath{\Upsilon}}
\DeclareUnicodeCharacter{03A6}{\ensuremath{\Phi}}
\DeclareUnicodeCharacter{03A8}{\ensuremath{\Psi}}
\DeclareUnicodeCharacter{03A9}{\ensuremath{\Omega}}
\DeclareUnicodeCharacter{00E1}{\'a}
\DeclareUnicodeCharacter{00E0}{\`a}
\DeclareUnicodeCharacter{00E2}{\^a}
\DeclareUnicodeCharacter{00E4}{\"a}
\DeclareUnicodeCharacter{00E5}{\aa{}}
\DeclareUnicodeCharacter{00E9}{\'e}
\DeclareUnicodeCharacter{00E8}{\`e}
\DeclareUnicodeCharacter{00EA}{\^e}
\DeclareUnicodeCharacter{00EB}{\"e}
\DeclareUnicodeCharacter{00ED}{\'{\i}}
\DeclareUnicodeCharacter{00EF}{\"{\i}}
\DeclareUnicodeCharacter{00F3}{\'o}
\DeclareUnicodeCharacter{00F2}{\`o}
\DeclareUnicodeCharacter{00F4}{\^o}
\DeclareUnicodeCharacter{00F6}{\"o}
\DeclareUnicodeCharacter{00F8}{\o{}}
\DeclareUnicodeCharacter{00FA}{\'u}
\DeclareUnicodeCharacter{00F9}{\`u}
\DeclareUnicodeCharacter{00FB}{\^u}
\DeclareUnicodeCharacter{00FC}{\"u}
\DeclareUnicodeCharacter{00F1}{\~n}
\DeclareUnicodeCharacter{00E7}{\c{c}}
\DeclareUnicodeCharacter{00DF}{\ss{}}
\DeclareUnicodeCharacter{00C4}{\"A}
\DeclareUnicodeCharacter{00C5}{\AA{}}
\DeclareUnicodeCharacter{00C9}{\'E}
\DeclareUnicodeCharacter{00D6}{\"O}
\DeclareUnicodeCharacter{00DC}{\"U}
\DeclareUnicodeCharacter{2013}{\textendash{}}
\DeclareUnicodeCharacter{2014}{\textemdash{}}
\DeclareUnicodeCharacter{2018}{`}
\DeclareUnicodeCharacter{2019}{'}
\DeclareUnicodeCharacter{201C}{``}
\DeclareUnicodeCharacter{201D}{''}
\DeclareUnicodeCharacter{2026}{\ldots{}}
\DeclareUnicodeCharacter{00A0}{~}
\DeclareUnicodeCharacter{00A7}{\S{}}
\DeclareUnicodeCharacter{2264}{\ensuremath{\le}}
\DeclareUnicodeCharacter{2265}{\ensuremath{\ge}}
\DeclareUnicodeCharacter{2260}{\ensuremath{\neq}}
\DeclareUnicodeCharacter{2248}{\ensuremath{\approx}}
\DeclareUnicodeCharacter{2261}{\ensuremath{\equiv}}
\DeclareUnicodeCharacter{223C}{\ensuremath{\sim}}
\DeclareUnicodeCharacter{2243}{\ensuremath{\simeq}}
\DeclareUnicodeCharacter{2245}{\ensuremath{\cong}}
\DeclareUnicodeCharacter{221D}{\ensuremath{\propto}}
\DeclareUnicodeCharacter{226A}{\ensuremath{\ll}}
\DeclareUnicodeCharacter{226B}{\ensuremath{\gg}}
\DeclareUnicodeCharacter{2192}{\ensuremath{\to}}
\DeclareUnicodeCharacter{2190}{\ensuremath{\leftarrow}}
\DeclareUnicodeCharacter{2194}{\ensuremath{\leftrightarrow}}
\DeclareUnicodeCharacter{21D2}{\ensuremath{\Rightarrow}}
\DeclareUnicodeCharacter{21D0}{\ensuremath{\Leftarrow}}
\DeclareUnicodeCharacter{21D4}{\ensuremath{\Leftrightarrow}}
\DeclareUnicodeCharacter{21A6}{\ensuremath{\mapsto}}
\DeclareUnicodeCharacter{2208}{\ensuremath{\in}}
\DeclareUnicodeCharacter{2209}{\ensuremath{\notin}}
\DeclareUnicodeCharacter{2286}{\ensuremath{\subseteq}}
\DeclareUnicodeCharacter{2282}{\ensuremath{\subset}}
\DeclareUnicodeCharacter{2287}{\ensuremath{\supseteq}}
\DeclareUnicodeCharacter{2283}{\ensuremath{\supset}}
\DeclareUnicodeCharacter{2229}{\ensuremath{\cap}}
\DeclareUnicodeCharacter{222A}{\ensuremath{\cup}}
\DeclareUnicodeCharacter{2216}{\ensuremath{\setminus}}
\DeclareUnicodeCharacter{2205}{\ensuremath{\emptyset}}
\DeclareUnicodeCharacter{2200}{\ensuremath{\forall}}
\DeclareUnicodeCharacter{2203}{\ensuremath{\exists}}
\DeclareUnicodeCharacter{2115}{\ensuremath{\mathbb{N}}}
\DeclareUnicodeCharacter{2124}{\ensuremath{\mathbb{Z}}}
\DeclareUnicodeCharacter{211A}{\ensuremath{\mathbb{Q}}}
\DeclareUnicodeCharacter{211D}{\ensuremath{\mathbb{R}}}
\DeclareUnicodeCharacter{2102}{\ensuremath{\mathbb{C}}}
\DeclareUnicodeCharacter{00B7}{\ensuremath{\cdot}}
\DeclareUnicodeCharacter{00D7}{\ensuremath{\times}}
\DeclareUnicodeCharacter{2212}{\ensuremath{-}}
\DeclareUnicodeCharacter{00B1}{\ensuremath{\pm}}
\DeclareUnicodeCharacter{2213}{\ensuremath{\mp}}
\DeclareUnicodeCharacter{221A}{\ensuremath{\surd}}
\DeclareUnicodeCharacter{2211}{\ensuremath{\sum}}
\DeclareUnicodeCharacter{220F}{\ensuremath{\prod}}
\DeclareUnicodeCharacter{222B}{\ensuremath{\int}}
\DeclareUnicodeCharacter{221E}{\ensuremath{\infty}}
\DeclareUnicodeCharacter{2202}{\ensuremath{\partial}}
\DeclareUnicodeCharacter{2207}{\ensuremath{\nabla}}
\DeclareUnicodeCharacter{2297}{\ensuremath{\otimes}}
\DeclareUnicodeCharacter{2295}{\ensuremath{\oplus}}
\DeclareUnicodeCharacter{2218}{\ensuremath{\circ}}
\DeclareUnicodeCharacter{2308}{\ensuremath{\lceil}}
\DeclareUnicodeCharacter{2309}{\ensuremath{\rceil}}
\DeclareUnicodeCharacter{230A}{\ensuremath{\lfloor}}
\DeclareUnicodeCharacter{230B}{\ensuremath{\rfloor}}
\DeclareUnicodeCharacter{27E8}{\ensuremath{\langle}}
\DeclareUnicodeCharacter{27E9}{\ensuremath{\rangle}}
\DeclareUnicodeCharacter{2016}{\ensuremath{\|}}
\DeclareUnicodeCharacter{2113}{\ensuremath{\ell}}
\DeclareUnicodeCharacter{00B0}{\ensuremath{^{\circ}}}
\DeclareUnicodeCharacter{00B2}{\ensuremath{^{2}}}
\DeclareUnicodeCharacter{00B3}{\ensuremath{^{3}}}
\DeclareUnicodeCharacter{00B9}{\ensuremath{^{1}}}
\DeclareUnicodeCharacter{2080}{\ensuremath{_{0}}}
\DeclareUnicodeCharacter{2081}{\ensuremath{_{1}}}
\DeclareUnicodeCharacter{2082}{\ensuremath{_{2}}}
\DeclareUnicodeCharacter{207F}{\ensuremath{^{n}}}
\DeclareUnicodeCharacter{2096}{\ensuremath{_{k}}}
%% --- 6. Numbering: one counter per section -----------------------
\numberwithin{equation}{section}
\setcounter{tocdepth}{1}
%% --- 7. The project macro set, then the claim environments -----------------
% spec/macros.tex, read relative to this module's directory (latexmk -cd).
% =============================================================================
%  spec/macros.tex -- the one macro set every generated contract inputs.
%
%  Two parts: part A, the macros spec/notation.md catalogues, and part B,
%  the macros local to one module, grouped by module.
%  harness/assemble_contract.py emits % =============================================================================
%  spec/macros.tex -- the one macro set every generated contract inputs.
%
%  Two parts: part A, the macros spec/notation.md catalogues, and part B,
%  the macros local to one module, grouped by module.
%  harness/assemble_contract.py emits % =============================================================================
%  spec/macros.tex -- the one macro set every generated contract inputs.
%
%  Two parts: part A, the macros spec/notation.md catalogues, and part B,
%  the macros local to one module, grouped by module.
%  harness/assemble_contract.py emits \input{../../macros} for it, so the
%  file is read from spec/modules/<id>/.
%
%  RULES
%    * A macro that spec/notation.md catalogues is a verbatim copy of it
%      (part A). Changing one is a notation change: bump spec/notation.md.
%    * Part B holds the module-local macros. A macro defined by two modules
%      with the same expansion appears once, under the first module (in the
%      order below) that defines it. Two macros with the same name and
%      DIFFERENT expansions are both kept under a module prefix, and the
%      claims write the prefixed names:
%        \Chain  -> \redChain (reductions, \mathfrak{C}: the reduction chain)
%                   \pcChain  (protocol-compiler, \mathfrak{R}: the two-segment
%                             chain; \Chsrc and \Chrec are built on it)
%        \Nmsm   -> \redNmsm  (reductions, \mathcal{N}: a POVM measurement
%                             channel)
%                   \pcNmsm   (\mathcal{N}_{\mathsf A}: Alice's measurement
%                             channel, written by
%                             transmitter/thm:certificate-sufficiency,
%                             transmitter/rem:xreg and
%                             reductions/lem:source-transfer)
%    * Deliberately omitted: a \provenance{tag}{details} macro. Provenance
%      is a header field of every claim file (spec/claim-format.md section
%      3); a claim body that calls such a macro fails to compile, which is
%      the intended signal.
%
%  REQUIRED PACKAGES, loaded by the assembler's preamble before this file:
%    amsmath, amssymb, amsfonts, mathtools, dsfont, braket, physics, xcolor
%    (braket before physics, so that physics wins for \ket, \bra, \tr, ...).
%
%  This file contains only definitions: no \usepackage, no document body.
% =============================================================================


% =============================================================================
%  PART A -- notation macros, the ones spec/notation.md catalogues, each a
%  verbatim copy of its catalogue row; grouped under numbered headings.
% =============================================================================

% -----------------------------------------------------------------------------
% 1. State spaces and operator sets (a "system label" -> spaces of operators)
% -----------------------------------------------------------------------------
\newcommand{\Hs}[1]{\mathcal{H}(#1)}        % Hilbert space of a system label
\newcommand{\Lin}[1]{\mathcal{L}(#1)}       % bounded linear operators on H(#1)
\newcommand{\Linn}[2]{\mathcal{L}(#1,#2)}   % linear maps H(#1) -> H(#2)
\newcommand{\Pos}[1]{\mathrm{Pos}(#1)}      % positive semidefinite operators on H(#1)
\newcommand{\Density}{\mathrm{D}}           % density operators: \Density(A) = D(A)
\newcommand{\Seq}[1]{S_{=}(#1)}             % normalized states (trace = 1)
\newcommand{\Sleq}[1]{S_{\leq}(#1)}         % subnormalized states (0 <= trace <= 1)
\newcommand{\id}{\mathds{1}}                % identity operator
\newcommand{\idon}[1]{\mathds{1}_{#1}}      % identity on a named system, e.g. \idon{C}

% DEPRECATED: applies \mathcal to both arguments itself, so it takes bare
% letters, not system labels. Prefer \Linn; do not use in new contracts.
\providecommand{\Lop}[2]{\mathcal{L}(\mathcal{#1}, \mathcal{#2})}


% -----------------------------------------------------------------------------
% 2. Composite systems
% -----------------------------------------------------------------------------
\newcommand{\dsum}{\oplus}                  % direct sum of system labels
% Tensor product: use \otimes directly.
% n-fold direct sum of an operator space: \Lin{A}^{\oplus n}


% -----------------------------------------------------------------------------
% 3. Events, conditioning and classical registers (for STATES)
%    Convention:  \wedge  -> subnormalized "and-event" component, e.g. \rho_{A\wedge x}
%                 |        -> normalized conditional state,        e.g. \rho_{A|x}
% (Write these directly in subscripts; no extra macros needed.)
% -----------------------------------------------------------------------------


% -----------------------------------------------------------------------------
% 4. Photon-number subspaces / direct-sum blocks
%    Convention:  operators carry the restriction "on top", in brackets;
%                 states carry it as a \wedge subscript (see Sec. 3).
% -----------------------------------------------------------------------------
\newcommand{\Pnum}[1]{P_{#1}}               % photon-number projector (system from context)
\newcommand{\Pnumon}[2]{P_{#1}^{#2}}        % with explicit system superscript, e.g. P_n^C
\newcommand{\blk}[2]{#1^{(#2)}}             % operator block, e.g. \blk{M}{\leqN} = M^{(<=N)}
\newcommand{\leqN}{\leq N}                  % use as \blk{M}{\leqN}
\newcommand{\gtN}{> N}                      % use as \blk{M}{\gtN}
\newcommand{\eqn}{= n}                      % use as \blk{M}{\eqn}


% -----------------------------------------------------------------------------
% 5. System labels used in the protocol
% -----------------------------------------------------------------------------
\newcommand{\Chan}{C}                       % quantum-channel (Eve) input to Bob
\newcommand{\Prb}{P}                        % probe input
\newcommand{\Flag}{{\mathrm{Flg}}}          % flag system (EXISTS ONLY AFTER squashing)
\newcommand{\Loss}{{\mathrm{Lss}}}          % loss system
\newcommand{\Vc}{{\mathrm{Vcm}}}            % vacuum system
\newcommand{\Nout}{N_{\mathrm{out}}}
\newcommand{\Nin}{N_{\mathrm{in}}}


% -----------------------------------------------------------------------------
% 6. Alphabets
% -----------------------------------------------------------------------------
\newcommand{\Yset}{\mathcal{Y}}             % Bob's (coarse-grained) detection alphabet
\newcommand{\Yraw}{\mathcal{Y}_{\mathrm{raw}}}% raw click patterns {0,1}^4
\newcommand{\Yscset}{\mathcal{Y}_{\singlec}}% single-click outcomes
\newcommand{\Ymcset}{\mathcal{Y}_{\mc}}     % multi-click outcomes
\newcommand{\Tset}{\mathcal{T}_{\mathrm{set}}}% private setting set (the probe settings in the optical instance)
\newcommand{\Xset}{\mathcal{X}}             % Alice's alphabet
\newcommand{\Dset}{\{H,V,+,-\}}             % polarization / detector labels
\newcommand{\pass}{\top}                    % pass sifting


% -----------------------------------------------------------------------------
% 7. Mode / outcome abbreviations
% -----------------------------------------------------------------------------
\newcommand{\on}{{\mathrm{on}}}
\newcommand{\off}{{\mathrm{off}}}
\newcommand{\dc}{{\mathrm{dc}}}             % dark count
\newcommand{\sq}{\mathrm{sq}}               % squashed
\newcommand{\mc}{\mathrm{mc}}               % multi-click
\newcommand{\nc}{\mathrm{nc}}               % no-click
\newcommand{\singlec}{\mathrm{sc}}          % single-click
\newcommand{\vac}{\mathrm{vac}}             % vacuum
\newcommand{\kout}{{k_\mathrm{out}}}
\newcommand{\kin}{{k_\mathrm{in}}}


% -----------------------------------------------------------------------------
% 8. Maps, channels and miscellaneous operators
% -----------------------------------------------------------------------------
\newcommand{\Prob}{\mathrm{Prob}}           % probability (use \Prob[\cdot] or \Prob(\cdot))
\newcommand{\lmin}{\lambda_{\min}}          % minimum eigenvalue
\newcommand{\Wflag}{W_{\Flag}}              % flag-space weight
\newcommand{\transpose}{\mathsf{T}}
\newcommand{\innerp}[2]{\left\langle #1 \vert #2 \right\rangle}
\renewcommand{\vec}[1]{\boldsymbol{#1}}     % NB: overrides the LaTeX arrow accent


% -----------------------------------------------------------------------------
% 9. Key-rate numerics
% -----------------------------------------------------------------------------
\newcommand{\cG}{\mathcal{G}}               % protocol map
\newcommand{\cZ}{\mathcal{Z}}               % pinching map
\newcommand{\Amap}{\mathcal{A}_{\mathrm{ann}}} % announcement + measurement map
\newcommand{\Amapnb}{\mathcal{A}_{\mathrm{ann}|\mathrm{nb}}}
\newcommand{\simpl}{}                       % empty suffix for simplified block maps
\newcommand{\acSet}{\mathcal{S}_{\mathrm{acc}}} % acceptance set
\newcommand{\testE}{\mathcal{E}_{\mathrm{test}}} % test-statistics map
\newcommand{\ecLeak}{\mathrm{leak}^{\mathrm{EC}}_{\mathrm{obs}}}
\newcommand{\pgen}{p_{\mathrm{gen}}}
\newcommand{\askeyr}{R_{\infty}}             % asymptotic key rate
\newcommand{\ppass}{p_{\mathrm{pass}}}       % sifting pass probability


% -----------------------------------------------------------------------------
% 10. Finite-size analysis / Renyi entropies
% -----------------------------------------------------------------------------

% --- Round types: test vs. key generation (register T, alphabet \Fset) --------
\newcommand{\tst}{\mathrm{test}}            % test round
\newcommand{\gen}{\mathrm{gen}}             % key-generation round
\newcommand{\Fset}{\mathcal{T}}             % round-type alphabet {test, gen} (not \Tset)
\newcommand{\Chat}{\widehat{C}}             % acceptance-test register
\newcommand{\Chatset}{\widehat{\mathcal{C}}}% its alphabet

% --- Basis / key extraction functions ----------------------------------------
\newcommand{\bA}{b_{A}}                     % Alice's basis-extraction function
\newcommand{\bB}{b_{B}}                     % Bob's basis-extraction function
\newcommand{\kB}{k_{B}}                     % Bob's key-extraction function

% --- Maps --------------------------------------------------------------------
\newcommand{\Mprot}{\mathcal{M}}            % single-round protocol map
\newcommand{\Phimap}{\boldsymbol{\Phi}}     % test-statistics map
\newcommand{\Pdeph}{\mathcal{P}}            % dephasing (pinching) map on \Blreg

% --- Isometries --------------------------------------------------------------
\newcommand{\Vmsm}{V_{\mathrm{msm}}}        % measurement isometry
\newcommand{\Vann}{V_{\mathrm{ann}}}        % announcement isometry
\newcommand{\VI}{V_{I}}                     % Eve's copy of the announcements
\newcommand{\VChat}{V_{\Chat}}              % acceptance-data isometry
\newcommand{\VR}{V_{R}}                     % sifting + key map -> raw key R
\newcommand{\VZ}{V_{Z}}                     % raw key -> key register Z
\newcommand{\Wprot}{W}                      % full protocol isometry

% --- Renyi entropies ---------------------------------------------------------
% SECURITY-CRITICAL: see spec/notation.md §4. Never change an arrow to
% make an inequality work out; never convert silently between the sandwiched
% (tilde) and Petz (overline) families.
\newcommand{\Hsandup}{\widetilde{H}^{\uparrow}_{\alpha}}   % sandwiched, up
\newcommand{\Hsanddown}{\widetilde{H}^{\downarrow}_{\alpha}}% sandwiched, down
\newcommand{\Hpetzup}{\overline{H}^{\uparrow}_{\alpha}}    % Petz, up
\newcommand{\Hpetzdown}{\overline{H}^{\downarrow}_{\alpha}}% Petz, down
\newcommand{\hup}{h^{\uparrow}_{\alpha}}                   % single-round rate
\newcommand{\hqkd}{h^{\uparrow,\mathrm{QKD}}_{\alpha}}     % single-round QKD rate
\newcommand{\hlow}{\underline{h}_{\alpha}}                 % lower bound on \hqkd
% --- spec/notation.md 3.5 and 4.5 -------------------------------------------
% \hlb  is the first-order linearization lower bound on \hlow: the headline
%       key rate (the `first-order` tier of foundations/def:tiers). NOT \hlow
%       itself.
% \tnum is the instance's one declared numerical tolerance. Its
%       near-collision with the declared marginal \tau_A is disclosed in
%       spec/notation.md 3.5; bare \tau is reserved and denotes neither.
\newcommand{\hlb}{\underline{h}^{\mathrm{FO}}_{\alpha}}    % first-order lower bound on \hlow
\newcommand{\tnum}{\tau_{\mathrm{num}}}                    % declared numerical tolerance


% =============================================================================
%  PART B -- module-local macros, grouped by module. Deprecated aliases,
%  defined so that a text written with them compiles, used in no new contract: \Lop,
%  \testE, \Pdeph, \pass, \ecLeak (spec/notation.md section 3.9).
% =============================================================================

% --- meat-core -----------------------------------------------------------------
\newcommand{\SecSet}{\mathbb{S}}            % sector index set of Q
\newcommand{\Qsys}{Q}                       % declared sector-carrying measured system
\newcommand{\Vset}{\mathcal{V}}             % declared joint outcome alphabet of the Q-POVM
\newcommand{\CAset}{\mathcal{C}_{A}}        % Alice announcement alphabet
\newcommand{\CBset}{\mathcal{C}_{B}}        % Bob announcement alphabet
\newcommand{\Party}{\mathfrak{P}}           % declared party partition
\newcommand{\Oset}[1]{\mathcal{O}_{#1}}     % readout alphabet of party #1
\newcommand{\rd}[1]{\pi_{#1}}               % readout map of party #1
\newcommand{\Sadm}{\mathcal{S}_{A}}         % admissible Alice subspace
\newcommand{\nomarg}{\mathsf{none}}         % "no trusted marginal declared"
\newcommand{\Zset}{\mathcal{Z}_{\perp}}     % raw-key alphabet incl. \perp
\newcommand{\Zval}{\mathcal{Z}}             % key alphabet without \perp
\newcommand{\Etil}{\widetilde{E}}           % single-round purifying register
\newcommand{\Ehat}{\widehat{E}}             % Eve's n-round register
\newcommand{\OmAT}{\Omega_{\mathrm{AT}}}    % acceptance event
\newcommand{\epa}{\varepsilon_{\mathrm{PA}}}
\newcommand{\eev}{\varepsilon_{\mathrm{EV}}}
\newcommand{\lEC}{\lambda_{\mathrm{EC}}}
\newcommand{\Decl}{\mathfrak{D}}            % protocol declaration
\newcommand{\CList}{\mathfrak{L}}           % declared constraint list
\newcommand{\PinS}{\mathcal{P}_{\SecSet}}   % sector pinching on Q
\newcommand{\Feas}{\mathcal{W}}             % feasible set of block states
\newcommand{\Chkey}{\Chatset_{\mathrm{key}}}% key-live acceptance values
\newcommand{\Kann}{\mathcal{K}}             % announcement-class alphabet
\newcommand{\hold}{\mathsf{h}}              % declared key holder (A or B)
\newcommand{\sfD}{D}                        % classical KL divergence (base 2)
\newcommand{\Dtwo}{D}                       % quantum relative entropy (base 2)
\newcommand{\Dsand}{\widetilde{D}_{\alpha}} % sandwiched Renyi divergence
\newcommand{\idmap}{\mathrm{id}}            % identity channel
\newcommand{\Aev}{\mathcal{A}}              % Eve's per-round attack channel
\newcommand{\AttSet}{\mathfrak{A}}          % licensed attack-channel set
\newcommand{\Chacc}{\Chatset_{\mathrm{acc}}}% acceptance support of \acSet
\newcommand{\epsstar}{\varepsilon^{\star}}      % distance to an exact POVM
\newcommand{\dfH}{\Delta_{\mathrm{herm}}}       % Hermiticity defect
\newcommand{\dfP}{\Delta_{\mathrm{pos}}}        % negative part
\newcommand{\dfC}{\Delta_{\mathrm{cpl}}}        % completeness defect
\newcommand{\lmoq}{g}                       % acceptance-side gradient block
\newcommand{\lmoJ}{G}                       % feasible-side gradient block
\newcommand{\bacc}{\beta_{\mathrm{acc}}}    % certified acceptance-side minimum
\newcommand{\Itset}{\mathfrak{I}}           % finite index set of iterates (C22)

% --- reductions ----------------------------------------------------------------
\newcommand{\Bsys}{B}                        % physical (pre-squashing) system
\newcommand{\Tcl}[1]{\mathcal{T}(#1)}        % trace-class operators on H(#1)
\newcommand{\Ins}{\mathcal{I}}               % generalized-instrument branch map
\newcommand{\redNmsm}{\mathcal{N}}              % POVM measurement channel
\newcommand{\Eprot}{\mathcal{E}}             % n-round protocol channel
\newcommand{\Grnd}{\mathsf{G}}               % one round map of a protocol
\newcommand{\Kfin}{\mathsf{K}}               % end-of-protocol processing
\newcommand{\Crit}{\mathsf{sec}}             % value of a security criterion
\newcommand{\Pflag}{\Pi_{\Flag}}             % flag-sector projector on Q
\newcommand{\Pleq}{\Pi_{\leqN}}              % preserved-block projector
\newcommand{\Pgt}{\Pi_{\gtN}}                % above-cutoff projector
\newcommand{\Sph}{S^{\mathrm{ph}}}           % family of physical inputs
\newcommand{\Srch}{S^{\mathrm{rch}}}         % reachable squashed states
\newcommand{\Ssq}{S^{\sq}}                   % squashed admissible set
\newcommand{\Ffss}{F^{\mathrm{FSS}}}         % flag-state-squasher target
\newcommand{\Fwp}{F^{\mathrm{WP}}}           % multi-bound WPFSS target
\newcommand{\Mwp}{M^{\mathrm{WP}}}           % renormalized above-cutoff POVM
\newcommand{\Mdec}{{M'}^{\mathrm{ph}}}       % declared physical family on B
\newcommand{\Ppvm}{P^{\Flag}}                % flag PVM
\newcommand{\Lmp}{{\Lambda_{\mathrm{mp}}}}   % measure-and-prepare channel
\newcommand{\Lwp}{{\Lambda^{\mathrm{WP}}}}   % multi-bound WPFSS channel
\newcommand{\Lfss}{{\Lambda^{\mathrm{FSS}}}} % flag-state-squasher channel
\newcommand{\lmax}{\lambda_{\max}}           % maximum eigenvalue
\newcommand{\Lnoise}{\widetilde{\Lambda}}    % flag-noise channel
\newcommand{\cg}{\varpi}                     % declared coarse-graining map
\newcommand{\MCK}{\mathsf{MC}}               % multiclick raw-pattern set
\newcommand{\redChain}{\mathfrak{C}}            % reduction chain
\newcommand{\Lnk}{\ell}                      % a single typed link
\newcommand{\dsq}{\preceq}                      % the squashing relation
\newcommand{\opn}[1]{\left\lVert #1\right\rVert}           % operator (spectral) norm
\newcommand{\mxn}[1]{\left\lVert #1\right\rVert_{\max}}    % entrywise max modulus
\newcommand{\Fld}{\mathbb{F}}                   % exact ordered substrate field
\newcommand{\Fldi}{\mathbb{F}_{i}}              % its Gaussian extension F + iF
\newcommand{\lat}[1]{\big\lfloor #1\big\rfloor_{\kappa}} % round DOWN onto the lattice

% --- protocol-compiler ---------------------------------------------------------
\newcommand{\Spec}{\mathfrak{S}}             % protocol specification (compiler input)
\newcommand{\pcChain}{\mathfrak{R}}            % reduction chain (two typed segments)
\newcommand{\Chsrc}{\pcChain_{\mathrm{src}}}   % its Alice-side source-map segment
\newcommand{\Chrec}{\pcChain_{\mathrm{rec}}}   % its receiver-side squashing segment
\newcommand{\rvirt}{\rho^{\mathrm{v}}}       % virtual (PRE-source-map) signal state
\newcommand{\Fph}{F^{\mathrm{ph}}}           % physical (pre-reduction) POVM element
\newcommand{\Vph}{\mathcal{V}^{\mathrm{ph}}} % physical outcome alphabet
\newcommand{\Tpub}{\mathcal{T}_{\mathrm{pub}}}% public-choice alphabet
\newcommand{\Uset}{\mathcal{U}}              % residual outcome alphabet
\newcommand{\idxm}{\mathsf{idx}}             % label-to-index map
\newcommand{\Asig}{A'}                       % transmitted (signal) system -- the REAL one
\newcommand{\Avirt}{A''}                     % virtual transmitted system (source descriptor of Xi_1)
\newcommand{\Abar}{\bar{A}}                  % Alice's measured register
\newcommand{\Ahat}{\hat{A}}                  % Alice's shield register
\newcommand{\pcNmsm}{\mathcal{N}_{\mathsf A}}  % Alice's measurement channel (see the RULES block)
\newcommand{\Phip}{\Phi^{+}}                 % maximally entangled vector label
\newcommand{\Smap}{\mathfrak{s}}             % the emitted sector map (fraktur: a compile-time object)

% --- receiver ------------------------------------------------------------------
\newcommand{\Modes}{\mathbb{M}}              % declared ordered mode set
\newcommand{\Inp}{\mathbb{I}}                % declared input modes
\newcommand{\InpC}{\mathbb{I}_{\Chan}}       % channel input modes
\newcommand{\InpS}{\mathbb{I}_{\Prb}}        % probe (source) input modes
\newcommand{\Env}{\mathbb{E}}                % dilation (environment) modes
\newcommand{\Port}[1]{\mathbb{P}_{#1}}       % mode set of detector port #1
\newcommand{\Emd}[1]{E_{#1}}                 % coordinate projector on modes #1
\newcommand{\Gmat}{G}                        % instrument (transfer) matrix
\newcommand{\Fock}{\mathfrak{F}}             % Fock / second-quantization functor
\newcommand{\Symn}[1]{\mathrm{Sym}^{#1}}     % symmetric tensor power
\newcommand{\normord}[1]{{:}\mspace{2mu}#1\mspace{2mu}{:}}   % normal ordering
\newcommand{\rk}{\operatorname{rank}}        % rank (physics.sty owns \rank)
\newcommand{\SCK}{\mathsf{SC}}               % single-click pattern set
\newcommand{\Recv}{\mathfrak{R}}             % receiver instrument tuple
\newcommand{\Cat}{\mathfrak{C}}              % component catalogue
\newcommand{\Grph}{\mathfrak{E}}             % declared element graph
\newcommand{\defic}{\mathbb{L}}              % loss deficiency operator
\newcommand{\Pvac}[1]{\Pi_{\vac}(#1)}        % joint no-click element on ports #1
\newcommand{\pmc}{p_{\mc}}                   % uniform high-photon mc weight
\newcommand{\Bnd}{\beta}                     % the scalar bound function beta(n)
\newcommand{\Rot}{R}                         % planar rotation matrix
\newcommand{\Wplate}{W}                      % retarder (wave-plate) matrix
\newcommand{\Bsp}{B}                         % beamsplitter core matrix

% --- statistics ----------------------------------------------------------------
\newcommand{\honest}{p^{\mathrm{hon}}}          % honest per-round distribution
\newcommand{\ecomp}{\varepsilon_{\mathrm{comp}}}% completeness budget
\newcommand{\simplex}[1]{\Delta(#1)}            % probability simplex
\newcommand{\Chatcon}{\Chatset_{\mathrm{con}}}  % constrained outcomes
\newcommand{\Ireg}[3]{\mathcal{I}_{#1}(#2,#3)}  % regularized incomplete beta
\newcommand{\Ncount}[1]{\mathsf{N}_{#1}}        % per-outcome count
\newcommand{\qlo}{\mathbf{q}^{-}}               % box lower bounds
\newcommand{\qhi}{\mathbf{q}^{+}}               % box upper bounds
 for it, so the
%  file is read from spec/modules/<id>/.
%
%  RULES
%    * A macro that spec/notation.md catalogues is a verbatim copy of it
%      (part A). Changing one is a notation change: bump spec/notation.md.
%    * Part B holds the module-local macros. A macro defined by two modules
%      with the same expansion appears once, under the first module (in the
%      order below) that defines it. Two macros with the same name and
%      DIFFERENT expansions are both kept under a module prefix, and the
%      claims write the prefixed names:
%        \Chain  -> \redChain (reductions, \mathfrak{C}: the reduction chain)
%                   \pcChain  (protocol-compiler, \mathfrak{R}: the two-segment
%                             chain; \Chsrc and \Chrec are built on it)
%        \Nmsm   -> \redNmsm  (reductions, \mathcal{N}: a POVM measurement
%                             channel)
%                   \pcNmsm   (\mathcal{N}_{\mathsf A}: Alice's measurement
%                             channel, written by
%                             transmitter/thm:certificate-sufficiency,
%                             transmitter/rem:xreg and
%                             reductions/lem:source-transfer)
%    * Deliberately omitted: a \provenance{tag}{details} macro. Provenance
%      is a header field of every claim file (spec/claim-format.md section
%      3); a claim body that calls such a macro fails to compile, which is
%      the intended signal.
%
%  REQUIRED PACKAGES, loaded by the assembler's preamble before this file:
%    amsmath, amssymb, amsfonts, mathtools, dsfont, braket, physics, xcolor
%    (braket before physics, so that physics wins for \ket, \bra, \tr, ...).
%
%  This file contains only definitions: no \usepackage, no document body.
% =============================================================================


% =============================================================================
%  PART A -- notation macros, the ones spec/notation.md catalogues, each a
%  verbatim copy of its catalogue row; grouped under numbered headings.
% =============================================================================

% -----------------------------------------------------------------------------
% 1. State spaces and operator sets (a "system label" -> spaces of operators)
% -----------------------------------------------------------------------------
\newcommand{\Hs}[1]{\mathcal{H}(#1)}        % Hilbert space of a system label
\newcommand{\Lin}[1]{\mathcal{L}(#1)}       % bounded linear operators on H(#1)
\newcommand{\Linn}[2]{\mathcal{L}(#1,#2)}   % linear maps H(#1) -> H(#2)
\newcommand{\Pos}[1]{\mathrm{Pos}(#1)}      % positive semidefinite operators on H(#1)
\newcommand{\Density}{\mathrm{D}}           % density operators: \Density(A) = D(A)
\newcommand{\Seq}[1]{S_{=}(#1)}             % normalized states (trace = 1)
\newcommand{\Sleq}[1]{S_{\leq}(#1)}         % subnormalized states (0 <= trace <= 1)
\newcommand{\id}{\mathds{1}}                % identity operator
\newcommand{\idon}[1]{\mathds{1}_{#1}}      % identity on a named system, e.g. \idon{C}

% DEPRECATED: applies \mathcal to both arguments itself, so it takes bare
% letters, not system labels. Prefer \Linn; do not use in new contracts.
\providecommand{\Lop}[2]{\mathcal{L}(\mathcal{#1}, \mathcal{#2})}


% -----------------------------------------------------------------------------
% 2. Composite systems
% -----------------------------------------------------------------------------
\newcommand{\dsum}{\oplus}                  % direct sum of system labels
% Tensor product: use \otimes directly.
% n-fold direct sum of an operator space: \Lin{A}^{\oplus n}


% -----------------------------------------------------------------------------
% 3. Events, conditioning and classical registers (for STATES)
%    Convention:  \wedge  -> subnormalized "and-event" component, e.g. \rho_{A\wedge x}
%                 |        -> normalized conditional state,        e.g. \rho_{A|x}
% (Write these directly in subscripts; no extra macros needed.)
% -----------------------------------------------------------------------------


% -----------------------------------------------------------------------------
% 4. Photon-number subspaces / direct-sum blocks
%    Convention:  operators carry the restriction "on top", in brackets;
%                 states carry it as a \wedge subscript (see Sec. 3).
% -----------------------------------------------------------------------------
\newcommand{\Pnum}[1]{P_{#1}}               % photon-number projector (system from context)
\newcommand{\Pnumon}[2]{P_{#1}^{#2}}        % with explicit system superscript, e.g. P_n^C
\newcommand{\blk}[2]{#1^{(#2)}}             % operator block, e.g. \blk{M}{\leqN} = M^{(<=N)}
\newcommand{\leqN}{\leq N}                  % use as \blk{M}{\leqN}
\newcommand{\gtN}{> N}                      % use as \blk{M}{\gtN}
\newcommand{\eqn}{= n}                      % use as \blk{M}{\eqn}


% -----------------------------------------------------------------------------
% 5. System labels used in the protocol
% -----------------------------------------------------------------------------
\newcommand{\Chan}{C}                       % quantum-channel (Eve) input to Bob
\newcommand{\Prb}{P}                        % probe input
\newcommand{\Flag}{{\mathrm{Flg}}}          % flag system (EXISTS ONLY AFTER squashing)
\newcommand{\Loss}{{\mathrm{Lss}}}          % loss system
\newcommand{\Vc}{{\mathrm{Vcm}}}            % vacuum system
\newcommand{\Nout}{N_{\mathrm{out}}}
\newcommand{\Nin}{N_{\mathrm{in}}}


% -----------------------------------------------------------------------------
% 6. Alphabets
% -----------------------------------------------------------------------------
\newcommand{\Yset}{\mathcal{Y}}             % Bob's (coarse-grained) detection alphabet
\newcommand{\Yraw}{\mathcal{Y}_{\mathrm{raw}}}% raw click patterns {0,1}^4
\newcommand{\Yscset}{\mathcal{Y}_{\singlec}}% single-click outcomes
\newcommand{\Ymcset}{\mathcal{Y}_{\mc}}     % multi-click outcomes
\newcommand{\Tset}{\mathcal{T}_{\mathrm{set}}}% private setting set (the probe settings in the optical instance)
\newcommand{\Xset}{\mathcal{X}}             % Alice's alphabet
\newcommand{\Dset}{\{H,V,+,-\}}             % polarization / detector labels
\newcommand{\pass}{\top}                    % pass sifting


% -----------------------------------------------------------------------------
% 7. Mode / outcome abbreviations
% -----------------------------------------------------------------------------
\newcommand{\on}{{\mathrm{on}}}
\newcommand{\off}{{\mathrm{off}}}
\newcommand{\dc}{{\mathrm{dc}}}             % dark count
\newcommand{\sq}{\mathrm{sq}}               % squashed
\newcommand{\mc}{\mathrm{mc}}               % multi-click
\newcommand{\nc}{\mathrm{nc}}               % no-click
\newcommand{\singlec}{\mathrm{sc}}          % single-click
\newcommand{\vac}{\mathrm{vac}}             % vacuum
\newcommand{\kout}{{k_\mathrm{out}}}
\newcommand{\kin}{{k_\mathrm{in}}}


% -----------------------------------------------------------------------------
% 8. Maps, channels and miscellaneous operators
% -----------------------------------------------------------------------------
\newcommand{\Prob}{\mathrm{Prob}}           % probability (use \Prob[\cdot] or \Prob(\cdot))
\newcommand{\lmin}{\lambda_{\min}}          % minimum eigenvalue
\newcommand{\Wflag}{W_{\Flag}}              % flag-space weight
\newcommand{\transpose}{\mathsf{T}}
\newcommand{\innerp}[2]{\left\langle #1 \vert #2 \right\rangle}
\renewcommand{\vec}[1]{\boldsymbol{#1}}     % NB: overrides the LaTeX arrow accent


% -----------------------------------------------------------------------------
% 9. Key-rate numerics
% -----------------------------------------------------------------------------
\newcommand{\cG}{\mathcal{G}}               % protocol map
\newcommand{\cZ}{\mathcal{Z}}               % pinching map
\newcommand{\Amap}{\mathcal{A}_{\mathrm{ann}}} % announcement + measurement map
\newcommand{\Amapnb}{\mathcal{A}_{\mathrm{ann}|\mathrm{nb}}}
\newcommand{\simpl}{}                       % empty suffix for simplified block maps
\newcommand{\acSet}{\mathcal{S}_{\mathrm{acc}}} % acceptance set
\newcommand{\testE}{\mathcal{E}_{\mathrm{test}}} % test-statistics map
\newcommand{\ecLeak}{\mathrm{leak}^{\mathrm{EC}}_{\mathrm{obs}}}
\newcommand{\pgen}{p_{\mathrm{gen}}}
\newcommand{\askeyr}{R_{\infty}}             % asymptotic key rate
\newcommand{\ppass}{p_{\mathrm{pass}}}       % sifting pass probability


% -----------------------------------------------------------------------------
% 10. Finite-size analysis / Renyi entropies
% -----------------------------------------------------------------------------

% --- Round types: test vs. key generation (register T, alphabet \Fset) --------
\newcommand{\tst}{\mathrm{test}}            % test round
\newcommand{\gen}{\mathrm{gen}}             % key-generation round
\newcommand{\Fset}{\mathcal{T}}             % round-type alphabet {test, gen} (not \Tset)
\newcommand{\Chat}{\widehat{C}}             % acceptance-test register
\newcommand{\Chatset}{\widehat{\mathcal{C}}}% its alphabet

% --- Basis / key extraction functions ----------------------------------------
\newcommand{\bA}{b_{A}}                     % Alice's basis-extraction function
\newcommand{\bB}{b_{B}}                     % Bob's basis-extraction function
\newcommand{\kB}{k_{B}}                     % Bob's key-extraction function

% --- Maps --------------------------------------------------------------------
\newcommand{\Mprot}{\mathcal{M}}            % single-round protocol map
\newcommand{\Phimap}{\boldsymbol{\Phi}}     % test-statistics map
\newcommand{\Pdeph}{\mathcal{P}}            % dephasing (pinching) map on \Blreg

% --- Isometries --------------------------------------------------------------
\newcommand{\Vmsm}{V_{\mathrm{msm}}}        % measurement isometry
\newcommand{\Vann}{V_{\mathrm{ann}}}        % announcement isometry
\newcommand{\VI}{V_{I}}                     % Eve's copy of the announcements
\newcommand{\VChat}{V_{\Chat}}              % acceptance-data isometry
\newcommand{\VR}{V_{R}}                     % sifting + key map -> raw key R
\newcommand{\VZ}{V_{Z}}                     % raw key -> key register Z
\newcommand{\Wprot}{W}                      % full protocol isometry

% --- Renyi entropies ---------------------------------------------------------
% SECURITY-CRITICAL: see spec/notation.md §4. Never change an arrow to
% make an inequality work out; never convert silently between the sandwiched
% (tilde) and Petz (overline) families.
\newcommand{\Hsandup}{\widetilde{H}^{\uparrow}_{\alpha}}   % sandwiched, up
\newcommand{\Hsanddown}{\widetilde{H}^{\downarrow}_{\alpha}}% sandwiched, down
\newcommand{\Hpetzup}{\overline{H}^{\uparrow}_{\alpha}}    % Petz, up
\newcommand{\Hpetzdown}{\overline{H}^{\downarrow}_{\alpha}}% Petz, down
\newcommand{\hup}{h^{\uparrow}_{\alpha}}                   % single-round rate
\newcommand{\hqkd}{h^{\uparrow,\mathrm{QKD}}_{\alpha}}     % single-round QKD rate
\newcommand{\hlow}{\underline{h}_{\alpha}}                 % lower bound on \hqkd
% --- spec/notation.md 3.5 and 4.5 -------------------------------------------
% \hlb  is the first-order linearization lower bound on \hlow: the headline
%       key rate (the `first-order` tier of foundations/def:tiers). NOT \hlow
%       itself.
% \tnum is the instance's one declared numerical tolerance. Its
%       near-collision with the declared marginal \tau_A is disclosed in
%       spec/notation.md 3.5; bare \tau is reserved and denotes neither.
\newcommand{\hlb}{\underline{h}^{\mathrm{FO}}_{\alpha}}    % first-order lower bound on \hlow
\newcommand{\tnum}{\tau_{\mathrm{num}}}                    % declared numerical tolerance


% =============================================================================
%  PART B -- module-local macros, grouped by module. Deprecated aliases,
%  defined so that a text written with them compiles, used in no new contract: \Lop,
%  \testE, \Pdeph, \pass, \ecLeak (spec/notation.md section 3.9).
% =============================================================================

% --- meat-core -----------------------------------------------------------------
\newcommand{\SecSet}{\mathbb{S}}            % sector index set of Q
\newcommand{\Qsys}{Q}                       % declared sector-carrying measured system
\newcommand{\Vset}{\mathcal{V}}             % declared joint outcome alphabet of the Q-POVM
\newcommand{\CAset}{\mathcal{C}_{A}}        % Alice announcement alphabet
\newcommand{\CBset}{\mathcal{C}_{B}}        % Bob announcement alphabet
\newcommand{\Party}{\mathfrak{P}}           % declared party partition
\newcommand{\Oset}[1]{\mathcal{O}_{#1}}     % readout alphabet of party #1
\newcommand{\rd}[1]{\pi_{#1}}               % readout map of party #1
\newcommand{\Sadm}{\mathcal{S}_{A}}         % admissible Alice subspace
\newcommand{\nomarg}{\mathsf{none}}         % "no trusted marginal declared"
\newcommand{\Zset}{\mathcal{Z}_{\perp}}     % raw-key alphabet incl. \perp
\newcommand{\Zval}{\mathcal{Z}}             % key alphabet without \perp
\newcommand{\Etil}{\widetilde{E}}           % single-round purifying register
\newcommand{\Ehat}{\widehat{E}}             % Eve's n-round register
\newcommand{\OmAT}{\Omega_{\mathrm{AT}}}    % acceptance event
\newcommand{\epa}{\varepsilon_{\mathrm{PA}}}
\newcommand{\eev}{\varepsilon_{\mathrm{EV}}}
\newcommand{\lEC}{\lambda_{\mathrm{EC}}}
\newcommand{\Decl}{\mathfrak{D}}            % protocol declaration
\newcommand{\CList}{\mathfrak{L}}           % declared constraint list
\newcommand{\PinS}{\mathcal{P}_{\SecSet}}   % sector pinching on Q
\newcommand{\Feas}{\mathcal{W}}             % feasible set of block states
\newcommand{\Chkey}{\Chatset_{\mathrm{key}}}% key-live acceptance values
\newcommand{\Kann}{\mathcal{K}}             % announcement-class alphabet
\newcommand{\hold}{\mathsf{h}}              % declared key holder (A or B)
\newcommand{\sfD}{D}                        % classical KL divergence (base 2)
\newcommand{\Dtwo}{D}                       % quantum relative entropy (base 2)
\newcommand{\Dsand}{\widetilde{D}_{\alpha}} % sandwiched Renyi divergence
\newcommand{\idmap}{\mathrm{id}}            % identity channel
\newcommand{\Aev}{\mathcal{A}}              % Eve's per-round attack channel
\newcommand{\AttSet}{\mathfrak{A}}          % licensed attack-channel set
\newcommand{\Chacc}{\Chatset_{\mathrm{acc}}}% acceptance support of \acSet
\newcommand{\epsstar}{\varepsilon^{\star}}      % distance to an exact POVM
\newcommand{\dfH}{\Delta_{\mathrm{herm}}}       % Hermiticity defect
\newcommand{\dfP}{\Delta_{\mathrm{pos}}}        % negative part
\newcommand{\dfC}{\Delta_{\mathrm{cpl}}}        % completeness defect
\newcommand{\lmoq}{g}                       % acceptance-side gradient block
\newcommand{\lmoJ}{G}                       % feasible-side gradient block
\newcommand{\bacc}{\beta_{\mathrm{acc}}}    % certified acceptance-side minimum
\newcommand{\Itset}{\mathfrak{I}}           % finite index set of iterates (C22)

% --- reductions ----------------------------------------------------------------
\newcommand{\Bsys}{B}                        % physical (pre-squashing) system
\newcommand{\Tcl}[1]{\mathcal{T}(#1)}        % trace-class operators on H(#1)
\newcommand{\Ins}{\mathcal{I}}               % generalized-instrument branch map
\newcommand{\redNmsm}{\mathcal{N}}              % POVM measurement channel
\newcommand{\Eprot}{\mathcal{E}}             % n-round protocol channel
\newcommand{\Grnd}{\mathsf{G}}               % one round map of a protocol
\newcommand{\Kfin}{\mathsf{K}}               % end-of-protocol processing
\newcommand{\Crit}{\mathsf{sec}}             % value of a security criterion
\newcommand{\Pflag}{\Pi_{\Flag}}             % flag-sector projector on Q
\newcommand{\Pleq}{\Pi_{\leqN}}              % preserved-block projector
\newcommand{\Pgt}{\Pi_{\gtN}}                % above-cutoff projector
\newcommand{\Sph}{S^{\mathrm{ph}}}           % family of physical inputs
\newcommand{\Srch}{S^{\mathrm{rch}}}         % reachable squashed states
\newcommand{\Ssq}{S^{\sq}}                   % squashed admissible set
\newcommand{\Ffss}{F^{\mathrm{FSS}}}         % flag-state-squasher target
\newcommand{\Fwp}{F^{\mathrm{WP}}}           % multi-bound WPFSS target
\newcommand{\Mwp}{M^{\mathrm{WP}}}           % renormalized above-cutoff POVM
\newcommand{\Mdec}{{M'}^{\mathrm{ph}}}       % declared physical family on B
\newcommand{\Ppvm}{P^{\Flag}}                % flag PVM
\newcommand{\Lmp}{{\Lambda_{\mathrm{mp}}}}   % measure-and-prepare channel
\newcommand{\Lwp}{{\Lambda^{\mathrm{WP}}}}   % multi-bound WPFSS channel
\newcommand{\Lfss}{{\Lambda^{\mathrm{FSS}}}} % flag-state-squasher channel
\newcommand{\lmax}{\lambda_{\max}}           % maximum eigenvalue
\newcommand{\Lnoise}{\widetilde{\Lambda}}    % flag-noise channel
\newcommand{\cg}{\varpi}                     % declared coarse-graining map
\newcommand{\MCK}{\mathsf{MC}}               % multiclick raw-pattern set
\newcommand{\redChain}{\mathfrak{C}}            % reduction chain
\newcommand{\Lnk}{\ell}                      % a single typed link
\newcommand{\dsq}{\preceq}                      % the squashing relation
\newcommand{\opn}[1]{\left\lVert #1\right\rVert}           % operator (spectral) norm
\newcommand{\mxn}[1]{\left\lVert #1\right\rVert_{\max}}    % entrywise max modulus
\newcommand{\Fld}{\mathbb{F}}                   % exact ordered substrate field
\newcommand{\Fldi}{\mathbb{F}_{i}}              % its Gaussian extension F + iF
\newcommand{\lat}[1]{\big\lfloor #1\big\rfloor_{\kappa}} % round DOWN onto the lattice

% --- protocol-compiler ---------------------------------------------------------
\newcommand{\Spec}{\mathfrak{S}}             % protocol specification (compiler input)
\newcommand{\pcChain}{\mathfrak{R}}            % reduction chain (two typed segments)
\newcommand{\Chsrc}{\pcChain_{\mathrm{src}}}   % its Alice-side source-map segment
\newcommand{\Chrec}{\pcChain_{\mathrm{rec}}}   % its receiver-side squashing segment
\newcommand{\rvirt}{\rho^{\mathrm{v}}}       % virtual (PRE-source-map) signal state
\newcommand{\Fph}{F^{\mathrm{ph}}}           % physical (pre-reduction) POVM element
\newcommand{\Vph}{\mathcal{V}^{\mathrm{ph}}} % physical outcome alphabet
\newcommand{\Tpub}{\mathcal{T}_{\mathrm{pub}}}% public-choice alphabet
\newcommand{\Uset}{\mathcal{U}}              % residual outcome alphabet
\newcommand{\idxm}{\mathsf{idx}}             % label-to-index map
\newcommand{\Asig}{A'}                       % transmitted (signal) system -- the REAL one
\newcommand{\Avirt}{A''}                     % virtual transmitted system (source descriptor of Xi_1)
\newcommand{\Abar}{\bar{A}}                  % Alice's measured register
\newcommand{\Ahat}{\hat{A}}                  % Alice's shield register
\newcommand{\pcNmsm}{\mathcal{N}_{\mathsf A}}  % Alice's measurement channel (see the RULES block)
\newcommand{\Phip}{\Phi^{+}}                 % maximally entangled vector label
\newcommand{\Smap}{\mathfrak{s}}             % the emitted sector map (fraktur: a compile-time object)

% --- receiver ------------------------------------------------------------------
\newcommand{\Modes}{\mathbb{M}}              % declared ordered mode set
\newcommand{\Inp}{\mathbb{I}}                % declared input modes
\newcommand{\InpC}{\mathbb{I}_{\Chan}}       % channel input modes
\newcommand{\InpS}{\mathbb{I}_{\Prb}}        % probe (source) input modes
\newcommand{\Env}{\mathbb{E}}                % dilation (environment) modes
\newcommand{\Port}[1]{\mathbb{P}_{#1}}       % mode set of detector port #1
\newcommand{\Emd}[1]{E_{#1}}                 % coordinate projector on modes #1
\newcommand{\Gmat}{G}                        % instrument (transfer) matrix
\newcommand{\Fock}{\mathfrak{F}}             % Fock / second-quantization functor
\newcommand{\Symn}[1]{\mathrm{Sym}^{#1}}     % symmetric tensor power
\newcommand{\normord}[1]{{:}\mspace{2mu}#1\mspace{2mu}{:}}   % normal ordering
\newcommand{\rk}{\operatorname{rank}}        % rank (physics.sty owns \rank)
\newcommand{\SCK}{\mathsf{SC}}               % single-click pattern set
\newcommand{\Recv}{\mathfrak{R}}             % receiver instrument tuple
\newcommand{\Cat}{\mathfrak{C}}              % component catalogue
\newcommand{\Grph}{\mathfrak{E}}             % declared element graph
\newcommand{\defic}{\mathbb{L}}              % loss deficiency operator
\newcommand{\Pvac}[1]{\Pi_{\vac}(#1)}        % joint no-click element on ports #1
\newcommand{\pmc}{p_{\mc}}                   % uniform high-photon mc weight
\newcommand{\Bnd}{\beta}                     % the scalar bound function beta(n)
\newcommand{\Rot}{R}                         % planar rotation matrix
\newcommand{\Wplate}{W}                      % retarder (wave-plate) matrix
\newcommand{\Bsp}{B}                         % beamsplitter core matrix

% --- statistics ----------------------------------------------------------------
\newcommand{\honest}{p^{\mathrm{hon}}}          % honest per-round distribution
\newcommand{\ecomp}{\varepsilon_{\mathrm{comp}}}% completeness budget
\newcommand{\simplex}[1]{\Delta(#1)}            % probability simplex
\newcommand{\Chatcon}{\Chatset_{\mathrm{con}}}  % constrained outcomes
\newcommand{\Ireg}[3]{\mathcal{I}_{#1}(#2,#3)}  % regularized incomplete beta
\newcommand{\Ncount}[1]{\mathsf{N}_{#1}}        % per-outcome count
\newcommand{\qlo}{\mathbf{q}^{-}}               % box lower bounds
\newcommand{\qhi}{\mathbf{q}^{+}}               % box upper bounds
 for it, so the
%  file is read from spec/modules/<id>/.
%
%  RULES
%    * A macro that spec/notation.md catalogues is a verbatim copy of it
%      (part A). Changing one is a notation change: bump spec/notation.md.
%    * Part B holds the module-local macros. A macro defined by two modules
%      with the same expansion appears once, under the first module (in the
%      order below) that defines it. Two macros with the same name and
%      DIFFERENT expansions are both kept under a module prefix, and the
%      claims write the prefixed names:
%        \Chain  -> \redChain (reductions, \mathfrak{C}: the reduction chain)
%                   \pcChain  (protocol-compiler, \mathfrak{R}: the two-segment
%                             chain; \Chsrc and \Chrec are built on it)
%        \Nmsm   -> \redNmsm  (reductions, \mathcal{N}: a POVM measurement
%                             channel)
%                   \pcNmsm   (\mathcal{N}_{\mathsf A}: Alice's measurement
%                             channel, written by
%                             transmitter/thm:certificate-sufficiency,
%                             transmitter/rem:xreg and
%                             reductions/lem:source-transfer)
%    * Deliberately omitted: a \provenance{tag}{details} macro. Provenance
%      is a header field of every claim file (spec/claim-format.md section
%      3); a claim body that calls such a macro fails to compile, which is
%      the intended signal.
%
%  REQUIRED PACKAGES, loaded by the assembler's preamble before this file:
%    amsmath, amssymb, amsfonts, mathtools, dsfont, braket, physics, xcolor
%    (braket before physics, so that physics wins for \ket, \bra, \tr, ...).
%
%  This file contains only definitions: no \usepackage, no document body.
% =============================================================================


% =============================================================================
%  PART A -- notation macros, the ones spec/notation.md catalogues, each a
%  verbatim copy of its catalogue row; grouped under numbered headings.
% =============================================================================

% -----------------------------------------------------------------------------
% 1. State spaces and operator sets (a "system label" -> spaces of operators)
% -----------------------------------------------------------------------------
\newcommand{\Hs}[1]{\mathcal{H}(#1)}        % Hilbert space of a system label
\newcommand{\Lin}[1]{\mathcal{L}(#1)}       % bounded linear operators on H(#1)
\newcommand{\Linn}[2]{\mathcal{L}(#1,#2)}   % linear maps H(#1) -> H(#2)
\newcommand{\Pos}[1]{\mathrm{Pos}(#1)}      % positive semidefinite operators on H(#1)
\newcommand{\Density}{\mathrm{D}}           % density operators: \Density(A) = D(A)
\newcommand{\Seq}[1]{S_{=}(#1)}             % normalized states (trace = 1)
\newcommand{\Sleq}[1]{S_{\leq}(#1)}         % subnormalized states (0 <= trace <= 1)
\newcommand{\id}{\mathds{1}}                % identity operator
\newcommand{\idon}[1]{\mathds{1}_{#1}}      % identity on a named system, e.g. \idon{C}

% DEPRECATED: applies \mathcal to both arguments itself, so it takes bare
% letters, not system labels. Prefer \Linn; do not use in new contracts.
\providecommand{\Lop}[2]{\mathcal{L}(\mathcal{#1}, \mathcal{#2})}


% -----------------------------------------------------------------------------
% 2. Composite systems
% -----------------------------------------------------------------------------
\newcommand{\dsum}{\oplus}                  % direct sum of system labels
% Tensor product: use \otimes directly.
% n-fold direct sum of an operator space: \Lin{A}^{\oplus n}


% -----------------------------------------------------------------------------
% 3. Events, conditioning and classical registers (for STATES)
%    Convention:  \wedge  -> subnormalized "and-event" component, e.g. \rho_{A\wedge x}
%                 |        -> normalized conditional state,        e.g. \rho_{A|x}
% (Write these directly in subscripts; no extra macros needed.)
% -----------------------------------------------------------------------------


% -----------------------------------------------------------------------------
% 4. Photon-number subspaces / direct-sum blocks
%    Convention:  operators carry the restriction "on top", in brackets;
%                 states carry it as a \wedge subscript (see Sec. 3).
% -----------------------------------------------------------------------------
\newcommand{\Pnum}[1]{P_{#1}}               % photon-number projector (system from context)
\newcommand{\Pnumon}[2]{P_{#1}^{#2}}        % with explicit system superscript, e.g. P_n^C
\newcommand{\blk}[2]{#1^{(#2)}}             % operator block, e.g. \blk{M}{\leqN} = M^{(<=N)}
\newcommand{\leqN}{\leq N}                  % use as \blk{M}{\leqN}
\newcommand{\gtN}{> N}                      % use as \blk{M}{\gtN}
\newcommand{\eqn}{= n}                      % use as \blk{M}{\eqn}


% -----------------------------------------------------------------------------
% 5. System labels used in the protocol
% -----------------------------------------------------------------------------
\newcommand{\Chan}{C}                       % quantum-channel (Eve) input to Bob
\newcommand{\Prb}{P}                        % probe input
\newcommand{\Flag}{{\mathrm{Flg}}}          % flag system (EXISTS ONLY AFTER squashing)
\newcommand{\Loss}{{\mathrm{Lss}}}          % loss system
\newcommand{\Vc}{{\mathrm{Vcm}}}            % vacuum system
\newcommand{\Nout}{N_{\mathrm{out}}}
\newcommand{\Nin}{N_{\mathrm{in}}}


% -----------------------------------------------------------------------------
% 6. Alphabets
% -----------------------------------------------------------------------------
\newcommand{\Yset}{\mathcal{Y}}             % Bob's (coarse-grained) detection alphabet
\newcommand{\Yraw}{\mathcal{Y}_{\mathrm{raw}}}% raw click patterns {0,1}^4
\newcommand{\Yscset}{\mathcal{Y}_{\singlec}}% single-click outcomes
\newcommand{\Ymcset}{\mathcal{Y}_{\mc}}     % multi-click outcomes
\newcommand{\Tset}{\mathcal{T}_{\mathrm{set}}}% private setting set (the probe settings in the optical instance)
\newcommand{\Xset}{\mathcal{X}}             % Alice's alphabet
\newcommand{\Dset}{\{H,V,+,-\}}             % polarization / detector labels
\newcommand{\pass}{\top}                    % pass sifting


% -----------------------------------------------------------------------------
% 7. Mode / outcome abbreviations
% -----------------------------------------------------------------------------
\newcommand{\on}{{\mathrm{on}}}
\newcommand{\off}{{\mathrm{off}}}
\newcommand{\dc}{{\mathrm{dc}}}             % dark count
\newcommand{\sq}{\mathrm{sq}}               % squashed
\newcommand{\mc}{\mathrm{mc}}               % multi-click
\newcommand{\nc}{\mathrm{nc}}               % no-click
\newcommand{\singlec}{\mathrm{sc}}          % single-click
\newcommand{\vac}{\mathrm{vac}}             % vacuum
\newcommand{\kout}{{k_\mathrm{out}}}
\newcommand{\kin}{{k_\mathrm{in}}}


% -----------------------------------------------------------------------------
% 8. Maps, channels and miscellaneous operators
% -----------------------------------------------------------------------------
\newcommand{\Prob}{\mathrm{Prob}}           % probability (use \Prob[\cdot] or \Prob(\cdot))
\newcommand{\lmin}{\lambda_{\min}}          % minimum eigenvalue
\newcommand{\Wflag}{W_{\Flag}}              % flag-space weight
\newcommand{\transpose}{\mathsf{T}}
\newcommand{\innerp}[2]{\left\langle #1 \vert #2 \right\rangle}
\renewcommand{\vec}[1]{\boldsymbol{#1}}     % NB: overrides the LaTeX arrow accent


% -----------------------------------------------------------------------------
% 9. Key-rate numerics
% -----------------------------------------------------------------------------
\newcommand{\cG}{\mathcal{G}}               % protocol map
\newcommand{\cZ}{\mathcal{Z}}               % pinching map
\newcommand{\Amap}{\mathcal{A}_{\mathrm{ann}}} % announcement + measurement map
\newcommand{\Amapnb}{\mathcal{A}_{\mathrm{ann}|\mathrm{nb}}}
\newcommand{\simpl}{}                       % empty suffix for simplified block maps
\newcommand{\acSet}{\mathcal{S}_{\mathrm{acc}}} % acceptance set
\newcommand{\testE}{\mathcal{E}_{\mathrm{test}}} % test-statistics map
\newcommand{\ecLeak}{\mathrm{leak}^{\mathrm{EC}}_{\mathrm{obs}}}
\newcommand{\pgen}{p_{\mathrm{gen}}}
\newcommand{\askeyr}{R_{\infty}}             % asymptotic key rate
\newcommand{\ppass}{p_{\mathrm{pass}}}       % sifting pass probability


% -----------------------------------------------------------------------------
% 10. Finite-size analysis / Renyi entropies
% -----------------------------------------------------------------------------

% --- Round types: test vs. key generation (register T, alphabet \Fset) --------
\newcommand{\tst}{\mathrm{test}}            % test round
\newcommand{\gen}{\mathrm{gen}}             % key-generation round
\newcommand{\Fset}{\mathcal{T}}             % round-type alphabet {test, gen} (not \Tset)
\newcommand{\Chat}{\widehat{C}}             % acceptance-test register
\newcommand{\Chatset}{\widehat{\mathcal{C}}}% its alphabet

% --- Basis / key extraction functions ----------------------------------------
\newcommand{\bA}{b_{A}}                     % Alice's basis-extraction function
\newcommand{\bB}{b_{B}}                     % Bob's basis-extraction function
\newcommand{\kB}{k_{B}}                     % Bob's key-extraction function

% --- Maps --------------------------------------------------------------------
\newcommand{\Mprot}{\mathcal{M}}            % single-round protocol map
\newcommand{\Phimap}{\boldsymbol{\Phi}}     % test-statistics map
\newcommand{\Pdeph}{\mathcal{P}}            % dephasing (pinching) map on \Blreg

% --- Isometries --------------------------------------------------------------
\newcommand{\Vmsm}{V_{\mathrm{msm}}}        % measurement isometry
\newcommand{\Vann}{V_{\mathrm{ann}}}        % announcement isometry
\newcommand{\VI}{V_{I}}                     % Eve's copy of the announcements
\newcommand{\VChat}{V_{\Chat}}              % acceptance-data isometry
\newcommand{\VR}{V_{R}}                     % sifting + key map -> raw key R
\newcommand{\VZ}{V_{Z}}                     % raw key -> key register Z
\newcommand{\Wprot}{W}                      % full protocol isometry

% --- Renyi entropies ---------------------------------------------------------
% SECURITY-CRITICAL: see spec/notation.md §4. Never change an arrow to
% make an inequality work out; never convert silently between the sandwiched
% (tilde) and Petz (overline) families.
\newcommand{\Hsandup}{\widetilde{H}^{\uparrow}_{\alpha}}   % sandwiched, up
\newcommand{\Hsanddown}{\widetilde{H}^{\downarrow}_{\alpha}}% sandwiched, down
\newcommand{\Hpetzup}{\overline{H}^{\uparrow}_{\alpha}}    % Petz, up
\newcommand{\Hpetzdown}{\overline{H}^{\downarrow}_{\alpha}}% Petz, down
\newcommand{\hup}{h^{\uparrow}_{\alpha}}                   % single-round rate
\newcommand{\hqkd}{h^{\uparrow,\mathrm{QKD}}_{\alpha}}     % single-round QKD rate
\newcommand{\hlow}{\underline{h}_{\alpha}}                 % lower bound on \hqkd
% --- spec/notation.md 3.5 and 4.5 -------------------------------------------
% \hlb  is the first-order linearization lower bound on \hlow: the headline
%       key rate (the `first-order` tier of foundations/def:tiers). NOT \hlow
%       itself.
% \tnum is the instance's one declared numerical tolerance. Its
%       near-collision with the declared marginal \tau_A is disclosed in
%       spec/notation.md 3.5; bare \tau is reserved and denotes neither.
\newcommand{\hlb}{\underline{h}^{\mathrm{FO}}_{\alpha}}    % first-order lower bound on \hlow
\newcommand{\tnum}{\tau_{\mathrm{num}}}                    % declared numerical tolerance


% =============================================================================
%  PART B -- module-local macros, grouped by module. Deprecated aliases,
%  defined so that a text written with them compiles, used in no new contract: \Lop,
%  \testE, \Pdeph, \pass, \ecLeak (spec/notation.md section 3.9).
% =============================================================================

% --- meat-core -----------------------------------------------------------------
\newcommand{\SecSet}{\mathbb{S}}            % sector index set of Q
\newcommand{\Qsys}{Q}                       % declared sector-carrying measured system
\newcommand{\Vset}{\mathcal{V}}             % declared joint outcome alphabet of the Q-POVM
\newcommand{\CAset}{\mathcal{C}_{A}}        % Alice announcement alphabet
\newcommand{\CBset}{\mathcal{C}_{B}}        % Bob announcement alphabet
\newcommand{\Party}{\mathfrak{P}}           % declared party partition
\newcommand{\Oset}[1]{\mathcal{O}_{#1}}     % readout alphabet of party #1
\newcommand{\rd}[1]{\pi_{#1}}               % readout map of party #1
\newcommand{\Sadm}{\mathcal{S}_{A}}         % admissible Alice subspace
\newcommand{\nomarg}{\mathsf{none}}         % "no trusted marginal declared"
\newcommand{\Zset}{\mathcal{Z}_{\perp}}     % raw-key alphabet incl. \perp
\newcommand{\Zval}{\mathcal{Z}}             % key alphabet without \perp
\newcommand{\Etil}{\widetilde{E}}           % single-round purifying register
\newcommand{\Ehat}{\widehat{E}}             % Eve's n-round register
\newcommand{\OmAT}{\Omega_{\mathrm{AT}}}    % acceptance event
\newcommand{\epa}{\varepsilon_{\mathrm{PA}}}
\newcommand{\eev}{\varepsilon_{\mathrm{EV}}}
\newcommand{\lEC}{\lambda_{\mathrm{EC}}}
\newcommand{\Decl}{\mathfrak{D}}            % protocol declaration
\newcommand{\CList}{\mathfrak{L}}           % declared constraint list
\newcommand{\PinS}{\mathcal{P}_{\SecSet}}   % sector pinching on Q
\newcommand{\Feas}{\mathcal{W}}             % feasible set of block states
\newcommand{\Chkey}{\Chatset_{\mathrm{key}}}% key-live acceptance values
\newcommand{\Kann}{\mathcal{K}}             % announcement-class alphabet
\newcommand{\hold}{\mathsf{h}}              % declared key holder (A or B)
\newcommand{\sfD}{D}                        % classical KL divergence (base 2)
\newcommand{\Dtwo}{D}                       % quantum relative entropy (base 2)
\newcommand{\Dsand}{\widetilde{D}_{\alpha}} % sandwiched Renyi divergence
\newcommand{\idmap}{\mathrm{id}}            % identity channel
\newcommand{\Aev}{\mathcal{A}}              % Eve's per-round attack channel
\newcommand{\AttSet}{\mathfrak{A}}          % licensed attack-channel set
\newcommand{\Chacc}{\Chatset_{\mathrm{acc}}}% acceptance support of \acSet
\newcommand{\epsstar}{\varepsilon^{\star}}      % distance to an exact POVM
\newcommand{\dfH}{\Delta_{\mathrm{herm}}}       % Hermiticity defect
\newcommand{\dfP}{\Delta_{\mathrm{pos}}}        % negative part
\newcommand{\dfC}{\Delta_{\mathrm{cpl}}}        % completeness defect
\newcommand{\lmoq}{g}                       % acceptance-side gradient block
\newcommand{\lmoJ}{G}                       % feasible-side gradient block
\newcommand{\bacc}{\beta_{\mathrm{acc}}}    % certified acceptance-side minimum
\newcommand{\Itset}{\mathfrak{I}}           % finite index set of iterates (C22)

% --- reductions ----------------------------------------------------------------
\newcommand{\Bsys}{B}                        % physical (pre-squashing) system
\newcommand{\Tcl}[1]{\mathcal{T}(#1)}        % trace-class operators on H(#1)
\newcommand{\Ins}{\mathcal{I}}               % generalized-instrument branch map
\newcommand{\redNmsm}{\mathcal{N}}              % POVM measurement channel
\newcommand{\Eprot}{\mathcal{E}}             % n-round protocol channel
\newcommand{\Grnd}{\mathsf{G}}               % one round map of a protocol
\newcommand{\Kfin}{\mathsf{K}}               % end-of-protocol processing
\newcommand{\Crit}{\mathsf{sec}}             % value of a security criterion
\newcommand{\Pflag}{\Pi_{\Flag}}             % flag-sector projector on Q
\newcommand{\Pleq}{\Pi_{\leqN}}              % preserved-block projector
\newcommand{\Pgt}{\Pi_{\gtN}}                % above-cutoff projector
\newcommand{\Sph}{S^{\mathrm{ph}}}           % family of physical inputs
\newcommand{\Srch}{S^{\mathrm{rch}}}         % reachable squashed states
\newcommand{\Ssq}{S^{\sq}}                   % squashed admissible set
\newcommand{\Ffss}{F^{\mathrm{FSS}}}         % flag-state-squasher target
\newcommand{\Fwp}{F^{\mathrm{WP}}}           % multi-bound WPFSS target
\newcommand{\Mwp}{M^{\mathrm{WP}}}           % renormalized above-cutoff POVM
\newcommand{\Mdec}{{M'}^{\mathrm{ph}}}       % declared physical family on B
\newcommand{\Ppvm}{P^{\Flag}}                % flag PVM
\newcommand{\Lmp}{{\Lambda_{\mathrm{mp}}}}   % measure-and-prepare channel
\newcommand{\Lwp}{{\Lambda^{\mathrm{WP}}}}   % multi-bound WPFSS channel
\newcommand{\Lfss}{{\Lambda^{\mathrm{FSS}}}} % flag-state-squasher channel
\newcommand{\lmax}{\lambda_{\max}}           % maximum eigenvalue
\newcommand{\Lnoise}{\widetilde{\Lambda}}    % flag-noise channel
\newcommand{\cg}{\varpi}                     % declared coarse-graining map
\newcommand{\MCK}{\mathsf{MC}}               % multiclick raw-pattern set
\newcommand{\redChain}{\mathfrak{C}}            % reduction chain
\newcommand{\Lnk}{\ell}                      % a single typed link
\newcommand{\dsq}{\preceq}                      % the squashing relation
\newcommand{\opn}[1]{\left\lVert #1\right\rVert}           % operator (spectral) norm
\newcommand{\mxn}[1]{\left\lVert #1\right\rVert_{\max}}    % entrywise max modulus
\newcommand{\Fld}{\mathbb{F}}                   % exact ordered substrate field
\newcommand{\Fldi}{\mathbb{F}_{i}}              % its Gaussian extension F + iF
\newcommand{\lat}[1]{\big\lfloor #1\big\rfloor_{\kappa}} % round DOWN onto the lattice

% --- protocol-compiler ---------------------------------------------------------
\newcommand{\Spec}{\mathfrak{S}}             % protocol specification (compiler input)
\newcommand{\pcChain}{\mathfrak{R}}            % reduction chain (two typed segments)
\newcommand{\Chsrc}{\pcChain_{\mathrm{src}}}   % its Alice-side source-map segment
\newcommand{\Chrec}{\pcChain_{\mathrm{rec}}}   % its receiver-side squashing segment
\newcommand{\rvirt}{\rho^{\mathrm{v}}}       % virtual (PRE-source-map) signal state
\newcommand{\Fph}{F^{\mathrm{ph}}}           % physical (pre-reduction) POVM element
\newcommand{\Vph}{\mathcal{V}^{\mathrm{ph}}} % physical outcome alphabet
\newcommand{\Tpub}{\mathcal{T}_{\mathrm{pub}}}% public-choice alphabet
\newcommand{\Uset}{\mathcal{U}}              % residual outcome alphabet
\newcommand{\idxm}{\mathsf{idx}}             % label-to-index map
\newcommand{\Asig}{A'}                       % transmitted (signal) system -- the REAL one
\newcommand{\Avirt}{A''}                     % virtual transmitted system (source descriptor of Xi_1)
\newcommand{\Abar}{\bar{A}}                  % Alice's measured register
\newcommand{\Ahat}{\hat{A}}                  % Alice's shield register
\newcommand{\pcNmsm}{\mathcal{N}_{\mathsf A}}  % Alice's measurement channel (see the RULES block)
\newcommand{\Phip}{\Phi^{+}}                 % maximally entangled vector label
\newcommand{\Smap}{\mathfrak{s}}             % the emitted sector map (fraktur: a compile-time object)

% --- receiver ------------------------------------------------------------------
\newcommand{\Modes}{\mathbb{M}}              % declared ordered mode set
\newcommand{\Inp}{\mathbb{I}}                % declared input modes
\newcommand{\InpC}{\mathbb{I}_{\Chan}}       % channel input modes
\newcommand{\InpS}{\mathbb{I}_{\Prb}}        % probe (source) input modes
\newcommand{\Env}{\mathbb{E}}                % dilation (environment) modes
\newcommand{\Port}[1]{\mathbb{P}_{#1}}       % mode set of detector port #1
\newcommand{\Emd}[1]{E_{#1}}                 % coordinate projector on modes #1
\newcommand{\Gmat}{G}                        % instrument (transfer) matrix
\newcommand{\Fock}{\mathfrak{F}}             % Fock / second-quantization functor
\newcommand{\Symn}[1]{\mathrm{Sym}^{#1}}     % symmetric tensor power
\newcommand{\normord}[1]{{:}\mspace{2mu}#1\mspace{2mu}{:}}   % normal ordering
\newcommand{\rk}{\operatorname{rank}}        % rank (physics.sty owns \rank)
\newcommand{\SCK}{\mathsf{SC}}               % single-click pattern set
\newcommand{\Recv}{\mathfrak{R}}             % receiver instrument tuple
\newcommand{\Cat}{\mathfrak{C}}              % component catalogue
\newcommand{\Grph}{\mathfrak{E}}             % declared element graph
\newcommand{\defic}{\mathbb{L}}              % loss deficiency operator
\newcommand{\Pvac}[1]{\Pi_{\vac}(#1)}        % joint no-click element on ports #1
\newcommand{\pmc}{p_{\mc}}                   % uniform high-photon mc weight
\newcommand{\Bnd}{\beta}                     % the scalar bound function beta(n)
\newcommand{\Rot}{R}                         % planar rotation matrix
\newcommand{\Wplate}{W}                      % retarder (wave-plate) matrix
\newcommand{\Bsp}{B}                         % beamsplitter core matrix

% --- statistics ----------------------------------------------------------------
\newcommand{\honest}{p^{\mathrm{hon}}}          % honest per-round distribution
\newcommand{\ecomp}{\varepsilon_{\mathrm{comp}}}% completeness budget
\newcommand{\simplex}[1]{\Delta(#1)}            % probability simplex
\newcommand{\Chatcon}{\Chatset_{\mathrm{con}}}  % constrained outcomes
\newcommand{\Ireg}[3]{\mathcal{I}_{#1}(#2,#3)}  % regularized incomplete beta
\newcommand{\Ncount}[1]{\mathsf{N}_{#1}}        % per-outcome count
\newcommand{\qlo}{\mathbf{q}^{-}}               % box lower bounds
\newcommand{\qhi}{\mathbf{q}^{+}}               % box upper bounds

\makeatletter
\newcommand*{\ClaimHeadFont}{\bfseries}   % switched by prominence
\newcommand*{\ClaimNoteFont}{\mdseries}
\newcommand*{\ClaimShortId}{}
\newcommand*{\SetShortId}[1]{\renewcommand*{\ClaimShortId}{#1}}
\newcommand*{\RestatedName}{}
% Upright bodies throughout: italic is reserved for non-normative text.
% \@endparpenalty\@M keeps the meta line on the statement's page.
\newtheoremstyle{claim}{\topsep}{\topsep}{\normalfont\@endparpenalty\@M}{}%
  {\ClaimHeadFont}{.}{.5em}{\thmname{#1}\thmnumber{ #2}\thmnote{ {\ClaimNoteFont(#3)}}}
\newtheoremstyle{claimass}{\topsep}{\topsep}{\normalfont\@endparpenalty\@M}{}%
  {\ClaimHeadFont}{.}{.5em}{\thmname{#1}\thmnumber{ #2}%
  \ifx\ClaimShortId\@empty\else\ [\textsf{\ClaimShortId}]\fi%
  \thmnote{ {\ClaimNoteFont(#3)}}}
\newtheoremstyle{claimremark}{.5\topsep}{.5\topsep}{\small\@endparpenalty\@M}{}%
  {\itshape}{.}{.5em}{\thmname{#1}\thmnumber{ #2}\thmnote{ {\upshape(#3)}}}
\newtheoremstyle{claimreq}{\topsep}{\topsep}{\normalfont\@endparpenalty\@M}{}%
  {\color{reqAccent}\scshape}{.}{.5em}%
  {\thmname{#1}\thmnumber{ #2}\thmnote{ {\upshape(#3)}}}
\makeatother
\theoremstyle{claim}
\newtheorem{theorem}{Theorem}[section]
\newtheorem{lemma}[theorem]{Lemma}
\newtheorem{proposition}[theorem]{Proposition}
\newtheorem{corollary}[theorem]{Corollary}
\newtheorem{fact}[theorem]{Fact}
\newtheorem{definition}[theorem]{Definition}
\newtheorem*{restatedclaim}{\RestatedName}
\theoremstyle{claimass}
\newtheorem{assumption}[theorem]{Assumption}
\newtheorem*{restatedassumption}{\RestatedName}
\theoremstyle{claimremark}
\newtheorem{remark}[theorem]{Remark}
\newtheorem{example}[theorem]{Example}
\theoremstyle{claimreq}
\newtheorem{requirement}[theorem]{Requirement}
\newtheorem*{restatedrequirement}{\RestatedName}
%% --- 8. Prominence: main results a margin rule and a bold heading ----------
\newtcolorbox{mainclaim}{enhanced jigsaw, breakable, blanker,
  borderline west={1.5pt}{-8pt}{mainRule}, left=0pt, right=0pt, top=0pt, bottom=0pt,
  boxsep=0pt, before skip=\topsep, after skip=\topsep, parbox=false,
  before upper={\renewcommand*{\ClaimNoteFont}{\bfseries}}}
\newenvironment{technicalclaim}{\renewcommand*{\ClaimHeadFont}{\mdseries\color{techGrey}}%
  \renewcommand*{\ClaimNoteFont}{}}{}
%% --- 9. Non-normative text, the meta line, navigation ------------------------
\makeatletter
% Tie what follows to this paragraph: no page break, no list penalty before
% the next theorem head (the \@nobreak route a section heading uses).
\newcommand{\keepwithnext}{\par\nobreak\vspace{3pt}\@nobreaktrue
  \everypar{\if@nobreak\@nobreakfalse\clubpenalty\@M\else\clubpenalty\@clubpenalty\everypar{}\fi}}
\newcommand{\claimsummary}[2]{% #1 claim id (named destination), #2 summary
  \par\addvspace{\bigskipamount}\Needspace{6\baselineskip}%
  {\small\itshape\color{nnGrey}\noindent\hypertarget{claim-#1}{}#2\par}\keepwithnext}
\newcommand{\claimanchor}[1]{% a claim without a summary: the destination alone
  \par\addvspace{\bigskipamount}\Needspace{6\baselineskip}%
  \phantomsection\hypertarget{claim-#1}{}\keepwithnext}
\newcommand{\claimleadin}[1]{{\color{nnGrey}\noindent#1\par}\keepwithnext}
\newcommand{\claimmeta}[1]{\par\nobreak\vspace{2pt}%
  {\footnotesize\color{metaGrey}\raggedright\interlinepenalty\@M\noindent#1\par}%
  \addvspace{\medskipamount}}
% After an assumption: the sentence generated from its verification field; after
% any claim with one: its plain-text reading (user_statement). Both
% non-normative, both grey.
\newcommand{\claimverification}[1]{\par\nobreak\vspace{2pt}%
  {\footnotesize\color{metaGrey}\noindent#1\par}\nobreak}
\newcommand{\claimreading}[1]{\par\nobreak\vspace{2pt}%
  {\footnotesize\color{nnGrey}\noindent\emph{Reading:}\enspace#1\par}\nobreak}
\makeatother
\newcommand{\claimbookmark}[4]{% #1 id, #2 kind word, #3 label, #4 plain title
  \bookmark[dest=claim-#1,level=2]{#2 \getrefnumber{#3} #4}}
\newcommand{\metasep}{\unskip\enspace\textperiodcentered\enspace\ignorespaces}
\newcommand{\metalabel}[1]{\emph{#1}\enspace}
\newcommand{\sectionintro}[1]{{\color{nnGrey}\noindent#1\par}\smallskip}
\newcommand{\sectionimports}[1]{{\footnotesize\color{metaGrey}\raggedright\noindent\metalabel{Imports}#1\par}\smallskip}
\newenvironment{overview}{\par\color{nnGrey}\itshape\noindent\ignorespaces}{\par}
% References: the whole "Def. 2.1" is the link; ids are the fallback text of
% every link into another PDF (a relative URI to a named destination).
\newcommand{\mref}[2]{\hyperref[#2]{#1~\ref*{#2}}}
\newcommand{\aref}[2]{\hyperref[#2]{\textsf{#1}}}
\DeclareUrlCommand{\claimid}{\urlstyle{sf}\def\UrlBreaks{\do\/\do\:\do\-\do\.\do\_}%
  \def\UrlBigBreaks{}\def\UrlNoBreaks{}}
\newcommand{\extref}[2]{\hyperref{#1}{}{claim-#2}{\claimid{#2}}}
\newcommand{\extaref}[3]{\hyperref{#1}{}{claim-#2}{\textsf{#3}}}
\newcommand{\coderef}[2]{\href{#1}{\claimid{#2}}}
% Claim-body commands (spec/claim-format.md section 4): \xref{<full id>} prints
% the kind, title and id of a claim of another contract, linked to its stable
% destination in that contract's PDF (one \DeclareXref per cited claim below);
% \term[gloss]{phrase} marks a defined phrase for the index (the gloss feeds
% the Terms table); \tref[phrase]{text} is a use of a defined term. The
% assembler rewrites both in the body into \termdef (the definition's named
% destination term-<slug>) and \termref / \exttermref (a link to it, in this
% or another contract's PDF); the fallbacks below serve text it copies
% elsewhere, such as the digest. \iface is an interface key.
\makeatletter
\newcommand{\DeclareXref}[4]{% id, kind word, title, contract PDF
  \@namedef{xref@#1}{\hyperref{#4}{}{claim-#1}{#2 ``#3'' ({\small\claimid{#1}})}}}
\newcommand{\xref}[1]{\@ifundefined{xref@#1}%
  {\textbf{[unresolved \nolinkurl{#1}]}\@latex@warning{Unresolved xref #1}}%
  {\@nameuse{xref@#1}}}
\makeatother
\DeclareUrlCommand{\iface}{\urlstyle{tt}}
\newcommand{\term}[2][]{\emph{#2}\index{#2}}
\newcommand{\tref}[2][]{#2}
\newcommand{\termdef}[2]{\hypertarget{#1}{}\emph{#2}\index{#2}}
\newcommand{\termref}[2]{\hyperlink{#1}{#2}}
\newcommand{\exttermref}[3]{\hyperref{#1}{}{#2}{#3}}
\hypersetup{pdftitle={Proof reductions by squashing (reductions) -- statements digest 0.22.0}}

% The digest: the contract's labels under the prefix C-, its numbers checked.
\externaldocument[C-]{contract}
\theoremstyle{claim}\newtheorem*{digestclaim}{\RestatedName}
\theoremstyle{claimass}\newtheorem*{digestassumption}{\RestatedName}
\theoremstyle{claimreq}\newtheorem*{digestrequirement}{\RestatedName}
\makeatletter
% #1 label, #2 the number the assembler computed: a local label carrying it, so
% that a reference between two digest items stays in the digest; the build
% stops when the contract's own number differs.
\newcommand{\digestlabel}[2]{\phantomsection
  \edef\digest@theirs{\getrefnumber{C-#1}}\def\digest@ours{#2}%
  \ifx\digest@theirs\digest@ours\else
    \PackageError{digest}{#1 is #2 here but \digest@theirs\space in the contract}{}\fi
  \def\@currentlabel{#2}\label{#1}}
\makeatother
% A display's local label (amsmath records the \tag, i.e. the contract's number).
\newcommand{\digesteqlabel}[1]{\label{#1}}

% Claims of other contracts cited here (\xref, [[...]]).
\DeclareXref{foundations/ass:snapped-measurement}{Assumption}{The device measures the snapped family where the run reads the snap}{../foundations/contract.pdf}
\DeclareXref{foundations/def:binary64}{Definition}{Binary64 numbers, declared values and dyadic lattices}{../foundations/contract.pdf}
\DeclareXref{foundations/def:declared-data}{Definition}{Declared data and the parties to a module boundary}{../foundations/contract.pdf}
\DeclareXref{foundations/def:domination}{Definition}{Domination of one operator family by another at block weights}{../foundations/contract.pdf}
\DeclareXref{foundations/def:evaluated-off-snap}{Definition}{A reference handed beside a snap; a block evaluated off its snap}{../foundations/contract.pdf}
\DeclareXref{foundations/def:exact-admissibility}{Definition}{Exact admissibility of a declared operator family}{../foundations/contract.pdf}
\DeclareXref{foundations/def:exact-island}{Definition}{The exact island}{../foundations/contract.pdf}
\DeclareXref{foundations/def:exact-reference}{Definition}{Reference of a measurement family, exact or enclosed}{../foundations/contract.pdf}
\DeclareXref{foundations/def:factor-enclosure}{Definition}{Factor enclosure of a block}{../foundations/contract.pdf}
\DeclareXref{foundations/def:pinching}{Definition}{Operator conventions and pinching}{../foundations/contract.pdf}
\DeclareXref{foundations/def:round-homogeneity}{Definition}{Round homogeneity of a per-round family}{../foundations/contract.pdf}
\DeclareXref{foundations/def:round-model}{Definition}{Round model, attack sequences and attack sets}{../foundations/contract.pdf}
\DeclareXref{foundations/def:sector-map}{Definition}{Sector maps and block-list operators}{../foundations/contract.pdf}
\DeclareXref{foundations/def:snap}{Definition}{Snapping declared data onto a lattice}{../foundations/contract.pdf}
\DeclareXref{foundations/lem:absorbing-domination}{Lemma}{Domination of a family that absorbs its residue in one element}{../foundations/contract.pdf}
\DeclareXref{foundations/lem:binary64-lattice}{Lemma}{Binary64 numbers on dyadic lattices}{../foundations/contract.pdf}
\DeclareXref{foundations/lem:domination}{Lemma}{Elementary properties of domination}{../foundations/contract.pdf}
\DeclareXref{foundations/lem:exact-accumulation}{Lemma}{Exactness of a sum that rounds nowhere}{../foundations/contract.pdf}
\DeclareXref{foundations/lem:rational-obstruction}{Lemma}{Rationality obstruction to the domination of a lattice family}{../foundations/contract.pdf}
\DeclareXref{foundations/lem:recentred-enclosure}{Lemma}{Zeroing the centres an enclosure cannot tell from zero}{../foundations/contract.pdf}
\DeclareXref{foundations/rem:computed-completeness}{Remark}{A complete computed sum from an incomplete family}{../foundations/contract.pdf}
\DeclareXref{foundations/rem:noise-floor}{Remark}{What domination controls, and why it needs a noise floor}{../foundations/contract.pdf}
\DeclareXref{foundations/req:boundary-admission}{Requirement}{Exact admissibility at every module boundary}{../foundations/contract.pdf}
\DeclareXref{foundations/req:consumer-decides-exactly}{Requirement}{Exact decision by every consumer}{../foundations/contract.pdf}
\DeclareXref{foundations/req:declared-lattice}{Requirement}{The declared lattice condition}{../foundations/contract.pdf}
\DeclareXref{foundations/req:declared-tolerance}{Requirement}{One declared tolerance per instance}{../foundations/contract.pdf}
\DeclareXref{foundations/req:exact-emission}{Requirement}{Exact verification before emission}{../foundations/contract.pdf}
\DeclareXref{foundations/req:norm-conversion}{Requirement}{Explicit conversion between norms}{../foundations/contract.pdf}
\DeclareXref{foundations/req:rounding-direction}{Requirement}{Rounding direction of emitted bounds}{../foundations/contract.pdf}
\DeclareXref{foundations/req:snap-reference}{Requirement}{Snaps of a measurement family: kept, or moved and disclosed}{../foundations/contract.pdf}
\DeclareXref{foundations/req:support-safe-snap}{Requirement}{Support-safe snaps of a measurement family}{../foundations/contract.pdf}
\DeclareXref{foundations/req:tolerance-scope}{Requirement}{Scope of the declared tolerance}{../foundations/contract.pdf}
\DeclareXref{foundations/req:who-snaps}{Requirement}{Snapping by the declarer only}{../foundations/contract.pdf}
\DeclareXref{meat-core/cor:composed-security}{Corollary}{Security of the protocol and its net key}{../meat-core/contract.pdf}
\DeclareXref{meat-core/def:admissible-declaration}{Definition}{Admissible declaration}{../meat-core/contract.pdf}
\DeclareXref{meat-core/def:declaration}{Definition}{The protocol declaration as the MEAT core reads it}{../meat-core/contract.pdf}
\DeclareXref{meat-core/def:hqkd-list}{Definition}{Constrained target set and single-round rate}{../meat-core/contract.pdf}
\DeclareXref{meat-core/def:roundmodel}{Definition}{Round model and licensed attack sets}{../meat-core/contract.pdf}
\DeclareXref{meat-core/frank-wolfe/req:bracket-on-reference}{Requirement}{The first-order bound evaluated on the reference companion}{../meat-core/contract.pdf}
\DeclareXref{meat-core/frank-wolfe/req:certified-evaluation}{Requirement}{The certified evaluation of the first-order bound}{../meat-core/contract.pdf}
\DeclareXref{meat-core/lem:memory-lift}{Lemma}{Approximating an attack by all-at-once attacks with finite memory}{../meat-core/contract.pdf}
\DeclareXref{meat-core/lem:range}{Lemma}{Constraint transport and range inclusion}{../meat-core/contract.pdf}
\DeclareXref{meat-core/prop:optol-nobound}{Proposition}{Instability of the rate under measurement perturbations}{../meat-core/contract.pdf}
\DeclareXref{meat-core/rem:cgate}{Remark}{Degeneracy of the licensed-class half of the gate}{../meat-core/contract.pdf}
\DeclareXref{meat-core/rem:modes-da1}{Remark}{The two source modes on the program}{../meat-core/contract.pdf}
\DeclareXref{meat-core/req:companion-admission}{Requirement}{Admission of the reference companion beside its declaration}{../meat-core/contract.pdf}
\DeclareXref{meat-core/thm:protocol-channel}{Theorem}{The declaration generates a protocol channel}{../meat-core/contract.pdf}
\DeclareXref{protocol-compiler/ass:atk}{Assumption}{Attack-set licence}{../protocol-compiler/contract.pdf}
\DeclareXref{protocol-compiler/def:chaincond}{Definition}{Interface conditions on a chain}{../protocol-compiler/contract.pdf}
\DeclareXref{protocol-compiler/def:declaration}{Definition}{Protocol declaration}{../protocol-compiler/contract.pdf}
\DeclareXref{protocol-compiler/def:homog}{Definition}{Licensed constraint list}{../protocol-compiler/contract.pdf}
\DeclareXref{protocol-compiler/def:spec}{Definition}{Protocol specification}{../protocol-compiler/contract.pdf}
\DeclareXref{protocol-compiler/req:chain-refusals}{Requirement}{Refusals of the interface conditions on a chain}{../protocol-compiler/contract.pdf}
\DeclareXref{protocol-compiler/req:link-row-provenance}{Requirement}{Omission of the link-soundness row only on a proved chain}{../protocol-compiler/contract.pdf}
\DeclareXref{receiver/ass:eig}{Assumption}{Declared slack on the computed block eigenvalues}{../receiver/contract.pdf}
\DeclareXref{receiver/ass:trust}{Assumption}{Trusted and exactly characterised receiver}{../receiver/contract.pdf}
\DeclareXref{receiver/component-catalogue/thm:catalogue-admissibility}{Theorem}{Admissibility of the catalogue entries}{../receiver/contract.pdf}
\DeclareXref{receiver/cor:compiled-instrument}{Corollary}{Receiver instrument from a well-formed graph declaration}{../receiver/contract.pdf}
\DeclareXref{receiver/cor:unconditioned-family}{Corollary}{The unconditioned receiver is a setting family}{../receiver/contract.pdf}
\DeclareXref{receiver/def:complex-enclosure}{Definition}{Complex enclosures, and enclosures of square roots and angles}{../receiver/contract.pdf}
\DeclareXref{receiver/def:evaluation-model}{Definition}{The evaluation model: floating point, enclosures and rigour}{../receiver/contract.pdf}
\DeclareXref{receiver/def:graph}{Definition}{Declared element graph}{../receiver/contract.pdf}
\DeclareXref{receiver/def:instrument}{Definition}{Receiver instrument}{../receiver/contract.pdf}
\DeclareXref{receiver/def:pmcex}{Definition}{Exact-window weight of a setting with a declared depth}{../receiver/contract.pdf}
\DeclareXref{receiver/def:residuals}{Definition}{The six residuals}{../receiver/contract.pdf}
\DeclareXref{receiver/def:setting-family}{Definition}{Setting family}{../receiver/contract.pdf}
\DeclareXref{receiver/def:slot}{Definition}{The strengthening slot}{../receiver/contract.pdf}
\DeclareXref{receiver/detector-imperfections/def:dark-composed}{Definition}{The dark-composed family}{../receiver/contract.pdf}
\DeclareXref{receiver/detector-imperfections/def:darkadm}{Definition}{Dark-count channel and its independent product form}{../receiver/contract.pdf}
\DeclareXref{receiver/detector-imperfections/def:darktol}{Definition}{Receiver tolerances deciding dark-channel admissibility}{../receiver/contract.pdf}
\DeclareXref{receiver/detector-imperfections/def:normalized-station}{Definition}{Normalized station and factored loss}{../receiver/contract.pdf}
\DeclareXref{receiver/detector-imperfections/def:pure-loss}{Definition}{Pure-loss channel on the channel modes}{../receiver/contract.pdf}
\DeclareXref{receiver/detector-imperfections/prop:dark-slack-bound}{Proposition}{Multiclick bound under a dark-count channel admitted with slack}{../receiver/contract.pdf}
\DeclareXref{receiver/detector-imperfections/prop:efficiency-normalization}{Proposition}{A receiver as uniform loss followed by its normalized station}{../receiver/contract.pdf}
\DeclareXref{receiver/detector-imperfections/prop:noise-floor}{Proposition}{Loss and dark counts make every dark-composed element definite}{../receiver/contract.pdf}
\DeclareXref{receiver/detector-imperfections/rem:d2slack}{Remark}{Cost of admitting a dark-count channel with slack}{../receiver/contract.pdf}
\DeclareXref{receiver/detector-imperfections/req:dark-channel-record}{Requirement}{The emitted dark-count channel and the scalars of its composed bound}{../receiver/contract.pdf}
\DeclareXref{receiver/detector-imperfections/req:dark-composed-reference}{Requirement}{The emitted blocks are dark-composed and carry their reference}{../receiver/contract.pdf}
\DeclareXref{receiver/detector-imperfections/thm:dark-counts}{Theorem}{Dark counts on a photon-number block-diagonal POVM}{../receiver/contract.pdf}
\DeclareXref{receiver/lem:enclosure-arithmetic}{Lemma}{Arithmetic on complex enclosures}{../receiver/contract.pdf}
\DeclareXref{receiver/prop:slotvalid}{Proposition}{Range of the strengthening-slot weights}{../receiver/contract.pdf}
\DeclareXref{receiver/req:evaluation-label}{Requirement}{The evaluation label on every artifact}{../receiver/contract.pdf}
\DeclareXref{receiver/req:exact-route-decision}{Requirement}{Deciding the exact route and enclosing the port Grams}{../receiver/contract.pdf}
\DeclareXref{receiver/req:rigorous-weight}{Requirement}{Rigorous evaluation of the multiclick weight}{../receiver/contract.pdf}
\DeclareXref{receiver/thm:click-povm}{Theorem}{Click-pattern POVM}{../receiver/contract.pdf}
\DeclareXref{receiver/thm:coarse-graining}{Theorem}{Coarse-graining of the click-pattern POVM}{../receiver/contract.pdf}
\DeclareXref{receiver/thm:compilation-well-defined}{Theorem}{Well-definedness of compilation}{../receiver/contract.pdf}
\DeclareXref{receiver/thm:emitted-bound}{Theorem}{Multiclick bounds carried by the emitted weights}{../receiver/contract.pdf}
\DeclareXref{receiver/thm:enclosed-blocks}{Theorem}{Rigorous enclosures of the photon-number blocks}{../receiver/contract.pdf}
\DeclareXref{receiver/thm:enclosed-weight}{Theorem}{The weight on enclosed bound data bounds the multiclick element}{../receiver/contract.pdf}
\DeclareXref{receiver/thm:exact-and-evaluated}{Theorem}{Exact statements, floating-point and rigorous evaluation}{../receiver/contract.pdf}
\DeclareXref{receiver/thm:representation-independence}{Theorem}{Representation independence of the click-pattern POVM}{../receiver/contract.pdf}
\DeclareXref{transmitter/ass:ens}{Assumption}{Emission of the specified ensemble}{../transmitter/contract.pdf}
\DeclareXref{transmitter/ass:seq}{Assumption}{Sequential round structure}{../transmitter/contract.pdf}
\DeclareXref{transmitter/ass:src}{Assumption}{Retention of the system A in trusted-source mode}{../transmitter/contract.pdf}
\DeclareXref{transmitter/decoy/ass:wcp}{Assumption}{Emission of the declared weak-coherent-pulse states}{../transmitter/contract.pdf}
\DeclareXref{transmitter/def:source-link}{Definition}{Source link}{../transmitter/contract.pdf}
\DeclareXref{transmitter/def:source-segment}{Definition}{Source segment and its interface conditions}{../transmitter/contract.pdf}
\DeclareXref{transmitter/def:source-spec}{Definition}{Source specification}{../transmitter/contract.pdf}
\DeclareXref{transmitter/def:srcert}{Definition}{Source-replacement certificate}{../transmitter/contract.pdf}
\DeclareXref{transmitter/thm:certificate-sufficiency}{Theorem}{Checkability and sufficiency of the source-replacement certificate}{../transmitter/contract.pdf}

\begin{document}
\begin{center}
{\LARGE\bfseries Proof reductions by squashing (reductions)}\\[4pt]
{\large Statements digest, version 0.22.0}
\end{center}
\noindent{\small Every definition, fact, assumption, result and requirement of the \href{contract.pdf}{contract}, with the contract's numbers; remarks, examples, technical steps and all proofs are left out. Each heading opens its statement in the contract.}

\part*{Mathematics}\pdfbookmark[0]{Mathematics}{digest-part-1}

\section*{1\quad Squashing relations and their statistics}\pdfbookmark[1]{1 Squashing relations and their statistics}{digest-group-reductions/def:system-setting}

% claim reductions/def:system-setting
\digestlabel{def:system-setting}{1.1}\pdfbookmark[2]{Definition 1.1 The system setting of the reductions}{digest-reductions/def:system-setting}
\renewcommand*{\RestatedName}{\hyperref{contract.pdf}{}{claim-reductions/def:system-setting}{Definition~1.1}}%
\begin{digestclaim}[The system setting of the reductions]
The
\emph{system setting}
of this contract is the following:
\begin{enumerate}[label=(\alph*),leftmargin=2.4em]
\item the Hilbert space $\Hs{\Bsys}$ of the system $\Bsys$ that Bob's
  physical receiver acts on is separable and may be infinite-dimensional;
\item every target space $\Hs{\Qsys}$ is finite-dimensional, with declared
  dimensions;
\item every side register other than the adversary's registers, that is,
  Alice's retained system and the classical registers, is finite-dimensional,
  with declared dimensions.
\end{enumerate}
\end{digestclaim}

% claim reductions/lem:sep
\digestlabel{lem:sep}{1.2}\pdfbookmark[2]{Lemma 1.2 Separation and tensor adjoints in the infinite-dimensional case}{digest-reductions/lem:sep}
\renewcommand*{\RestatedName}{\hyperref{contract.pdf}{}{claim-reductions/lem:sep}{Lemma~1.2}}%
\begin{digestclaim}[Separation and tensor adjoints in the infinite-dimensional case]
\emph{Setting.} $\Bsys$ is the system Bob's physical receiver acts on and
$\Qsys$ the system it is squashed to. $\Hs{\Bsys}$ is separable and may be
infinite-dimensional, while $\Hs{\Qsys}$ and every side register are
finite-dimensional (Definition~\ref{def:system-setting}, and, for the
adversary's registers, the round model of
\xref{foundations/def:round-model}, every register of which is
finite-dimensional). $\Tcl{\Bsys}$ denotes the
trace-class operators on $\Hs{\Bsys}$; a \emph{channel}
$\Lambda\colon\Tcl{\Bsys}\to\Tcl{\Qsys}$ is a completely positive
trace-preserving map, and its \emph{adjoint}
$\Lambda^{\dagger}\colon\Lin{\Qsys}\to\Lin{\Bsys}$ is the unital completely
positive map fixed by
\begin{equation*}\tag{\getrefnumber{C-eq:pairing}}\digesteqlabel{eq:pairing}
  \tr\big[X\,\Lambda(\rho)\big]=\tr\big[\Lambda^{\dagger}(X)\,\rho\big]
  \qquad\text{for all }X\in\Lin{\Qsys},\ \rho\in\Tcl{\Bsys}.
\end{equation*}
More generally, if $\Lambda\colon\Tcl{\Bsys}\to\Tcl{\Bsys'}$, for a Hilbert
space $\Hs{\Bsys'}$ of any dimension, is only completely positive and trace
non-increasing ($\tr\Lambda(\rho)\le\tr\rho$ for $\rho\succeq0$), then
\eqref{eq:pairing} for all $X\in\Lin{\Bsys'}$ and $\rho\in\Tcl{\Bsys}$ fixes
a completely positive map $\Lambda^{\dagger}\colon\Lin{\Bsys'}\to\Lin{\Bsys}$,
again called its \emph{adjoint}, and $\Lambda^{\dagger}$ is sub-unital:
$\idon{\Bsys}-\Lambda^{\dagger}(\idon{\Bsys'})\succeq0$. Where such a
$\Lambda$ is moreover trace preserving, $\Lambda^{\dagger}$ is unital,
$\Lambda^{\dagger}(\idon{\Bsys'})=\idon{\Bsys}$, whatever the dimension of
$\Hs{\Bsys'}$.

Let $\Hs{\Bsys_1},\dots,\Hs{\Bsys_n}$ be separable and
$\Hs{\Qsys_1},\dots,\Hs{\Qsys_n}$ finite-dimensional.
\begin{enumerate}[label=(\roman*),leftmargin=3.0em]
\item \emph{(Separation.)} If $X\in\Lin{\Bsys_1}$ satisfies $\tr[X\rho]=0$ for
  every $\rho\in\Seq{\Bsys_1}$, then $X=0$.
\item \emph{(Tensor channels and their adjoints.)} If $\Lambda_i\colon
  \Tcl{\Bsys_i}\to\Tcl{\Qsys_i}$ are channels, then exactly one trace-norm
  continuous linear map $\Lambda_1\otimes\dots\otimes\Lambda_n$ from
  $\Tcl{\Bsys_1\cdots\Bsys_n}$, the trace class of the completed tensor
  product of the $\Hs{\Bsys_i}$, to $\Tcl{\Qsys_1\cdots\Qsys_n}$ takes every
  $\rho_1\otimes\dots\otimes\rho_n$ with $\rho_i\in\Tcl{\Bsys_i}$ to
  $\Lambda_1(\rho_1)\otimes\dots\otimes\Lambda_n(\rho_n)$; it is a channel,
  and $(\Lambda_1\otimes\dots\otimes\Lambda_n)^{\dagger}
  =\Lambda_1^{\dagger}\otimes\dots\otimes\Lambda_n^{\dagger}$ on all of
  $\Lin{\Qsys_1\cdots\Qsys_n}$. The same holds with any of the $\Lambda_i$
  replaced by $\idmap_{R_i}$ for a finite-dimensional $R_i$.
\item \emph{(Tensor channels from finite inputs.)} Let
  $\Hs{R_1},\dots,\Hs{R_n}$ be finite-dimensional,
  $\Hs{\Bsys'_1},\dots,\Hs{\Bsys'_n}$ separable, of any dimension, and
  $\Phi_i\colon\Lin{R_i}\to\Tcl{\Bsys'_i}$ completely positive and trace
  preserving. Then exactly one linear map $\Phi_1\otimes\dots\otimes\Phi_n$
  from $\Lin{R_1\cdots R_n}$ to $\Tcl{\Bsys'_1\cdots\Bsys'_n}$, the trace
  class of the completed tensor product of the $\Hs{\Bsys'_i}$, takes every
  $X_1\otimes\dots\otimes X_n$ with $X_i\in\Lin{R_i}$ to
  $\Phi_1(X_1)\otimes\dots\otimes\Phi_n(X_n)$, and it is completely positive
  and trace preserving. The identity map $\idmap_{R}$ of a finite-dimensional
  $R$ is such a $\Phi_i$.
\end{enumerate}
\end{digestclaim}

% claim reductions/def:squashing
\digestlabel{def:squashing}{1.3}\pdfbookmark[2]{Definition 1.3 Squashing relation}{digest-reductions/def:squashing}
\renewcommand*{\RestatedName}{\hyperref{contract.pdf}{}{claim-reductions/def:squashing}{Definition~1.3}}%
\begin{digestclaim}[Squashing relation]
Let $M=\{M_v\}_{v\in\Vset}$ be a POVM on $\Bsys$ and $F=\{F_v\}_{v\in\Vset}$ a
POVM on $\Qsys$ over the same finite alphabet $\Vset$, and let
$\hat O=(\hat O_k)_{k=1}^{m}$ and $O=(O_k)_{k=1}^{m}$ be tuples of Hermitian
operators $\hat O_k\in\Lin{\Qsys}$, $O_k\in\Lin{\Bsys}$. For a channel
$\Lambda\colon\Tcl{\Bsys}\to\Tcl{\Qsys}$, with adjoint $\Lambda^{\dagger}$ as
in Lemma~\ref{lem:sep}, write
\begin{equation*}\tag{\getrefnumber{C-eq:squashing}}\digesteqlabel{eq:squashing}
  (M;O)\;\dsq_{\Lambda}\;(F;\hat O)
  \quad:\Longleftrightarrow\quad
  \Lambda^{\dagger}(F_v)=M_v\ \ \forall v\in\Vset
  \ \ \text{and}\ \
  \Lambda^{\dagger}(\hat O_k)=O_k\ \ \forall k\in\{1,\dots,m\},
\end{equation*}
and call \eqref{eq:squashing} a
\emph{squashing relation}:
$M$ \emph{squashes to} the
\emph{target POVM}
$F$ via the
\emph{squashing channel}
$\Lambda$, which \emph{transports} $O$ to $\hat O$. The operators
$\hat O_k$ are the relation's
\emph{carried operators}.
The empty case $m=0$ is written $M\dsq_{\Lambda}F$, and $\dsq$ replaces
$\dsq_{\Lambda}$ when only the existence of a squashing channel is asserted.
We call $\Lambda^{\dagger}(\hat O_k)$ the
\emph{pull-back}
of $\hat O_k$.
\end{digestclaim}

% claim reductions/lem:statistics
\digestlabel{lem:statistics}{1.4}\pdfbookmark[2]{Lemma 1.4 Squashing as statistics equivalence}{digest-reductions/lem:statistics}
\renewcommand*{\RestatedName}{\hyperref{contract.pdf}{}{claim-reductions/lem:statistics}{Lemma~1.4}}%
\begin{digestclaim}[Squashing as statistics equivalence]
Let $M,F,O,\hat O$ be as in Definition~\ref{def:squashing} and let
$\Lambda\colon\Tcl{\Bsys}\to\Tcl{\Qsys}$ be a channel. The following are
equivalent.
\begin{enumerate}[label=(\roman*),leftmargin=3.0em]
\item $(M;O)\dsq_{\Lambda}(F;\hat O)$.
\item For every $\rho\in\Seq{\Bsys}$, every $v$ and every $k$,
  $\tr[M_v\rho]=\tr[F_v\Lambda(\rho)]$ and
  $\tr[O_k\rho]=\tr[\hat O_k\Lambda(\rho)]$.
\item For every finite-dimensional side system $R$, every
  $\rho\in\Seq{\Bsys R}$, every $X\in\Lin{R}$, every $v$ and every $k$, writing
  $\sigma\coloneqq(\Lambda\otimes\idmap_R)(\rho)$,
  \begin{equation*}
    \tr\big[(M_v\otimes X)\rho\big]=\tr\big[(F_v\otimes X)\sigma\big]
    \qquad\text{and}\qquad
    \tr\big[(O_k\otimes X)\rho\big]=\tr\big[(\hat O_k\otimes X)\sigma\big].
  \end{equation*}
\end{enumerate}
Moreover, two channels $\Lambda,\Lambda'$ that both realize the relation
satisfy $\Lambda^{\dagger}=\Lambda'^{\dagger}$ on
$\operatorname{span}\big(\{F_v\}_{v\in\Vset}\cup\{\hat O_k\}_{k=1}^{m}\big)$.
\end{digestclaim}

% claim reductions/cor:pullback
\digestlabel{cor:pullback}{1.7}\pdfbookmark[2]{Corollary 1.7 Pull-back of operator bounds}{digest-reductions/cor:pullback}
\renewcommand*{\RestatedName}{\hyperref{contract.pdf}{}{claim-reductions/cor:pullback}{Corollary~1.7}}%
\begin{digestclaim}[Pull-back of operator bounds]
Let $(M;O)\dsq_{\Lambda}(F;\hat O)$, let $R$ be a finite-dimensional side
system, let $\rho\in\Seq{\Bsys R}$ and let
$\sigma=(\Lambda\otimes\idmap_R)(\rho)$. Then for every $k$,
\begin{equation*}\tag{\getrefnumber{C-eq:transport}}\digesteqlabel{eq:transport}
  \tr_{\Qsys}\big[(\hat O_k\otimes\idon{R})\,\sigma\big]
  =\tr_{\Bsys}\big[(O_k\otimes\idon{R})\,\rho\big]
\end{equation*}
as operators on $R$. Consequently, on every such $\sigma$:
\begin{enumerate}[label=(\roman*),leftmargin=3.0em]
\item \emph{Exact values.} If $O_k=c\,\idon{\Bsys}$ then
  $\tr_{\Qsys}[(\hat O_k\otimes\idon{R})\sigma]=c\,\sigma_R$; in particular
  $\tr[\hat O_k\sigma]=c$, and at $c=0$,
  $\tr[(\hat O_k\otimes X)\sigma]=0$ for every $X\in\Lin{R}$.
\item \emph{Positivity.} If $O_k\succeq0$ then
  $\tr_{\Qsys}[(\hat O_k\otimes\idon{R})\sigma]\succeq0$; in particular
  $\tr[\hat O_k\sigma]\ge0$.
\item \emph{Spectral bounds.} If
  $a\,\idon{\Bsys}\preceq O_k\preceq b\,\idon{\Bsys}$ then
  $a\,\sigma_R\preceq\tr_{\Qsys}[(\hat O_k\otimes\idon{R})\sigma]
  \preceq b\,\sigma_R$; in particular
  $a\le\tr[\hat O_k\sigma]\le b$.
\end{enumerate}
\end{digestclaim}

% claim reductions/def:certificate
\digestlabel{def:certificate}{1.8}\pdfbookmark[2]{Definition 1.8 Certificates of transported constraints}{digest-reductions/def:certificate}
\renewcommand*{\RestatedName}{\hyperref{contract.pdf}{}{claim-reductions/def:certificate}{Definition~1.8}}%
\begin{digestclaim}[Certificates of transported constraints]
The
\emph{certificate}
of a transported constraint is the hypothesis that makes it hold: for a
constraint from Corollary~\ref{cor:pullback}, the hypothesis of the case that
yields it, an operator identity in its case~(i) and an operator inequality in
its cases (ii)--(iii), both facts about the device and the channel, valid
whatever the input; and, for an assumed constraint of
Corollary~\ref{cor:transported}\,(ii), an assumption on the input states.
\end{digestclaim}

% claim reductions/def:msm
\digestlabel{def:msm}{1.10}\pdfbookmark[2]{Definition 1.10 Measurement channel}{digest-reductions/def:msm}
\renewcommand*{\RestatedName}{\hyperref{contract.pdf}{}{claim-reductions/def:msm}{Definition~1.10}}%
\begin{digestclaim}[Measurement channel]
For a POVM $M=\{M_v\}_{v\in\Vset}$ on $\Bsys$ and a classical register $Y$ with
alphabet $\Vset$, the
\emph{measurement channel}
of $M$ is
\begin{equation*}\tag{\getrefnumber{C-eq:msm}}\digesteqlabel{eq:msm}
  \redNmsm_M\colon\Tcl{\Bsys}\to\Tcl{Y},\qquad
  \redNmsm_M(\rho)\coloneqq\sum_{v\in\Vset}\tr\big[M_v\rho\big]\,\dyad{v}_{Y}.
\end{equation*}
\end{digestclaim}

% claim reductions/lem:msmfact
\digestlabel{lem:msmfact}{1.11}\pdfbookmark[2]{Lemma 1.11 Measurement-channel factorization}{digest-reductions/lem:msmfact}
\renewcommand*{\RestatedName}{\hyperref{contract.pdf}{}{claim-reductions/lem:msmfact}{Lemma~1.11}}%
\begin{digestclaim}[Measurement-channel factorization]
Let $M,F$ be POVMs on $\Bsys,\Qsys$ over $\Vset$ and let
$\Lambda\colon\Tcl{\Bsys}\to\Tcl{\Qsys}$ be a channel. Then
\begin{equation*}\tag{\getrefnumber{C-eq:msmfact}}\digesteqlabel{eq:msmfact}
  M\dsq_{\Lambda}F
  \qquad\Longleftrightarrow\qquad
  \redNmsm_M=\redNmsm_F\circ\Lambda .
\end{equation*}
\end{digestclaim}

\section*{2\quad Protocols, security criteria and exact operational reductions}\pdfbookmark[1]{2 Protocols, security criteria and exact operational reductions}{digest-group-reductions/def:protocol}

% claim reductions/def:protocol
\digestlabel{def:protocol}{2.1}\pdfbookmark[2]{Definition 2.1 Measurement-factorizing n-round protocol}{digest-reductions/def:protocol}
\renewcommand*{\RestatedName}{\hyperref{contract.pdf}{}{claim-reductions/def:protocol}{Definition~2.1}}%
\begin{digestclaim}[Measurement-factorizing $n$-round protocol]
An
\emph{$n$-round protocol}
with receiver POVM $M$ is a channel
\begin{equation*}\tag{\getrefnumber{C-eq:protocol}}\digesteqlabel{eq:protocol}
  \Eprot_M\;=\;\Kfin\circ\Grnd_n\circ\dots\circ\Grnd_1,
  \qquad
  \Grnd_j\;=\;\bar{\Grnd}_j\circ\big(\idmap\otimes\redNmsm_M\big),
\end{equation*}
in which the measurement channel $\redNmsm_M$ of
Definition~\ref{def:msm} acts on the round-$j$
\emph{measured system}
$\Bsys_j$, the quantum system that the
\emph{measuring parties},
the parties the declaration names as measuring, receive in round $j$, and no
other system, and the identity acts on everything else. $\bar{\Grnd}_j$ is the
round's remaining processing (Alice's measurement on $A_j$, the round's
announcements appended to the public record, and the retained per-round data),
and $\Kfin$ is the end-of-protocol processing: sifting, parameter estimation
and the accept/abort decision, error correction, verification, privacy
amplification and the key maps. The
\emph{squashed protocol}
$\Eprot_F$ is \eqref{eq:protocol} with $\Bsys_j$ replaced by $\Qsys_j$, $M$ by
$F$, and $\bar{\Grnd}_j$, $\Kfin$ unchanged. One POVM and one channel serve
all $n$ rounds: the same $M$, and in the squashed protocol the same $F$ and
the same squashing channel, stands in every $\Grnd_j$, with no
dependence on $j$. That the measuring devices measure by the same $M$ in
every round, with no memory between rounds, is Assumption~\ref{ass:round};
$F$ and the squashing channel are objects of the analysis, chosen once for all
rounds. That is what makes the composite of the per-round channels a tensor
power, such as $\Lambda^{\otimes n}$ for the squashing channel, rather than a
product of $n$ different maps.

\emph{One measuring party, or several.} $M$ is one POVM on the measured
system, over the alphabet of the measuring parties' joint outcome. Where one
party measures, which is the one-party case, $\Bsys_j$ is that party's
system, Bob's in a prepare-and-measure declaration. Where several parties
measure, $\Bsys_j=\bigotimes_f\Bsys_{f,j}$ is the product of their
party-scoped parts, each factor $\Bsys_{f,j}$ belonging to one of them, and
\eqref{eq:protocol} places \emph{all} of their measurements at the one step
$\redNmsm_M$, before the round's remaining processing $\bar{\Grnd}_j$ and so
before every public message of the round.
The declaration says on which side Alice's measurement sits, and each reading
is covered by exactly one clause of the first paragraph: the measurement on
$A_j$ listed in $\bar{\Grnd}_j$ is that of a party the declaration does
\emph{not} name as measuring (Alice's, in a prepare-and-measure declaration),
while where she \emph{is} named as measuring, her system is a factor
$\Bsys_{f,j}$ of the measured system and her measurement is a factor of
$\redNmsm_M$, not a step of $\bar{\Grnd}_j$. Either way $\bar{\Grnd}_j$ and
$\Kfin$ touch no $\Bsys_j$, which is all the theorems of this contract read.
A declaration realizes the second reading by an embedding: it declares
Alice's system $A$
one-dimensional, $d_A=1$, folds her measured system into the sectors of
$\Qsys$ and her outcome components into $v$, and names the components she
holds in its party partition (\xref{protocol-compiler/def:declaration}). That the
physical protocol has that order, every public message of round $j$ sent after
every measurement of round $j$ whichever party sends it and whichever party
measured, is Assumption~A-SEQ of the transmitter
(\xref{transmitter/ass:seq}). The definition says nothing about which party
sends which message or about a message's dependence on earlier messages of
its own round, and a quantum operation of round $j$ that depended on a
message of round $j$ is outside it, as it is outside A-SEQ.
\end{digestclaim}

% claim reductions/def:criterion
\digestlabel{def:criterion}{2.2}\pdfbookmark[2]{Definition 2.2 Security criterion}{digest-reductions/def:criterion}
\renewcommand*{\RestatedName}{\hyperref{contract.pdf}{}{claim-reductions/def:criterion}{Definition~2.2}}%
\begin{digestclaim}[Security criterion]
A
\emph{security criterion}
is a family $f=(f_E)_E$ of functions, one for every admissible adversary
system $E$, defined on states of the protocol's output registers together with
$E$; an \emph{admissible adversary system} is a finite-dimensional system,
the form the round model of \xref{foundations/def:round-model} gives
the adversary's registers. Its value
on the protocol $\Eprot_M$ at an input $\rho$, a state of the systems the
protocol acts on together with $E$, is
\begin{equation*}\tag{\getrefnumber{C-eq:criterion}}\digesteqlabel{eq:criterion}
  \Crit(\Eprot_M,\rho)\coloneqq
  f_E\big((\Eprot_M\otimes\idmap_E)(\rho)\big).
\end{equation*}
\end{digestclaim}

% claim reductions/def:opreduction
\digestlabel{def:opreduction}{2.3}\pdfbookmark[2]{Definition 2.3 Exact operational reduction}{digest-reductions/def:opreduction}
\renewcommand*{\RestatedName}{\hyperref{contract.pdf}{}{claim-reductions/def:opreduction}{Definition~2.3}}%
\begin{digestclaim}[Exact operational reduction]
A
\emph{protocol presentation}
is a triple $\mathsf{P}=(\mathbb{A},W,\mathsf{out})$: a set $\mathbb{A}$, each
element of which is called an
\emph{attack strategy},
a tuple $W$ of
\emph{output registers},
and a map $\mathsf{out}\colon\mathbb{A}\to\Sleq{W}$ assigning to each
strategy the
\emph{full output}
the protocol produces under it. The full output is the state on \emph{every}
register a security criterion may read (Definition~\ref{def:criterion}): both
parties' key registers, the public transcript with the acceptance and abort
data, and the adversary's final register, and no marginal of them.

Let $\mathsf{P}=(\mathbb{A},W,\mathsf{out})$ and
$\mathsf{P}'=(\mathbb{A}',W,\mathsf{out}')$ be presentations \emph{over the
same output registers} $W$. An
\emph{exact operational reduction}
from $\mathsf{P}$ to $\mathsf{P}'$ is a map $r\colon\mathbb{A}\to\mathbb{A}'$
with
\begin{equation*}\tag{\getrefnumber{C-eq:opreduction}}\digesteqlabel{eq:opreduction}
  \mathsf{out}(a)\;=\;\mathsf{out}'\big(r(a)\big)
  \qquad\text{for every }a\in\mathbb{A} .
\end{equation*}
We call $\mathbb{A}$ the
\emph{physical strategy set}
and $\mathbb{A}'$ the
\emph{analysed strategy set},
equivalently the \emph{admissible} strategy set, and
$r(\mathbb{A})\subseteq\mathbb{A}'$ the reduction's
\emph{image}.

\emph{The conditions, in full.}
\begin{enumerate}[label=(R\arabic*),leftmargin=3.0em]
\item \emph{(one output tuple)} the two presentations carry literally the same
  output registers, and which register of $\mathsf{P}'$ is which register of
  $\mathsf{P}$ is part of the data rather than inferred. Where a model
  protocol's adversary register is a \emph{different} system, an
  identification of it with the physical one has to be declared and the
  criterion has to be invariant under that identification; this contract
  declares none and needs none, the adversary's registers being literally the
  same system in both of its instances;
\item \emph{(exact equality)} \eqref{eq:opreduction} is equality of the two
  states on the nose: not of a marginal, not up to a distance, and not up to
  any map on $W$;
\item \emph{(totality and typing)} $r$ is defined at every strategy of
  $\mathbb{A}$ and takes its values in $\mathbb{A}'$. That its values lie in
  $\mathbb{A}'$ is a condition on $r$ and on $\mathbb{A}'$ together, and it is
  the only place where a restriction of the analysed set has to be earned.
\end{enumerate}

\emph{What is not required, each item being used by an instance in this
contract.} $r$ need not be injective, need not be surjective onto
$\mathbb{A}'$, and need not preserve any structure of $\mathbb{A}$; the
analysed set $\mathbb{A}'$ may be, and in this contract typically is, strictly
larger than the image, because what a numerical analysis can optimize over is
a set cut out by finitely many conditions and not an image
(Definition~\ref{def:attset}); no metric, topology, order, linear or convex
structure on $\mathbb{A}$, on $\mathbb{A}'$ or on the output states is used;
$\mathsf{out}$ and $\mathsf{out}'$ need not be linear, affine, continuous or
given by channels, and a strategy need not be a channel or a tuple of
channels (in Theorem~\ref{thm:transfer} a strategy is a state); and nothing
is asserted in the other direction, a bound over $\mathbb{A}$ transferring to
$\mathbb{A}'$ under no hypothesis of this definition.

\emph{Two presentations, one per adversary system.} A security criterion in
the sense of Definition~\ref{def:criterion} is a \emph{family} $(f_E)_E$, one
function per admissible adversary system. Fixing $E$ fixes $W$ and turns each
$f_E$ into one function on $\Sleq{W}$, so the quantification over $E$ in that
definition is a quantification over presentations: a statement of this
contract that holds ``for every admissible $E$'' exhibits one presentation
pair and one reduction per $E$, and the transfer of
Lemma~\ref{lem:opreduction-transfer} is applied to each.
\end{digestclaim}

% claim reductions/lem:opreduction-transfer
\digestlabel{lem:opreduction-transfer}{2.4}\pdfbookmark[2]{Lemma 2.4 Transfer and composition of exact operational reductions}{digest-reductions/lem:opreduction-transfer}
\renewcommand*{\RestatedName}{\hyperref{contract.pdf}{}{claim-reductions/lem:opreduction-transfer}{Lemma~2.4}}%
\begin{digestclaim}[Transfer and composition of exact operational reductions]
Let $r$ be an exact operational reduction from
$\mathsf{P}=(\mathbb{A},W,\mathsf{out})$ to
$\mathsf{P}'=(\mathbb{A}',W,\mathsf{out}')$ in the sense of
Definition~\ref{def:opreduction}. Write $\mathcal{F}(W)$ for the set of all
functions $\Sleq{W}\to[-\infty,+\infty]$; a security criterion in the sense of
Definition~\ref{def:criterion}, at a fixed adversary system, is one such
function.
\begin{enumerate}[label=(\roman*),leftmargin=3.0em]
\item \emph{(transfer, with the same $\varepsilon$ and for every criterion at
  once)} for every $f\in\mathcal{F}(W)$ and every $a\in\mathbb{A}$,
  $f(\mathsf{out}(a))=f(\mathsf{out}'(r(a)))$, and consequently
  \begin{equation*}\tag{\getrefnumber{C-eq:opsup}}\digesteqlabel{eq:opsup}
    \sup_{a\in\mathbb{A}}f\big(\mathsf{out}(a)\big)
    \;=\;\sup_{b\in r(\mathbb{A})}f\big(\mathsf{out}'(b)\big)
    \;\le\;\sup_{b\in\mathbb{A}'}f\big(\mathsf{out}'(b)\big) .
  \end{equation*}
  In particular, if $f(\mathsf{out}'(b))\le\varepsilon$ for every
  $b\in\mathbb{A}'$, then $f(\mathsf{out}(a))\le\varepsilon$ for every
  $a\in\mathbb{A}$, with the \emph{same} $\varepsilon$, and simultaneously
  for every $f\in\mathcal{F}(W)$, no property of $f$ being used;
\item \emph{(enlargement of the analysed set)} if
  $\mathbb{A}'\subseteq\mathbb{A}''$ and
  $\mathsf{out}''|_{\mathbb{A}'}=\mathsf{out}'$, then the same $r$ is an
  exact operational reduction from $\mathsf{P}$ to
  $\mathsf{P}''=(\mathbb{A}'',W,\mathsf{out}'')$. So it is enough that the
  image $r(\mathbb{A})$ lie in the analysed set, and the last inequality of
  \eqref{eq:opsup} is where that slack is spent;
\item \emph{(restriction of the physical set)} if
  $\mathbb{A}_0\subseteq\mathbb{A}$, then $r|_{\mathbb{A}_0}$ is an exact
  operational reduction from $(\mathbb{A}_0,W,\mathsf{out}|_{\mathbb{A}_0})$ to
  $\mathsf{P}'$;
\item \emph{(one map, many analysed sets)} let $(\mathbb{A}'_k)_{k\in K}$ be a
  family of subsets of $\mathbb{A}'$ such that \emph{one and the same} $r$
  satisfies $r(\mathbb{A})\subseteq\mathbb{A}'_k$ for every $k\in K$. Then $r$
  is an exact operational reduction from $\mathsf{P}$ to
  $\big(\bigcap_{k\in K}\mathbb{A}'_k,W,\mathsf{out}'\big)$, and by (i) a bound
  over that intersection transfers. The hypothesis cannot be weakened to the
  existence, for each $k$, of \emph{some} exact operational reduction into
  $\mathbb{A}'_k$: let $W$ be one classical bit, $\mathbb{A}=\{a\}$,
  $\mathbb{A}'=\{b_1,b_2\}$ with $b_1\ne b_2$ and
  $\mathsf{out}(a)=\mathsf{out}'(b_1)=\mathsf{out}'(b_2)=\dyad{0}$, and
  $\mathbb{A}'_k=\{b_k\}$. Each $r_k\colon a\mapsto b_k$ is an exact
  operational reduction into $\mathbb{A}'_k$, the intersection is empty, so
  ``$f(\mathsf{out}'(b))\le0$ for every $b$ in the intersection'' holds
  vacuously; yet $f(\mathsf{out}(a))=1$ for
  $f(\rho)\coloneqq\tr[\dyad{0}\rho]$, so the conclusion of (i) is false for
  the intersection;
\item \emph{(composition, and how the analysed sets compose)} let
  $\mathbb{A}'_0\subseteq\mathbb{A}'$ contain $r(\mathbb{A})$ and let $r'$ be an
  exact operational reduction from
  $\big(\mathbb{A}'_0,W,\mathsf{out}'|_{\mathbb{A}'_0}\big)$ to
  $\mathsf{P}''=(\mathbb{A}'',W,\mathsf{out}'')$. Then $r'\circ r$ is an exact
  operational reduction from $\mathsf{P}$ to $\mathsf{P}''$, with image
  $r'(r(\mathbb{A}))$. The analysed set the composite needs is the second
  reduction's, or any subset of $\mathbb{A}''$ containing
  $r'(r(\mathbb{A}))$; the first reduction's analysed set enters only as a
  domain the second must cover, which is what
  $r(\mathbb{A})\subseteq\mathbb{A}'_0$ says. Composition is associative and
  the identity map on $\mathbb{A}$ is an exact operational reduction from
  $\mathsf{P}$ to itself;
\item \emph{(the direction, and what the slack costs)} nothing above bounds
  $\sup_{\mathbb{A}'}f\circ\mathsf{out}'$ by
  $\sup_{\mathbb{A}}f\circ\mathsf{out}$, and the first can be strictly larger:
  take $\mathbb{A}=\{a\}$, $\mathbb{A}'=\{r(a),b\}$ with
  $\mathsf{out}'(b)\ne\mathsf{out}(a)$ and an $f$ separating the two states.
  An analysed set larger than the image therefore costs \emph{tightness} and
  never soundness.
\end{enumerate}
\end{digestclaim}

\section*{3\quad Output-state equality and security transfer}\pdfbookmark[1]{3 Output-state equality and security transfer}{digest-group-reductions/thm:factorization}

% claim reductions/thm:factorization
\digestlabel{thm:factorization}{3.1}\pdfbookmark[2]{Theorem 3.1 Equality of the physical and squashed output states}{digest-reductions/thm:factorization}
\renewcommand*{\RestatedName}{\hyperref{contract.pdf}{}{claim-reductions/thm:factorization}{Theorem~3.1}}%
\begin{digestclaim}[Equality of the physical and squashed output states]
Let $M\dsq_{\Lambda}F$ over $\Vset$ and let $\Eprot_M,\Eprot_F$ be as in
Definition~\ref{def:protocol}. Then
\begin{equation*}\tag{\getrefnumber{C-eq:factorization}}\digesteqlabel{eq:factorization}
  \Eprot_M\;=\;\Eprot_F\circ\big(\idmap\otimes\Lambda^{\otimes n}\big),
\end{equation*}
the identity acting on every register other than $\Bsys_1\dots\Bsys_n$; and
consequently, for every finite-dimensional adversary system $E$ and every input
$\rho$,
\begin{equation*}\tag{\getrefnumber{C-eq:sameoutput}}\digesteqlabel{eq:sameoutput}
  (\Eprot_M\otimes\idmap_E)(\rho)
  =(\Eprot_F\otimes\idmap_E)\big(\sigma\big),
  \qquad
  \sigma\coloneqq\big(\idmap\otimes\Lambda^{\otimes n}\otimes\idmap_E\big)(\rho).
\end{equation*}
In particular $\Crit(\Eprot_M,\rho)=\Crit(\Eprot_F,\sigma)$ for every
security criterion in the sense of Definition~\ref{def:criterion}.
\end{digestclaim}

% claim reductions/def:roundmodel
\digestlabel{def:roundmodel}{3.2}\pdfbookmark[2]{Definition 3.2 Round model and attack sequences, typed by source mode}{digest-reductions/def:roundmodel}
\renewcommand*{\RestatedName}{\hyperref{contract.pdf}{}{claim-reductions/def:roundmodel}{Definition~3.2}}%
\begin{digestclaim}[Round model and attack sequences, typed by source mode]
The round model of the protocol \eqref{eq:protocol} is that of
\xref{foundations/def:round-model} with per-round system $\Bsys$, Alice's
system $A$ and the declared source mode. An attack sequence for the
protocol is thus a family $(\Aev_j)_{j=1}^{n}$ of channels whose output carries
the systems the adversary supplies in the declared source mode, the display
being that definition's at $S=\Bsys$:
\begin{equation*}\tag{\getrefnumber{C-eq:attackchannel}}\digesteqlabel{eq:attackchannel}
  \Aev_j\in
  \begin{cases}
    \mathrm{CPTP}\big(E_{j-1},\ \Bsys_j\,E'_j\big),
      & \text{trusted-source mode},\\[0.3ex]
    \mathrm{CPTP}\big(E_{j-1},\ A_j\,\Bsys_j\,E'_j\big),
      & \text{untrusted-source mode},
  \end{cases}
\end{equation*}
where $E_0$ is Eve's initial register, $E'_j$ her outgoing memory, and $E_j$
collects $E'_j$ together with whatever public data of round $j$ she retains. In
trusted-source mode $A_j$ is \emph{not} an output of $\Aev_j$: it is
supplied by Alice and carried into the round as a separate input. The
protocol's output under the sequence is obtained from \eqref{eq:protocol} by
supplying $\Bsys_j$ --- and, in untrusted-source mode, $A_j$ as well
--- as the output of $\Aev_j$ in round $j$. The corresponding
\emph{squashed sequence} is
\begin{equation*}\tag{\getrefnumber{C-eq:composite}}\digesteqlabel{eq:composite}
  \Aev^{\sq}_j\;\coloneqq\;
  \begin{cases}
    \big(\Lambda\otimes\idmap_{E'_j}\big)\circ\Aev_j
      \in\mathrm{CPTP}\big(E_{j-1},\Qsys_jE'_j\big),
      & \text{trusted-source mode},\\[0.3ex]
    \big(\idmap_{A}\otimes\Lambda\otimes\idmap_{E'_j}\big)\circ\Aev_j
      \in\mathrm{CPTP}\big(E_{j-1},A_j\Qsys_jE'_j\big),
      & \text{untrusted-source mode},
  \end{cases}
\end{equation*}
the squashing channel acting on the $\Bsys_j$ factor in both cases and
on nothing else.

\emph{Several measuring parties.} Where the measured system of a
round is the product $\Bsys_j=\bigotimes_f\Bsys_{f,j}$ of party-scoped parts
(Definition~\ref{def:protocol}), the squashing channel is the product
$\Lambda=\bigotimes_f\Lambda_f$ of one channel per part,
$\Qsys_j=\bigotimes_f\Qsys_{f,j}$, and the alphabet $\Vset$ of $M$ and $F$ is
the product of the parts' alphabets. Neither line of \eqref{eq:attackchannel} changes:
$\Bsys_j$ is one output system of $\Aev_j$ whatever its factorisation, so the
adversary supplies all of its parts in one channel and is given the same
freedom as with one measuring party; in particular the parts may be
arbitrarily correlated, and the adversary is not required to act on them
independently. Nor does \eqref{eq:composite} change: the product channel
still acts on the $\Bsys_j$ factor alone, leaving $A_j$ and $E'_j$ untouched.
That the product channel realizes a squashing relation for the product family
is Proposition~\ref{prop:tensor}, which this definition does not assert. With
one measuring party the product has one factor.
\end{digestclaim}

% claim reductions/cor:sequential
\digestlabel{cor:sequential}{3.3}\pdfbookmark[2]{Corollary 3.3 Output-state equality against adaptive attack sequences}{digest-reductions/cor:sequential}
\renewcommand*{\RestatedName}{\hyperref{contract.pdf}{}{claim-reductions/cor:sequential}{Corollary~3.3}}%
\begin{digestclaim}[Output-state equality against adaptive attack sequences]
Let $M\dsq_{\Lambda}F$ and fix a source mode. For every attack sequence
$(\Aev_j)_{j=1}^{n}$ of the corresponding type \eqref{eq:attackchannel}, the
output state of the physical protocol $\Eprot_M$ under $(\Aev_j)_j$ equals the
output state of the squashed protocol $\Eprot_F$ under $(\Aev^{\sq}_j)_j$ of
the same branch of \eqref{eq:composite}. Consequently every security criterion
takes the same value on the two, and if $\Eprot_F$ is $\varepsilon$-secure
against every sequence drawn from families $\AttSet^{\sq}_j$ with
\begin{equation*}\tag{\getrefnumber{C-eq:seqcover}}\digesteqlabel{eq:seqcover}
  \AttSet^{\sq}_j\;\supseteq\;
  \begin{cases}
    \big(\Lambda\otimes\idmap\big)\circ
      \mathrm{CPTP}\big(E_{j-1},\Bsys_jE'_j\big),
      & \text{trusted-source mode},\\[0.3ex]
    \big(\idmap_{A}\otimes\Lambda\otimes\idmap\big)\circ
      \mathrm{CPTP}\big(E_{j-1},A_j\Bsys_jE'_j\big),
      & \text{untrusted-source mode},
  \end{cases}
\end{equation*}
then $\Eprot_M$ is $\varepsilon$-secure against \emph{all} attack sequences of
that type.
\end{digestclaim}

% claim reductions/prop:platform
\digestlabel{prop:platform}{3.4}\pdfbookmark[2]{Proposition 3.4 Measurement factorization of the platform's protocol channel}{digest-reductions/prop:platform}
\renewcommand*{\RestatedName}{\hyperref{contract.pdf}{}{claim-reductions/prop:platform}{Proposition~3.4}}%
\begin{digestclaim}[Measurement factorization of the platform's protocol channel]
Let $\Decl$ be a protocol declaration in the sense of
\xref{meat-core/def:declaration} that satisfies the conditions of
\xref{meat-core/def:admissible-declaration}, in the format in which $A$ carries one sector, one
system is measured, Alice's measurement is declared as a POVM
$\{M_x\}_{x\in\Xset}$, the public discussion is two messages neither of which
reads the other, and no entry is scoped to a party. Write $a$ for its joint
announcement map, $\phi$ for its acceptance-statistic map and $r$ for its
raw-key map, and let $\Mprot^{\sq}$ be the single-round protocol channel it
generates, on $A\Qsys$ with the declared POVM $F=\{F_v\}_{v\in\Vset}$ on the
measured system $\Qsys$, the system the declaration's measuring parties hold
(Definition~\ref{def:protocol}). Let $M=\{M_v\}_{v\in\Vset}$ be a POVM on a
separable $\Bsys$, and let $\Mprot^{\mathrm{ph}}$ be the map that the same
components define with $(\Bsys,M)$ in place of $(\Qsys,F)$,
\begin{equation*}\tag{\getrefnumber{C-eq:platformph}}\digesteqlabel{eq:platformph}
  \Mprot^{\mathrm{ph}}(\omega)\coloneqq\sum_{x,v}
  \dyad{r(x,v)}_{Z}\otimes\dyad{\phi(a(x,v))}_{\Chat}\otimes
  \dyad{a(x,v)}_{I}\;\tr\big[(M_x\otimes M_v)\,\omega\big],
  \qquad\omega\in\Tcl{A\Bsys}.
\end{equation*}
Then:
\begin{enumerate}[label=(\roman*),leftmargin=3.0em]
\item both maps satisfy the hypothesis of Definition~\ref{def:protocol} with
  one remaining processing $\bar{\Grnd}$,
  \begin{equation*}
    \Mprot^{\sq}=\bar{\Grnd}\circ\big(\idmap_A\otimes\redNmsm_F\big),
    \qquad
    \Mprot^{\mathrm{ph}}=\bar{\Grnd}\circ\big(\idmap_A\otimes\redNmsm_M\big),
  \end{equation*}
  where $\bar{\Grnd}$ measures $A$ with Alice's POVM $\{M_x\}$ and applies to
  the pair $(x,v)$ the classical map $\dyad{x,v}\mapsto
  \dyad{r(x,v)}_{Z}\otimes\dyad{\phi(a(x,v))}_{\Chat}\otimes\dyad{a(x,v)}_{I}$.
  Alice's measurement is thus a step of $\bar{\Grnd}$, where
  Definition~\ref{def:protocol} places the measurement of a party the
  declaration does not name as measuring, and not of the measurement channel;
\item if $M\dsq_{\Lambda}F$, then for every finite-dimensional $E$ and every
  $\omega\in\Sleq{A\Bsys E}$
  \begin{equation*}\tag{\getrefnumber{C-eq:platformfact}}\digesteqlabel{eq:platformfact}
    \big(\Mprot^{\mathrm{ph}}\otimes\idmap_E\big)(\omega)
    =\big(\Mprot^{\sq}\otimes\idmap_E\big)
     \big((\idmap_A\otimes\Lambda\otimes\idmap_E)(\omega)\big).
  \end{equation*}
\end{enumerate}
\end{digestclaim}

% claim reductions/def:sound
\digestlabel{def:sound}{3.10}\pdfbookmark[2]{Definition 3.10 Reachable states and sound constraints}{digest-reductions/def:sound}
\renewcommand*{\RestatedName}{\hyperref{contract.pdf}{}{claim-reductions/def:sound}{Definition~3.10}}%
\begin{digestclaim}[Reachable states and sound constraints]
Let $M\dsq_{\Lambda}F$ and let $\Sph=(\Sph_E)_E$ be a family of
\emph{physical inputs}, one set
$\Sph_E\subseteq\Seq{A^{n}\Bsys^{n}E}$ per admissible adversary system $E$. A
state $\sigma\in\Seq{A^{n}\Qsys^{n}E}$ is \emph{reachable} if
\begin{equation*}\tag{\getrefnumber{C-eq:reachable}}\digesteqlabel{eq:reachable}
  \sigma=\big(\idmap_{A^{n}}\otimes\Lambda^{\otimes n}\otimes\idmap_E\big)(\rho)
  \qquad\text{for some }\rho\in\Sph_E ,
\end{equation*}
and $\Srch_E$ denotes the set of reachable states. A family $C=(C_E)_E$ of
subsets $C_E\subseteq\Seq{A^{n}\Qsys^{n}E}$, and any constraint whose solution
set it is, is \emph{sound} if $\Srch_E\subseteq C_E$ for every admissible $E$. Both notions
depend on $\Lambda$ and on $\Sph$, which are suppressed in the notation and
fixed throughout any one analysis.
\end{digestclaim}

% claim reductions/thm:transfer
\digestlabel{thm:transfer}{3.11}\pdfbookmark[2]{Theorem 3.11 Security transfer from a constrained squashed model}{digest-reductions/thm:transfer}
\begin{mainclaim}
\renewcommand*{\RestatedName}{\hyperref{contract.pdf}{}{claim-reductions/thm:transfer}{Theorem~3.11}}%
\begin{digestclaim}[Security transfer from a constrained squashed model]
Let $M\dsq_{\Lambda}F$ and let $\Eprot_M,\Eprot_F$ be as in
Definition~\ref{def:protocol}. For every admissible adversary system $E$ let
$\Sph_E\subseteq\Seq{A^{n}\Bsys^{n}E}$ and
$\Ssq_E\subseteq\Seq{A^{n}\Qsys^{n}E}$ satisfy
\begin{equation*}\tag{\getrefnumber{C-eq:inclusion}}\digesteqlabel{eq:inclusion}
  \big(\idmap_{A^{n}}\otimes\Lambda^{\otimes n}\otimes\idmap_E\big)
  \big(\Sph_E\big)\;\subseteq\;\Ssq_E .
\end{equation*}
If $\Crit(\Eprot_F,\sigma)\le\varepsilon$ for all $\sigma\in\Ssq_E$ and all
$E$, then $\Crit(\Eprot_M,\rho)\le\varepsilon$ for all $\rho\in\Sph_E$ and all
$E$ --- with the \emph{same} $\varepsilon$, and simultaneously for every
security criterion in the sense of Definition~\ref{def:criterion}.
\end{digestclaim}
\end{mainclaim}

% claim reductions/cor:sound
\digestlabel{cor:sound}{3.12}\pdfbookmark[2]{Corollary 3.12 Imposing and intersecting sound constraints}{digest-reductions/cor:sound}
\renewcommand*{\RestatedName}{\hyperref{contract.pdf}{}{claim-reductions/cor:sound}{Corollary~3.12}}%
\begin{digestclaim}[Imposing and intersecting sound constraints]
Let $M\dsq_{\Lambda}F$ and fix a family $\Sph$ of physical inputs.
\begin{enumerate}[label=(\roman*),leftmargin=3.0em]
\item A family $\Ssq$ satisfies \eqref{eq:inclusion} if and only if it is
  sound.
\item An arbitrary intersection of sound families is sound.
\item Consequently, if $(C^{(i)})_{i\in I}$ is any collection of sound families
  and $\Crit(\Eprot_F,\sigma)\le\varepsilon$ for every admissible $E$ and every
  $\sigma\in\bigcap_{i\in I}C^{(i)}_E$, then $\Crit(\Eprot_M,\rho)\le
  \varepsilon$ for every admissible $E$ and every $\rho\in\Sph_E$, with the
  same $\varepsilon$ and for every criterion of the form \eqref{eq:criterion}.
\item The tightest sound family is $\Srch$ itself and the loosest is
  $\Ssq_E=\Seq{A^{n}\Qsys^{n}E}$; enlarging a sound family costs tightness and
  never soundness.
\end{enumerate}
\end{digestclaim}

% claim reductions/cor:transported
\digestlabel{cor:transported}{3.13}\pdfbookmark[2]{Corollary 3.13 Sound constraints from transported operators}{digest-reductions/cor:transported}
\renewcommand*{\RestatedName}{\hyperref{contract.pdf}{}{claim-reductions/cor:transported}{Corollary~3.13}}%
\begin{digestclaim}[Sound constraints from transported operators]
Let $(M;O)\dsq_{\Lambda}(F;\hat O)$, fix a family $\Sph$ of physical inputs and
write $R\coloneqq A^{n}E$. Let $\hat O^{(n)}\in\Lin{\Qsys^{n}}$ be Hermitian
and let $O^{(n)}\coloneqq(\Lambda^{\dagger})^{\otimes n}(\hat O^{(n)})$ be its
pull-back. Then
\begin{equation*}\tag{\getrefnumber{C-eq:transportn}}\digesteqlabel{eq:transportn}
  \tr_{\Qsys^{n}}\big[(\hat O^{(n)}\otimes\idon{R})\,\sigma\big]
  =\tr_{\Bsys^{n}}\big[(O^{(n)}\otimes\idon{R})\,\rho\big]
\end{equation*}
for every $\rho\in\Sph_E$ and its image $\sigma$, and each of the following
families of squashed states is sound:
\begin{enumerate}[label=(\roman*),leftmargin=3.0em]
\item \emph{(certified)} for each certificate case of
  Corollary~\ref{cor:pullback} --- identity, positivity, spectral bounds ---
  read with $\Bsys^{n},\Qsys^{n},\Lambda^{\otimes n},R$ in place of
  $\Bsys,\Qsys,\Lambda,R$, the family of squashed states satisfying its
  conclusion;
\item \emph{(assumed)}
  $\{\sigma:\tr[(\hat O^{(n)}\otimes X)\sigma]\bowtie c\}$ for a fixed
  $\bowtie\,\in\{\le,=,\ge\}$, a fixed $X\in\Lin{R}$, Hermitian unless
  $\bowtie$ is $=$, and a fixed $c$, provided
  $\tr[(O^{(n)}\otimes X)\rho]\bowtie c$ holds for every $\rho\in\Sph_E$.
\end{enumerate}
Operators with a prescribed pull-back are supplied round by round: if
$\hat O^{(n)}$ is $\hat O_k$ on $\Qsys_j$ and $\idon{\Qsys}$ on the remaining
rounds, then $O^{(n)}$ is $O_k$ on $\Bsys_j$ and $\idon{\Bsys}$ elsewhere; and
a product $\hat O_{k_1}\otimes\dots\otimes\hat O_{k_n}$ pulls back to
$O_{k_1}\otimes\dots\otimes O_{k_n}$.
\end{digestclaim}

\section*{4\quad Operations on squashing relations}\pdfbookmark[1]{4 Operations on squashing relations}{digest-group-reductions/prop:preorder}

% claim reductions/prop:preorder
\digestlabel{prop:preorder}{4.1}\pdfbookmark[2]{Proposition 4.1 The squashing relation as a pre-order}{digest-reductions/prop:preorder}
\renewcommand*{\RestatedName}{\hyperref{contract.pdf}{}{claim-reductions/prop:preorder}{Proposition~4.1}}%
\begin{digestclaim}[The squashing relation as a pre-order]
$M\dsq_{\idmap}M$, and $M\dsq_{\Lambda_1}F_1$ together with
$F_1\dsq_{\Lambda_2}F_2$ imply $M\dsq_{\Lambda_2\circ\Lambda_1}F_2$. The same
holds for the relation with carried operators, the operator tuples
composing the same way: $(M;O)\dsq_{\Lambda_1}(F_1;\hat O^{1})$ and
$(F_1;\hat O^{1})\dsq_{\Lambda_2}(F_2;\hat O^{2})$ give
$(M;O)\dsq_{\Lambda_2\circ\Lambda_1}(F_2;\hat O^{2})$.
\end{digestclaim}

% claim reductions/cor:monotone
\digestlabel{cor:monotone}{4.2}\pdfbookmark[2]{Corollary 4.2 Monotonicity of the worst-case criterion along the pre-order}{digest-reductions/cor:monotone}
\renewcommand*{\RestatedName}{\hyperref{contract.pdf}{}{claim-reductions/cor:monotone}{Corollary~4.2}}%
\begin{digestclaim}[Monotonicity of the worst-case criterion along the pre-order]
For a POVM $G$ and a family $S$ of inputs write
$\varepsilon^{*}(G,S)\coloneqq\sup_E\sup_{\sigma\in S_E}
\Crit(\Eprot_G,\sigma)$. Let $F_2\dsq_{\Lambda_{12}}F_1$ and let both
squashed protocols use the same processing $\bar{\Grnd}_j$ and $\Kfin$. Then for
every family $S$,
\begin{equation*}
  \varepsilon^{*}(F_2,S)
  =\varepsilon^{*}\big(F_1,(\idmap\otimes\Lambda_{12}^{\otimes n}\otimes\idmap_E)(S)\big)
  \;\le\;\varepsilon^{*}(F_1,S_1)
  \quad\text{for every }S_1\supseteq
  (\idmap\otimes\Lambda_{12}^{\otimes n}\otimes\idmap_E)(S).
\end{equation*}
\end{digestclaim}

% claim reductions/lem:linear
\digestlabel{lem:linear}{4.3}\pdfbookmark[2]{Lemma 4.3 Linearity: post-processing, private settings and mixtures}{digest-reductions/lem:linear}
\renewcommand*{\RestatedName}{\hyperref{contract.pdf}{}{claim-reductions/lem:linear}{Lemma~4.3}}%
\begin{digestclaim}[Linearity: post-processing, private settings and mixtures]
Let $\Lambda\colon\Tcl{\Bsys}\to\Tcl{\Qsys}$ be a channel.
\begin{enumerate}[label=(\roman*),leftmargin=3.0em]
\item \emph{(post-processing)} If $M\dsq_{\Lambda}F$ and $P(z|v)$ is a
  stochastic matrix from $\Vset$ to a finite alphabet $\mathcal{U}$, then the
  post-processed families $(PM)_z\coloneqq\sum_vP(z|v)M_v$ and
  $(PF)_z\coloneqq\sum_vP(z|v)F_v$ are POVMs and
  $PM\dsq_{\Lambda}PF$.
\item \emph{(a private classical variable)} For each $t\in\Tset$ let
  $\{M_{y|t}\}_{y\in\Yset}$ be a \emph{conditional} POVM on $\Bsys$ ---
  complete for each fixed $t$, $\sum_yM_{y|t}=\idon{\Bsys}$ --- and let
  $\{F_{y|t}\}_{y\in\Yset}$ be conditional POVMs
  on the same $\Qsys$ with $M_{|t}\dsq_{\Lambda}F_{|t}$ for \emph{one common}
  $\Lambda$. Let $(q_t)_{t\in\Tset}$ be a probability distribution. Then, over
  the \emph{joint} alphabet $\Vset=\Yset\times\Tset$,
  \begin{equation*}\tag{\getrefnumber{C-eq:private}}\digesteqlabel{eq:private}
    \big\{q_tM_{y|t}\big\}_{(y,t)}\;\dsq_{\Lambda}\;
    \big\{q_tF_{y|t}\big\}_{(y,t)} ,
  \end{equation*}
  which in the comma notation reads
  $\{M_{y,t}\}\dsq_{\Lambda}\{F_{y,t}\}$ with $M_{y,t}=q_tM_{y|t}$ and
  $F_{y,t}=q_tF_{y|t}$.
\item \emph{(mixtures)} In the setting of (ii),
  $\big\{\sum_tq_tM_{y|t}\big\}_y\dsq_{\Lambda}
   \big\{\sum_tq_tF_{y|t}\big\}_y$.
\end{enumerate}
\end{digestclaim}

% claim reductions/def:joint-family
\digestlabel{def:joint-family}{4.4}\pdfbookmark[2]{Definition 4.4 Joint alphabet, classes and the joint family of a private setting}{digest-reductions/def:joint-family}
\renewcommand*{\RestatedName}{\hyperref{contract.pdf}{}{claim-reductions/def:joint-family}{Definition~4.4}}%
\begin{digestclaim}[Joint alphabet, classes and the joint family of a private setting]
Let $\Yset$ be a finite non-empty ordered
\emph{outcome alphabet},
let $\Tset$ be a finite non-empty ordered set of
\emph{settings},
and let $(q_t)_{t\in\Tset}$ be
\emph{class weights},
each a declared dyadic rational on the declared lattice
(Definition~\ref{def:exactadm}\,(b)), with $q_t>0$ and
$\sum_{t\in\Tset}q_t=1$ \emph{exactly}.
\begin{enumerate}[label=(\alph*),leftmargin=2.4em]
\item \emph{(the joint alphabet, its classes and weights)} The
  \emph{joint alphabet}
  is $\Vset\coloneqq\Yset\times\Tset$, ordered with $y$ varying
  fastest, which is its declared order; its
  \emph{classes}
  are $\Vset_t\coloneqq\Yset\times\{t\}$, $t\in\Tset$, a partition
  $\Vset=\bigsqcup_{t\in\Tset}\Vset_t$, each class ordered as $\Yset$ is; and
  the triple $\big(\Vset,\{\Vset_t\}_{t\in\Tset},\{q_t\}_{t\in\Tset}\big)$ is
  the \emph{joint alphabet with classes and weights}. When no private setting
  is declared, $\Tset$ is a singleton $\{\ast\}$ with $q_{\ast}=1$, and
  $\Vset=\Yset\times\{\ast\}$ is identified with $\Yset$: one class, of weight
  one.
\item \emph{(the joint family)} Let $\Bsys$ be a system and, for each
  $t\in\Tset$, let $\{\widetilde M_{y|t}\}_{y\in\Yset}$ be a
  \emph{conditional} POVM on $\Bsys$ --- complete at each fixed $t$,
  $\sum_{y}\widetilde M_{y|t}=\idon{\Bsys}$. Their
  \emph{joint family}
  is
  \begin{equation*}\tag{\getrefnumber{C-eq:jointfamily}}\digesteqlabel{eq:jointfamily}
    M_{(y,t)}\;\coloneqq\;q_t\,\widetilde M_{y|t},
    \qquad (y,t)\in\Vset ,
  \end{equation*}
  formed from the conditional elements and the weights by one product each,
  exactly. It is a POVM on $\Bsys$ over $\Vset$ by
  Lemma~\ref{lem:linear}\,(ii) --- $\sum_{(y,t)}q_t\widetilde M_{y|t}
  =\sum_tq_t\idon{\Bsys}=\idon{\Bsys}$ --- and it satisfies the
  \emph{class-marginal identity}
  \begin{equation*}\tag{\getrefnumber{C-eq:classmarginal}}\digesteqlabel{eq:classmarginal}
    \sum_{v\in\Vset_t}M_v\;=\;q_t\,\idon{\Bsys}
    \qquad\text{for every }t\in\Tset ,
  \end{equation*}
  an operator identity on $\Bsys$, by completeness of the conditional POVM at
  $t$.
\end{enumerate}
Two consequences are immediate and are used throughout. First, given
conditional relations $\widetilde M_{|t}\dsq_{\Lambda}\widetilde F_{|t}$,
$t\in\Tset$, with one common channel $\Lambda$, Lemma~\ref{lem:linear}\,(ii)
gives the relation $\{q_t\widetilde M_{y|t}\}\dsq_{\Lambda}\{q_t\widetilde
F_{y|t}\}$ for the joint families: the setting never enters the target space,
no block is opened for it, and the weights are carried by the target elements
themselves. A relation for a joint family need not have a target of this
form. Second, where every $\widetilde M_{y|t}$ is block
diagonal for an orthogonal decomposition of $\Hs{\Bsys}$, so is every
$M_{(y,t)}$, and \eqref{eq:classmarginal} holds block by block; where such a
joint family is itself a target, the decomposition being finite and the
carried operators block diagonal for it as well,
Corollary~\ref{cor:blockdiag} applies to it.
\end{digestclaim}

% claim reductions/prop:pinching
\digestlabel{prop:pinching}{4.6}\pdfbookmark[2]{Proposition 4.6 Sector pinching as a squashing}{digest-reductions/prop:pinching}
\renewcommand*{\RestatedName}{\hyperref{contract.pdf}{}{claim-reductions/prop:pinching}{Proposition~4.6}}%
\begin{digestclaim}[Sector pinching as a squashing]
Let $\{\Pi_s\}_{s\in\SecSet}$ be finitely many mutually orthogonal projections
on $\Hs{\Qsys}$ with $\sum_s\Pi_s=\idon{\Qsys}$, and let
$\PinS(\sigma)\coloneqq\sum_s\Pi_s\sigma\Pi_s$. Then $\PinS$ is a channel with
$\PinS^{\dagger}=\PinS$, and if every $F_v$ is block diagonal,
$F_v=\sum_s\Pi_sF_v\Pi_s$, then $F\dsq_{\PinS}F$.
\end{digestclaim}

% claim reductions/cor:blockdiag
\digestlabel{cor:blockdiag}{4.7}\pdfbookmark[2]{Corollary 4.7 Restriction to block-diagonal squashed states}{digest-reductions/cor:blockdiag}
\renewcommand*{\RestatedName}{\hyperref{contract.pdf}{}{claim-reductions/cor:blockdiag}{Corollary~4.7}}%
\begin{digestclaim}[Restriction to block-diagonal squashed states]
Let $(M;O)\dsq_{\Lambda}(F;\hat O)$ and suppose every $F_v$ and every
$\hat O_k$ is block diagonal for a finite orthogonal decomposition
$\{\Pi_s\}_{s\in\SecSet}$ of $\Hs{\Qsys}$. Then
\begin{enumerate}[label=(\roman*),leftmargin=3.0em]
\item $(M;O)\dsq_{\PinS\circ\Lambda}(F;\hat O)$, with the same target, the same
  tuple and the same pull-backs;
\item every state reachable through $\PinS\circ\Lambda$ is block diagonal in
  each round, so ``$\sigma$ is block diagonal in every round's sector
  decomposition'' is a sound constraint for that channel
  (Definition~\ref{def:sound});
\item consequently the squashed analysis may be restricted to block-diagonal
  states without excluding a reachable one.
\end{enumerate}
\end{digestclaim}

% claim reductions/prop:tensor
\digestlabel{prop:tensor}{4.8}\pdfbookmark[2]{Proposition 4.8 Tensor products of squashing relations}{digest-reductions/prop:tensor}
\renewcommand*{\RestatedName}{\hyperref{contract.pdf}{}{claim-reductions/prop:tensor}{Proposition~4.8}}%
\begin{digestclaim}[Tensor products of squashing relations]
If $M\dsq_{\Lambda}F$ and $M'\dsq_{\Lambda'}F'$ then
$\{M_v\otimes M'_{v'}\}_{(v,v')}\dsq_{\Lambda\otimes\Lambda'}
\{F_v\otimes F'_{v'}\}_{(v,v')}$.
\end{digestclaim}

% claim reductions/lem:obstructions
\digestlabel{lem:obstructions}{4.9}\pdfbookmark[2]{Lemma 4.9 Necessary conditions for a squashing relation}{digest-reductions/lem:obstructions}
\renewcommand*{\RestatedName}{\hyperref{contract.pdf}{}{claim-reductions/lem:obstructions}{Lemma~4.9}}%
\begin{digestclaim}[Necessary conditions for a squashing relation]
Let $M\dsq F$, let $X=\sum_{v}a_{v}F_{v}$ be a real linear combination of
target effects and let $X'\coloneqq\sum_{v}a_{v}M_{v}$. Then:
\begin{enumerate}[label=(\roman*),leftmargin=3.0em]
\item \emph{(spectral inclusion)}
  $\operatorname{spec}(X')\subseteq[\inf\operatorname{spec}X,\
  \sup\operatorname{spec}X]$; in particular
  $\operatorname{spec}(M_v)\subseteq[\inf\operatorname{spec}F_v,\
  \sup\operatorname{spec}F_v]$ and $\|M_v\|\le\|F_v\|$ for every $v$;
\item \emph{(span monotonicity)}
  $\dim\operatorname{span}\{M_v\}\le\dim\operatorname{span}\{F_v\}
  \le(\dim\Hs{\Qsys})^{2}$, and every linear relation
  $\sum_va_vF_v=c\,\idon{\Qsys}$ forces $\sum_va_vM_v=c\,\idon{\Bsys}$.
\end{enumerate}
\end{digestclaim}

\section*{5\quad Branching and flag targets}\pdfbookmark[1]{5 Branching and flag targets}{digest-group-reductions/lem:top}

% claim reductions/lem:top
\digestlabel{lem:top}{5.1}\pdfbookmark[2]{Lemma 5.1 Flag PVM as a top element}{digest-reductions/lem:top}
\renewcommand*{\RestatedName}{\hyperref{contract.pdf}{}{claim-reductions/lem:top}{Lemma~5.1}}%
\begin{digestclaim}[Flag PVM as a top element]
Let $M$ be any POVM on any $\Bsys$ with finite alphabet $\Vset$, and let
$\Ppvm=\{\dyad{v}\}_{v\in\Vset}$ be the \emph{flag PVM} on
$\Hs{\Flag}=\mathbb{C}^{|\Vset|}$ with orthonormal basis
$\{\ket{v}\}_{v\in\Vset}$. Then
\begin{equation*}\tag{\getrefnumber{C-eq:mpsquasher}}\digesteqlabel{eq:mpsquasher}
  M\;\dsq_{\Lmp}\;\Ppvm,
  \qquad
  \Lmp(\rho)\coloneqq\sum_{v\in\Vset}\tr[M_v\rho]\,\dyad{v} .
\end{equation*}
More generally, if $P$ is any PVM on a space $\Qsys$ with the same alphabet and
$P_v=0\Rightarrow M_v=0$, then $M\dsq P$.
\end{digestclaim}

% claim reductions/def:generalized-instrument
\digestlabel{def:generalized-instrument}{5.2}\pdfbookmark[2]{Definition 5.2 Generalized instrument}{digest-reductions/def:generalized-instrument}
\renewcommand*{\RestatedName}{\hyperref{contract.pdf}{}{claim-reductions/def:generalized-instrument}{Definition~5.2}}%
\begin{digestclaim}[Generalized instrument]
A \emph{generalized instrument} on $\Bsys$ is a finite family
$\{\Ins_j\}_{j\in J}$ of CP maps $\Ins_j\colon\Tcl{\Bsys}\to\Tcl{\Bsys'_j}$
whose traces add up to the trace, $\sum_j\tr\circ\Ins_j=\tr$: branch $j$
occurs with probability $\tr\Ins_j(\rho)$ and leaves the subnormalized state
$\Ins_j(\rho)$ on the system $\Bsys'_j$, which may differ from branch to
branch.
\end{digestclaim}

% claim reductions/lem:master
\digestlabel{lem:master}{5.3}\pdfbookmark[2]{Lemma 5.3 Instrument-conditioned squashing}{digest-reductions/lem:master}
\renewcommand*{\RestatedName}{\hyperref{contract.pdf}{}{claim-reductions/lem:master}{Lemma~5.3}}%
\begin{digestclaim}[Instrument-conditioned squashing]
Let $J$ be a finite set and let $\Ins_j\colon\Tcl{\Bsys}\to\Tcl{\Bsys'_j}$ be
CP maps with $\sum_{j}\tr\circ\Ins_j=\tr$, that is, a generalized instrument
(Definition~\ref{def:generalized-instrument}). For each $j$ let $M^{(j)}$ be a
POVM on $\Bsys'_j$ over a common finite alphabet $\Vset$ with
$(M^{(j)};O^{(j)})\dsq_{\Lambda_j}(F^{(j)};\hat O^{(j)})$, where $F^{(j)}$ is a
POVM on a finite-dimensional $\Qsys_j$. Each $\Ins_j$ is trace
non-increasing; write $\Ins_j^{\dagger}$ for its adjoint in the sense of
Lemma~\ref{lem:sep}, the completely positive map fixed by
$\tr[X\Ins_j(\rho)]=\tr[\Ins_j^{\dagger}(X)\rho]$ for $X\in\Lin{\Bsys'_j}$ and
$\rho\in\Tcl{\Bsys}$. Then:
\begin{enumerate}[label=(\roman*),leftmargin=3.0em]
\item $M_{v,j}\coloneqq\Ins_j^{\dagger}\big(M^{(j)}_v\big)$ defines a POVM on
  $\Bsys$ with joint outcome $(v,j)\in\Vset\times J$;
\item with $\Qsys\coloneqq\bigoplus_{j\in J}\Qsys_j$ and $F_{v,j}$ the operator
  equal to $F^{(j)}_v$ in the $j$-th block and zero elsewhere,
  \begin{equation*}\tag{\getrefnumber{C-eq:master}}\digesteqlabel{eq:master}
    \{M_{v,j}\}\;\dsq_{\Lambda}\;\{F_{v,j}\},
    \qquad
    \Lambda(\rho)\coloneqq\bigoplus_{j\in J}\Lambda_j\big(\Ins_j(\rho)\big);
  \end{equation*}
\item \emph{(coarse graining)} with the same channel,
  $\big\{\sum_jM_{v,j}\big\}_v\dsq_{\Lambda}
   \big\{\bigoplus_jF^{(j)}_v\big\}_v$;
\item \emph{(branch weights)} $\Lambda^{\dagger}(\Pi_j)
  =\Ins_j^{\dagger}(\idon{\Bsys'_j})$, where $\Pi_j$ projects onto the $j$-th
  block. In particular, if $\Ins_j^{\dagger}(\idon{\Bsys'_j})=q_j\,\idon{\Bsys}$
  for constants $q_j$, then for every finite-dimensional side system $R$ and
  every reachable $\sigma$,
  $\tr_{\Qsys}[(\Pi_j\otimes\idon{R})\sigma]=q_j\,\sigma_R$, and in particular
  $\tr[\Pi_j\sigma]=q_j$;
\item \emph{(transported operators)} the composite carries the operator equal
  to $\hat O^{(j)}_k$ in block $j$ and zero elsewhere, with pull-back
  $\Ins_j^{\dagger}\big(O^{(j)}_k\big)$.
\end{enumerate}
\end{digestclaim}

% claim reductions/cor:instances
\digestlabel{cor:instances}{5.8}\pdfbookmark[2]{Corollary 5.8 Three special cases of instrument-conditioned squashing}{digest-reductions/cor:instances}
\renewcommand*{\RestatedName}{\hyperref{contract.pdf}{}{claim-reductions/cor:instances}{Corollary~5.8}}%
\begin{digestclaim}[Three special cases of instrument-conditioned squashing]
Let $\Lambda_j$ and $F^{(j)}$ be branch channels and branch targets as in
Lemma~\ref{lem:master}.
\begin{enumerate}[label=(\roman*),leftmargin=3.0em]
\item \emph{(convex mixture with a known weight)} If
  $M_v=(1-q)G_v+qR_v$ for POVMs $G,R$ on $\Bsys$ and $q\in[0,1]$, with
  $G\dsq_{\Lambda_G}F^{(G)}$ and $R$ \emph{arbitrary}, then
  $M\dsq\big\{F^{(G)}_v\dsum\dyad{v}_{\Flag}\big\}_v$, with the sound branch
  constraints $\tr[\Pi_1\sigma]=1-q$ and $\tr[\Pflag\sigma]=q$, in the operator
  form of Lemma~\ref{lem:master}\,(iv).
\item \emph{(block-wise squashing)} Let $\Hs{\Bsys}=\bigoplus_j\Hs{\Bsys_j}$ be
  a finite orthogonal decomposition and let every $M_v$ be block diagonal,
  $M_v=\bigoplus_jM^{(j)}_v$, with $M^{(j)}\dsq_{\Lambda_j}F^{(j)}$. Then
  $M\dsq\bigoplus_jF^{(j)}$.
\item \emph{(POVM-conditioned squashing)} If $\Bsys=\Bsys_0D$ with a POVM
  $\{N_j\}_{j\in J}$ on $D$ and $M_v=\sum_jM_{v|j}\otimes N_j$ with
  $M_{|j}\dsq_{\Lambda_j}F_{|j}$ on $\Qsys_j$, then
  $\{M_v\}_v\dsq_{\Lambda}\big\{\bigoplus_jF_{v|j}\big\}_v$ with
  \begin{equation*}\tag{\getrefnumber{C-eq:povmcond}}\digesteqlabel{eq:povmcond}
    \Lambda(\rho)=\bigoplus_{j\in J}\Lambda_j\Big(
      \tr_{D}\big[(\idon{\Bsys_0}\otimes N_j^{1/2})\,\rho\,
      (\idon{\Bsys_0}\otimes N_j^{1/2})\big]\Big).
  \end{equation*}
\end{enumerate}
\end{digestclaim}

% claim reductions/prop:common
\digestlabel{prop:common}{5.9}\pdfbookmark[2]{Proposition 5.9 Common and dominating targets}{digest-reductions/prop:common}
\renewcommand*{\RestatedName}{\hyperref{contract.pdf}{}{claim-reductions/prop:common}{Proposition~5.9}}%
\begin{digestclaim}[Common and dominating targets]
In the setting of Lemma~\ref{lem:master}, suppose additionally
$F^{(j)}\dsq_{\Psi_j}F$ for a single POVM $F$ on a space $\Qsys_{*}$ (a
\emph{dominating target}; $\Psi_j=\idmap$ if all branch targets already
coincide). Then
\begin{equation*}
  \Big\{\sum_{j}M_{v,j}\Big\}_v\;\dsq_{\Lambda'}\;F,
  \qquad
  \Lambda'\coloneqq\sum_{j\in J}\Psi_j\circ\Lambda_j\circ\Ins_j .
\end{equation*}
\end{digestclaim}

% claim reductions/cor:flagsenough
\digestlabel{cor:flagsenough}{5.11}\pdfbookmark[2]{Corollary 5.11 Arbitrary unpopulated block in a flag-extended target}{digest-reductions/cor:flagsenough}
\renewcommand*{\RestatedName}{\hyperref{contract.pdf}{}{claim-reductions/cor:flagsenough}{Corollary~5.11}}%
\begin{digestclaim}[Arbitrary unpopulated block in a flag-extended target]
Let $\{G_v\}_{v\in\Vset}$ be any POVM on a finite-dimensional space
$\Bsys_0$ and let $\{H_v\}_{v\in\Vset}$ be any POVM on $\Bsys_0$. Then, on
$\Hs{\Bsys_0}\dsum\Hs{\Flag}$,
\begin{equation*}\tag{\getrefnumber{C-eq:flagsenough}}\digesteqlabel{eq:flagsenough}
  \big\{G_v\dsum\dyad{v}_{\Flag}\big\}_v
  \;\dsq_{\Lambda}\;
  \big\{H_v\dsum\dyad{v}_{\Flag}\big\}_v,
  \qquad
  \Lambda(\rho)\coloneqq
   0_{\Bsys_0}\dsum\sum_v\tr\big[(G_v\dsum\dyad{v})\rho\big]\,\dyad{v}_{\Flag}.
\end{equation*}
In particular $H_v=\tfrac{1}{|\Vset|}\idon{\Bsys_0}$ is admissible, and the
squashing channel depends on $G$ while the target does not.
\end{digestclaim}

% claim reductions/thm:mixture
\digestlabel{thm:mixture}{5.13}\pdfbookmark[2]{Theorem 5.13 Register-conditioned squashing}{digest-reductions/thm:mixture}
\renewcommand*{\RestatedName}{\hyperref{contract.pdf}{}{claim-reductions/thm:mixture}{Theorem~5.13}}%
\begin{digestclaim}[Register-conditioned squashing]
Let $\Bsys=\Bsys_0D$, let $\{N_j\}_{j\in J}$ be a POVM on $D$ with $J$ finite,
and let
\begin{equation*}\tag{\getrefnumber{C-eq:regmodel}}\digesteqlabel{eq:regmodel}
  M_v=\sum_{j\in J}M_{v|j}\otimes N_j ,
  \qquad v\in\Vset ,
\end{equation*}
where $\{M_{v|j}\}_{v}$ is a POVM on $\Bsys_0$ for each $j$. Suppose
$(M_{|j};O^{(j)})\dsq_{\Lambda_j}(F_{|j};\hat O^{(j)})$ for each $j$, with
$F_{|j}$ a POVM on a finite-dimensional $\Qsys_j$. Put
$\Hs{\Qsys}=\bigoplus_{j\in J}\Hs{\Qsys_j}$ with block projections $\Pi_j$ and
\begin{equation*}\tag{\getrefnumber{C-eq:regchannel}}\digesteqlabel{eq:regchannel}
  \Lambda(\rho)\coloneqq\bigoplus_{j\in J}\Lambda_j\Big(
    \tr_{D}\big[(\idon{\Bsys_0}\otimes N_j^{1/2})\,\rho\,
    (\idon{\Bsys_0}\otimes N_j^{1/2})\big]\Big).
\end{equation*}
Then $\Lambda$ is a channel and:
\begin{enumerate}[label=(\roman*),leftmargin=3.0em]
\item $\{M_v\}_v\dsq_{\Lambda}\big\{\bigoplus_jF_{v|j}\big\}_v$; in particular
  $\{M_v\}$ is a POVM on $\Bsys$;
\item the composite transports the operator $\hat O^{(j)}_k$ placed in block
  $j$ to $O^{(j)}_k\otimes N_j$;
\item $\Lambda^{\dagger}(\Pi_j)=\idon{\Bsys_0}\otimes N_j$, so on every
  reachable $\sigma$ and every finite-dimensional side system $R$,
  \begin{equation*}\tag{\getrefnumber{C-eq:branchweight}}\digesteqlabel{eq:branchweight}
    \tr_{\Qsys}\big[(\Pi_j\otimes\idon{R})\,\sigma\big]
    =\tr_{\Bsys}\big[\big((\idon{\Bsys_0}\otimes N_j)\otimes\idon{R}\big)\,\rho\big];
  \end{equation*}
  this is an identity, not a constraint, unless the right-hand side is a
  declared multiple of $\rho_R$;
\item if every $F_{|j}$ and every $\hat O^{(j)}_k$ is block diagonal for a
  finite orthogonal decomposition of $\Hs{\Qsys_j}$, then every
  $\bigoplus_jF_{v|j}$ and every carried operator is block diagonal for the
  union decomposition of $\Hs{\Qsys}$, and $\PinS\circ\Lambda$ realizes the
  same relation with block-diagonal reachable states
  (Corollary~\ref{cor:blockdiag}).
\end{enumerate}
\end{digestclaim}

\section*{6\quad Event-level device facts, attack sets and the licence}\pdfbookmark[1]{6 Event-level device facts, attack sets and the licence}{digest-group-reductions/prop:obstruction}

% claim reductions/prop:obstruction
\digestlabel{prop:obstruction}{6.1}\pdfbookmark[2]{Proposition 6.1 Obstruction to encoding event bounds in classical targets}{digest-reductions/prop:obstruction}
\renewcommand*{\RestatedName}{\hyperref{contract.pdf}{}{claim-reductions/prop:obstruction}{Proposition~6.1}}%
\begin{digestclaim}[Obstruction to encoding event bounds in classical targets]
Call a target POVM $F$ on $\Qsys$ \emph{classical} if every $F_v$ is diagonal
in one fixed orthonormal basis $\{\ket{f}\}_f$ of $\Hs{\Qsys}$. There exist a
POVM $M$ on $\Hs{\Bsys}=\mathbb{C}^{2}$, an event
$\mathcal{O}_{*}\subseteq\Vset$ and a constant $p_{*}>0$ with
$\sum_{v\in\mathcal{O}_{*}}M_v\succeq p_{*}\idon{\Bsys}$ such that \emph{no}
pair $(F,\Lambda)$ exists with $F$ classical, $M\dsq_{\Lambda}F$ and
$\sum_{v\in\mathcal{O}_{*}}F_v\succeq p_{*}\idon{\Qsys}$. Concretely,
\begin{equation*}\tag{\getrefnumber{C-eq:obstruction}}\digesteqlabel{eq:obstruction}
  M_1=\tfrac12\dyad{0},\quad
  M_2=\tfrac12\dyad{+},\quad
  M_3=\idon{\Bsys}-M_1-M_2,\quad
  \mathcal{O}_{*}=\{1,2\},\quad
  p_{*}=\tfrac12\Big(1-\tfrac{1}{\sqrt2}\Big),
\end{equation*}
with $\ket{+}=(\ket{0}+\ket{1})/\sqrt2$.
\end{digestclaim}

% claim reductions/def:attset
\digestlabel{def:attset}{6.3}\pdfbookmark[2]{Definition 6.3 Admissible attack sets in the consumer's format}{digest-reductions/def:attset}
\renewcommand*{\RestatedName}{\hyperref{contract.pdf}{}{claim-reductions/def:attset}{Definition~6.3}}%
\begin{digestclaim}[Admissible attack sets in the consumer's format]
Let $\CList=\{(\hat O_k,\bowtie_k,\Theta_k)\}_{k=1}^{m}$ be a declared list of
Hermitian $\hat O_k\in\Lin{\Qsys}$, relations
$\bowtie_k\in\{\preceq,=,\succeq\}$ and right-hand sides typed by the declared
source mode. For a list whose entries carry real constants $c_k$, as every
list this module emits does (Definition~\ref{def:licensed}), the admissible
sets are the attack sets of $\CList$ of
\xref{foundations/def:round-model} at the per-round system $S=\Qsys$, which the
consuming contract states at the same system
(\xref{meat-core/def:roundmodel}); their displays at that system are,
verbatim,
\begin{equation*}\tag{\getrefnumber{C-eq:attset}}\digesteqlabel{eq:attset}
  \AttSet_{\CList}(E',\,\Qsys E'')\coloneqq
  \Big\{\Aev\in\mathrm{CPTP}(E',\Qsys E''):\
  \Aev^{\dagger}\big(\hat O_k\otimes\idon{E''}\big)
  \bowtie_k c_k\,\idon{E'}\ \ \forall k\Big\}
  \qquad\text{(trusted-source mode)},
\end{equation*}
and, with the constraint operator lifted to the larger supplied system,
\begin{equation*}\tag{\getrefnumber{C-eq:attsetun}}\digesteqlabel{eq:attsetun}
  \AttSet_{\CList}(E',\,A\Qsys E'')\coloneqq
  \Big\{\Aev\in\mathrm{CPTP}(E',A\Qsys E''):\
  \Aev^{\dagger}\big(\idon{A}\otimes\hat O_k\otimes\idon{E''}\big)
  \bowtie_k c_k\,\idon{E'}\ \ \forall k\Big\}
\end{equation*}
in untrusted-source mode, the two distinguished by the output system
named in the argument. This contract writes $\bowtie_k$ in the operator
spelling of \eqref{eq:attset} and \eqref{eq:attsetun} in both source modes,
and everything it emits is admitted under either spelling. An entry that declares a Hermitian
right-hand side $\Theta_k$ on Alice's system, rather than a real constant, is
a \emph{general} entry: it is admissible input to the consuming contract but
is never licensed (\xref{protocol-compiler/def:homog}). Every entry this module emits
carries a real scalar $c_k$ and is written $(\hat O_k,\bowtie_k,c_k)$
(Definition~\ref{def:licensed}, Theorem~\ref{thm:licence}\,(iv)).

\emph{Party-local entries.} Where the measured system is a product
$\Qsys=\bigotimes_f\Qsys_f$ and $\CList$ carries entries tagged party-local
(Definition~\ref{def:licensed}), the admissible set is cut out, for each such
entry $(\hat O_k,\bowtie_k,c_k;f)$, by the condition
\begin{equation*}\tag{\getrefnumber{C-eq:attsetpl}}\digesteqlabel{eq:attsetpl}
  \big(\idmap_{R'}\otimes\Aev\big)^{\dagger}
  \Big(X\otimes\big(\hat O_k-c_k\,\idon{\Qsys_f}\big)\otimes\idon{E''}\Big)
  \;\bowtie_k\;0
  \qquad\text{for every finite-dimensional }R'\text{ and every }
  X\in\Pos{R'R\,\Qsys_{f^c}} ,
\end{equation*}
in place of the scalar condition of \eqref{eq:attset} or
\eqref{eq:attsetun}, where $\Qsys_{f^c}$ is the product of the other
factors and $R$ is $A$ in untrusted-source mode and trivial in
trusted-source mode; the untagged entries keep their conditions. The
reference system $R'$ is what carries Alice's retained system in
trusted-source mode, where it is no output of $\Aev$. Taking $R'$ trivial
and $X=\idon{}$ in \eqref{eq:attsetpl} gives the scalar condition for the
lifted operator $\idon{\Qsys_{f^c}}\otimes\hat O_k$, so the set with
party-local entries is contained in the set with the same entries read as
scalars. The condition \eqref{eq:attsetpl} is that of an entry tagged with the
factor $f$ in \xref{foundations/def:round-model}\,(iii) at $S=\Qsys$, which the
consuming contract states at the same system. An attack sequence is
$\CList$-admissible (an admissible attack sequence for $\CList$) if
$\Aev_j$ lies in the corresponding set with
$E'=E_{j-1}$ and $E''=E'_j$ for every $j$.
\end{digestclaim}

% claim reductions/def:licensed
\digestlabel{def:licensed}{6.5}\pdfbookmark[2]{Definition 6.5 Licensed entries and the party-local form}{digest-reductions/def:licensed}
\renewcommand*{\RestatedName}{\hyperref{contract.pdf}{}{claim-reductions/def:licensed}{Definition~6.5}}%
\begin{digestclaim}[Licensed entries and the party-local form]
The list $\CList$ is
\emph{licensed}
for the declared source mode when each of its entries either is untagged
with its right-hand side in the form that mode admits, or is in the
party-local form below. The untagged forms are $\Theta_k=c_k\,\tau_A$ with
$c_k\in\mathbb{R}$ in trusted-source mode, the state-level entry reading
$\tr_{\Qsys}[(\idon{A}\otimes\hat O_k)\,\omega]\bowtie_kc_k\,\tau_A$; and
$\Theta_k=c_k\in\mathbb{R}$ in untrusted-source mode, the entry reading
$\tr[(\idon{A}\otimes\hat O_k)\,\omega]\bowtie_kc_k$, a comparison of real
numbers (\xref{protocol-compiler/def:homog}). In both cases the untagged
entries are written $(\hat O_k,\bowtie_k,c_k)$.

The
\emph{party-local form}.
Where the measured system is declared as a product
$\Qsys=\bigotimes_{f}\Qsys_f$ of party-scoped factors, an entry may in
addition be
\emph{scoped to a factor}
$f$. It is written $(\hat O_k,\bowtie_k,c_k;f)$ with
$\hat O_k\in\Lin{\Qsys_f}$ Hermitian and $c_k\in\mathbb{R}$, it carries the
party-local tag, and its state-level reading, in either source mode, is
\begin{equation*}\tag{\getrefnumber{C-eq:partylocal-entry}}\digesteqlabel{eq:partylocal-entry}
  \tr_{\Qsys_f}\!\Big[\big(\idon{}\otimes(\hat O_k-c_k\,\idon{\Qsys_f})\big)
  \,\omega\Big]\;\bowtie_k\;0
  \qquad\text{as an operator on the systems of }\omega\text{ other than }
  \Qsys_f ,
\end{equation*}
those systems being $A$, where the declaration carries one, and the factors
$\Qsys_{f'}$, $f'\ne f$. Equivalently
$\tr_{\Qsys_f}[(\idon{}\otimes\hat O_k)\omega]\bowtie_kc_k\tr_{\Qsys_f}\omega$:
the entry is homogeneous with respect to the marginal that is actually
there. The attack-set licence of Theorem~\ref{thm:licence} covers a
party-local entry only through a certificate on the factor's own
squashing channel (its part (v)); the product structure that makes it
sound is Lemma~\ref{lem:party-local}.

\emph{What the tag does not change.} An entry without the tag is read in
the form of its source mode, exactly as above. With one factor, the
party-local reading in trusted-source mode is
$\tr_{\Qsys}[(\idon{A}\otimes\hat O_k)\omega]\bowtie_kc_k\,\omega_A$, which
under the consumer's marginal condition $\omega_A=\tau_A$ is the
trusted-source form; with one factor and $d_A=1$ it is the untrusted-source
comparison of real numbers. \emph{With one factor, in untrusted-source mode
and $d_A>1$ it is neither}: it reads
$\tr_{\Qsys}[(\idon{A}\otimes\hat O_k)\omega]\bowtie_kc_k\,\omega_A$, an
operator inequality on $A$ against the marginal that is actually there, and
the untrusted-source scalar entry is its trace
(Lemma~\ref{lem:party-local}\,(iii)), implied by it and not equivalent to
it. With more than one factor it implies the untagged reading of the same
$(\hat O_k,\bowtie_k,c_k)$ for the lifted operator $\idon{}\otimes\hat O_k$,
which is obtained from it by the partial trace over the other factors (and
over $A$ in untrusted-source mode; in trusted-source mode under the
consumer's marginal condition $\omega_A=\tau_A$), and in general it is not
implied by it.
\end{digestclaim}

% claim reductions/lem:party-local
\digestlabel{lem:party-local}{6.6}\pdfbookmark[2]{Lemma 6.6 Party-local transport}{digest-reductions/lem:party-local}
\renewcommand*{\RestatedName}{\hyperref{contract.pdf}{}{claim-reductions/lem:party-local}{Lemma~6.6}}%
\begin{digestclaim}[Party-local transport]
Let the measured system be a tensor product $\Bsys=\Bsys_f\otimes\Bsys_{f^c}$
of a \emph{part} $\Bsys_f$ and the product $\Bsys_{f^c}$ of all other parts,
let $\Lambda_f\colon\Tcl{\Bsys_f}\to\Tcl{\Qsys_f}$ and
$\Lambda_{f^c}\colon\Tcl{\Bsys_{f^c}}\to\Tcl{\Qsys_{f^c}}$ be channels, and
put $\Lambda\coloneqq\Lambda_f\otimes\Lambda_{f^c}$ and
$\Qsys\coloneqq\Qsys_f\otimes\Qsys_{f^c}$. Let $\hat O\in\Lin{\Qsys_f}$ be
Hermitian with a spectral pull-back certificate \emph{on the part},
\begin{equation*}\tag{\getrefnumber{C-eq:partcert}}\digesteqlabel{eq:partcert}
  \Lambda_f^{\dagger}(\hat O)\;\bowtie\;c\,\idon{\Bsys_f},
  \qquad\bowtie\,\in\{\preceq,=,\succeq\},\quad c\in\mathbb{R} .
\end{equation*}
Let $R$ and $R'$ be finite-dimensional systems, $E'$, $E''$
finite-dimensional, and $\Aev\in\mathrm{CPTP}(E',R\,\Bsys\,E'')$ \emph{any}
channel, with
$\Aev^{\sq}\coloneqq(\idmap_R\otimes\Lambda\otimes\idmap_{E''})\circ\Aev$.
Write $\hat D\coloneqq\hat O-c\,\idon{\Qsys_f}$.
\begin{enumerate}[label=(\roman*),leftmargin=2.2em]
\item \emph{(channel level)} for every $X\in\Pos{R'R\,\Qsys_{f^c}}$,
\begin{equation*}\tag{\getrefnumber{C-eq:partylocal-channel}}\digesteqlabel{eq:partylocal-channel}
  \big(\idmap_{R'}\otimes\Aev^{\sq}\big)^{\dagger}
  \big(X\otimes\hat D\otimes\idon{E''}\big)\;\bowtie\;0
  \qquad\text{as an operator on }R'E' ,
\end{equation*}
the tensor factors written in the order $R'R\,\Qsys_{f^c}$, $\Qsys_f$, $E''$.
\item \emph{(state level)} for every $\varrho\in\Pos{R'E'}$, the operator
$\omega\coloneqq\tr_{E''}\big[(\idmap_{R'}\otimes\Aev^{\sq})(\varrho)\big]$
on $R'R\,\Qsys$ satisfies
\begin{equation*}\tag{\getrefnumber{C-eq:partylocal-state}}\digesteqlabel{eq:partylocal-state}
  \tr_{\Qsys_f}\!\big[(\idon{R'R\,\Qsys_{f^c}}\otimes\hat D)\,\omega\big]
  \;\bowtie\;0
  \qquad\text{as an operator on }R'R\,\Qsys_{f^c},
\end{equation*}
equivalently
$\tr_{\Qsys_f}[(\idon{}\otimes\hat O)\,\omega]\bowtie
c\,\tr_{\Qsys_f}\omega$: the right-hand side is $c$ times the marginal of
$\omega$ \emph{that is actually there}, and no marginal is assumed.
\item \emph{(the scalar entry is its trace)} taking the trace of
\eqref{eq:partylocal-state} gives
$\tr[(\idon{}\otimes\hat O)\,\omega]\bowtie c\,\tr\omega$, the scalar form.
With $R'$ trivial and $R=A$, $\Qsys_{f^c}$ trivial (one factor), \eqref{eq:partylocal-state}
reads $\tr_{\Qsys}[(\idon{A}\otimes\hat O)\omega]\bowtie c\,\omega_A$.
\item \emph{(the product structure is needed)} if $\Lambda$ is not a product
the conclusion can fail while the scalar one holds. Let $\Qsys_f$ and
$\Qsys_{f^c}$ be qubits, $R$, $R'$ trivial, $\hat O=\dyad{0}$, $c=\tfrac12$,
and let $\Lambda$ discard its input and prepare
$\tfrac12(\dyad{00}+\dyad{11})$. Then
$\Lambda^{\dagger}(\idon{}\otimes\hat O)=\tfrac12\idon{\Bsys}\succeq c\,\idon{}$,
so every squashed attack satisfies the scalar entry with $\succeq$; but
$\tr_{\Qsys_f}[(\idon{}\otimes\hat D)\omega]
=\tfrac14\dyad{0}-\tfrac14\dyad{1}$, which is not positive semidefinite.
\end{enumerate}
\end{digestclaim}

% claim reductions/thm:licence
\digestlabel{thm:licence}{6.7}\pdfbookmark[2]{Theorem 6.7 The attack-set licence}{digest-reductions/thm:licence}
\begin{mainclaim}
\renewcommand*{\RestatedName}{\hyperref{contract.pdf}{}{claim-reductions/thm:licence}{Theorem~6.7}}%
\begin{digestclaim}[The attack-set licence]
Let $(M;O)\dsq_{\Lambda}(F;\hat O)$ with $\Lambda\colon\Tcl{\Bsys}\to
\Tcl{\Qsys}$, and suppose each carried operator has a \emph{spectral pull-back
certificate}: for each $k$ there are $\bowtie_k\in\{\preceq,=,\succeq\}$ and
$c_k\in\mathbb{R}$ with
\begin{equation*}\tag{\getrefnumber{C-eq:certificate}}\digesteqlabel{eq:certificate}
  O_k\;=\;\Lambda^{\dagger}(\hat O_k)\;\bowtie_k\;c_k\,\idon{\Bsys}
  \qquad\text{as operators on }\Hs{\Bsys}.
\end{equation*}
Then, with $\CList\coloneqq\{(\hat O_k,\bowtie_k,c_k)\}_{k=1}^{m}$:
\begin{enumerate}[label=(\roman*),leftmargin=3.0em]
\item \emph{(trusted-source mode)} for all finite-dimensional $E',E''$ and
  \emph{every} $\Aev\in\mathrm{CPTP}(E',\Bsys E'')$, the composite
  $\Aev^{\sq}\coloneqq(\Lambda\otimes\idmap_{E''})\circ\Aev$ lies in
  $\AttSet_{\CList}(E',\Qsys E'')$ of \eqref{eq:attset};
\item \emph{(untrusted-source mode)} for every
  $\Aev\in\mathrm{CPTP}(E',A\Bsys E'')$, the composite
  $\Aev^{\sq}\coloneqq(\idmap_A\otimes\Lambda\otimes\idmap_{E''})\circ\Aev$
  lies in $\AttSet_{\CList}(E',A\Qsys E'')$ of \eqref{eq:attsetun};
\item \emph{(the licensed reduction, in the declared mode)} consequently, for
  every attack sequence $(\Aev_j)_{j=1}^{n}$ against the physical protocol of
  the type \eqref{eq:attackchannel} for the declared source mode, the sequence
  $(\Aev^{\sq}_j)_j$ of the \emph{same} branch of \eqref{eq:composite} is
  $\CList$-admissible --- by (i) with the trusted branch of
  \eqref{eq:attackchannel}, by (ii) with the untrusted branch --- and produces
  the \emph{same} protocol output state, by the correspondingly named branch of
  Corollary~\ref{cor:sequential}; so security of the squashed round model
  against $\CList$-admissible sequences implies security of the physical
  protocol against all sequences of that type, with the same $\varepsilon$;
\item \emph{(licensed entries)} $\CList$ is licensed for either source mode in
  the sense of Definition~\ref{def:licensed}, with the same constants $c_k$;
  on every reachable state $\sigma$ the transported entries hold in the form
  with the marginal that is present,
  $\tr_{\Qsys}[(\idon{A}\otimes\hat O_k)\sigma]\bowtie_kc_k\,\sigma_A$, and
  so in the untrusted-source form
  $\tr[(\idon{A}\otimes\hat O_k)\sigma]\bowtie_kc_k$; and the trusted-source
  form $\tr_{\Qsys}[(\idon{A}\otimes\hat O_k)\sigma]\bowtie_kc_k\,\tau_A$
  holds on every reachable state whose marginal on $A$ is $\tau_A$.
\item \emph{(party-local entries, product measured system)} \emph{in place
  of the hypotheses above}, which type $\hat O_k$ on $\Qsys$ and certify it
  against $\Lambda$, let
  $\Bsys=\bigotimes_f\Bsys_f$, let $M^{f}\dsq_{\Lambda_f}F^{f}$ for every
  factor with $\Lambda_f\colon\Tcl{\Bsys_f}\to\Tcl{\Qsys_f}$, and put
  $M\coloneqq\{\bigotimes_fM^{f}_{v_f}\}$, $F\coloneqq\{\bigotimes_fF^{f}_{v_f}\}$
  and $\Lambda\coloneqq\bigotimes_f\Lambda_f$, so that $M\dsq_{\Lambda}F$
  (Proposition~\ref{prop:tensor}). Let each entry of $\CList$ be scoped to a
  factor $f_k$ with $\hat O_k\in\Lin{\Qsys_{f_k}}$ and a certificate
  \emph{on that factor's channel},
  $\Lambda_{f_k}^{\dagger}(\hat O_k)\bowtie_kc_k\,\idon{\Bsys_{f_k}}$. Then
  for every attack $\Aev$ of either type of \eqref{eq:attackchannel}, the
  squashed attack of the same branch of \eqref{eq:composite} satisfies
  \eqref{eq:attsetpl} for every $k$, and also the scalar condition of
  \eqref{eq:attset} or \eqref{eq:attsetun} for the lifted operator
  $\idon{\Qsys_{f_k^c}}\otimes\hat O_k$ with the same $c_k$; every
  squashed sequence is admissible for $\CList$ read with any subset of its entries
  tagged party-local, it produces the same protocol output state
  as the physical sequence, and the security implication of (iii) holds
  with the same $\varepsilon$. On every reachable state the entries hold in
  the form \eqref{eq:partylocal-entry}.
\end{enumerate}
\end{digestclaim}
\end{mainclaim}

\section*{7\quad Typed links and reduction chains}\pdfbookmark[1]{7 Typed links and reduction chains}{digest-group-reductions/def:link}

% claim reductions/def:link
\digestlabel{def:link}{7.1}\pdfbookmark[2]{Definition 7.1 Typed link}{digest-reductions/def:link}
\renewcommand*{\RestatedName}{\hyperref{contract.pdf}{}{claim-reductions/def:link}{Definition~7.1}}%
\begin{digestclaim}[Typed link]
A
\emph{typed link}
carries a stage $\textsf{stage}\in\{\textsf{source},\textsf{receiver}\}$.
A link at the stage \textsf{source} is a source link in the sense of
\xref{transmitter/def:source-link}: its descriptors are Alice-side, a
transmitted system with a signal ensemble on either end and a channel
witness between them, and no condition on it mentions the receiver's
systems, outcome alphabets or measurements.
\emph{This module reads that type as the transmitter states it}: the channel
witness is supplied either
\emph{as data}, in which case the source identity is decided on declared
numbers, or \emph{as a theorem constructor}, in which case the link
carries an obligation record, namely the triple of the kind, the id of
the soundness claim and the assumptions the obligation carries, which
\xref{transmitter/def:source-segment}\,(A1) defines and which this module
neither computes nor re-decides, together with the member of its
declared descriptor family that its descriptors are; and a source link
carries in addition that family, an optional
retained-register sector map, its
transported constraint entries with its licence slot, and its status.
Which kinds are admitted at all, and which components a consumer decides
\emph{on data} (the emitted soundness registry, the triple compared
against it entry by entry with those assumptions demanded in the ledger, and
the provenance rule for a non-empty segment), is
\xref{transmitter/def:source-segment}, conditions (A1) to (A3); this module
re-decides none of them, adds none, and adds no source-stage kind. A
\emph{receiver-stage typed link}
$\Lnk$ is a tuple
\begin{equation*}\tag{\getrefnumber{C-eq:link}}\digesteqlabel{eq:link}
  \Lnk=\Big(\textsf{stage},\ \textsf{kind},\
  \big(\Bsys^{\mathrm{in}},M^{\mathrm{in}}\big),\
  \big(\Bsys^{\mathrm{out}},M^{\mathrm{out}}\big),\
  \Lambda_{\Lnk},\
  \big(\hat O^{\mathrm{out}},O^{\mathrm{in}}\big),\
  \Srch_{\Lnk},\
  \CList_{\Lnk},\
  \textsf{licence},\ \textsf{status}\Big)
\end{equation*}
in which: $\textsf{stage}=\textsf{receiver}$;
\textsf{kind} is one of the receiver-stage link kinds the module's interface
declares (the core's kinds, and those a submodule's interface adds);
$(\Bsys^{\mathrm{in}},M^{\mathrm{in}})$ is the
\emph{source family}
--- the system and POVM the link starts from --- and
$(\Bsys^{\mathrm{out}},M^{\mathrm{out}})$ its \emph{target};
$\Lambda_{\Lnk}$ is the
\emph{channel witness},
with
$\big(M^{\mathrm{in}};O^{\mathrm{in}}\big)\dsq_{\Lambda_{\Lnk}}
\big(M^{\mathrm{out}};\hat O^{\mathrm{out}}\big)$;
$\Srch_{\Lnk}$ is the
\emph{reachable-set statement}
of Definition~\ref{def:sound} for $\Lambda_{\Lnk}$ and the declared family of
inputs; $\CList_{\Lnk}$ is the list of transported constraint entries
$(\hat O_k,\bowtie_k,c_k)$ with, per entry, the certificate case of
Corollary~\ref{cor:pullback} that discharges it; \textsf{licence} is the
\emph{licence record}
of the module's outputs, present exactly when $\CList_{\Lnk}\ne\varnothing$;
and $\textsf{status}\in\{\textsf{exact},\textsf{certified\_outer},
\textsf{exploratory\_numeric}\}$, with \textsf{exact} meaning that
$\Lambda_{\Lnk}^{\dagger}(M^{\mathrm{out}}_v)=M^{\mathrm{in}}_v$ holds with no
error term.

\emph{The count.} A receiver-stage link has ten slots. Six of
them are the required fields --- source family, target, channel witness,
reachable-set statement, transported constraint entries, licence record ---
and the remaining four are \textsf{stage}, \textsf{kind}, the carried
operator pair, and \textsf{status}. Wherever this contract says ``the six
required fields'' it means the first group; the tuple \eqref{eq:link} is the
whole of what an implementer binds.
\end{digestclaim}

% claim reductions/def:chain
\digestlabel{def:chain}{7.2}\pdfbookmark[2]{Definition 7.2 Reduction chain}{digest-reductions/def:chain}
\renewcommand*{\RestatedName}{\hyperref{contract.pdf}{}{claim-reductions/def:chain}{Definition~7.2}}%
\begin{digestclaim}[Reduction chain]
A
\emph{reduction chain}
is a pair $\redChain=(\redChain_{\mathrm{src}},\redChain_{\mathrm{rec}})$ of
a source segment and a
\emph{receiver segment}.
The source segment
$\redChain_{\mathrm{src}}=(\Xi_1,\dots,\Xi_{L_{\mathrm s}})$, $L_{\mathrm s}\ge0$,
is a source segment in the sense of \xref{transmitter/def:source-segment}
(written $\Chsrc$ there): a list of source links
(\xref{transmitter/def:source-link}) satisfying that definition's
interface conditions (S1)--(S4) --- descriptor match, index
preservation, the source identity where checkable, and the slots --- every
descriptor of which is Alice-side. The receiver segment
$\redChain_{\mathrm{rec}}=(\Lnk_1,\dots,\Lnk_k)$, $k\ge1$, is a list of
receiver-stage typed links (Definition~\ref{def:link}) that is
\begin{enumerate}[label=(C\arabic*),leftmargin=3.0em]
\item
  \emph{composable}:
  $(\Bsys^{\mathrm{out}}_i,M^{\mathrm{out}}_i)
  =(\Bsys^{\mathrm{in}}_{i+1},M^{\mathrm{in}}_{i+1})$ for $1\le i<k$;
\item
  \emph{tuple-chained}:
  $\hat O^{\mathrm{out}}_i=O^{\mathrm{in}}_{i+1}$ for
  $1\le i<k$, so that each declared operator on the final target has a
  pull-back at every intermediate stage;
\item
  \emph{stage-ordered}:
  every \textsf{source}-stage link precedes every \textsf{receiver}-stage
  link, which the pair form makes structural.
\end{enumerate}
Its
\emph{composite}
is the \emph{pair}
\begin{equation*}\tag{\getrefnumber{C-eq:chaincomposite}}\digesteqlabel{eq:chaincomposite}
  \big(\Xi_{\redChain},\ \Lambda_{\redChain}\big),
  \qquad
  \Xi_{\redChain}\coloneqq\Xi_{L_{\mathrm s}}\circ\dots\circ\Xi_1 ,
  \qquad
  \Lambda_{\redChain}\coloneqq\Lambda_{\Lnk_k}\circ\dots\circ\Lambda_{\Lnk_1},
\end{equation*}
of the
\emph{source composite}
--- the composite the transmitter's (S3) writes $\Xi$, the identity when
$L_{\mathrm s}=0$ --- and the
\emph{receiver composite}.
The two act on different systems and no condition above relates them:
$\Xi_{\redChain}$ carries the virtual transmitted system to the
transmitted one, on Alice's side, and $\Lambda_{\redChain}$ the
measured system to the target, on the receiver's. What the pair
composes to \emph{as a reduction} is Theorem~\ref{thm:two-sided}, and
Theorem~\ref{thm:chain} is what the receiver composite alone delivers.

\emph{Which chains are admitted.} The rule is a condition on a
\emph{declaration} that carries both a non-empty source segment and a
constraint list: such a declaration is
\emph{admitted}
only where every entry of the list is licensed in the sense of
Definition~\ref{def:licensed} --- a real constant $c_k$ as the right-hand
side, the entry tagged or untagged --- \emph{and only in
trusted-source mode}, Theorem~\ref{thm:two-sided} being stated there
alone, so that a declaration pairing a non-empty source segment with
untrusted-source mode is refused rather than admitted in silence. An
entry declaring a Hermitian right-hand side on Alice's retained system, a
\emph{general} entry in the sense of Definition~\ref{def:attset}, is
\emph{refused} beside one. The rule
is decided by the consumer that assembles those two.

\emph{Where the source segment is empty}, $L_{\mathrm s}=0$, the source
composite is the identity, $\Avirt=\Asig$, the virtual ensemble is the
specified one, and the pair \eqref{eq:chaincomposite} carries exactly the
receiver composite: Alice's declared states are the ones the security
statement is about, and the receiver segment is the whole of the chain's
content (Theorem~\ref{thm:chain}).
\end{digestclaim}

% claim reductions/thm:chain
\digestlabel{thm:chain}{7.4}\pdfbookmark[2]{Theorem 7.4 Composition of a chain}{digest-reductions/thm:chain}
\renewcommand*{\RestatedName}{\hyperref{contract.pdf}{}{claim-reductions/thm:chain}{Theorem~7.4}}%
\begin{digestclaim}[Composition of a chain]
Let $\redChain=(\redChain_{\mathrm{src}},\redChain_{\mathrm{rec}})$ be a
reduction chain with receiver segment
$\redChain_{\mathrm{rec}}=(\Lnk_1,\dots,\Lnk_k)$ and receiver composite
$\Lambda_{\redChain}$ of \eqref{eq:chaincomposite}. Then:
\begin{enumerate}[label=(\roman*),leftmargin=3.0em]
\item $\big(M^{\mathrm{in}}_1;O^{\mathrm{in}}_1\big)
  \dsq_{\Lambda_{\redChain}}
  \big(M^{\mathrm{out}}_k;\hat O^{\mathrm{out}}_k\big)$, and
  $\Lambda_{\redChain}^{\dagger}
  =\Lambda_{\Lnk_1}^{\dagger}\circ\dots\circ\Lambda_{\Lnk_k}^{\dagger}$;
\item Theorem~\ref{thm:factorization}, Theorem~\ref{thm:transfer} and
  Corollary~\ref{cor:sequential} apply to the composite: the physical protocol
  with $M^{\mathrm{in}}_1$ and the final squashed protocol with
  $M^{\mathrm{out}}_k$ have the same output state, round by round and
  adaptively;
\item an entry $(\hat O_k,\bowtie_k,c_k)$ declared on the \emph{final} target
  is sound for the receiver composite when the certificate
  \eqref{eq:certificate} holds for its pull-back
  $\Lambda_{\redChain}^{\dagger}(\hat O_k)$, which (C2) computes link by link; and
  Theorem~\ref{thm:licence} then applies with $\Lambda_{\redChain}$;
\item the composite status is \textsf{exact} whenever every link's is, and no
  other composite value is asserted.
\end{enumerate}
Nothing above reads the source segment: (i)--(iv) are statements about the
receiver composite $\Lambda_{\redChain}$ and hold at every length
$L_{\mathrm s}\ge0$. What the \emph{pair} \eqref{eq:chaincomposite} composes to
as a reduction is Theorem~\ref{thm:two-sided}, under the source segment's own
interface conditions and the admission rule of Definition~\ref{def:chain}, and
this theorem does not restate it. Where the source segment is empty the source
composite is the identity and the receiver composite is the whole of the
chain's reduction.
\end{digestclaim}

\section*{8\quad The source-side reduction and the two-sided theorem}\pdfbookmark[1]{8 The source-side reduction and the two-sided theorem}{digest-group-reductions/lem:precomposition}

% claim reductions/lem:precomposition
\digestlabel{lem:precomposition}{8.1}\pdfbookmark[2]{Lemma 8.1 Licensed admissibility survives precomposition with a source channel}{digest-reductions/lem:precomposition}
\renewcommand*{\RestatedName}{\hyperref{contract.pdf}{}{claim-reductions/lem:precomposition}{Lemma~8.1}}%
\begin{digestclaim}[Licensed admissibility survives precomposition with a source channel]
Let $\CList=\{(\hat O_k,\bowtie_k,c_k)\}_{k=1}^{m}$ be a list of Hermitian
$\hat O_k$, relations $\bowtie_k\in\{\preceq,=,\succeq\}$ and \emph{real}
constants $c_k$. Let $\Hs{E_{\mathrm{pre}}}$ and $\Hs{E'}$ be separable, let
$\Qsys$, $E''$, $A$, $R$ and $R'$ be finite-dimensional, and let
$\Gamma\in\mathrm{CPTP}(E_{\mathrm{pre}},E')$ be \emph{any} channel. The sets
below are the ones the displays \eqref{eq:attset}, \eqref{eq:attsetun} and
\eqref{eq:attsetpl} of Definition~\ref{def:attset} cut out, read at these
registers with the adjoints of Lemma~\ref{lem:sep}; where $E_{\mathrm{pre}}$
and $E'$ are finite-dimensional they are the attack sets that definition
names.
\begin{enumerate}[label=(\roman*),leftmargin=3.0em]
\item \emph{(trusted-source mode)} if
  $\Aev\in\AttSet_{\CList}(E',\Qsys E'')$ of \eqref{eq:attset}, then
  $\Aev\circ\Gamma\in\AttSet_{\CList}(E_{\mathrm{pre}},\Qsys E'')$;
\item \emph{(untrusted-source mode)} if
  $\Aev\in\AttSet_{\CList}(E',A\Qsys E'')$ of \eqref{eq:attsetun}, then
  $\Aev\circ\Gamma\in\AttSet_{\CList}(E_{\mathrm{pre}},A\Qsys E'')$;
\item \emph{(party-local entries)} let $\Qsys=\bigotimes_f\Qsys_f$ and let the
  entry $(\hat O_k,\bowtie_k,c_k;f_k)$ carry the party-local tag
  (Definition~\ref{def:licensed}), so that its condition on $\Aev$ is
  \eqref{eq:attsetpl}. If $\Aev$ satisfies \eqref{eq:attsetpl} for every
  finite-dimensional $R'$ and every $X\in\Pos{R'R\,\Qsys_{f_k^c}}$, then so does
  $\Aev\circ\Gamma$, with $E_{\mathrm{pre}}$ in place of $E'$. Consequently the
  admissible set of Definition~\ref{def:attset}, with \emph{any} subset of its
  entries tagged and the rest read as scalars, is closed under precomposition
  with a channel into its input register;
\item \emph{(the real constant is what carries it)} replace the right-hand side
  $c_k\,\idon{E'}$ of \eqref{eq:attset} by a Hermitian $\Theta\in\Lin{E'}$ that
  is not a real multiple of $\idon{E'}$. Then the condition
  $\Aev^{\dagger}(\hat O\otimes\idon{E''})\succeq\Theta$ is not preserved by
  precomposition: take $E_{\mathrm{pre}}=E'=\Qsys=\mathbb{C}^{2}$ with declared
  basis $\{\ket{0},\ket{1}\}$, $E''$ one-dimensional, $\Aev=\idmap$ the
  identity channel $E'\to\Qsys$, $\hat O=\Theta=\dyad{0}$, and
  $\Gamma(\rho)\coloneqq\tr[\rho]\,\dyad{1}$. The condition holds for $\Aev$
  with equality, and
  $(\Aev\circ\Gamma)^{\dagger}(\hat O)=0\not\succeq\dyad{0}=\Theta$.
\end{enumerate}
\end{digestclaim}

% claim reductions/lem:source-transfer
\digestlabel{lem:source-transfer}{8.3}\pdfbookmark[2]{Lemma 8.3 Exactness of the source map in trusted-source mode}{digest-reductions/lem:source-transfer}
\renewcommand*{\RestatedName}{\hyperref{contract.pdf}{}{claim-reductions/lem:source-transfer}{Lemma~8.3}}%
\begin{digestclaim}[Exactness of the source map in trusted-source mode]
Let a source specification in \emph{trusted-source mode}
(\xref{transmitter/def:source-spec}) declare an alphabet $\Xset$ with weights
$(p_x)_{x\in\Xset}$, the \emph{specified} ensemble $\{p_x,\rho_x\}$ on the
transmitted system $\Asig$ and the virtual ensemble
$\{p_x,\rvirt_x\}$ on the virtual transmitted system $\Avirt$, with
$\Hs{\Asig}$ separable and $\Hs{\Avirt}$ finite-dimensional, as
\xref{transmitter/def:source-spec} declares the virtual transmitted system,
and let its source segment satisfy
the interface conditions (S1)--(S4) of
\xref{transmitter/def:source-segment}, so that its composite
$\Xi\colon\Tcl{\Avirt}\to\Tcl{\Asig}$ satisfies the \emph{source identity}
\begin{equation*}\tag{\getrefnumber{C-eq:sourceidentity}}\digesteqlabel{eq:sourceidentity}
  \Xi\big(\rvirt_x\big)=\rho_x\qquad\text{for every }x\in\Xset .
\end{equation*}
Fix $n$, a finite-dimensional adversary register $E_0$ and a protocol in the
sense of Definition~\ref{def:protocol}, and put
\begin{equation*}\tag{\getrefnumber{C-eq:prepstates}}\digesteqlabel{eq:prepstates}
  \rho^{\mathrm{sp}}\coloneqq\bigotimes_{j=1}^{n}
    \Big(\sum_{x\in\Xset}p_x\,\dyad{x}_{X_j}\otimes\rho_x\Big),
  \qquad
  \rho^{\mathrm{v}}\coloneqq\bigotimes_{j=1}^{n}
    \Big(\sum_{x\in\Xset}p_x\,\dyad{x}_{X_j}\otimes\rvirt_x\Big),
\end{equation*}
the $j$-th factor living on $X_j\Asig_j$ and on $X_j\Avirt_j$ respectively.

\emph{The two presentations.} A \emph{strategy}
is a pair
$a=(\zeta,(\Aev_j)_{j=1}^{n})$ of an initial adversary state
$\zeta\in\Seq{E_0}$ and channels
\begin{equation*}\tag{\getrefnumber{C-eq:srcstrategy}}\digesteqlabel{eq:srcstrategy}
  \Aev_1\in\mathrm{CPTP}\big(\Asig^{n}E_0,\ \Bsys_1E'_1\big),
  \qquad
  \Aev_j\in\mathrm{CPTP}\big(E_{j-1},\ \Bsys_jE'_j\big)\ \ (2\le j\le n),
\end{equation*}
with $E'_j$, $E_j$ as in Definition~\ref{def:roundmodel}: the \emph{trusted}
branch of \eqref{eq:attackchannel}, the only branch this lemma is about, in
which $A_j$ is Alice's and is no output of $\Aev_j$. A strategy of the
\emph{virtual} type is the same with $\Avirt^{n}$ in place of $\Asig^{n}$ in
\eqref{eq:srcstrategy}. The \emph{specified-source} presentation
$\mathsf{P}^{\mathrm{sp}}$ has these strategies and the output map
$\mathsf{out}^{\mathrm{sp}}(a)\coloneqq$ the state \eqref{eq:protocol} produces
on the output registers of $\Kfin$ together with $E_n$, from the input
$\rho^{\mathrm{sp}}\otimes\zeta$, with $\Bsys_j$ supplied by $\Aev_j$ and the
classical $X_j$ carried into $\bar{\Grnd}_j$; the \emph{virtual-source}
presentation $\mathsf{P}^{\mathrm{v}}$ has the strategies of the virtual type
and the output map built the same way from $\rho^{\mathrm{v}}\otimes\zeta$, the
two protocols agreeing in \emph{every} other component: the
measured system and its POVM, $\bar{\Grnd}_j$ and $\Kfin$. The output registers are literally
the same, which is condition (R1) of Definition~\ref{def:opreduction}.
\begin{enumerate}[label=(\roman*),leftmargin=3.0em]
\item \emph{(the source map)} define
  $r_{\mathrm{src}}(\zeta,(\Aev_j)_j)\coloneqq(\zeta,(\Aev^{\mathrm{v}}_j)_j)$
  by
  \begin{equation*}\tag{\getrefnumber{C-eq:srcmap-strategy}}\digesteqlabel{eq:srcmap-strategy}
    \Aev^{\mathrm{v}}_1\coloneqq
      \Aev_1\circ\big(\Xi^{\otimes n}\otimes\idmap_{E_0}\big),
    \qquad
    \Aev^{\mathrm{v}}_j\coloneqq\Aev_j\ \ (2\le j\le n).
  \end{equation*}
  Then $r_{\mathrm{src}}$ is an exact operational reduction from
  $\mathsf{P}^{\mathrm{sp}}$ to $\mathsf{P}^{\mathrm{v}}$ in the sense of
  Definition~\ref{def:opreduction}, with the whole of the virtual type as the
  admissible strategy set;
\item \emph{(the two descriptions of the virtual protocol)} let the virtual
  protocol be described \emph{entanglement-based} instead: Alice's retained
  systems $A_1,\dots,A_n$ carry, jointly with the virtual transmitted systems,
  the preparation $\Psi^{\otimes n}$, where $\Psi$ is a source-replacement
  certificate for the declared data $(A,\mathbb{T},\tau_A,\{M_x\})$
  against $\{p_x,\rvirt_x\}$ in any of the declared forms of
  \xref{transmitter/def:srcert}, satisfying (SR1)--(SR3); and $x_j$ is
  produced by her declared measurement channel $\pcNmsm$, the
  measurement map of part (iii) of \xref{transmitter/thm:certificate-sufficiency}, applied to
  $A_j$ at its declared time in round $j$. Write $\mathsf{out}^{\mathrm{eb}}$
  for the resulting output map on the virtual type. Then, under A-SEQ
  (\xref{transmitter/ass:seq}),
  $\mathsf{out}^{\mathrm{eb}}(a)=\mathsf{out}^{\mathrm{v}}(a)$ for every
  strategy $a$ of that type, so the identity map on strategies is an exact
  operational reduction between the two presentations in either direction;
\item \emph{(the composite)} consequently $r_{\mathrm{src}}$ is an exact
  operational reduction from $\mathsf{P}^{\mathrm{sp}}$ to the virtual protocol
  in its entanglement-based description, the object the rest of this
  contract and the consuming engine analyse;
\item \emph{(one channel for every label is essential)} \eqref{eq:sourceidentity}
  with \emph{one} $\Xi$ cannot be weakened to a family $(\Xi_x)_x$ with
  $\Xi_x(\rvirt_x)=\rho_x$. Take $n=1$, $\Xset=\{0,1\}$,
  $p_0=p_1=\tfrac12$, $\Asig=\Avirt=\mathbb{C}^{2}$,
  $\rvirt_0=\rvirt_1=\dyad{0}$, $\rho_0=\dyad{0}$, $\rho_1=\dyad{1}$, and
  $\Xi_x(\cdot)\coloneqq\tr[\cdot]\,\rho_x$. Every identity
  $\Xi_x(\rvirt_x)=\rho_x$ holds, yet for the attack that measures the received
  signal in the basis $\{\ket{0},\ket{1}\}$ and keeps the outcome in $E'_1$,
  the marginal on $X_1E'_1$ is perfectly correlated in the specified-source
  protocol and, for \emph{every} attack of the virtual type, a product in the
  virtual one, so no map of strategies can make the two
  full outputs equal; the protocol exhibited in the proof announces nothing in round one, so that
  $E_1$ is $E'_1$ there.
\end{enumerate}
\end{digestclaim}

% claim reductions/thm:two-sided
\digestlabel{thm:two-sided}{8.6}\pdfbookmark[2]{Theorem 8.6 The two-sided reduction in trusted-source mode}{digest-reductions/thm:two-sided}
\begin{mainclaim}
\renewcommand*{\RestatedName}{\hyperref{contract.pdf}{}{claim-reductions/thm:two-sided}{Theorem~8.6}}%
\begin{digestclaim}[The two-sided reduction in trusted-source mode]
Assume the hypotheses of Lemma~\ref{lem:source-transfer}: a trusted-source
specification with virtual ensemble $\{p_x,\rvirt_x\}$, a source segment
whose composite $\Xi$ satisfies the source identity
\eqref{eq:sourceidentity}, a source-replacement certificate satisfying
(SR1)--(SR3), and A-SEQ. Assume in addition $M\dsq_{\Lambda}F$ together with
a spectral pull-back certificate \eqref{eq:certificate} for the list
$\CList=\{(\hat O_k,\bowtie_k,c_k)\}$, or the factorwise hypotheses of
Theorem~\ref{thm:licence}\,(v) for a list with entries carrying the
party-local tag. Let $\mathsf{P}^{\mathrm{sp}}$ be the specified-source
presentation of Lemma~\ref{lem:source-transfer}, its receiver measuring $M$
on $\Bsys_j$, and let $\mathsf{P}^{\mathrm{v},\sq}_{\CList}$ be the
presentation whose protocol has the virtual source and the target $F$ on
$\Qsys_j$ and whose strategies are the $\CList$-admissible sequences of the
virtual type (Definition~\ref{def:attset}). Define
\begin{equation*}\tag{\getrefnumber{C-eq:twosidedmap}}\digesteqlabel{eq:twosidedmap}
  r_{\mathrm{two}}\big(\zeta,(\Aev_j)_j\big)
  \coloneqq\big(\zeta,(\Aev^{\mathrm{v},\sq}_j)_j\big),
  \qquad
  \Aev^{\mathrm{v},\sq}_1
    \coloneqq\lambda_1\circ\Aev_1\circ
      \big(\Xi^{\otimes n}\otimes\idmap_{E_0}\big),
  \quad
  \Aev^{\mathrm{v},\sq}_j\coloneqq\lambda_j\circ\Aev_j\ \ (j\ge2),
\end{equation*}
where $\lambda_j\coloneqq\Lambda\otimes\idmap_{E'_j}$ is the trusted-source
branch of \eqref{eq:composite}, the only branch the strategies
\eqref{eq:srcstrategy} of $\mathsf{P}^{\mathrm{sp}}$ are typed in.
\begin{enumerate}[label=(\roman*),leftmargin=3.0em]
\item \emph{(one map, and the order of the two replacements is immaterial)}
  $r_{\mathrm{two}}=r_{\sq}\circ r_{\mathrm{src}}
  =r_{\mathrm{src}}\circ r_{\sq}$, where $r_{\mathrm{src}}$ is the source map
  \eqref{eq:srcmap-strategy} and $r_{\sq}$ the squashing map
  $(\Aev_j)_j\mapsto(\Aev^{\sq}_j)_j$ of \eqref{eq:composite}; both composites
  are defined on the same strategies and agree there. And $r_{\mathrm{two}}$ is
  an exact operational reduction from $\mathsf{P}^{\mathrm{sp}}$ to the
  virtual, squashed protocol taken with \emph{all} strategies of its type;
\item \emph{(the image is admissible)} let $E_0$, every $E'_j$ and $\Qsys$ be
  finite-dimensional, as the round model of Definition~\ref{def:roundmodel}
  makes the adversary's registers, every register of it being
  finite-dimensional, and Definition~\ref{def:system-setting} makes the
  target space, and let
  $\Hs{\Avirt}$ be finite-dimensional, as it is in both
  cases of \xref{transmitter/def:source-spec}: trivially for an ensemble
  declared as matrices, and by its clause (b) for a separable specified
  system. Then every $\Aev^{\mathrm{v},\sq}_j$ of \eqref{eq:twosidedmap}
  lies in the set Definition~\ref{def:attset} assigns to its round, so
  $r_{\mathrm{two}}$ has its image in the strategy set of
  $\mathsf{P}^{\mathrm{v},\sq}_{\CList}$ and is an exact operational reduction
  into it. \emph{No hypothesis on the specified transmitted system is needed}:
  $\Hs{\Asig}$ may be separable and infinite-dimensional, as the transmitted
  system of a source model is, and it is not a register
  Definition~\ref{def:system-setting} or the round model ranges over
  but a system in transit;
\item \emph{(transfer, with the same $\varepsilon$)} consequently, if the
  virtual, squashed protocol is $\varepsilon$-secure against every
  $\CList$-admissible sequence of the virtual type, a security criterion in
  the sense of Definition~\ref{def:criterion} taking a value at most
  $\varepsilon$ at each of them, then the specified-source physical protocol
  is $\varepsilon$-secure against \emph{every} attack sequence of its
  type, with the \emph{same} $\varepsilon$ and simultaneously for every such
  criterion;
\item \emph{(in the description the certificate presents)} (i) to (iii) hold
  with the virtual protocol read entanglement-based, Alice retaining $A_j$,
  which trusted-source mode is the mode that gives her, and measuring
  $\{M_x\}$, this being the description the consuming engine analyses;
\item \emph{(the empty source segment)} where the source segment is empty,
  $\Xi$ is the identity, $\Avirt=\Asig$, $\rvirt_x=\rho_x$,
  $r_{\mathrm{src}}$ is the identity map on strategies and $r_{\mathrm{two}}$ is
  $r_{\sq}$; (i) to (iii) are then Corollary~\ref{cor:sequential} and
  Theorem~\ref{thm:licence}\,(iii) word for word, and the theorem states
  nothing new about the case this module emits.
\end{enumerate}
\end{digestclaim}
\end{mainclaim}

\section*{9\quad The number model and exact decisions}\pdfbookmark[1]{9 The number model and exact decisions}{digest-group-reductions/def:exactadm}

% claim reductions/def:exactadm
\digestlabel{def:exactadm}{9.1}\pdfbookmark[2]{Definition 9.1 The number model, the declared lattice and exact admissibility}{digest-reductions/def:exactadm}
\renewcommand*{\RestatedName}{\hyperref{contract.pdf}{}{claim-reductions/def:exactadm}{Definition~9.1}}%
\begin{digestclaim}[The number model, the declared lattice and exact admissibility]
\emph{(a) The model.} Every numerical datum this module reads or emits,
the displacement record below aside, is a binary64 value on the wire and
\emph{is} that value: no exact object is posited behind it and none is
approximated by it. A binary64 value is a dyadic
rational, and \emph{this contract reads it as that rational}
(\xref{foundations/def:binary64}\,(iii)). Every predicate this contract
decides is decided, and every emitted number but one is formed, by finitely
many operations in the substrate of Definition~\ref{def:substrate} and
Lemma~\ref{lem:exactpsd}. \emph{The exception is the emission displacement
$\Delta_s$ of \eqref{eq:displacement}}: an entrywise maximum \emph{modulus},
hence the square root of an exact rational and in general irrational, emitted
as a binary64 rendering of a record of what the normalizer of
Definition~\ref{def:emitnorm} did, in a direction nothing records; it decides
nothing, and no statement of this contract reads it. That is the reading the
exact island (\xref{foundations/def:exact-island}) gives, which
contains both the admission decisions on declared blocks and a producer's
construction of what it emits. The only non-numerical data are combinatorial
(alphabets, orders, block identifiers, a coarse-graining map, a branch index
set, a cutoff) and carry no arithmetic.

\emph{(b) The declared lattice.} The request declares one integer exponent
$\kappa$, and every declared and every emitted operator entry, together with
every \emph{real scalar} the request declares of the following kinds --- the
class weights $q_t$, any certified weight, the entries of a declared
stochastic matrix, and any declared device constant --- satisfies
\begin{equation*}\tag{\getrefnumber{C-eq:latinv}}\digesteqlabel{eq:latinv}
  2^{\kappa}x\in\mathbb{Z},
  \qquad |x|\le1,
  \qquad 1\le\kappa\le53 .
\end{equation*}
Such an $x$ is exactly a binary64 number
(\xref{foundations/lem:binary64-lattice}\,(i)), its significand carrying at
most $\kappa\le53$ bits. The list is exhaustive: the constant a caller declares for a
requested constraint entry is not among them, and where it is off the
lattice it is rounded onto it or refused (Corollary~\ref{cor:entries}\,(iii)). It is a condition on a \emph{set} of values: it says which numbers
may be declared and nothing about how near one number is to another. A
complex entry satisfies \eqref{eq:latinv} when its real and imaginary parts
do.

\emph{(c) Exact admissibility on the declared lattice.} Let
$\{A_v\}_{v\in I}$ be a finite family of matrices presented by its blocks
over declared bases, indexed by the finite set $I$ the invoking rule names
(the outcome alphabet at a fixed setting for a declared block family, a class
$\Vset_t$ of the datum of a reduction instance
(Definition~\ref{def:reduction-instance}), $\Vset$ for the family this module
emits), and let $c$ be the constant the rule names for its completeness
identity: $1$ for a POVM, the class weight $q_t$ for a
class of a datum. The family is
\emph{exactly admissible on the declared lattice}
when
\begin{enumerate}[label=(\alph*),leftmargin=2.4em]
\item \emph{(the sector map)} the block kind, the block identifiers and their
  order are the ones the consumer holds, a comparison of declared labels,
  decided and never measured;
\item \emph{(Hermiticity)} $A_v=A_v^{\dagger}$, entry for entry;
\item \emph{(positivity)} $A_v\succeq0$, decided by \eqref{eq:psdtest};
\item \emph{(completeness)} $\sum_{v\in I}A_v=c\,\idon{}$ on each block, as an
  identity of the entries;
\item \emph{(the sum rounds nowhere)} every entrywise binary64
  addition of that accumulation, taken in the declared order of $I$,
  is exact, so that the \emph{computed} sum \emph{is} the sum of (d);
\end{enumerate}
and every entry satisfies \eqref{eq:latinv}. Conditions (a)--(e) say that
the family is exactly admissible with constant $c$ for the sector map the
consumer holds (\xref{foundations/def:exact-admissibility}), each block being
presented at the dimension the consumer's map declares for it. Condition (e)
is a condition in its own right and is stated on the \emph{additions}: a
vanishing completeness \emph{defect} says that the computed sum is the
identity, not that $\sum_vA_v=c\,\idon{}$ holds of the numbers.
Lemma~\ref{C-lem:latsum} shows that (e) is satisfiable on this
contract's own emission, rather than a condition that would refuse every
family.

\emph{(d) Scope, stated once and binding wherever this definition is
invoked.} This definition and the rules that invoke it are about data this
contract requires to be presented as a finite list of matrices of declared
dimension: the data sectors of a reduction instance
(Definition~\ref{def:reduction-instance}\,(I1)) and the declared blocks they
are formed from; the declared classical scalars and matrices a submodule's
route reads (a stochastic matrix and its descended form, the class weights,
each certified weight, any declared device constant); and the family this
module emits. They are \emph{not} available for the objects an instance's
hypotheses (H) are about: operators on a part of $\Hs{\Bsys}$ that no
finite list of matrices presents, an above-cutoff family, a register POVM, an
uncharacterised branch. There is nothing on those to decide, and their
conditions remain hypotheses of the instance
(Definition~\ref{def:reduction-instance}\,(I1)).
\end{digestclaim}

% claim reductions/def:substrate
\digestlabel{def:substrate}{9.2}\pdfbookmark[2]{Definition 9.2 The arithmetic substrate}{digest-reductions/def:substrate}
\renewcommand*{\RestatedName}{\hyperref{contract.pdf}{}{claim-reductions/def:substrate}{Definition~9.2}}%
\begin{digestclaim}[The arithmetic substrate]
The
\emph{substrate}
of this contract is an ordered field $\Fld$ with
$\mathbb{Q}\subseteq\Fld\subseteq\mathbb{R}$ together with its Gaussian
extension $\Fldi\coloneqq\Fld+i\Fld$, which is closed under complex
conjugation. What is fixed is the \emph{predicate} each rule decides and the
requirement that it be decided exactly; which such field carries the
arithmetic is not. The choice cannot change any decision. Binary64 values are
dyadic rationals (Definition~\ref{def:exactadm}\,(a)), and every number a rule
reads is formed from the declared entries and from rational constants by
field operations, floors of rationals and roundings of rationals to a
lattice, so it is rational; since $\mathbb{Q}$ is an ordered subfield of every
admissible $\Fld$ with the induced order, and $\mathbb{Q}+i\mathbb{Q}$ a
subfield of $\Fldi$, each decision, a comparison or an equality of such
numbers, is the same decision taken over $\mathbb{Q}$. In particular
$\Fld=\mathbb{Q}$ always suffices. The latitude is bounded: it is a choice of
arithmetic substrate, not of mathematics.
\end{digestclaim}

% claim reductions/lem:exactpsd
\digestlabel{lem:exactpsd}{9.3}\pdfbookmark[2]{Lemma 9.3 Exact decision of positivity and of an operator-norm bound}{digest-reductions/lem:exactpsd}
\renewcommand*{\RestatedName}{\hyperref{contract.pdf}{}{claim-reductions/lem:exactpsd}{Lemma~9.3}}%
\begin{digestclaim}[Exact decision of positivity and of an operator-norm bound]
Let $\Fld,\Fldi$ be the substrate of Definition~\ref{def:substrate} and let
$d,d'\ge1$.
\begin{enumerate}[label=(\roman*),leftmargin=3.0em]
\item Let $X=X^{\dagger}\in\Lin{\mathbb{C}^{d}}$ have all entries in $\Fldi$ and
  write
  \begin{equation*}\tag{\getrefnumber{C-eq:charpoly}}\digesteqlabel{eq:charpoly}
    \det\big(\lambda\,\idon{d}+X\big)=\sum_{k=0}^{d}e_k(X)\,\lambda^{d-k},
    \qquad e_0(X)=1 ,
  \end{equation*}
  so that $e_k(X)$ is the sum of the $k\times k$ principal minors of $X$. Then
  every $e_k(X)$ lies in $\Fld$, and
  \begin{equation*}\tag{\getrefnumber{C-eq:psdtest}}\digesteqlabel{eq:psdtest}
    X\succeq0
    \quad\Longleftrightarrow\quad
    e_k(X)\ge0\ \text{ for every }k=1,\dots,d .
  \end{equation*}
\item Let $Z\in\Linn{\mathbb{C}^{d'}}{\mathbb{C}^{d}}$ have all entries in $\Fldi$ and
  let $\delta\in\Fld$, $\delta\ge0$. Then
  \begin{equation*}\tag{\getrefnumber{C-eq:normtest}}\digesteqlabel{eq:normtest}
    \opn{Z}\le\delta
    \quad\Longleftrightarrow\quad
    \delta^{2}\,\idon{d'}-Z^{\dagger}Z\;\succeq\;0 ,
  \end{equation*}
  the right-hand side being a test of the form (i) on a Hermitian matrix with
  entries in $\Fldi$.
\end{enumerate}
Both predicates are therefore decided by finitely many field operations and
sign comparisons in $\Fld$: no rounding, no threshold, and \emph{no
eigenvalue is computed}.
\end{digestclaim}

\section*{10\quad Reduction instances}\pdfbookmark[1]{10 Reduction instances}{digest-group-reductions/def:reduction-instance}

% claim reductions/def:reduction-instance
\digestlabel{def:reduction-instance}{10.1}\pdfbookmark[2]{Definition 10.1 Reduction instance}{digest-reductions/def:reduction-instance}
\renewcommand*{\RestatedName}{\hyperref{contract.pdf}{}{claim-reductions/def:reduction-instance}{Definition~10.1}}%
\begin{digestclaim}[Reduction instance]
Fix the declared lattice exponent $\kappa$ of
Definition~\ref{def:exactadm}\,(b) and a joint alphabet with
classes and class weights $\big(\Vset,\{\Vset_t\}_{t\in\Tset},\{q_t\}\big)$
(Definition~\ref{def:joint-family}\,(a)). Let $\SecSet_{\mathrm{dat}}$ be a
finite ordered list of
\emph{data sector}s,
each $s$ with a declared dimension $d_s\ge1$ and a declared
\emph{support}
$\Vset^{(s)}\subseteq\Vset$ that meets every class. A
\emph{datum}
for these sectors is a family
$A=\{A^{(s)}_v\}_{s\in\SecSet_{\mathrm{dat}},\,v\in\Vset}$ of $d_s\times d_s$
matrices with entries in $\Fldi$ and $A^{(s)}_v=0$ for
$v\notin\Vset^{(s)}$: these elements lie
\emph{outside the sector support},
and no step of this contract reads, rounds or repairs them. It is an
\emph{admissible datum},
written $A\in\mathcal{D}_{\mathrm{adm}}$, when for every $s$:
(a) $A^{(s)}_v=(A^{(s)}_v)^{\dagger}$ entry for entry; (b)
$A^{(s)}_v\succeq0$, decided by \eqref{eq:psdtest}; (c) it is
\emph{class complete},
\begin{equation*}\tag{\getrefnumber{C-eq:classcomplete}}\digesteqlabel{eq:classcomplete}
  \sum_{v\in\Vset_t}A^{(s)}_v\;=\;q_t\,\idon{d_s}
  \qquad\text{exactly, for every }t\in\Tset ;
\end{equation*}
and it is an
\emph{exactly admissible datum},
written $A\in\mathcal{D}_{\kappa}$, when in addition (d) every entry
satisfies \eqref{eq:latinv}. These are conditions (b)--(d) of
Definition~\ref{def:exactadm}\,(c) and the lattice invariant, read per data
sector with the class weight $q_t$ as the constant of the completeness
identity.

A
\emph{reduction instance}
$\mathcal{J}$ over $\big(\Vset,\{\Vset_t\},\{q_t\}\big)$ consists of the
following four parts.
\begin{enumerate}[label=(I\arabic*),leftmargin=3.0em]
\item \emph{(the handed model and its datum)} A system $\Bsys$ with
  $\Hs{\Bsys}$ separable (Definition~\ref{def:system-setting}); data sectors
  $\SecSet_{\mathrm{dat}}$ with their declared dimensions and supports; a
  \emph{handed datum}
  $A\in\mathcal{D}_{\mathrm{adm}}$, every entry of which is formed from the
  numbers the instance was handed by finitely many operations in $\Fldi$; a
  finite list of
  \emph{hypotheses (H)}
  about the rest of the model (operators without a finite presentation,
  facts about a part of $\Hs{\Bsys}$ the datum does not describe, a register
  POVM), each of which is either decided exactly on declared data, the
  request being refused where it fails, or carried as an assumption of the
  emitted ledger; and a
  \emph{model map}
  $A'\mapsto M[A']=\{M_v[A']\}_{v\in\Vset}$ that assigns to every admissible
  datum $A'\in\mathcal{D}_{\mathrm{adm}}$ a POVM on $\Bsys$ over $\Vset$,
  under (H). The
  \emph{handed model}
  is $M[A]$, the model map at the handed datum. The
  \emph{stability clause}
  requires that no hypothesis in (H) read an entry of the datum, so that (H)
  holds of $M[A']$ for every $A'\in\mathcal{D}_{\mathrm{adm}}$ as soon as it
  holds of the handed model; and that parts (I2)--(I4) be established
  for \emph{every} $A'\in\mathcal{D}_{\mathrm{adm}}$ under (H), not only at
  the handed datum.
\item \emph{(the target)} A finite ordered
  \emph{sector set}
  $\SecSet=\SecSet_{\mathrm{dat}}\sqcup\SecSet_{\mathrm{sc}}$, the second
  summand the
  \emph{scalar sector}s,
  each $s\in\SecSet$ with a declared dimension $d_s\ge1$; for each scalar
  sector $s$ and each $v\in\Vset$ a
  \emph{scalar value}
  $\phi_s(v)\in[0,1]$ satisfying \eqref{eq:latinv} with
  $\sum_{v\in\Vset}\phi_s(v)=1$ exactly, each formed from the numbers the
  instance was handed by finitely many operations in $\Fld$ and, where a
  rounding onto the lattice is needed, by the \emph{downward} rounding of a
  quantity the value is a certified lower bound of, any complementary value
  being recomputed from the rounded ones and never rounded on its own; the
  target space $\Hs{\Qsys}\coloneqq\bigoplus_{s\in\SecSet}\mathbb{C}^{d_s}$
  with sector projectors $\Pi_s$, so that $\dim\Hs{\Qsys}=\sum_sd_s$; and the
  \emph{target family}
  \begin{equation*}\tag{\getrefnumber{C-eq:instance-target}}\digesteqlabel{eq:instance-target}
    F_v[A']\;\coloneqq\;\bigoplus_{s\in\SecSet}F^{(s)}_v[A'],
    \qquad
    F^{(s)}_v[A']\;\coloneqq\;
    \begin{cases}
      A'^{(s)}_v & s\in\SecSet_{\mathrm{dat}},\\[2pt]
      \phi_s(v)\,\idon{d_s} & s\in\SecSet_{\mathrm{sc}},
    \end{cases}
  \end{equation*}
  assigned one sector block at a time, so that $F_v[A']$ is
  sector-block-diagonal \emph{by construction}: no off-sector entry is
  written and none is to be verified. The sector map of the instance
  is the list $\SecSet$ with its order, each sector's dimension and its kind
  (data or scalar), and it is what a consumer compares
  (Definition~\ref{def:exactadm}\,(c), condition (a)). The
  \emph{default party partition}
  of the target assigns every factor of every $\Hs{\Qsys_s}$,
  $s\in\SecSet$, to \emph{Bob}: \emph{every sector of the target is Bob's}.
  That is the assignment the physical model makes: the target describes what
  reaches Bob's receiver, the handed model living on his system $\Bsys$. An
  instance may in addition be declared with \emph{another party's label},
  one label for the whole instance, and then the same assignment is made
  with that label in its place: every factor of every $\Hs{\Qsys_s}$ is that
  party's, because the handed model lives on the system \emph{that} party
  receives. Where no label is declared the label is Bob's. Nothing else in
  the instance depends on the label: it is a declared identifier, compared
  and never computed, and no condition of (I1)--(I4) reads it.
\item \emph{(the existence theorem)} A channel
  $\Lambda[A']=\PinS\circ\Lambda_0[A']\colon\Tcl{\Bsys}\to\Tcl{\Qsys}$,
  given by a constructor and recorded as one, with $\PinS$ the sector
  pinching of Proposition~\ref{prop:pinching} for $\{\Pi_s\}_{s\in\SecSet}$;
  and a finite
  \emph{transported set}
  $\hat O[A']$ of sector-block-diagonal Hermitian operators on $\Qsys$
  containing every sector projector $\Pi_s$, with
  pull-backs $O[A']$; such that
  \begin{equation*}\tag{\getrefnumber{C-eq:instance-relation}}\digesteqlabel{eq:instance-relation}
    \big(M[A'];\,O[A']\big)\ \dsq_{\Lambda[A']}\ \big(F[A'];\,\hat O[A']\big)
    \qquad\text{for every }A'\in\mathcal{D}_{\mathrm{adm}}\text{ under (H)}.
  \end{equation*}
  By linearity of $\Lambda[A']^{\dagger}$, every real linear combination of
  transported operators is transported to the same combination of their
  pull-backs.
\item \emph{(the certified entries)} A finite list
  $\CList[A']=\{(\hat O_k,\bowtie_k,c_k)\}_{k}$ of
  \emph{certified entries},
  each $\hat O_k$ a real linear combination of transported operators of the
  form
  \begin{equation*}\tag{\getrefnumber{C-eq:instance-entry}}\digesteqlabel{eq:instance-entry}
    \hat O_k\;=\;\sum_{v\in\Vset}\sum_{s\in\SecSet}
      \alpha_{k,v,s}\,\Pi_sF_v[A']\Pi_s\;+\;Z_k ,
  \end{equation*}
  with real coefficients $\alpha_{k,v,s}$ \emph{on the lattice}, each an
  integer multiple of $2^{-\kappa}$, and $Z_k$ a real combination of sector
  projectors with coefficients of the same kind, $\bowtie_k\in\{\preceq,=,\succeq\}$
  and $c_k$ a declared constant on the lattice, whose certificate
  \eqref{eq:certificate},
  $\Lambda[A']^{\dagger}(\hat O_k)\bowtie_kc_k\,\idon{\Bsys}$, is
  \emph{discharged} for every $A'\in\mathcal{D}_{\mathrm{adm}}$ under (H); and
  a
  \emph{class-marginal flag}
  recording whether the class-marginal identity
  \begin{equation*}\tag{\getrefnumber{C-eq:instance-classmarginal}}\digesteqlabel{eq:instance-classmarginal}
    \sum_{v\in\Vset_t}M_v[A']\;=\;q_t\,\idon{\Bsys}
    \qquad\text{for every }t\in\Tset
  \end{equation*}
  is \emph{discharged}, that is, established by the construction for every
  $A'\in\mathcal{D}_{\mathrm{adm}}$ under (H), or not.
  The modulus bound of \eqref{eq:latinv} is not imposed on a coefficient:
  $|x|\le1$ is a condition on a declared or emitted \emph{entry} (a value of
  the datum, a scalar value, a declared scalar or constant), where a POVM
  element and a completeness identity give it, and not on the coefficients
  of a linear combination, which a branching instance forms as the
  difference $\beta_s-c_k$ of two lattice multiples
  (Theorem~\ref{thm:instrument-instance}), again a lattice multiple;
  \emph{being on the lattice is all that is asked of a coefficient}. What
  \eqref{eq:latinv} does bind is the entries of the operator the combination
  produces (Corollary~\ref{cor:entries}\,(iii)).
\end{enumerate}
The instance \emph{at} a datum $A'\in\mathcal{D}_{\mathrm{adm}}$ is the same
$\Bsys$, (H), sectors, supports, scalar values, constructor and entry rule
with $A'$ in place of the handed datum; by the stability clause it is a
reduction instance whenever $\mathcal{J}$ is.

\emph{Five consequences, each immediate.} \emph{First}, for every
$A'\in\mathcal{D}_{\mathrm{adm}}$ the target $\{F_v[A']\}_{v\in\Vset}$ is a
POVM on $\Qsys$: every sector block is Hermitian and positive semidefinite,
by (a) and (b) on a data sector and by $\phi_s(v)\in[0,1]$ on a scalar one;
and it is complete, on a data sector because
$\sum_{v\in\Vset}A'^{(s)}_v=\sum_{t}\sum_{v\in\Vset_t}A'^{(s)}_v
=\sum_tq_t\idon{d_s}=\idon{d_s}$ by \eqref{eq:classcomplete} and
$\sum_tq_t=1$, on a scalar sector because $\sum_v\phi_s(v)=1$. Its data
sectors \emph{are} the datum, entry by entry, so the emitted target carries
the datum and nothing else of the model. \emph{Second}, every $F_v[A']$,
every $\Pi_s$ and every transported operator is sector-block-diagonal, so
Corollary~\ref{cor:blockdiag}\,(i) applies: a relation established for
$\Lambda_0[A']$ holds for $\PinS\circ\Lambda_0[A']$ with the same target,
the same transported set and the same pull-backs, which is why (I3) may be
stated with the pinching in front, and the squashed analysis may be
restricted to block-diagonal states without excluding a reachable one
(Corollary~\ref{cor:blockdiag}\,(iii)). \emph{Third}, the pull-back of every
sector projector is positive semidefinite, $\Lambda[A']^{\dagger}(\Pi_s)
\succeq0$, because the adjoint of a channel is a positive map, and
$\sum_s\Lambda[A']^{\dagger}(\Pi_s)=\idon{\Bsys}$ by unitality; so
$(\Pi_s,\succeq,0)$ carries the certificate of
Corollary~\ref{cor:pullback}\,(ii) for \emph{every} instance and needs no
part of (I4) to say so. \emph{That certificate travels with the transported
set, whatever the request carries}, which is where (I3) puts every sector
projector and where the artefact records the certificate case per entry; it
becomes an \emph{entry} of $\CList$ only where a caller requests it \emph{as
an entry} (Corollary~\ref{cor:entries}\,(i)), the emitted list being empty
where the request carries no entry. \emph{Fourth}, the relation pulls a class
sum of the target back to the class sum of the model,
$\Lambda[A']^{\dagger}\big(\sum_{v\in\Vset_t}F_v[A']\big)
=\sum_{v\in\Vset_t}M_v[A']$, by linearity; so
\eqref{eq:instance-classmarginal} holds of the model exactly when the
pull-back identity $\Lambda[A']^{\dagger}\big(\sum_{v\in\Vset_t}F_v[A']\big)
=q_t\idon{\Bsys}$ does, which is the expectation-form entry
Corollary~\ref{cor:entries}\,(ii) emits where the flag says it is discharged,
and never the operator identity
$\sum_{v\in\Vset_t}F_v[A']=q_t\idon{\Qsys}$, which fails on every scalar
sector whose values do not sum to $q_t$ over $\Vset_t$. \emph{Fifth},
$\dim\Hs{\Qsys}=\sum_sd_s$ is a sum of declared dimensions, finite, and the
sector map is a comparison of declared labels; both are read off the
declaration with no arithmetic performed.

The link of Definition~\ref{def:link} an instance produces is the relation
\eqref{eq:instance-relation} at the declared datum, with the transported set
and the entry list as its slots, and its chain (Definition~\ref{def:chain})
is that link preceded by whatever post-processing the instance's submodule
applies before it.
\end{digestclaim}

% claim reductions/cor:identity-instance
\digestlabel{cor:identity-instance}{10.2}\pdfbookmark[2]{Corollary 10.2 The identity instance}{digest-reductions/cor:identity-instance}
\renewcommand*{\RestatedName}{\hyperref{contract.pdf}{}{claim-reductions/cor:identity-instance}{Corollary~10.2}}%
\begin{digestclaim}[The identity instance]
Let $\big(\Vset,\{\Vset_t\},\{q_t\}\big)$ be a joint alphabet with classes
and weights, let $\SecSet_0$ be a declared finite ordered sector set with
dimensions $d_s$, and let $\{M_v\}_{v\in\Vset}$ be a POVM on
$\Hs{\Bsys}=\bigoplus_{s\in\SecSet_0}\mathbb{C}^{d_s}$ presented as the
sector-block list $M_v=\bigoplus_sA^{(s)}_v$ of an admissible datum
$A\in\mathcal{D}_{\mathrm{adm}}$ for the sectors $\SecSet_0$ with the default
supports $\Vset^{(s)}=\Vset$, that is, a declared family every block of which
is decided exactly admissible on the declared lattice on entry with
the class sums at $q_t$ (Definition~\ref{def:exactadm}\,(c)). Then the
\emph{identity instance} $\mathcal{J}_{\mathrm{id}}$, given by
\begin{enumerate}[label=(I\arabic*),leftmargin=3.0em]
\item $\Bsys$ as given; $\SecSet_{\mathrm{dat}}=\SecSet_0$; the handed datum
  $A$; (H) empty; and the model map $M_v[A']\coloneqq\bigoplus_{s}A'^{(s)}_v$;
\item $\SecSet_{\mathrm{sc}}=\varnothing$, so that $\Qsys=\Bsys$ and
  $F[A']=M[A']$;
\item $\Lambda=\PinS$ for the decomposition $\SecSet_0$ (that is,
  $\Lambda_0=\idmap$), independent of the datum, with the transported set
  $\{\Pi_s\}_{s\in\SecSet_0}$ and pull-backs $\Pi_s$;
\item an empty entry list, and the class-marginal identity
  discharged,
\end{enumerate}
is a reduction instance. Its chain is the identity link: one receiver-stage
link of kind
\textsf{identity} with $\Lambda_{\Lnk}=\PinS$ (the identity on every state
block diagonal for $\SecSet_0$, and $\idmap$ outright when $\SecSet_0$ has one
sector), with $M^{\mathrm{out}}=M^{\mathrm{in}}$, the sector projectors
carried to themselves, the reachable-set statement the declared input
family, pinched, no entries, no licence and
$\textsf{status}=\textsf{exact}$.
\end{digestclaim}

% claim reductions/cor:flag-instance
\digestlabel{cor:flag-instance}{10.3}\pdfbookmark[2]{Corollary 10.3 The flag instance of an uncharacterised branch}{digest-reductions/cor:flag-instance}
\renewcommand*{\RestatedName}{\hyperref{contract.pdf}{}{claim-reductions/cor:flag-instance}{Corollary~10.3}}%
\begin{digestclaim}[The flag instance of an uncharacterised branch]
Let $\big(\Vset,\{\Vset_t\},\{q_t\}\big)$ be a joint alphabet with classes
and weights and let $\{M_v\}_{v\in\Vset}$ be any POVM on any system $\Bsys$
with $\Hs{\Bsys}$ separable, about which nothing beyond that is declared.
Then the \emph{flag instance} $\mathcal{J}_{\mathrm{flag}}[M]$, given by
\begin{enumerate}[label=(I\arabic*),leftmargin=3.0em]
\item $\Bsys$ as given; no data sector, $\SecSet_{\mathrm{dat}}=\varnothing$,
  so that the datum is empty and $\mathcal{D}_{\mathrm{adm}}=\{\varnothing\}$;
  (H) the single hypothesis that $\{M_v\}$ is a POVM on $\Bsys$; and the
  constant model map $M[\varnothing]\coloneqq M$;
\item one scalar sector $(f)$ of dimension $1$ for each $f\in\Vset$, in the
  order of $\Vset$, with $\phi_f(v)\coloneqq\delta_{vf}$, so that
  $\Hs{\Qsys}=\Hs{\Flag}=\mathbb{C}^{|\Vset|}$ and $F_v=\dyad{v}_{\Flag}$,
  the flag PVM $\Ppvm$ of Lemma~\ref{lem:top};
\item $\Lambda=\PinS\circ\Lmp=\Lmp$, the measure-and-prepare channel
  \eqref{eq:mpsquasher} of $M$, with the transported set
  $\{\Pi_f=\dyad{f}\}_{f\in\Vset}$ and pull-backs $M_f$;
\item an empty entry list, and the class-marginal identity
  discharged exactly when (H) is extended by the declared fact
  $\sum_{v\in\Vset_t}M_v=q_t\idon{\Bsys}$ for every $t$, as it is when the
  branch is declared to be a joint family
  (Definition~\ref{def:joint-family}\,(b)), and not otherwise,
\end{enumerate}
is a reduction instance. The normalizer of Definition~\ref{def:emitnorm} has
nothing to act on: the declared datum is the handed one, no $\Delta_s$
exists, and the declared model is $M$ itself.
\end{digestclaim}

% claim reductions/thm:instrument-instance
\digestlabel{thm:instrument-instance}{10.4}\pdfbookmark[2]{Theorem 10.4 The branching instance: an instrument followed by branch instances}{digest-reductions/thm:instrument-instance}
\renewcommand*{\RestatedName}{\hyperref{contract.pdf}{}{claim-reductions/thm:instrument-instance}{Theorem~10.4}}%
\begin{digestclaim}[The branching instance: an instrument followed by branch instances]
Let $J$ be a finite ordered set, let $\{\Ins_j\}_{j\in J}$,
$\Ins_j\colon\Tcl{\Bsys}\to\Tcl{\Bsys'_j}$, be a
generalized instrument (Definition~\ref{def:generalized-instrument}),
$\sum_j\tr\circ\Ins_j=\tr$, given independently of any datum, and for each
$j$ let $\mathcal{J}_j$ be a reduction instance on $\Bsys'_j$ over one common
joint alphabet with classes and weights $\big(\Vset,\{\Vset_t\},\{q_t\}\big)$,
with handed datum $A_j$, data sectors $\SecSet_{j,\mathrm{dat}}$ with supports
$\Vset^{(j,s)}$, scalar sectors $\SecSet_{j,\mathrm{sc}}$ with values
$\phi_{j,s}$, hypotheses (H$_j$), channel $\Lambda_j[\cdot]$, transported set
$\hat O^{(j)}[\cdot]$ and entry list $\CList_j[\cdot]$. Write
$\SecSet\coloneqq\bigsqcup_{j\in J}\{j\}\times\SecSet_j$, ordered by $j$ and
then as $\SecSet_j$ is, $\Pi_{(j,s)}$ for its sector projectors,
$\Pi_j\coloneqq\sum_{s\in\SecSet_j}\Pi_{(j,s)}$ for the \emph{branch
projectors}, $\hat X|_j$ for the operator equal to $\hat X$ in block $j$ and
zero elsewhere, and $A\coloneqq(A_j)_{j\in J}$ for the joint datum. Then each
of the two constructions (i) and (ii) is a reduction instance, with
\begin{itemize}[leftmargin=2.2em]
\item the channel
  \begin{equation*}\tag{\getrefnumber{C-eq:instrument-channel}}\digesteqlabel{eq:instrument-channel}
    \Lambda[A']\;\coloneqq\;\PinS\circ\bigoplus_{j\in J}\Lambda_j[A'_j]\circ\Ins_j ,
  \end{equation*}
\item the sector set $\SecSet$, the hypotheses (H)~$=\bigcup_j$(H$_j$)
  together with the hypothesis that $\{\Ins_j\}$ is a generalized
  instrument and, for every $j$ in a declared set $J_{\mathrm{w}}\subseteq J$
  of branches with declared lattice constants $q_j$, the
  \emph{branch-weight hypothesis}
  $\Ins_j^{\dagger}(\idon{\Bsys'_j})=q_j\idon{\Bsys}$,
\item the transported set consisting of the sector projectors and, for every
  $j$ and every $\hat X\in\hat O^{(j)}$, the operator $\hat X|_j$, with
  pull-back $\Ins_j^{\dagger}\big(\Lambda_j^{\dagger}(\hat X)\big)$, in
  particular
  \begin{equation*}\tag{\getrefnumber{C-eq:instrument-transport}}\digesteqlabel{eq:instrument-transport}
    \Lambda^{\dagger}\big(\Pi_{(j,s)}\big)
    =\Ins_j^{\dagger}\big(\Lambda_j^{\dagger}(\Pi_s)\big),
    \qquad
    \Lambda^{\dagger}\big(\Pi_j\big)=\Ins_j^{\dagger}\big(\idon{\Bsys'_j}\big) ,
  \end{equation*}
\item the entry list consisting of, for every branch $j$ and every entry
  $(\hat O_k,\bowtie_k,c_k)\in\CList_j$, the entry
  \begin{equation*}\tag{\getrefnumber{C-eq:instrument-entry}}\digesteqlabel{eq:instrument-entry}
    \big(\hat O_k|_j-c_k\,\Pi_j,\ \bowtie_k,\ 0\big),
  \end{equation*}
  together with the \emph{branch-weight entry} $(\Pi_j,=,q_j)$ for every
  $j\in J_{\mathrm{w}}$, where that set is not empty:
\end{itemize}
\begin{enumerate}[label=(\roman*),leftmargin=3.0em]
\item \emph{(the joint-outcome form)} over the joint alphabet
  $\Vset\times J$ of Definition~\ref{def:joint-family}\,(a) whose
  outcome alphabet is $\Yset\times J$, ordered with the branch label
  major and the outcome minor, and whose settings are $\Tset$: writing its
  element $(y,j,t)$ as $(v,j)$ with $v=(y,t)\in\Vset$, its
  declared order has
  $(v,j)<(v',j')$ iff $t<t'$, or $t=t'$ and $j<j'$, or $t=t'$, $j=j'$ and
  $y<y'$, so that every class $\Vset_t\times J$ is a contiguous run and the
  restriction of the order to $\Vset\times\{j\}$ \emph{is} the order of
  $\Vset$; its classes are $\Vset_t\times J$ and its weights $q_t$. On it, the model map
  $M_{(v,j)}[A']\coloneqq\Ins_j^{\dagger}\big(M^{(j)}_v[A'_j]\big)$; data
  sectors $(j,s)$, $s\in\SecSet_{j,\mathrm{dat}}$, of dimension $d_{j,s}$ and
  support $\Vset^{(j,s)}\times\{j\}$, with
  $A'^{(j,s)}_{(v,j')}\coloneqq\delta_{jj'}A'^{(s)}_{j,v}$; scalar sectors
  $(j,s)$, $s\in\SecSet_{j,\mathrm{sc}}$, with
  $\phi_{(j,s)}(v,j')\coloneqq\delta_{jj'}\phi_{j,s}(v)$; and the target
  $F_{(v,j)}[A']=F^{(j)}_v[A'_j]\big|_j$, the branch target in block $j$ and
  zero elsewhere;
\item \emph{(the coarse-grained form)} over $\big(\Vset,\{\Vset_t\},\{q_t\}\big)$
  itself: the model map
  $M_v[A']\coloneqq\sum_{j}\Ins_j^{\dagger}\big(M^{(j)}_v[A'_j]\big)$; the same
  data and scalar sectors with supports $\Vset^{(j,s)}$,
  $A'^{(j,s)}_v\coloneqq A'^{(s)}_{j,v}$ and $\phi_{(j,s)}(v)\coloneqq
  \phi_{j,s}(v)$; and the target $F_v[A']=\bigoplus_jF^{(j)}_v[A'_j]$.
\end{enumerate}
In both forms the class-marginal identity is discharged when every branch
discharges its own, and the flag is set in exactly that case; and the
normalizer of Definition~\ref{def:emitnorm} acts on the composite datum
\emph{branch by branch}: its output at the data sectors of branch $j$ is the
output of the normalizer of $\mathcal{J}_j$ at $A_j$, so that the declared
datum is $A'=(A'_j)_j$ with $A'_j$ the branches' declared data and the
declared model is $M[A']$ at those.
\end{digestclaim}

% claim reductions/cor:register-instance
\digestlabel{cor:register-instance}{10.5}\pdfbookmark[2]{Corollary 10.5 The register instance: register-conditioned squashing of instances}{digest-reductions/cor:register-instance}
\renewcommand*{\RestatedName}{\hyperref{contract.pdf}{}{claim-reductions/cor:register-instance}{Corollary~10.5}}%
\begin{digestclaim}[The register instance: register-conditioned squashing of instances]
Let $\Bsys=\Bsys_0D$, let $\{N_j\}_{j\in J}$ be a POVM on the register $D$
with $J$ finite and ordered, and for each $j$ let $\mathcal{J}_j$ be a
reduction instance on $\Bsys_0$ over one common joint alphabet with classes
and weights, with handed model $\{M_{v|j}\}_{v\in\Vset}$, channel
$\Lambda_j$, transported set $\hat O^{(j)}$ and entry list $\CList_j$. Let
$\Ins_j$ be the instrument of Theorem~\ref{thm:mixture},
$\Ins_j(\rho)=\tr_{D}\big[(\idon{\Bsys_0}\otimes N_j^{1/2})\rho
(\idon{\Bsys_0}\otimes N_j^{1/2})\big]$. Then the branching instance of
Theorem~\ref{thm:instrument-instance} for $\{\Ins_j\}$ and
$\{\mathcal{J}_j\}$, the \emph{register instance}, is a reduction
instance on $\Bsys$ in both forms, with the handed model
\begin{equation*}\tag{\getrefnumber{C-eq:register-model}}\digesteqlabel{eq:register-model}
  M_v=\sum_{j\in J}M_{v|j}\otimes N_j
  \quad\text{(coarse-grained)},
  \qquad
  M_{(v,j)}=M_{v|j}\otimes N_j
  \quad\text{(joint-outcome)},
\end{equation*}
the channel \eqref{eq:regchannel} pinched,
$\Lambda=\PinS\circ\bigoplus_j\Lambda_j\circ\Ins_j$, the target
$\bigoplus_jF^{(j)}_v$ resp.\ $F^{(j)}_v|_j$, the hypotheses
$\bigcup_j$(H$_j$) together with ``$\{N_j\}$ is a POVM on $D$'' and, for every
$j$ in a declared set $J_{\mathrm{w}}\subseteq J$ with declared lattice
constants $q_j$, ``$N_j=q_j\idon{D}$'', and:
\begin{enumerate}[label=(\roman*),leftmargin=3.0em]
\item \emph{(transport)} for every branch $j$ and sector $s$ of
  $\mathcal{J}_j$, and for every $\hat X\in\hat O^{(j)}$,
  \begin{equation*}\tag{\getrefnumber{C-eq:register-transport}}\digesteqlabel{eq:register-transport}
    \Lambda^{\dagger}\big(\Pi_{(j,s)}\big)=\Lambda_j^{\dagger}(\Pi_s)\otimes N_j,
    \qquad
    \Lambda^{\dagger}\big(\Pi_j\big)=\idon{\Bsys_0}\otimes N_j,
    \qquad
    \Lambda^{\dagger}\big(\hat X|_j\big)=\Lambda_j^{\dagger}(\hat X)\otimes N_j ;
  \end{equation*}
\item \emph{(the entry clause)} every entry $(\hat O,\bowtie,c)$ of a branch
  $j$ is carried as $\big(\hat O|_j-c\,\Pi_j,\bowtie,0\big)$, with pull-back
  \begin{equation*}\tag{\getrefnumber{C-eq:register-entry}}\digesteqlabel{eq:register-entry}
    \Lambda^{\dagger}\big(\hat O|_j-c\,\Pi_j\big)
    =\big(\Lambda_j^{\dagger}(\hat O)-c\,\idon{\Bsys_0}\big)\otimes N_j ,
  \end{equation*}
  a tensor product of the branch's certified operator with $N_j\succeq0$, so
  that the certificate holds at $0$ with the branch's relation; the
  branch-weight entry $(\Pi_j,=,q_j)$ is listed for every $j\in J_{\mathrm{w}}$;
\item \emph{(the class-marginal identity)} it is discharged when every
  branch discharges it, and the flag is set in exactly that case;
\item \emph{(the declared model)} the normalizer acts branch by branch, and
  the declared model is
  \begin{equation*}\tag{\getrefnumber{C-eq:register-declared}}\digesteqlabel{eq:register-declared}
    \Mdec_v=\sum_{j\in J}{\Mdec}^{(j)}_v\otimes N_j
    \quad\text{resp.}\quad
    \Mdec_{(v,j)}={\Mdec}^{(j)}_v\otimes N_j ,
  \end{equation*}
  with ${\Mdec}^{(j)}$ the declared model of branch $j$.
\end{enumerate}
\end{digestclaim}

% claim reductions/lem:precomposed-instance
\digestlabel{lem:precomposed-instance}{10.6}\pdfbookmark[2]{Lemma 10.6 Precomposition of a reduction instance with a channel}{digest-reductions/lem:precomposed-instance}
\renewcommand*{\RestatedName}{\hyperref{contract.pdf}{}{claim-reductions/lem:precomposed-instance}{Lemma~10.6}}%
\begin{digestclaim}[Precomposition of a reduction instance with a channel]
Let $\Bsys$ and $\Bsys'$ be systems with separable Hilbert spaces, and let
$\Phi\colon\Tcl{\Bsys}\to\Tcl{\Bsys'}$ be a channel with adjoint
$\Phi^{\dagger}\colon\Lin{\Bsys'}\to\Lin{\Bsys}$ (Lemma~\ref{lem:sep}).
\begin{enumerate}[label=(\roman*),leftmargin=3.0em]
\item \emph{(relations)} If $(M';O')\dsq_{\Lambda}(F;\hat O)$ with
  $\Lambda\colon\Tcl{\Bsys'}\to\Tcl{\Qsys}$, then
  $\Phi^{\dagger}(M')\coloneqq\{\Phi^{\dagger}(M'_v)\}_{v\in\Vset}$ is a POVM
  on $\Bsys$ and
  \begin{equation*}\tag{\getrefnumber{C-eq:precomposed-relation}}\digesteqlabel{eq:precomposed-relation}
    \big(\Phi^{\dagger}(M');\Phi^{\dagger}(O')\big)\;\dsq_{\Lambda\circ\Phi}\;
    \big(F;\hat O\big),\qquad
    (\Lambda\circ\Phi)^{\dagger}=\Phi^{\dagger}\circ\Lambda^{\dagger}.
  \end{equation*}
  If $\Lambda^{\dagger}(\hat O_k)\bowtie c\,\idon{\Bsys'}$ for a relation
  $\bowtie\,\in\{\preceq,=,\succeq\}$ and a real $c$, then
  $(\Lambda\circ\Phi)^{\dagger}(\hat O_k)\bowtie c\,\idon{\Bsys}$, with the
  same relation and the same constant.
\item \emph{(instances)} Let $\mathcal{J}$ be a reduction instance on
  $\Bsys'$ with model map $M[\cdot]$, channel
  $\Lambda[\cdot]=\PinS\circ\Lambda_0[\cdot]$ and pull-backs $O[\cdot]$.
  The \emph{precomposed instance} $\mathcal{J}\circ\Phi$ has the system
  $\Bsys$; the joint alphabet, data sectors, supports,
  handed datum, hypotheses (H), scalar values,
  sector set, target family, transported set, entry list and
  class-marginal flag of $\mathcal{J}$; the model map
  $A'\mapsto\Phi^{\dagger}(M[A'])$; the channel
  $A'\mapsto\Lambda[A']\circ\Phi=\PinS\circ(\Lambda_0[A']\circ\Phi)$; and the
  pull-backs $A'\mapsto\Phi^{\dagger}(O[A'])$. It is a reduction instance
  on $\Bsys$.
\item \emph{(at a datum, and the emission)} For every admissible datum $A'$,
  $\mathcal{J}\circ\Phi$ at $A'$ is $\mathcal{J}$ at $A'$ precomposed with
  $\Phi$. The emission normalizer produces the same declared datum from
  the handed datum of either, so the emitted target, scalar values, entries
  and constants of $\mathcal{J}\circ\Phi$ are those of $\mathcal{J}$; its
  emitted channel is the emitted channel of $\mathcal{J}$ precomposed with
  $\Phi$, and its declared model is $\Phi^{\dagger}$ of the declared model of
  $\mathcal{J}$.
\end{enumerate}
\end{digestclaim}

\section*{11\quad Exact emission}\pdfbookmark[1]{11 Exact emission}{digest-group-reductions/def:emitnorm}

% claim reductions/def:emitnorm
\digestlabel{def:emitnorm}{11.1}\pdfbookmark[2]{Definition 11.1 The emission normalizer, the declared datum and the declared model}{digest-reductions/def:emitnorm}
\renewcommand*{\RestatedName}{\hyperref{contract.pdf}{}{claim-reductions/def:emitnorm}{Definition~11.1}}%
\begin{digestclaim}[The emission normalizer, the declared datum and the declared model]
Let $\mathcal{J}$ be a reduction instance
(Definition~\ref{def:reduction-instance}) with handed datum
$A\in\mathcal{D}_{\mathrm{adm}}$, formed exactly. Those numbers \emph{are} the
data sectors of the emitted target
(Definition~\ref{def:reduction-instance}\,(I2)), so they are this module's own
declaration to its consumer. Where the exact arithmetic that formed
them has left the declared lattice \eqref{eq:latinv}, they
are re-established on it, per data sector $s$ and per class $t$, by the
following steps, the
\emph{emission normalizer}.
The steps act on the elements $A^{(s)}_v$ with
$v\in\Vset_t\cap\Vset^{(s)}$, the elements of the class $t$ in the sector's
support; the elements outside the sector support are exactly zero,
lie on the lattice, and are touched by no step. Every positive
semidefiniteness below is decided exactly by \eqref{eq:psdtest}, on matrices
with entries in $\Fldi$. The
structural support of a matrix is the set of indices $i$ whose row $i$
or column $i$ holds a non-zero entry (\xref{foundations/def:snap}). The
\emph{designation order} ranks the elements of a class in the sector's
support by largest exact trace and, among equal traces, by the
declared order of the class. \emph{A tie is exact and ordinary here}:
the traces compared are exact rationals, and a sector symmetric under a
relabelling of the class ties every element of it, so ``largest trace''
alone names no element; the declared order is the order in which this
contract reads $\Vset$, the outcome varying fastest, and it is emitted with
the artefact, so the order is a function of the declaration and of nothing
else. \emph{Rounding to nearest} a matrix at $\kappa$ rounds every real and
imaginary part of every entry to the nearest integer multiple of
$2^{-\kappa}$, \emph{ties to even}: a part lying exactly halfway between two
multiples goes to the one whose integer numerator is even. \emph{The tie is
exact}: the rounding is carried out in the substrate of
Lemma~\ref{lem:exactpsd} on exact rationals, so a halfway part is an exact
equality rather than a floating-point near-miss. \emph{Ties to even is
odd-symmetric}: negating a part negates both of the multiples nearest it and
their numerators, which preserves evenness, so the rounding commutes with
$X\mapsto X^{\dagger}$ --- mirror entries have equal real parts and opposite
imaginary parts before it and therefore after it --- and a matrix rounded
from a Hermitian one is Hermitian; a zero entry stays zero.
\begin{enumerate}[label=(\arabic*),leftmargin=2.4em]
\item \emph{Keep.} A class $t$ at a sector $s$ every element of which
  satisfies \eqref{eq:latinv} entry by entry is \emph{kept}:
  $A'^{(s)}_v\coloneqq A^{(s)}_v$ for every $v\in\Vset_t\cap\Vset^{(s)}$, and
  no further step acts on it.
\item \emph{Snap for fidelity.} At every other class and sector, for $w=0$
  and then $w=2^{-\kappa},2^{-\kappa+1},\dots,2^{-1}$ in this order, let
  $P^{w}_v$ be $A^{(s)}_v$ with every off-diagonal entry multiplied by
  $1-w$, that is $(1-w)A^{(s)}_v+w\,T_v$ with $T_v$ the diagonal part of
  $A^{(s)}_v$. Every element of the class in the sector's support receives
  $P^{w}_v$ rounded to nearest at $\kappa$, and the \emph{class residue}
  $E^{w}\coloneqq q_t\idon{d_s}-\sum_{v\in\Vset_t\cap\Vset^{(s)}}A'^{(s)}_v$,
  a Hermitian matrix of lattice values, is split by its \emph{components}:
  the classes of the equivalence relation on the indices $i$ whose row $i$ of
  $E^{w}$ is non-zero generated by $i\sim l$ whenever $E^{w}_{il}\ne0$. For a
  component $C$ write $E_C$ for the matrix that agrees with $E^{w}$ on the
  entries whose row and column both lie in $C$ and is zero elsewhere. Each
  $E_C$ is added to the first element $A'^{(s)}_u$ in the designation order
  whose structural support in the handed datum, that of $A^{(s)}_u$,
  contains $C$, its \emph{host}. The first $w$ at which every component has
  a host and every element so formed is decided positive semidefinite gives
  the family, and $w$ is recorded. Where no $w$ does, the class is snapped
  by step~(3).
\item \emph{The last resort.} The support-safe steps of (2) produced no snap
  of the class, and the class is snapped by a uniform lift, applied to the
  class with constant $q_t$ and starting from the handed datum. Let $n$ be
  the number of elements of the class in the sector's support, $v_1$ the
  first of them in the designation order, and, for $w\in(0,1)$,
  \begin{equation*}\tag{\getrefnumber{C-eq:lastresort}}\digesteqlabel{eq:lastresort}
    B^{w}_v\;\coloneqq\;(1-w)\,A^{(s)}_v+w\,\frac{q_t}{n}\,\idon{d_s},
    \qquad v\in\Vset_t\cap\Vset^{(s)},
  \end{equation*}
  a family that is Hermitian, complete at $q_t$, with entries in $\Fldi$, and
  positive definite in every element, since $q_t>0$ at a class that is not
  kept (at $q_t=0$ completeness and positive semidefiniteness make every
  element zero, which is on the lattice). Its designation order is that of
  the class, $\tr B^{w}_v$ being an increasing affine function of
  $\tr A^{(s)}_v$. For $w=2^{-\kappa},2^{-\kappa+1},\dots,2^{-1}$ in this
  order, every element other than $v_1$ receives $B^{w}_v$ rounded to
  nearest at $\kappa$, and $v_1$ receives $q_t\idon{d_s}$ minus their sum;
  the first $w$ at which every element so formed is decided positive
  semidefinite gives the family. Where no $w$ does, every element $v\ne v_1$
  receives the first matrix $A'^{(s)}_v$ of the following list that is
  Hermitian, satisfies \eqref{eq:latinv} entry by entry, and satisfies
  \begin{equation*}\tag{\getrefnumber{C-eq:rounddown}}\digesteqlabel{eq:rounddown}
    0\;\preceq\;A'^{(s)}_v\;\preceq\;B^{1/2}_v :
  \end{equation*}
  $B^{1/2}_v$ itself; then, for $r=\kappa,\kappa-1,\dots,1$ in this order,
  the shift $B^{1/2}_v-2^{-r}\idon{d_s}$ and the scaling
  $(1-2^{-r})B^{1/2}_v$, each rounded to nearest at $\kappa$; finally the
  zero matrix, which always qualifies, $B^{1/2}_v$ being positive
  semidefinite; and $v_1$ receives $q_t\idon{d_s}$ minus their sum. The
  weight $w$, or the fall-back, is recorded, and the class is flagged as one
  whose support is not preserved: an element may then hold a non-zero entry
  in a row or column in which the handed element holds none, that is,
  outside its structural support (\xref{foundations/def:snap}).
\item \emph{All six conditions are decided last}, after (2) or (3), each of
  which changes partial sums from the touched element onward: (a)--(e) of
  Definition~\ref{def:exactadm} together with the lattice invariant
  \eqref{eq:latinv} of every entry, read on the target family
  \eqref{eq:instance-target} the datum enters. They are the six
  Theorem~\ref{thm:emit}\,(ii) decides; a family for which one of them fails
  is refused.
\end{enumerate}
A class at a sector that step~(1) does not keep is \emph{moved}, whichever
of (2) and (3) snapped it. Write $A'$ for the resulting datum --- the
\emph{declared datum},
an element of $\mathcal{D}_{\kappa}$ by (1) or the step of (2) and (3) that
snapped the class, each of which makes it Hermitian, on the lattice and
class complete exactly, positive semidefinite by an exact decision in
(2) and in (3) at a weight of its ladder, and by construction in the fall-back of (3), and
by (4), which decides all six --- and
\begin{equation*}\tag{\getrefnumber{C-eq:declmodel}}\digesteqlabel{eq:declmodel}
  \Mdec\;\coloneqq\;M[A'],
  \qquad\text{that is}\qquad
  \Mdec_v\;=\;M_v[A']\quad(v\in\Vset),
\end{equation*}
for the
\emph{declared model}:
the handed model with the datum replaced by $A'$, as the instance's
model map defines it. Per data sector $s$, those being the sectors the
normalizer touches and the only ones both data have,
\begin{equation*}\tag{\getrefnumber{C-eq:displacement}}\digesteqlabel{eq:displacement}
  \Delta_s\;\coloneqq\;\max_{v\in\Vset}\
  \mxn{A'^{(s)}_v-A^{(s)}_v}
\end{equation*}
is the
\emph{emission displacement};
it is recorded per data sector, together with the kept and the moved classes,
for a moved class the step that snapped it, the weight $w$ of (2) or (3) or
the fall-back of (3), the hosts of the components of (2), and the flag of a
class snapped by (3).

The normalizer has five properties, each of which Theorem~\ref{thm:emit}, or
a claim that instantiates it, uses.
\emph{First, $A'=A$ and every $\Delta_s=0$ whenever the exact construction
already satisfies \eqref{eq:latinv}}: every class is then kept, the
normalizer is the identity and nothing below has an antecedent. That is the
ordinary case at a nominal declaration and is not the general one: every
exact step of a construction multiplies, a product of dyadics adds
significand widths, and a lattice established at the instance is closed
under none of them; that is the case in which a producer re-establishes its
lattice at an emission.
\emph{Second, the normalizer reads data sectors and nothing else.} No object
the hypotheses (H) of the instance are about --- an operator without a finite
presentation, a family on a part of $\Hs{\Bsys}$ the datum does not describe,
a register POVM --- is read, and none is repaired; so (H) holds of $\Mdec$
by the stability clause of
Definition~\ref{def:reduction-instance}\,(I1), $\Mdec$ is a POVM on $\Bsys$
by (I1) at $A'$, and every statement the instance makes about its model under
(H) holds of $\Mdec$. For an instance whose datum is a family of preserved
blocks below a cutoff, this is the statement that the declared model agrees
with the handed one above the cutoff and satisfies every above-cutoff fact the
handed one does.
\emph{Third, what the normalizer produces is a datum, and the family the
emitted identity is about is the model the instance builds from it.} Two
objects, two symbols, and every statement below names which it is on: $A'$ is
the declared datum, a finite list of sector blocks carried entry by entry in
the emitted target; $\Mdec$ is the declared model on $\Bsys$, the object an
adjoint landing on $\Bsys$ can equal (Theorem~\ref{thm:emit}\,(iii)). Where
the model map is the identity on a finite $\Bsys$ the two coincide; where the
model has parts the datum does not describe --- an above-cutoff family,
further branches --- they do not, and the instance's model map says how $A'$
sits inside $\Mdec$. For a branching instance
(Theorem~\ref{thm:instrument-instance}) the normalizer acts branch by branch,
so the declared model is the branches' declared models composed.
\emph{Fourth, this is a snap of what this module declares and never of what it
was handed}, which is the condition under
which a producer may snap at all.
\emph{Fifth, no step revives a zero row or column, save the last resort
(3)}: for every data sector $s$ and every $v\in\Vset$ whose class at $s$ is
not snapped by (3) the structural support of $A'^{(s)}_v$ is contained in
that of $A^{(s)}_v$. In (2) the family $\{T_v\}$ of diagonal parts is
complete at $q_t$, the diagonal of $\sum_vA^{(s)}_v=q_t\idon{d_s}$ being
$q_t$ in every place, and $T_v$ is diagonal, with non-negative entries,
non-zero only at indices of the structural support of $A^{(s)}_v$, so
$P^{w}$ is a mixture with a complete diagonal family that adds only on
coordinates of each element's own structural support; $P^{w}_v$ has the
structural support of $A^{(s)}_v$, and a rounding to nearest keeps a zero
entry zero. Each $E_C$ has its non-zero rows and columns in $C$, which the
host's structural support contains, and the entries of $E^{w}$ between two
components are zero by their definition, so $E^{w}=\sum_CE_C$ and every
non-zero entry of the residue goes to a host. In (2) the residue need not be
positive semidefinite, and the family is taken only where every element is
decided positive semidefinite. No element is mixed with a multiple of the
identity outside (3). In (3) the elements of the class in the sector's
support may gain rows and columns, which is what its flag discloses; the
elements outside the support stay zero there too. And (3) always yields a
family whose elements are positive semidefinite: at a weight of its ladder by decision, and
in the fall-back because every element of $B^{1/2}$ is positive definite, the
rounded elements satisfy \eqref{eq:rounddown}, and
$A'^{(s)}_{v_1}=B^{1/2}_{v_1}+\sum_{v\ne v_1}\big(B^{1/2}_v-A'^{(s)}_v\big)
\succeq B^{1/2}_{v_1}\succ0$ by the completeness of $B^{1/2}$.
\end{digestclaim}

% claim reductions/thm:emit
\digestlabel{thm:emit}{11.2}\pdfbookmark[2]{Theorem 11.2 Exact emission of a reduction instance}{digest-reductions/thm:emit}
\begin{mainclaim}
\renewcommand*{\RestatedName}{\hyperref{contract.pdf}{}{claim-reductions/thm:emit}{Theorem~11.2}}%
\begin{digestclaim}[Exact emission of a reduction instance]
Let the declared data pass every refusal rule of this module's
interface, so that every declared block and scalar it reads is
exactly admissible on the declared lattice
(Definition~\ref{def:exactadm}). Let $\mathcal{J}$ be a reduction instance
(Definition~\ref{def:reduction-instance}) with handed datum $A$, and let $A'$
be the declared datum the normalizer of Definition~\ref{def:emitnorm}
produces from it (where the normalizer refuses, nothing is emitted). Let
$\{F_v\}\coloneqq F[A']$, $\Lambda\coloneqq\Lambda[A']$ and
$\CList\coloneqq\CList[A']$ be the target, the channel and the entry list of
the instance at $A'$, and $\Mdec=M[A']$ the declared model
\eqref{eq:declmodel}. Then:
\begin{enumerate}[label=(\roman*),leftmargin=3.0em]
\item \emph{(everything on the list below is computed exactly, and the
  displacement $\Delta_s$ is the one emitted number that is not on it)} every
  number on the following list is obtained from the numbers the instance was
  handed by finitely many operations in $\Fldi$, floors and comparisons: the
  data-sector values $A'^{(s)}_v$; the scalar values
  $\phi_s(v)$; the sector projectors; and every emitted constraint operator
  and constant of Corollary~\ref{cor:entries}. \emph{That list is everything
  this module emits with exactly one exception}: the emission displacement
  $\Delta_s$ of \eqref{eq:displacement}, an entrywise maximum \emph{modulus}
  and hence the square root of an exact rational, in general irrational, which
  is emitted only as a record of what the emission normalizer did and which
  no statement of this contract reads. \emph{No floating-point
  comparison, no eigenvalue routine, no enclosure and no interval occurs
  anywhere in that list}; the intermediate numbers an instance forms on the
  way to its datum and its scalar values are listed by the corollary that
  instantiates this theorem for it;
\item \emph{(the emitted family is exactly admissible, and is decided so
  before it is emitted)} the family $\{F_v\}$, presented as the sector-block list
  \eqref{eq:instance-target}, satisfies conditions (a)--(e) of
  Definition~\ref{def:exactadm} together with \eqref{eq:latinv} of every
  entry --- six conditions --- and each of the six is \emph{decided}, exactly, on the numbers
  about to be emitted: (a) by comparison of declared labels against the
  instance's sector map; (b) and (d) by entrywise equality and (c) by
  Lemma~\ref{lem:exactpsd}, per sector block; (e) by Lemma~\ref{C-lem:latsum}, whose hypothesis
  \eqref{C-eq:latsumcond} is a declared input condition; and \eqref{eq:latinv}
  entry by entry. \emph{A family for which any of the six cannot be
  established is not emitted}.
  \emph{And the same numbers discharge what the consuming declaration asks of
  a party's side}: the tuple $\big(\Qsys,\SecSet,\{F_v\}_{v\in\Vset}\big)$
  \emph{together with the party partition} of
  Definition~\ref{def:reduction-instance}\,(I2) --- the
  default party partition, every factor Bob's, unless the instance
  declares another party's label, in which case every factor is that party's
  --- satisfies the well-formedness conditions the consuming declaration
  places on \emph{that party's} side (\xref{meat-core/def:admissible-declaration}\,(i) and
  (iii), with the party labels per tensor factor demanded through its (viii))
  ---
  \begin{equation*}\tag{\getrefnumber{C-eq:emitwellformed}}\digesteqlabel{eq:emitwellformed}
    F_v=F_v^{\dagger},\qquad
    F_v\succeq0,\qquad
    \sum_{v\in\Vset}F_v=\idon{\Qsys},
  \end{equation*}
  with every $F_v$ block diagonal for $\SecSet$ and
  $\Hs{\Qsys}=\bigoplus_{s\in\SecSet}\Hs{\Qsys_s}$ finite-dimensional, every
  factor of every $\Hs{\Qsys_s}$ the instance's declared party's and Bob's by
  default. \emph{Those are conditions on a party partition and not on a
  defect}, and they are a \emph{different} condition from the conditions of
  \xref{foundations/def:exact-admissibility} that (a)--(e) discharge, which
  say nothing about who owns which factor;
\item \emph{(the defining identity is about the numbers emitted, and it is an
  identity on $\Bsys$)} for every $v\in\Vset$,
  \begin{equation*}\tag{\getrefnumber{C-eq:emitidentity}}\digesteqlabel{eq:emitidentity}
    \Lambda^{\dagger}(F_v)\;=\;\Mdec_v
    \qquad\text{exactly},
  \end{equation*}
  \emph{both sides operators on $\Bsys$}, which is where $\Lambda^{\dagger}$
  lands, $\Lambda$ being a channel $\Tcl{\Bsys}\to\Tcl{\Qsys}$. \emph{Only the
  data sectors carry anything the normalizer produced}: $\Mdec$ is the handed
  model with the datum replaced by $A'$ as the instance's model map defines
  it, and everything the hypotheses (H) are about is the declared one,
  unchanged. Moreover every claim of the instance holds verbatim of the
  declared model --- the relation \eqref{eq:instance-relation} with its
  transported set and pull-backs, every certificate of
  \eqref{eq:instance-entry}, and the class-marginal identity
  \eqref{eq:instance-classmarginal} where the instance discharges it --- and
  Theorem~\ref{thm:transfer} and Theorem~\ref{thm:chain} hold of the one-link
  receiver segment (Definition~\ref{def:link}) whose
  source family is $(\Bsys,\Mdec)$, whose target is $(\Qsys,F)$ and
  whose channel witness is $\Lambda$, with
  $\textsf{status}=\textsf{exact}$;
\item \emph{(the emitted bounds, and the direction each is rounded in)} every
  scalar value that was rounded to reach the lattice was rounded
  \emph{down}, as the lower bound of the quantity it certifies, and every
  complementary value was \emph{recomputed} from the rounded ones and never
  rounded on its own (Definition~\ref{def:reduction-instance}\,(I2)); and a
  caller-declared constant of modulus at most one off the lattice is rounded
  in the direction that enlarges the feasible set for its declared relation
  (Corollary~\ref{cor:entries}\,(iii)), an equality constant off the lattice
  and any constant of modulus above one being refused. \emph{Every other
  number this module emits is exact, the displacement record $\Delta_s$ of
  (i) aside}; these are all the rounding directions there are among the
  numbers a statement of this contract reads, and the direction of each is
  recorded;
\item \emph{(what relates the declared model to the handed one)} where some
  $\Delta_s>0$ the emitted numbers are a declaration of a family near the
  constructed one, and nearness is not what relates them: \emph{no bound
  relates a key length computed from $\Mdec$ alone to one computed from the
  handed model $M[A]$, or from the supplier's own output}, and this contract
  has no theorem converting an operator perturbation of the target into a
  security penalty, nor any other relation between $A'$ and $A$ on a moved
  class: the consuming contract's impossibility
  (\xref{meat-core/prop:optol-nobound}) is the whole of what is proved
  there; what a run rests on there instead is~(vi);
\item \emph{(what a run on the emission rests on where it reads a moved
  block, and only there)} let $B=B_{\mathrm{snap}}\sqcup B_{\mathrm{em}}$
  be the blocks that the record of the run names, for this artefact, as
  resting on \xref{foundations/ass:snapped-measurement}, kept as two
  disjoint parts, since they are blocks of two different snaps and their
  labels may coincide: $B_{\mathrm{em}}$ holds the data sectors of the
  emission so named, and $B_{\mathrm{snap}}$, where $A$ is formed from
  measurement blocks of a snap (\xref{foundations/def:snap}) that the
  declaring party made of the measurement it states, the blocks of that
  snap so named; the declaring party is the instance that states the
  measurement of the device, or a reference of it, from which $A$ is formed.
  Where $B$ is empty, this part asserts nothing, and no part of this theorem
  rests on that assumption. Where $B$ is not empty, let that assumption
  hold. Then on every block in $B_{\mathrm{snap}}$, the block of the
  measurement the device performs is the snapped block, with the entries
  recorded for that snap, from which $A$ is formed, in place of the block of
  the model the declaring party states; and every conclusion about the
  device that the claims of the run draw from the emitted target, under
  their hypothesis that its data sectors carry the handed datum $A$, holds
  when they carry the declared datum $A'$ on every data sector in
  $B_{\mathrm{em}}$, which by~(iii) is the datum whose model $\Mdec$ the
  emitted target squashes exactly. The run carries no bound relating its key
  length to the declaring party's model on the blocks in
  $B_{\mathrm{snap}}$, nor to $M[A]$ on the data sectors in
  $B_{\mathrm{em}}$: no claim relates that model to the snapped block on a
  block in $B_{\mathrm{snap}}$, nor $A$ to $A'$ on a data sector in
  $B_{\mathrm{em}}$.
\end{enumerate}
\end{digestclaim}
\end{mainclaim}

% claim reductions/cor:entries
\digestlabel{cor:entries}{11.3}\pdfbookmark[2]{Corollary 11.3 The emitted constraint entries}{digest-reductions/cor:entries}
\renewcommand*{\RestatedName}{\hyperref{contract.pdf}{}{claim-reductions/cor:entries}{Corollary~11.3}}%
\begin{digestclaim}[The emitted constraint entries]
Assume the hypotheses of Theorem~\ref{thm:emit} for a reduction instance
$\mathcal{J}$ with declared datum $A'$, target $\{F_v\}$, channel $\Lambda$
and declared model $\Mdec$. A request \emph{carries entries} when it selects
the case in which the instance's own entries are emitted: in the flag-state
submodule that case is the licensed route; in a core-only artefact it is a
caller's request for an entry --- for a default carried operator as an entry,
or for the class-marginal entry. \emph{What is emitted is decided by the
request}: \emph{where the request carries no entry the emitted constraint
list is empty} --- whatever entries the instance's own list $\CList[A']$
contains --- and the artefact carries no attack-set licence; \emph{where the
request carries entries the emitted list is the instance's own entries, the
class-marginal entry on request and any further requested entry that (iii)
admits}, so that
the list is non-empty and the licence of Theorem~\ref{thm:licence} is
present. In \emph{both} cases the
certificate $(\Pi_s,\succeq,0)$ of every sector projector of every instance
is carried \emph{per entry in the transported set}, with the certificate case
of Corollary~\ref{cor:pullback} that makes each usable, and not as a
constraint entry. Then
\emph{(i) and the identity \eqref{eq:entryclass} of (ii) of every request,
and what each of (ii) and (iii) emits of a request that carries entries}:
\begin{enumerate}[label=(\roman*),leftmargin=3.0em]
\item \emph{(the default carried operators)} every operator of the
  instance's \emph{default carried set} --- every sector projector $\Pi_s$,
  $s\in\SecSet$, which Definition~\ref{def:reduction-instance}\,(I3) puts in
  the transported set of every instance, together with the sums of sector
  projectors that instance transports besides (for a target split at a
  cutoff, the projector onto the blocks below it and the projector onto the
  flag blocks; for a branching instance, the branch projector $\Pi_j$) ---
  is emitted at $c=0$ with $\bowtie\,=\,\succeq$: its entries
  are $0$ and $1$, hence on the lattice at every $\kappa\ge1$, and its
  pull-back
  \begin{equation*}\tag{\getrefnumber{C-eq:entrysector}}\digesteqlabel{eq:entrysector}
    \Lambda^{\dagger}(\Pi_s)\;\succeq\;0
  \end{equation*}
  is exactly positive for every instance, as therefore is the pull-back of
  every sum of them. \emph{It is emitted with the transported set}, per entry
  and whatever the request carries, as the certificate case of
  Corollary~\ref{cor:pullback}\,(ii) that entry records, and it becomes an
  entry of $\CList$ only where a caller requests it \emph{as an entry} ---
  a request that makes the emitted list non-empty and so brings the licence
  with it;
\item \emph{(the class-marginal entry)} For $\hat O=\sum_{v\in\Vset_t}F_v$ at
  a declared $t\in\Tset$ --- the expectation form of the physical
  class-marginal identity \eqref{eq:classmarginal} ---
  \begin{equation*}\tag{\getrefnumber{C-eq:entryclass}}\digesteqlabel{eq:entryclass}
    \Lambda^{\dagger}\big(\hat O\big)\;=\;q_t\,\idon{\Bsys}
    \qquad\text{exactly}
  \end{equation*}
  holds \emph{whenever} the instance discharges the class-marginal identity
  \eqref{eq:instance-classmarginal}, and is established by this contract
  \emph{only} there --- \emph{a fact about the instance and the channel, which
  this part asserts of every request}. \emph{Where the request carries
  entries}, so that the emitted list is non-empty and the licence is present,
  the entry is then emitted, on request, as the single equality entry
  $(\hat O,=,q_t)$ at the declared $q_t$; where the instance does not
  discharge it the request is refused, and the entry is never emitted
  on the strength of a branch that does not establish it;
\item \emph{(any further entry --- the general rule)} \emph{under the same
  hypothesis as (ii)'s emission sentence --- the request carries entries},
  let
  $\hat O=\sum_{v\in\Vset}\sum_{s\in\SecSet}\alpha_{v,s}\Pi_sF_v\Pi_s+Z$ with
  declared real lattice coefficients $\alpha_{v,s}$ and $Z$ a fixed operator
  of the declared sector decomposition (a sector projector or a combination
  of them whose coefficients are declared lattice values, as
  Definition~\ref{def:reduction-instance}\,(I4) requires), the form
  \eqref{eq:instance-entry}, with a relation
  $\bowtie$ and a declared real constant $c$.
  \emph{Where a declared constant of modulus at most one does not satisfy
  \eqref{eq:latinv} it is first rounded onto the lattice in the direction that
  enlarges the feasible set for $\bowtie$} --- down for $\succeq$, up for
  $\preceq$ --- and \emph{an equality entry whose constant is off the lattice
  is refused rather than split into a pair}, an equality having no loosening
  direction; \emph{a declared constant of modulus above one is refused under
  every relation}. Suppose the certificate
  \begin{equation*}\tag{\getrefnumber{C-eq:entrygeneral}}\digesteqlabel{eq:entrygeneral}
    \Lambda^{\dagger}(\hat O)\ \bowtie\ c\,\idon{\Bsys}
  \end{equation*}
  is established for the emitted numbers, at the constant $c$ so obtained, in
  one of two ways: for an entry of the instance's own list $\CList[A']$, by
  (I4) of the instance; for a requested entry, by a \emph{decided
  certificate} --- a condition on the emitted numbers and the request's
  declared constants, decided exactly from them, that a claim stated for the
  instance proves sufficient for \eqref{eq:entrygeneral} under the instance's
  hypotheses (H) --- that holds there. Then the entry is emitted at that $c$.
  \emph{No other requested entry is emitted}: a requested operator of this
  form for which no claim stated for the instance gives a decided
  certificate, or whose decided certificate fails at the emitted numbers, is
  refused, and no certificate is taken from the request on trust.
\end{enumerate}
Every entry that (i)--(iii) emit is \emph{untagged}, $(\hat O,\bowtie,c)$,
and is read in the form of the declared source mode
(Definition~\ref{def:licensed}).

\emph{Each certificate above is discharged for the emitted numbers
themselves}, those numbers being the family every part is read on
(Theorem~\ref{thm:emit}\,(iii)), so no constant is shifted to pay for a
displacement between the family a certificate is asserted of and the numbers
its operator is assembled from. The rounding rule of (iii) applies only to a
constant a caller declares, never to one this module or an instance
computes. \emph{Together, (i)--(iii) cover every entry this module can
emit}: the default carried operators of every instance, the class-marginal
entry a caller may request, and every entry of the form
\eqref{eq:entrygeneral} --- which includes every entry an instance carries in
its own list, such as a flag-sector entry of a weight-preserving squasher,
and the branch entries a branching instance inherits
(Theorem~\ref{thm:instrument-instance},
Corollary~\ref{cor:register-instance}). \emph{A requested operator outside
that form, or of that form without a decided certificate that holds, is not
emitted at some other constant; it is refused}. So the
three parts are
exhaustive over what is \emph{emitted}, which is what the artefact carries.
\emph{Where the request carries no entry there is nothing for (ii)--(iii) to
be exhaustive about}: $\CList=\varnothing$ there, whatever entries the
instance's own list contains, and the certificate of every default carried
operator travels with the transported set alone.
\end{digestclaim}

% claim reductions/thm:emitted-product
\digestlabel{thm:emitted-product}{11.5}\pdfbookmark[2]{Theorem 11.5 The emission of several measuring parties}{digest-reductions/thm:emitted-product}
\renewcommand*{\RestatedName}{\hyperref{contract.pdf}{}{claim-reductions/thm:emitted-product}{Theorem~11.5}}%
\begin{digestclaim}[The emission of several measuring parties]
Let $\mathbb{F}$ be a finite ordered set of \emph{factors}, each $f$ with a
party label, and for each $f$ let $\mathcal{J}^{f}$ be a
reduction instance (Definition~\ref{def:reduction-instance}) on a
system $\Bsys_f$ over
$\big(\Vset^{f},\{\Vset^{f}_{t}\}_{t\in\Tset^{f}},\{q^{f}_{t}\}\big)$ to
which Theorem~\ref{thm:emit} applies, with emitted datum $A'_f$, emitted
target $F^{f}\coloneqq F^{f}[A'_f]$ on
$\Qsys_f=\bigoplus_{s\in\SecSet^{f}}\Qsys_{f,s}$, channel
$\Lambda_f\coloneqq\Lambda_f[A'_f]$, declared model ${M'}^{\mathrm{ph},f}=M^{f}[A'_f]$
and entry list $\CList^{f}$. Put $\Bsys\coloneqq\bigotimes_f\Bsys_f$,
$\Qsys\coloneqq\bigotimes_f\Qsys_f$, $\Vset\coloneqq\prod_f\Vset^{f}$,
$\Lambda\coloneqq\bigotimes_f\Lambda_f$, all products in the order of
$\mathbb{F}$, and for $v=(v_f)_f\in\Vset$
\begin{equation*}\tag{\getrefnumber{C-eq:producttarget}}\digesteqlabel{eq:producttarget}
  F_v\coloneqq\bigotimes_{f}F^{f}_{v_f},\qquad
  \Mdec_v\coloneqq\bigotimes_{f}{M'}^{\mathrm{ph},f}_{v_f} .
\end{equation*}
Suppose that \emph{what is emitted} is the ordered list of the factors'
artefacts, each with its party label, and the rule \eqref{eq:producttarget},
and that no entry of any $F_v$ is emitted, declared or rounded. Then:
\begin{enumerate}[label=(\roman*),leftmargin=3.0em]
\item \emph{(a measurement, block diagonal for the product sector map)}
  $\{F_v\}_{v\in\Vset}$ is a POVM on $\Qsys$: every $F_v$ is Hermitian and
  positive semidefinite and $\sum_{v\in\Vset}F_v=\idon{\Qsys}$, as identities
  of the exact rational numbers the factors emitted. With the \emph{product
  sector map} --- index set $\prod_f\SecSet^{f}$ in the product order, the
  last factor fastest, identifiers the factors' identifiers joined in that
  order, dimensions the products, projectors
  $\Pi_{s}=\bigotimes_f\Pi^{f}_{s_f}$ --- every $F_v$ is block diagonal, its
  block at $s=(s_f)_f$ being the Kronecker product
  $\bigotimes_fF^{f,(s_f)}_{v_f}$ of the factors' blocks. The party
  partition assigns every tensor factor of $\Hs{\Qsys_{f,s_f}}$ to the party
  of $f$.
\item \emph{(the squashing relation and the declared model)}
  $\Lambda^{\dagger}(F_v)=\Mdec_v$ for every $v$, exactly, so
  $\Mdec\dsq_{\Lambda}F$; the declared joint model is the \emph{product} of
  the factors' declared models. $\Lambda$ is of the form the artefact
  records for one factor: with $\PinS$ the pinching for the product
  projectors, $\PinS\circ\Lambda=\Lambda$, and the squashed
  analysis may be restricted to states block diagonal for the product sector
  map without excluding a reachable one. The transported set contains every
  product projector $\Pi_s$, every lifted projector
  $\idon{}\otimes\Pi^{f}_{s_f}\otimes\idon{}$ and every lifted transported
  operator $\idon{}\otimes\hat O\otimes\idon{}$, $\hat O$ transported by
  $\mathcal{J}^{f}$, with pull-back the lift of $\hat O$'s pull-back.
\item \emph{(entries keep their constants)} for every factor $f$ and every
  entry $(\hat O_k,\bowtie_k,c_k)\in\CList^{f}$, the certificate
  $\Lambda_f^{\dagger}(\hat O_k)\bowtie_kc_k\,\idon{\Bsys_f}$ which
  $\mathcal{J}^{f}$ discharges at $A'_f$ is a certificate on that factor's
  channel in the sense of Theorem~\ref{thm:licence}\,(v). So the joint list
  $\CList\coloneqq\bigsqcup_f\{(\hat O_k,\bowtie_k,c_k;f)\}$, each entry
  scoped to its factor and carrying the same operator, the same relation and
  the same constant, is licensed with its entries read party-locally
  (Definition~\ref{def:licensed}), or as scalars on the lifted operators
  $\idon{}\otimes\hat O_k\otimes\idon{}$, or any mixture; every squashed
  attack sequence against $\Mdec$ is admissible for it and reproduces
  the protocol's output state; and the lifted operators are block diagonal
  for the product sector map with entries that are the factor's emitted
  entries and zeros and ones. No constant is re-derived and none is rounded.
\item \emph{(class-marginal identities)} with classes
  $\Vset_{t}=\prod_f\Vset^{f}_{t_f}$ and weights $q_t=\prod_fq^{f}_{t_f}$,
  $\sum_{v\in\Vset_t}\Mdec_v=q_t\,\idon{\Bsys}$ holds for every $t$ if every
  factor's class-marginal identity is discharged. Where it is
  discharged for the factors in a subset $\mathbb{G}\subseteq\mathbb{F}$ only,
  the identity holds for the classes that are full in every factor outside
  $\mathbb{G}$:
  $\sum_{v:\,v_f\in\Vset^{f}_{t_f}\,(f\in\mathbb{G})}\Mdec_v
  =\prod_{f\in\mathbb{G}}q^{f}_{t_f}\,\idon{\Bsys}$. What is emitted for it
  is the expectation-form
  entry on the pulled-back class sum and never an operator identity on
  $\Qsys$. The weight $q_t$ is a product of declared numbers and is formed
  exactly by whoever reads it; it is not emitted as a number of this module,
  and it could not be: each $q^{f}_{t_f}$ lies on its factor's declared
  lattice (Definition~\ref{def:exactadm}\,(b),
  Definition~\ref{def:joint-family}\,(a)), and a product of two lattice
  numbers lies in general on a finer lattice only and need not be a binary64
  number, exactly as for the blocks of \eqref{eq:producttarget}.
\item \emph{(what was rounded)} the numbers behind \eqref{eq:producttarget}
  are the factors' emitted numbers and nothing else. Every displacement,
  kept or moved class and rounding direction is a factor's and is
  recorded with that factor's emission,
  and Theorem~\ref{thm:emit}\,(v) applies to each factor as it stands: no
  bound relates a key length computed from $\Mdec$ alone to one computed
  from the product of the handed models. The product introduces no further
  displacement.
\item \emph{(one factor)} for $|\mathbb{F}|=1$ parts (i), (ii), (iv) and (v)
  are the corresponding statements of Theorem~\ref{thm:emit} for that
  instance, and what is emitted is that instance's artefact, every entry
  untagged. Part (iii) at one factor is Theorem~\ref{thm:licence}\,(v) at one
  factor, which says \emph{more} than (i)--(iv) of that theorem in
  untrusted-source mode with $d_A>1$ and the same otherwise
  (Definition~\ref{def:licensed}); the untagged entry that is emitted is
  implied by it (Lemma~\ref{lem:party-local}\,(iii)), so the one-factor
  emission asserts nothing beyond Theorem~\ref{thm:licence}\,(i)--(iv).
\end{enumerate}
These statements are about the factors' instance links and their emitted
families.
\end{digestclaim}

% claim reductions/cor:emitted-identity
\digestlabel{cor:emitted-identity}{11.6}\pdfbookmark[2]{Corollary 11.6 The emitted family of the identity instance}{digest-reductions/cor:emitted-identity}
\renewcommand*{\RestatedName}{\hyperref{contract.pdf}{}{claim-reductions/cor:emitted-identity}{Corollary~11.6}}%
\begin{digestclaim}[The emitted family of the identity instance]
Let the declared data pass every refusal rule of this module's interface, let
$\mathcal{J}_{\mathrm{id}}$ be the identity instance of
Corollary~\ref{cor:identity-instance} with handed datum $A$, and let $A'$ be
the datum the normalizer of Definition~\ref{def:emitnorm} declares. Then
Theorem~\ref{thm:emit} holds of $\mathcal{J}_{\mathrm{id}}$; the normalizer
moves nothing, $A'=A$ and $\Delta_s=0$ for every data sector $s$; and, for
every $v\in\Vset$,
\begin{equation*}\tag{\getrefnumber{C-eq:emitted-identity}}\digesteqlabel{eq:emitted-identity}
  \Lambda^{\dagger}(F_v)\;=\;\PinS(F_v)\;=\;F_v\;=\;\Mdec_v
  \;=\;\bigoplus_{s\in\SecSet_0}A'^{(s)}_v
  \qquad\text{exactly, on }\Bsys=\Qsys ,
\end{equation*}
so the emitted family \emph{is} the declared datum, which is the handed
datum, exactly admissible and decided so before emission; the emitted
constraint list is \emph{empty where the request carries no entry} --- this
instance's own list being empty in any case
(Corollary~\ref{cor:identity-instance}\,(I4)) --- and where the request
carries entries it carries, on request, the class-marginal entries
$\big(\sum_{v\in\Vset_t}F_v,=,q_t\big)$, which the instance discharges, and
a further requested entry $(\hat O,\bowtie,c)$ of the form
\eqref{eq:instance-entry}, which Corollary~\ref{cor:entries}\,(iii) emits
exactly where $\hat O\bowtie c\,\idon{\Bsys}$ holds at the constant $c$ it
emits, this instance carrying none of its own: $\Lambda^{\dagger}(\hat O)=\hat O$, so the
certificate \eqref{eq:entrygeneral} of such an entry is that operator
inequality on the emitted numbers, decided exactly by finitely many field
operations and sign comparisons, a decided certificate in the sense of
Corollary~\ref{cor:entries}\,(iii) for every operator of that form; its
default carried operators --- here the sector projectors and nothing else,
the transported set being $\{\Pi_s\}_{s\in\SecSet_0}$ --- are emitted with
the transported set at $c=0$ with $\succeq$ in either case
(Corollary~\ref{cor:entries}\,(i)); no scalar value is emitted and no bound
of this module's own is rounded; the chain is the identity link with
$\textsf{status}=\textsf{exact}$. These statements are about the instance's
link and its emitted family.
The instance is selectable only once a declared input carries an
already-finite POVM.
\end{digestclaim}

\section*{12\quad The reference companion}\pdfbookmark[1]{12 The reference companion}{digest-group-reductions/cor:reference-companion}

% claim reductions/cor:reference-companion
\digestlabel{cor:reference-companion}{12.1}\pdfbookmark[2]{Corollary 12.1 The reference companion of a reduction instance}{digest-reductions/cor:reference-companion}
\begin{mainclaim}
\renewcommand*{\RestatedName}{\hyperref{contract.pdf}{}{claim-reductions/cor:reference-companion}{Corollary~12.1}}%
\begin{digestclaim}[The reference companion of a reduction instance]
Let $\mathcal{J}$ be a reduction instance
(Definition~\ref{def:reduction-instance}) whose handed datum $A$ is
formed, by the rule of its part (I1), from a family $K=\{K^{(b)}_j\}$ of
measurement blocks, over blocks $b$ of a sector map and outcomes $j$ of an
alphabet (one such family per setting where the blocks are given per
setting), that a declaring instance states as a snap
(\xref{foundations/def:snap}) of the measurement of a device. Let
$G=\{G^{(b)}_j\}$ be a family over the same blocks, alphabet and dimensions,
every block of which is Hermitian and positive semidefinite and satisfies
exactly every completeness identity that the declaration of $K$ states of
$K$ on that block. In the use this corollary is for, $G^{(b)}$ is the block
of the handed reference (\xref{foundations/def:evaluated-off-snap})
where the declaring instance hands one, and $K^{(b)}$ elsewhere, so that
$G=K$ where it hands none. Suppose
\begin{itemize}[leftmargin=4.2em]
\item[(H-REF)] $\mathcal{J}$ is a reduction instance over the substrate
  $\Fld=\mathbb{R}$ (Definition~\ref{def:substrate}), and the rule of its
  part (I1) forms, from every family with the properties just listed for $G$,
  a datum in $\mathcal{D}_{\mathrm{adm}}$.
\end{itemize}
Let $A^{\mathrm{ref}}$ be the datum that rule forms from $G$ in place of
$K$, and call the instance at $A^{\mathrm{ref}}$ --- the same $\Bsys$, (H),
sectors, supports, scalar values, constructor and entry rule, with
$A^{\mathrm{ref}}$ in place of the handed datum --- the
\emph{reference companion} of $\mathcal{J}$ at $G$. Then:
\begin{enumerate}[label=(\roman*),leftmargin=3.0em]
\item \emph{(an admissible target over the emitted shape)}
  $A^{\mathrm{ref}}\in\mathcal{D}_{\mathrm{adm}}$, the reference companion
  is a reduction instance, and its target $\{F_v[A^{\mathrm{ref}}]\}_{v\in\Vset}$
  is a POVM on $\Qsys$, sector-block-diagonal by construction. Its alphabet
  with classes and weights, its sector map (the sector list, its order, each
  sector's dimension and kind), its scalar values and its party partition are
  those of the target the instance has at every admissible datum, the
  declared datum the normalizer produces included, and it differs
  from such a target on the data sectors alone, where it
  carries $A^{\mathrm{ref}}$. Where $G=K$, $A^{\mathrm{ref}}=A$: the
  reference companion is $\mathcal{J}$ at its handed datum;
\item \emph{(an exact squash of the companion's model)}
  $\big(M[A^{\mathrm{ref}}];O[A^{\mathrm{ref}}]\big)\dsq_{\Lambda[A^{\mathrm{ref}}]}
  \big(F[A^{\mathrm{ref}}];\hat O[A^{\mathrm{ref}}]\big)$, so that
  $\Lambda[A^{\mathrm{ref}}]^{\dagger}(F_v[A^{\mathrm{ref}}])=M_v[A^{\mathrm{ref}}]$
  exactly for every $v$; every certificate of the entry list
  $\CList[A^{\mathrm{ref}}]$ is discharged, and the class-marginal identity
  \eqref{eq:instance-classmarginal} holds at $A^{\mathrm{ref}}$ where the
  instance's flag says it is discharged; and Theorem~\ref{thm:transfer} and
  Theorem~\ref{thm:chain} hold of the one-link receiver segment
  (Definition~\ref{def:link}) whose source family is
  $(\Bsys,M[A^{\mathrm{ref}}])$, whose target is
  $(\Qsys,F[A^{\mathrm{ref}}])$ and whose channel witness is
  $\Lambda[A^{\mathrm{ref}}]$, with $\textsf{status}=\textsf{exact}$;
\item \emph{(the companion is known through enclosures)} let every block of
  $G$ be read with an entrywise enclosure
  (\xref{foundations/def:exact-reference}), with radius zero on every block
  where $G^{(b)}=K^{(b)}$. Where the rule of (I1) forms each entry of the
  datum as a finite expression in entries of the blocks and in numbers the
  instance was handed, and every reciprocal node of that expression receives
  an enclosure by part (iv) of \xref{receiver/lem:enclosure-arithmetic}, as
  it does where the expression has none, the rules of that lemma applied to
  that expression give a complex enclosure of every element of
  $A^{\mathrm{ref}}$ whose centre is formed by the same expression from the
  centres and whose radius is zero at every entry whose inputs are all
  enclosed with radius zero; part (v) of that lemma, applied element by
  element, makes it an entrywise enclosure of $A^{\mathrm{ref}}$, whose
  centre is the Hermitian part of that centre and whose radius is zero at an
  entry $(i,l)$ where the radii at $(i,l)$ and $(l,i)$ both are. Every entry operator
  $\hat O_k[A^{\mathrm{ref}}]$ of \eqref{eq:instance-entry}, a combination
  of the blocks $\Pi_sF_v[A^{\mathrm{ref}}]\Pi_s$ with real coefficients
  $\alpha_{k,v,s}$ plus a combination $Z_k$ of sector projectors, is then
  enclosed entrywise by the same combination of the centres, with radius
  $\sum_{v}|\alpha_{k,v,s}|\,R^{(s)}_{v,il}$ at the entry $(i,l)$ of each
  data sector $s$, where $(C,R)$ is any entrywise enclosure of the datum and
  $R^{(s)}_{v,il}$ is the radius of its element $v$ at that entry, and
  radius zero on every scalar sector;
\item \emph{(where (H-REF) holds)} (H-REF) holds for the identity instance of
  Corollary~\ref{cor:identity-instance}; and it holds for the branching
  instance of Theorem~\ref{thm:instrument-instance}, and for the register
  instance of Corollary~\ref{cor:register-instance}, whenever it holds for
  every branch instance whose datum is formed from blocks of $K$, the data of
  the other branches being formed from no block of $K$;
\item \emph{(what the companion rests on, and only there)} no part of
  (i)--(iv) rests on \xref{foundations/ass:snapped-measurement}. Let $B$ be
  the blocks on which $G^{(b)}=K^{(b)}$ that the record of the run names as
  resting on that assumption. Where $B$ is empty this part asserts nothing.
  Where it is not, let the assumption hold: then on every block in $B$ the
  block of the measurement the device performs is the snapped block
  $K^{(b)}$, with the entries recorded for the snap, in place of the block
  of the declaring instance's model, so that the companion is formed there
  from the block the device then measures; and no claim relates a key length
  computed from the companion to that model on the blocks in $B$;
\item \emph{(entries requested beside the instance's own)} let
  $(\hat O,\bowtie,c)$ be an entry emitted at a request that carries entries
  (Corollary~\ref{cor:entries}), written in the form
  \eqref{eq:instance-entry} with coefficients $\alpha_{v,s}$ and a
  combination $Z$ of sector projectors,
  $\hat O=\sum_{v\in\Vset}\sum_{s\in\SecSet}\alpha_{v,s}\Pi_sF_v[A']\Pi_s+Z$
  at the declared datum $A'$, and call
  $(\hat O^{\mathrm{ref}},\bowtie,c)$, with
  \begin{equation*}\tag{\getrefnumber{C-eq:companion-form}}\digesteqlabel{eq:companion-form}
    \hat O^{\mathrm{ref}}\;\coloneqq\;\sum_{v\in\Vset}\sum_{s\in\SecSet}
      \alpha_{v,s}\,\Pi_sF_v[A^{\mathrm{ref}}]\Pi_s\;+\;Z ,
  \end{equation*}
  its \emph{companion form} at that writing. Part~(iii) encloses
  $\hat O^{\mathrm{ref}}$. Where $\hat O$ is a sum of sector projectors
  (Corollary~\ref{cor:entries}\,(i)), written with every $\alpha_{v,s}=0$,
  $\hat O^{\mathrm{ref}}=\hat O$ and
  $\Lambda[A^{\mathrm{ref}}]^{\dagger}(\hat O)\succeq0$. Where it is the
  class-marginal entry $\big(\sum_{v\in\Vset_t}F_v,=,q_t\big)$ of
  Corollary~\ref{cor:entries}\,(ii), written with $\alpha_{v,s}=1$ for
  $v\in\Vset_t$, $0$ otherwise, and $Z=0$, $\hat O^{\mathrm{ref}}=\hat O$,
  and where the instance discharges \eqref{eq:instance-classmarginal},
  $\Lambda[A^{\mathrm{ref}}]^{\dagger}(\hat O^{\mathrm{ref}})=q_t\,\idon{\Bsys}$
  exactly. For any other entry this part asserts no certificate at
  $A^{\mathrm{ref}}$. In every case, where
  $\Lambda[A^{\mathrm{ref}}]^{\dagger}(\hat O^{\mathrm{ref}})\bowtie
  c\,\idon{\Bsys}$ holds,
  $\big(M[A^{\mathrm{ref}}];O[A^{\mathrm{ref}}],
  \Lambda[A^{\mathrm{ref}}]^{\dagger}(\hat O^{\mathrm{ref}})\big)
  \dsq_{\Lambda[A^{\mathrm{ref}}]}
  \big(F[A^{\mathrm{ref}}];\hat O[A^{\mathrm{ref}}],\hat O^{\mathrm{ref}}\big)$,
  and the companion form is an entry carrying the spectral certificate
  \eqref{eq:certificate} at its constant.
\end{enumerate}
\end{digestclaim}
\end{mainclaim}

% claim reductions/cor:companion-factor-enclosure
\digestlabel{cor:companion-factor-enclosure}{12.2}\pdfbookmark[2]{Corollary 12.2 Factor enclosures carried to the reference companion}{digest-reductions/cor:companion-factor-enclosure}
\renewcommand*{\RestatedName}{\hyperref{contract.pdf}{}{claim-reductions/cor:companion-factor-enclosure}{Corollary~12.2}}%
\begin{digestclaim}[Factor enclosures carried to the reference companion]
Let $\mathcal{J}$, $K$, $G$ and $A^{\mathrm{ref}}$ be as in
Corollary~\ref{cor:reference-companion}, the blocks $K^{(b)}_j$ indexed by a
block $b$ and an outcome $j$, together with the setting where the blocks are
given per setting. Call the rule of part (I1) of $\mathcal{J}$
(Definition~\ref{def:reduction-instance}) a \emph{non-negative combination}
when for every data sector $s$ and every $v\in\Vset$ there are
non-negative rationals $\lambda^{(s,v)}_{b,j}$ such that, for every family
$K'$ with the properties that Corollary~\ref{cor:reference-companion} lists
for $G$, the datum $A'$ the rule forms from $K'$ satisfies
\begin{equation*}\tag{\getrefnumber{C-eq:datum-combination}}\digesteqlabel{eq:datum-combination}
  A'^{(s)}_v\;=\;\sum_{b,j}\lambda^{(s,v)}_{b,j}\,K'^{(b)}_j ,
\end{equation*}
every block $K'^{(b)}_j$ with $\lambda^{(s,v)}_{b,j}>0$ being of dimension
$d_s$. Let the rule be a non-negative combination, and let some blocks
$G^{(b)}_j$ carry factor enclosures
$\big(L^{(b)}_{j},\varrho^{(b)}_{j},W^{(b)}_{j}\big)$
(\xref{foundations/def:factor-enclosure}).
\begin{enumerate}[label=(\roman*),leftmargin=3.0em]
\item Let $s$ be a data sector and $v\in\Vset$ such that every block
  $G^{(b)}_j$ with $\lambda^{(s,v)}_{b,j}>0$ carries one, and list the pairs
  $(b,j)$ with $\lambda^{(s,v)}_{b,j}>0$ in a fixed order. The triple whose
  matrix box has as its centre and its radii those of the boxes
  $\big(L^{(b)}_{j},\varrho^{(b)}_{j}\big)$ side by side in that order, and
  whose core is the diagonal matrix of the
  $\lambda^{(s,v)}_{b,j}W^{(b)}_{j}$ in that order, is a factor enclosure of
  $A^{\mathrm{ref},(s)}_v$, its number of columns the sum of theirs; where
  every $\lambda^{(s,v)}_{b,j}$ is zero, it has no column.
\item The block of the companion's target $F_v[A^{\mathrm{ref}}]$ on the data
  sector $s$ is $A^{\mathrm{ref},(s)}_v$, so the triple of (i) is a factor
  enclosure of that block.
\item The rule of (I1) of the identity instance of
  Corollary~\ref{cor:identity-instance} is a non-negative combination: each
  element of its datum is the block of the family it is read from at the
  same sector and outcome, with the coefficient $1$, every other coefficient
  being $0$.
\end{enumerate}
\end{digestclaim}

% claim reductions/lem:row-relaxation
\digestlabel{lem:row-relaxation}{12.3}\pdfbookmark[2]{Lemma 12.3 A rounded constraint entry with its constant relaxed}{digest-reductions/lem:row-relaxation}
\renewcommand*{\RestatedName}{\hyperref{contract.pdf}{}{claim-reductions/lem:row-relaxation}{Lemma~12.3}}%
\begin{digestclaim}[A rounded constraint entry with its constant relaxed]
Let $\Lambda\colon\Tcl{\Bsys}\to\Tcl{\Qsys}$ be a channel with
$(M;O)\dsq_{\Lambda}(F;\hat O)$ (Definition~\ref{def:squashing}), and let
$(\hat O_k,\bowtie_k,c_k)$ be an entry with $\hat O_k$ a carried operator,
$\bowtie_k\in\{\preceq,=,\succeq\}$, $c_k\in\mathbb{R}$ and the certificate
$\Lambda^{\dagger}(\hat O_k)\bowtie_kc_k\,\idon{\Bsys}$ of
\eqref{eq:certificate}. Let $\hat O'_k$ be a Hermitian operator on $\Qsys$,
let $E_k\coloneqq\hat O'_k-\hat O_k$, and let $\rho_{il}\ge0$ be numbers with
$|(E_k)_{il}|\le\rho_{il}$ for every entry in a fixed orthonormal basis of
$\Hs{\Qsys}$, and $\delta_k\ge\max_i\sum_l\rho_{il}$. Define the
\emph{relaxed entries} of $(\hat O_k,\bowtie_k,c_k)$ as
$(\hat O'_k,\succeq,c'_k)$ with $c'_k\le c_k-\delta_k$ where $\bowtie_k$ is
$\succeq$; $(\hat O'_k,\preceq,c''_k)$ with $c''_k\ge c_k+\delta_k$ where
$\bowtie_k$ is $\preceq$; and both of these where $\bowtie_k$ is $=$. Then:
\begin{enumerate}[label=(\roman*),leftmargin=2.4em]
\item $-\delta_k\idon{\Qsys}\preceq E_k\preceq\delta_k\idon{\Qsys}$;
\item \emph{(certificate)} $(M;O,\Lambda^{\dagger}(\hat O'_k))\dsq_{\Lambda}
  (F;\hat O,\hat O'_k)$, and every relaxed entry carries the spectral
  certificate \eqref{eq:certificate} at its own constant:
  $\Lambda^{\dagger}(\hat O'_k)\succeq c'_k\idon{\Bsys}$, respectively
  $\Lambda^{\dagger}(\hat O'_k)\preceq c''_k\idon{\Bsys}$;
\item \emph{(attack sets)} for all finite-dimensional $E',E''$, every
  $\Aev$ in the set \eqref{eq:attset}, \eqref{eq:attsetun} or, for an entry
  tagged party-local with the factor $f$, \eqref{eq:attsetpl}, of the list
  holding the exact entry lies in the set of the list in which the exact
  entry is replaced by its relaxed entries, the other entries unchanged;
\item \emph{(states)} for every finite-dimensional system $X$ and every
  positive semidefinite $\omega$ on $\Hs{X}\otimes\Hs{\Qsys}$,
  $\tr_{\Qsys}[(\idon{X}\otimes\hat O_k)\,\omega]\succeq c_k\tr_{\Qsys}\omega$
  implies $\tr_{\Qsys}[(\idon{X}\otimes\hat O'_k)\,\omega]\succeq
  c'_k\tr_{\Qsys}\omega$, and the same with $\preceq$, $c''_k$ in place of
  $\succeq$, $c'_k$; with $X$ trivial these are the scalar readings
  $\tr[\hat O_k\omega]\ge c_k\tr\omega$ and the like.
\end{enumerate}
\end{digestclaim}

\section*{13\quad Assumptions}\pdfbookmark[1]{13 Assumptions}{digest-group-reductions/ass:meas}

% claim reductions/ass:meas
\digestlabel{ass:meas}{13.1}\pdfbookmark[2]{Assumption 13.1 Each measuring party's device as one POVM per round}{digest-reductions/ass:meas}
\SetShortId{A-MEAS}
\renewcommand*{\RestatedName}{\hyperref{contract.pdf}{}{claim-reductions/ass:meas}{Assumption~13.1}}%
\begin{digestassumption}[Each measuring party's device as one POVM per round]
In every round, each measuring party's interaction with the quantum system it
receives satisfies:
(i) it is described by a single POVM on that system;
(ii) it emits the classical outcome and nothing else, and no
post-measurement quantum system is retained or reused;
(iii) all downstream processing (announcements, sifting, acceptance, key map,
error correction, privacy amplification) acts on classical data;
(iv) the round's measurement channel is the product of one
measurement channel per factor, each acting on its own factor alone, for each
of the following factorizations of the measured system $\Bsys_j$:
where the declaration declares several measured factors
(\xref{protocol-compiler/def:spec}), the factorization
$\Bsys_j=\bigotimes_f\Bsys_{f,j}$ with one factor per measured factor,
whichever parties carry their labels and two factors of one party included;
and where several parties measure, the factorization into the parties'
parts of Definition~\ref{def:protocol}.
\end{digestassumption}

% claim reductions/ass:round
\digestlabel{ass:round}{13.2}\pdfbookmark[2]{Assumption 13.2 Round homogeneity of the physical measurement}{digest-reductions/ass:round}
\SetShortId{A-ROUND-SQ}
\renewcommand*{\RestatedName}{\hyperref{contract.pdf}{}{claim-reductions/ass:round}{Assumption~13.2}}%
\begin{digestassumption}[Round homogeneity of the physical measurement]
Let $M$ be the physical POVM: the POVM on the measured system of a
round that the analysis takes as the measurement the devices of the
measuring parties perform on it.
The family $\{M\}$ is round-homogeneous
(\xref{foundations/def:round-homogeneity}) in both clauses of that
definition:
(i) the same $M$ describes the measurement of every round;
(ii) a round's measurement acts as $M$ describes, whatever the inputs and
outputs of the other rounds, with no detector memory across rounds.
\end{digestassumption}

% claim reductions/ass:recv
\digestlabel{ass:recv}{13.3}\pdfbookmark[2]{Assumption 13.3 Supplier artefact whose claims hold of the device}{digest-reductions/ass:recv}
\SetShortId{A-RECV}
\renewcommand*{\RestatedName}{\hyperref{contract.pdf}{}{claim-reductions/ass:recv}{Assumption~13.3}}%
\begin{digestassumption}[Supplier artefact whose claims hold of the device]
Every declared block, weight, matrix and scalar that a
reduction instance of this module reads from a supplier's artefact
satisfies:
(i) the artefact is the output of the supplying module, at the module version
it declares, which is an \emph{accepted} version of the supplying module, and
on the inputs it declares, and the validation results it carries are that
module's own;
(ii) the facts the supplier's claims state hold of those numbers as declared,
under that module's own assumption ledger.
\end{digestassumption}

% claim reductions/ass:probe
\digestlabel{ass:probe}{13.5}\pdfbookmark[2]{Assumption 13.5 The station's private setting drawn from its own declared randomness}{digest-reductions/ass:probe}
\SetShortId{A-PROBE}
\renewcommand*{\RestatedName}{\hyperref{contract.pdf}{}{claim-reductions/ass:probe}{Assumption~13.5}}%
\begin{digestassumption}[The station's private setting drawn from its own declared randomness]
The private per-round setting $t\in\Tset$ of a measuring party's device
satisfies, in every round:
(i) it is drawn from that party's own randomness with the
declared class weights $q_t$;
(ii) it is drawn independently of the adversary's actions and of the round's
quantum input;
(iii) it is not inferred from any outcome.
\end{digestassumption}

% claim reductions/ass:tsm
\digestlabel{ass:tsm}{13.6}\pdfbookmark[2]{Assumption 13.6 Output-state test statistics}{digest-reductions/ass:tsm}
\SetShortId{A-TSM}
\renewcommand*{\RestatedName}{\hyperref{contract.pdf}{}{claim-reductions/ass:tsm}{Assumption~13.6}}%
\begin{digestassumption}[Output-state test statistics]
Every per-round public choice of known bias in the physical protocol that is
applied at the measurement the emitted family $\{F_v\}$ describes, a private
per-round setting of that measurement and any round-type designation the
consumer adds by relabelling its outcomes included,
(i) is a component of the emitted outcome $v$, and
(ii) enters the analysis only through the statistics of $\{F_v\}$.
\end{digestassumption}

% claim reductions/ass:constdecl
\digestlabel{ass:constdecl}{13.7}\pdfbookmark[2]{Assumption 13.7 Declared constants of the constraint list}{digest-reductions/ass:constdecl}
\SetShortId{A-CONST}
\renewcommand*{\RestatedName}{\hyperref{contract.pdf}{}{claim-reductions/ass:constdecl}{Assumption~13.7}}%
\begin{digestassumption}[Declared constants of the constraint list]
Every constant $c_k$ in an emitted constraint entry
(i) is fixed before the protocol run by the device fact that certifies it,
and
(ii) is not a function of any statistic observed during the run.
\end{digestassumption}

\part*{Engine requirements}\pdfbookmark[0]{Engine requirements}{digest-part-2}

\section*{14\quad Requirements: channel witnesses, post-processing and relations}\pdfbookmark[1]{14 Requirements: channel witnesses, post-processing and relations}{digest-group-reductions/req:channel-witness}

% claim reductions/req:channel-witness
\digestlabel{req:channel-witness}{14.1}\pdfbookmark[2]{Requirement 14.1 A channel witness for every emitted link}{digest-reductions/req:channel-witness}
\renewcommand*{\RestatedName}{\hyperref{contract.pdf}{}{claim-reductions/req:channel-witness}{Requirement~14.1}}%
\begin{digestrequirement}[A channel witness for every emitted link]
The engine emits every link with an explicit channel witness and
refuses a request that names a reduction coming with none.
\end{digestrequirement}

% claim reductions/req:pinched-channel-witness
\digestlabel{req:pinched-channel-witness}{14.2}\pdfbookmark[2]{Requirement 14.2 The pinched squashing channel as the emitted witness}{digest-reductions/req:pinched-channel-witness}
\renewcommand*{\RestatedName}{\hyperref{contract.pdf}{}{claim-reductions/req:pinched-channel-witness}{Requirement~14.2}}%
\begin{digestrequirement}[The pinched squashing channel as the emitted witness]
The engine emits $\PinS\circ\Lambda$, the squashing channel followed by the
pinching onto the sectors of Corollary~\ref{cor:blockdiag}, as the
artefact's channel witness (output \iface{squashing_channel}), so that the
block-diagonality restriction of that corollary's part (iii) and every other
emitted constraint are sound for \emph{one} channel.
\end{digestrequirement}

% claim reductions/req:carried-tuple-emitted
\digestlabel{req:carried-tuple-emitted}{14.3}\pdfbookmark[2]{Requirement 14.3 Emission of the carried tuple, and constraints built from it}{digest-reductions/req:carried-tuple-emitted}
\renewcommand*{\RestatedName}{\hyperref{contract.pdf}{}{claim-reductions/req:carried-tuple-emitted}{Requirement~14.3}}%
\begin{digestrequirement}[Emission of the carried tuple, and constraints built from it]
Beside the channel witness of each link it emits, the engine emits the
carried tuple of the link's squashing relation
(Definition~\ref{def:squashing}; \iface{transported_operators}), and it
builds every constraint it emits from an operator in that tuple.
\end{digestrequirement}

% claim reductions/req:no-approximate-relations
\digestlabel{req:no-approximate-relations}{14.4}\pdfbookmark[2]{Requirement 14.4 Refusal of approximate squashing relations}{digest-reductions/req:no-approximate-relations}
\renewcommand*{\RestatedName}{\hyperref{contract.pdf}{}{claim-reductions/req:no-approximate-relations}{Requirement~14.4}}%
\begin{digestrequirement}[Refusal of approximate squashing relations]
The engine refuses a declared squashing relation that carries an error
term, with the reason \iface{approximate_squashing_unsupported}, and emits no
link with the status \iface{certified_outer}, which is reserved for
approximate squashing.
\end{digestrequirement}

% claim reductions/req:round-homogeneous-emission
\digestlabel{req:round-homogeneous-emission}{14.5}\pdfbookmark[2]{Requirement 14.5 Emission of the round-homogeneous case only}{digest-reductions/req:round-homogeneous-emission}
\renewcommand*{\RestatedName}{\hyperref{contract.pdf}{}{claim-reductions/req:round-homogeneous-emission}{Requirement~14.5}}%
\begin{digestrequirement}[Emission of the round-homogeneous case only]
The engine emits one declared model and one squashing channel for
every round, both built from the one physical POVM that
Assumption~\ref{ass:round} holds fixed (the interface convention
\iface{round_homogeneity}). It emits no round-dependent squasher.
\end{digestrequirement}

% claim reductions/req:post-processing-order
\digestlabel{req:post-processing-order}{14.6}\pdfbookmark[2]{Requirement 14.6 Noise and coarse-graining before the first squashing link}{digest-reductions/req:post-processing-order}
\renewcommand*{\RestatedName}{\hyperref{contract.pdf}{}{claim-reductions/req:post-processing-order}{Requirement~14.6}}%
\begin{digestrequirement}[Noise and coarse-graining before the first squashing link]
The engine applies a noise or coarse-graining map on the physical POVM
before the first squashing link, and refuses a request to apply one after a
squashing link (input \iface{post_processing_order_declaration}, refusal
\iface{post_processing_order_violated}). The refusal reaches exactly those
two kinds of map: a stochastic relabelling applied to both sides after a
link, depending on the private component of the outcome only, is not
refused by it (Remark~\ref{C-rem:intertwining}, \eqref{C-eq:relabel}).
\end{digestrequirement}

% claim reductions/req:declared-relation-check
\digestlabel{req:declared-relation-check}{14.7}\pdfbookmark[2]{Requirement 14.7 Exact check of a declared relation}{digest-reductions/req:declared-relation-check}
\renewcommand*{\RestatedName}{\hyperref{contract.pdf}{}{claim-reductions/req:declared-relation-check}{Requirement~14.7}}%
\begin{digestrequirement}[Exact check of a declared relation]
The engine tests every declared relation (input \iface{declared_relations})
against both necessary conditions of Lemma~\ref{lem:obstructions} and
refuses one that fails either, with \iface{declared_relation_impossible}.
With $X$ and $X'$ as in that lemma, it decides (i) in operator form on the
rational enclosures $a\le b$ the declaration exhibits: the pair $a\idon{\Qsys}\preceq X\preceq b\idon{\Qsys}$
exactly by \eqref{eq:psdtest} and, where it holds, the same pair for $X'$;
and (ii) as the exact rank over $\Fldi$ of the matrix whose columns are the
vectorized declared elements, by elimination in the field, together with its
linear relations: an exact basis of the solutions $(a,c)$ of
$\sum_va_vF_v=c\,\idon{\Qsys}$, obtained by elimination in the field, each
basis solution tested on the declared physical elements as
$\sum_va_vM_v=c\,\idon{\Bsys}$. No eigenvalue is computed and no threshold is
used.
\end{digestrequirement}

% claim reductions/req:reserved-link-kinds
\digestlabel{req:reserved-link-kinds}{14.8}\pdfbookmark[2]{Requirement 14.8 Refusal of the reserved receiver-stage link kind}{digest-reductions/req:reserved-link-kinds}
\renewcommand*{\RestatedName}{\hyperref{contract.pdf}{}{claim-reductions/req:reserved-link-kinds}{Requirement~14.8}}%
\begin{digestrequirement}[Refusal of the reserved receiver-stage link kind]
The engine emits no receiver-stage link of the reserved kind
\iface{label-space}, and refuses a request that declares a link of that kind
with \iface{link_kind_deferred}.
\end{digestrequirement}

% claim reductions/req:unpadded-branch-blocks
\digestlabel{req:unpadded-branch-blocks}{14.9}\pdfbookmark[2]{Requirement 14.9 Branch blocks emitted as supplied}{digest-reductions/req:unpadded-branch-blocks}
\renewcommand*{\RestatedName}{\hyperref{contract.pdf}{}{claim-reductions/req:unpadded-branch-blocks}{Requirement~14.9}}%
\begin{digestrequirement}[Branch blocks emitted as supplied]
The engine does not pad the branch blocks of a composite target: a
reduction instance emits the blocks its branches supply, and a padded variant a
submodule proves from Corollary~\ref{cor:flagsenough} is not emitted unless
that submodule's interface declares it.
\end{digestrequirement}

\section*{15\quad Requirements: chains, constraint slots and the licence}\pdfbookmark[1]{15 Requirements: chains, constraint slots and the licence}{digest-group-reductions/req:constraint-and-licence-slots}

% claim reductions/req:constraint-and-licence-slots
\digestlabel{req:constraint-and-licence-slots}{15.1}\pdfbookmark[2]{Requirement 15.1 Constraint and licence slots on every artefact}{digest-reductions/req:constraint-and-licence-slots}
\renewcommand*{\RestatedName}{\hyperref{contract.pdf}{}{claim-reductions/req:constraint-and-licence-slots}{Requirement~15.1}}%
\begin{digestrequirement}[Constraint and licence slots on every artefact]
The engine emits every artefact with a constraint list
\iface{constraint_list} and a licence slot \iface{attack_set_license}, even
when both are empty: a record without them excludes a class of device
knowledge that no classical target can carry
(Proposition~\ref{prop:obstruction}, Remark~\ref{C-rem:obstruction-scope}\,(a)),
and the licence slot is where the licence of Theorem~\ref{thm:licence} is
carried.
\end{digestrequirement}

% claim reductions/req:licence-record
\digestlabel{req:licence-record}{15.2}\pdfbookmark[2]{Requirement 15.2 Content and presence of the licence record}{digest-reductions/req:licence-record}
\renewcommand*{\RestatedName}{\hyperref{contract.pdf}{}{claim-reductions/req:licence-record}{Requirement~15.2}}%
\begin{digestrequirement}[Content and presence of the licence record]
The engine emits the licence record \iface{attack_set_license} naming
the theorem that licenses the constraint list (Theorem~\ref{thm:licence});
the hypothesis it was discharged on, which is the certificate
\eqref{eq:certificate} of each entry; and, per entry, what discharges that
hypothesis: the claim that establishes the certificate under
Corollary~\ref{cor:entries} --- the instance's own claim for an entry of its
list, that corollary for a default carried operator and for the
class-marginal entry, the claim giving the decided certificate for a
requested entry --- and, where that claim reads a fact a supplier's claims
state, the supplier's artefact under Assumption~\ref{ass:recv}, named by its
version. It emits the record only when every entry has its certificate and
what discharges it. The artefact names the proof technique the licence was
issued for, since under another technique the licence must be re-established
(Remark~\ref{C-rem:notlicensed}\,(e)).
\end{digestrequirement}

% claim reductions/req:declared-constants-only
\digestlabel{req:declared-constants-only}{15.3}\pdfbookmark[2]{Requirement 15.3 Refusal of data-dependent constants}{digest-reductions/req:declared-constants-only}
\renewcommand*{\RestatedName}{\hyperref{contract.pdf}{}{claim-reductions/req:declared-constants-only}{Requirement~15.3}}%
\begin{digestrequirement}[Refusal of data-dependent constants]
The engine builds no constraint entry whose constant $c_k$ is a function of a
statistic observed during the run; a request for one is refused with
\iface{data_dependent_constant}. Such an entry would be unsound regardless of
provenance (Remark~\ref{C-rem:notlicensed}\,(c)), the constants being declared
before the run (Assumption~\ref{ass:constdecl}).
\end{digestrequirement}

% claim reductions/req:mixed-status-chain
\digestlabel{req:mixed-status-chain}{15.4}\pdfbookmark[2]{Requirement 15.4 Refusal of a chain with a link that is not exact}{digest-reductions/req:mixed-status-chain}
\renewcommand*{\RestatedName}{\hyperref{contract.pdf}{}{claim-reductions/req:mixed-status-chain}{Requirement~15.4}}%
\begin{digestrequirement}[Refusal of a chain with a link that is not exact]
The engine assigns a chain no composite status other than the one of
Theorem~\ref{thm:chain}\,(iv): a chain containing a link that is not
\textsf{exact} is refused, with
\iface{approximate_squashing_unsupported}, rather than given a composite
status, because no theorem of this contract converts an error term into a
security penalty. The refusal is a rule of the engine and not a converse
this contract proves.
\end{digestrequirement}

% claim reductions/req:identity-link-recorded
\digestlabel{req:identity-link-recorded}{15.5}\pdfbookmark[2]{Requirement 15.5 Identity link recorded for an already-finite receiver}{digest-reductions/req:identity-link-recorded}
\renewcommand*{\RestatedName}{\hyperref{contract.pdf}{}{claim-reductions/req:identity-link-recorded}{Requirement~15.5}}%
\begin{digestrequirement}[Identity link recorded for an already-finite receiver]
The engine records the reduction of an already-finite receiver as the
identity link of Example~\ref{C-ex:identity}, never as an empty chain: an empty
chain would leave the artefact unable to say whether a reduction was performed
and found trivial or never considered, and the composition order the chain
carries as data (Definition~\ref{def:chain}) has to be carried even when there
is nothing to order.
\end{digestrequirement}

\section*{16\quad Requirements: constraint entries}\pdfbookmark[1]{16 Requirements: constraint entries}{digest-group-reductions/req:certified-entries-only}

% claim reductions/req:certified-entries-only
\digestlabel{req:certified-entries-only}{16.1}\pdfbookmark[2]{Requirement 16.1 Emission of certified constraint entries only}{digest-reductions/req:certified-entries-only}
\renewcommand*{\RestatedName}{\hyperref{contract.pdf}{}{claim-reductions/req:certified-entries-only}{Requirement~16.1}}%
\begin{digestrequirement}[Emission of certified constraint entries only]
Every constraint entry the engine emits carries exactly one
certificate, named by its case of Corollary~\ref{cor:pullback}
(identity, positivity or spectral bounds), and is of the certified kind, case
(i) of Corollary~\ref{cor:transported}; the engine emits no entry of the
assumed kind, case (ii). Each emitted entry's certificate case is recorded
with the transported set, as
Requirement~\ref{req:transported-set-certificates} states.
\end{digestrequirement}

% claim reductions/req:class-marginal-operator-form
\digestlabel{req:class-marginal-operator-form}{16.2}\pdfbookmark[2]{Requirement 16.2 Refusal of the operator form of the class-marginal identity}{digest-reductions/req:class-marginal-operator-form}
\renewcommand*{\RestatedName}{\hyperref{contract.pdf}{}{claim-reductions/req:class-marginal-operator-form}{Requirement~16.2}}%
\begin{digestrequirement}[Refusal of the operator form of the class-marginal identity]
The engine refuses, with the reason \iface{unsound_constraint_requested}, a
request to impose on the squashed states the operator form
$\sum_{v\in\Vset_t}F_v=q_t\,\idon{\Qsys}$ of the class-marginal identity
(Remark~\ref{C-rem:unsound}).
\end{digestrequirement}

% claim reductions/req:constraint-list-by-request
\digestlabel{req:constraint-list-by-request}{16.3}\pdfbookmark[2]{Requirement 16.3 The constraint list and the licence decided by the request}{digest-reductions/req:constraint-list-by-request}
\renewcommand*{\RestatedName}{\hyperref{contract.pdf}{}{claim-reductions/req:constraint-list-by-request}{Requirement~16.3}}%
\begin{digestrequirement}[The constraint list and the licence decided by the request]
Where the request carries no entry (Corollary~\ref{cor:entries}), the engine
emits an empty \iface{constraint_list}, whatever entries the instance's own
list contains, and no \iface{attack_set_license}. Where the request carries
entries, it emits in \iface{constraint_list} the instance's own entries, the
class-marginal entry on request and any further requested entry that
Corollary~\ref{cor:entries}\,(iii) admits
(Requirement~\ref{req:uncertified-operator-refused}), and it emits
the licence of Theorem~\ref{thm:licence} in \iface{attack_set_license}. The
licence is present if and only if the list is non-empty
(\iface{licence_presence}): a non-empty list without a licence is refused
with \iface{licence_missing}, and a licence without a list with
\iface{licence_without_constraints}.
\end{digestrequirement}

% claim reductions/req:transported-set-certificates
\digestlabel{req:transported-set-certificates}{16.4}\pdfbookmark[2]{Requirement 16.4 Certificates carried with the transported set}{digest-reductions/req:transported-set-certificates}
\renewcommand*{\RestatedName}{\hyperref{contract.pdf}{}{claim-reductions/req:transported-set-certificates}{Requirement~16.4}}%
\begin{digestrequirement}[Certificates carried with the transported set]
Whatever the request carries, the engine emits in
\iface{transported_operators}, per operator of the transported set,
the certificate case of Corollary~\ref{cor:pullback} that makes it usable. In
particular the certificate $(\Pi_s,\succeq,0)$ of every sector projector of
every instance travels there and not as a constraint entry; such an operator
becomes an entry of \iface{constraint_list} only where a caller requests it
as an entry (Corollary~\ref{cor:entries}\,(i)).
\end{digestrequirement}

% claim reductions/req:class-marginal-entry
\digestlabel{req:class-marginal-entry}{16.5}\pdfbookmark[2]{Requirement 16.5 The class-marginal entry only where the instance discharges it}{digest-reductions/req:class-marginal-entry}
\renewcommand*{\RestatedName}{\hyperref{contract.pdf}{}{claim-reductions/req:class-marginal-entry}{Requirement~16.5}}%
\begin{digestrequirement}[The class-marginal entry only where the instance discharges it]
The engine emits the class-marginal entry $\big(\sum_{v\in\Vset_t}F_v,=,q_t\big)$
at a declared $t$, in the expectation form of
Corollary~\ref{cor:entries}\,(ii), only on request and only where the
instance discharges the class-marginal identity
\eqref{eq:instance-classmarginal}. A request for it of an instance that does
not discharge the identity is refused with
\iface{constraint_without_certificate}, and the entry is never emitted on the
strength of a branch that does not establish it.
\end{digestrequirement}

% claim reductions/req:caller-constant-rounding
\digestlabel{req:caller-constant-rounding}{16.6}\pdfbookmark[2]{Requirement 16.6 Outward rounding of a caller's constant off the lattice}{digest-reductions/req:caller-constant-rounding}
\renewcommand*{\RestatedName}{\hyperref{contract.pdf}{}{claim-reductions/req:caller-constant-rounding}{Requirement~16.6}}%
\begin{digestrequirement}[Outward rounding of a caller's constant off the lattice]
Where the declared constant of a requested entry does not satisfy
\eqref{eq:latinv}, the engine first rounds it onto the lattice in the
direction that enlarges the feasible set for the entry's relation, down for
$\succeq$ and up for $\preceq$, and records the direction taken per entry in
\iface{emission_record} (\xref{foundations/req:rounding-direction}). An
equality entry whose constant is off the lattice is refused with
\iface{constraint_without_certificate} rather than split into a pair, and a
declared constant of modulus above one is refused with
\iface{constraint_without_certificate} under every relation, before any
rounding.
\end{digestrequirement}

% claim reductions/req:uncertified-operator-refused
\digestlabel{req:uncertified-operator-refused}{16.7}\pdfbookmark[2]{Requirement 16.7 Refusal of a requested operator without a decided certificate}{digest-reductions/req:uncertified-operator-refused}
\renewcommand*{\RestatedName}{\hyperref{contract.pdf}{}{claim-reductions/req:uncertified-operator-refused}{Requirement~16.7}}%
\begin{digestrequirement}[Refusal of a requested operator without a decided certificate]
The engine emits an operator that a request carries as an entry only where
Corollary~\ref{cor:entries} covers it: as a default carried operator (part
(i)), as the class-marginal entry of an instance that discharges the
class-marginal identity (part (ii)), or, under part (iii), where a claim
stated for the instance gives a decided certificate for the operator's form
and the engine decides that condition exactly on the emitted numbers and the
request's declared constants, finding that it holds. It refuses, with
\iface{constraint_without_certificate}, every other requested operator: one
outside the form of Corollary~\ref{cor:entries}\,(i)--(iii), one of that form
for which no claim stated for the instance gives a decided certificate, and
one whose decided certificate fails at the emitted numbers. It takes no
certificate \eqref{eq:certificate} from the request on trust, and it never
emits a requested operator at some other constant.
\end{digestrequirement}

% claim reductions/req:party-local-tag
\digestlabel{req:party-local-tag}{16.8}\pdfbookmark[2]{Requirement 16.8 The party-local tag and its scope}{digest-reductions/req:party-local-tag}
\renewcommand*{\RestatedName}{\hyperref{contract.pdf}{}{claim-reductions/req:party-local-tag}{Requirement~16.8}}%
\begin{digestrequirement}[The party-local tag and its scope]
The engine emits every constraint entry with the reading it is to be read in.
An entry is emitted untagged, as $(\hat O,\bowtie,c)$, and is read in the
form of the declared source mode (Definition~\ref{def:licensed}); or, only
where the measured system is a declared product
$\Qsys=\bigotimes_f\Qsys_f$ of party-scoped factors and the entry's
certificate is discharged on that factor's own channel $\Lambda_f$
(Theorem~\ref{thm:licence}\,(v)), it is emitted with the tag
\iface{party_local} and a scope, as $(\hat O,\bowtie,c;f)$, with $\hat O$ an
operator on $\Qsys_f$ alone carried as a sector-block list against that
factor's sector map. In this module the tag is emitted only with the
entries of several measuring parties
(Theorem~\ref{thm:emitted-product}\,(iii)).
A request for the party-local form on a one-factor emission is answered
with the untagged entry, at the same operator, relation and constant, and is
not refused.
\end{digestrequirement}

% claim reductions/req:licensed-entry-lattice
\digestlabel{req:licensed-entry-lattice}{16.9}\pdfbookmark[2]{Requirement 16.9 Lattice condition on the entries of emitted constraint operators}{digest-reductions/req:licensed-entry-lattice}
\renewcommand*{\RestatedName}{\hyperref{contract.pdf}{}{claim-reductions/req:licensed-entry-lattice}{Requirement~16.9}}%
\begin{digestrequirement}[Lattice condition on the entries of emitted constraint operators]
Where the licence is present, the engine decides, on the emitted numbers and
before the licence is issued, that every entry of every emitted constraint
operator satisfies \eqref{eq:latinv}, and refuses the request
otherwise, with the reason \iface{constraint_without_certificate} (the lattice
half of the rule \iface{licensed_hermiticity};
Corollary~\ref{cor:entries}\,(iii)).
\end{digestrequirement}

\section*{17\quad Requirements: exact admission and emission}\pdfbookmark[1]{17 Requirements: exact admission and emission}{digest-group-reductions/req:exact-admission-refusal}

% claim reductions/req:exact-admission-refusal
\digestlabel{req:exact-admission-refusal}{17.1}\pdfbookmark[2]{Requirement 17.1 Refusal of a declaration outside the substrate}{digest-reductions/req:exact-admission-refusal}
\renewcommand*{\RestatedName}{\hyperref{contract.pdf}{}{claim-reductions/req:exact-admission-refusal}{Requirement~17.1}}%
\begin{digestrequirement}[Refusal of a declaration outside the substrate]
The engine refuses a declaration whose entries are not exactly representable
in the field $\Fld$ chosen for the arithmetic
(Definition~\ref{def:substrate}), with the reason
\iface{declaration_not_exactly_admissible}, rather than rounding the entries
into that field. No declaration whose entries satisfy \eqref{eq:latinv} is
refused on this ground.
\end{digestrequirement}

% claim reductions/req:emission-refusal
\digestlabel{req:emission-refusal}{17.2}\pdfbookmark[2]{Requirement 17.2 Refusal of an emitted family that is not exactly admissible}{digest-reductions/req:emission-refusal}
\renewcommand*{\RestatedName}{\hyperref{contract.pdf}{}{claim-reductions/req:emission-refusal}{Requirement~17.2}}%
\begin{digestrequirement}[Refusal of an emitted family that is not exactly admissible]
The engine emits a family only if the six conditions of
Theorem~\ref{thm:emit}\,(ii) --- (a)--(e) of
Definition~\ref{def:exactadm} together with \eqref{eq:latinv} of every entry
--- hold of the numbers about to be emitted, decided exactly after the
emission normalizer where that acted; otherwise it refuses with the
reason \iface{emitted_family_not_admissible}, naming the predicate that
failed, and repairs nothing it was handed. Under the same reason, where the
run declares the strict option \iface{strict_snaps}
(\xref{foundations/req:snap-reference}\,(iv)) and no reference companion of
Requirement~\ref{req:reference-companion} is formed beside the artefact, it
refuses every moved class of Definition~\ref{def:emitnorm}, one snapped by
the last resort (3) included, naming the data sector, the class, the
cause and the fix; where a companion is formed, whether the class rests is
decided by the run after its search, and the run refuses there under the
same option. Where the run does not declare the option, no class is refused
for being moved. The cause of a moved class that would rest is that no
reference companion bypasses the normalizer, which, by
Requirement~\ref{req:reference-companion}, is where no claim of this module
establishes (H-REF) for the instance; for a class snapped by (3) it also
names that a component of its residue has no host. The fix named is
declared numbers on which the exact construction stays on the lattice, so
that the class is kept; for a class
snapped by (3) also a model that includes the device's loss and dark counts,
so that the elements of the class are positive definite on the sector. A
moved class is mixed only with the family of its diagonal parts in (2), or
with a multiple of the identity in the last resort (3).
\end{digestrequirement}

% claim reductions/req:emission-record
\digestlabel{req:emission-record}{17.3}\pdfbookmark[2]{Requirement 17.3 The record of an emission}{digest-reductions/req:emission-record}
\renewcommand*{\RestatedName}{\hyperref{contract.pdf}{}{claim-reductions/req:emission-record}{Requirement~17.3}}%
\begin{digestrequirement}[The record of an emission]
The engine emits with every artefact, in \iface{emission_record}: the
declared lattice exponent $\kappa$; whether the run declared the strict
option of \xref{foundations/req:snap-reference}\,(iv); the reference kind of
every snap the normalizer makes, \emph{exact}, every radius zero, the
reference being the handed datum itself
(Definition~\ref{def:emitnorm}), so that the enclosure setting, recorded as
\xref{foundations/req:who-snaps} asks of the record of every snap, does not
bear on it; per data sector, the emission displacement
$\Delta_s$ of \eqref{eq:displacement}; per data sector and class, whether the
class was kept (Definition~\ref{def:emitnorm}\,(1)) or moved, and for a moved
class the step (2) or (3) of Definition~\ref{def:emitnorm} that snapped it,
the weight $w$ taken or the fall-back of (3), the hosts of the components of
the residue under (2), and for (3) the flag \iface{support_not_preserved}; and
the direction of every bound it rounds, each scalar value rounded
down and each complementary value recomputed (Theorem~\ref{thm:emit}\,(iv)).
A data sector \emph{rests on} \xref{foundations/ass:snapped-measurement}
where some class of it is moved, save where it is
evaluated off its snap (\xref{foundations/def:evaluated-off-snap}):
in a run whose key length is formed from a number on the reference
companion of Requirement~\ref{req:reference-companion}, its first-order
bound or a number certified on it, which carries no data of this emission's
normalizer, every data sector that holds a moved class is, a kept sector
being named as kept and not counted among them, and none rests on the
assumption. Where the key length is formed instead from a number certified
on the snapped declaration in the fallback of
\xref{meat-core/frank-wolfe/req:certified-evaluation}\,(i), none is, and
every data sector that holds a moved class rests on the assumption, with the
cause \iface{companion_not_certified}
(\xref{foundations/req:snap-reference}\,(iii)(b)). The record names every moved sector, with the flag of any class
of it snapped by (3), and carries
\xref{foundations/ass:snapped-measurement} among the assumptions the
artefact may rest on exactly where there is one. This record is emitted with
the artefact, before the run knows whether the number its key length is
formed from is formed on the reference companion; what it names is therefore provisional, and the
record of the run decides which sectors rest on the assumption, each with
its cause and group (\xref{foundations/req:snap-reference}\,(iii);
Theorem~\ref{thm:emit}\,(vi)).
\end{digestrequirement}

% claim reductions/req:reference-companion
\digestlabel{req:reference-companion}{17.4}\pdfbookmark[2]{Requirement 17.4 Emitting the reference companion beside the artefact}{digest-reductions/req:reference-companion}
\renewcommand*{\RestatedName}{\hyperref{contract.pdf}{}{claim-reductions/req:reference-companion}{Requirement~17.4}}%
\begin{digestrequirement}[Emitting the reference companion beside the artefact]
Where a claim of this module establishes (H-REF) of
Corollary~\ref{cor:reference-companion} for the reduction instance the
artefact is emitted from, and either the declaring instance hands the
reference beside its snap of the measurement it states
(\xref{foundations/def:evaluated-off-snap}\,(i)) on at least one block or
the emission normalizer moves at least one class
(Definition~\ref{def:emitnorm}), the engine emits beside the artefact, in
\iface{reference_companion}, the reference companion of that instance at the
family $G$ that carries the handed reference on every handed block and the
snapped block the instance declares on every other block, so that, where no
block is handed, $G$ is the instance's snapped family and the companion's
datum is the handed datum
(Corollary~\ref{cor:reference-companion}\,(i)):
\begin{enumerate}[label=(\roman*),leftmargin=2.4em]
\item \emph{The datum.} Every enclosure of a handed block is first
  recentred by \xref{foundations/lem:recentred-enclosure}, so that every
  entry whose enclosure contains $0$ has centre $0$. The engine then forms,
  in exact arithmetic on the exact island
  (\xref{foundations/def:exact-island}), the enclosure of the companion's
  datum that Corollary~\ref{cor:reference-companion}\,(iii) gives, from these
  enclosures and from the snapped blocks with radius zero, by the
  construction it applies to the snapped family: the same rule of (I1), the
  same weight route, scalar values, sectors and supports. It then
  recentres that enclosure again, element by element and block by block,
  by \xref{foundations/lem:recentred-enclosure}, after making its centre
  Hermitian and its radii symmetric by part (iv) of that lemma where they
  are not, every block of the datum being Hermitian
  (Corollary~\ref{cor:reference-companion}\,(i)); where the rule of (I1)
  combines entries of several blocks, a sum that is exactly zero may so
  have a centre that is not, and this second recentring sets it to zero
  wherever its enclosure contains zero. The entries of (iii) are formed
  from the datum's enclosure so recentred. It applies no
  normalizer (Definition~\ref{def:emitnorm}), decides no lattice condition,
  rounds nothing, and decides nothing of the companion's admissibility,
  which is Corollary~\ref{cor:reference-companion}\,(i) on the claims the
  handed reference names (\xref{foundations/req:snap-reference}\,(i)). It
  emits the centres, entries of $\mathbb{Q}(i)$, and the radii, rationals,
  entry by entry, so that an entry the reference knows exactly is emitted
  exactly, with radius zero.
\item \emph{The shape.} The companion is emitted with the sector map, the
  alphabet with its classes and weights, the scalar values and the party
  partition of the artefact's target, which are equal by
  Corollary~\ref{cor:reference-companion}\,(i), and with the artefact's
  channel witness and transported set read at the companion's datum.
\item \emph{The entries.} On a route whose request carries entries
  (Corollary~\ref{cor:entries}), the companion's entry list holds the
  instance's own entries at the companion's datum and, for every entry
  requested beside them that the artefact carries, its companion form
  (Corollary~\ref{cor:reference-companion}\,(vi)) at the writing by which
  the engine recognised it, where its certificate at the companion's datum
  is established: for a sum of sector projectors and for the class-marginal
  entry, by Corollary~\ref{cor:reference-companion}\,(vi), the operator
  being the artefact's, enclosed with radius zero; for another entry, where
  a claim of this module states for the instance a condition on an
  enclosure of the companion's datum under which its certificate holds,
  by deciding that condition exactly on the datum's enclosure recentred as
  (i) states (for the operator \eqref{C-eq:two-sector-entry} on the flat instance's
  licensed route, \eqref{eq:companion-two-sector} of
  Corollary~\ref{cor:flat-companion}). A requested entry whose certificate
  at the companion's datum is not so established is not formed on the
  companion and has no entry there in its place. For each entry
  $(\hat O_k,\bowtie_k,c_k)$ of that list, with the enclosure $(C_k,R_k)$
  of its operator that Corollary~\ref{cor:reference-companion}\,(iii)
  gives, or radius zero where the operator is the artefact's: the operator
  $\hat O'_k$ whose every real and imaginary part is the binary64 number
  nearest to that of $C_k$, ties to even, which is Hermitian and has the
  blocks of $\hat O_k$; the rationals
  $\rho_{il}\coloneqq|\operatorname{Re}(\hat O'_k-C_k)_{il}|
  +|\operatorname{Im}(\hat O'_k-C_k)_{il}|+(R_k)_{il}$ and
  $\delta_k\coloneqq\max_i\sum_l\rho_{il}$, computed exactly; and the
  relaxed entries of Lemma~\ref{lem:row-relaxation} at $\delta_k$, with
  $c'_k$ the largest binary64 number at most $c_k-\delta_k$ and $c''_k$ the
  smallest at least $c_k+\delta_k$. Where $\delta_k=0$ the entry is
  $\hat O_k$ itself and is emitted as it is, an equality included. Each
  emitted entry carries the certificate case of the entry it relaxes, its
  $\delta_k$, and the mark that it is relaxed, its certificate being
  Lemma~\ref{lem:row-relaxation}\,(ii); its operator is added to the
  companion's transported set with its pull-back. The lattice condition of
  Requirement~\ref{req:licensed-entry-lattice} binds the artefact's own
  entries and not these.
\item \emph{No normalizer.} No step of the emission normalizer
  (Definition~\ref{def:emitnorm}) acts on the companion's datum: the
  companion carries the reference, not a snap of it.
\item \emph{The record.} The companion carries, per block of $G$, whether
  it is the handed reference or the snapped block, the kind of every handed
  block (\emph{exact}, every radius zero, or \emph{enclosed}, with the
  largest radius of the block after recentring), the entries recentred in
  the handed blocks and in the datum, the largest radius of the datum's
  enclosure after its recentring, each $\delta_k$, every requested entry
  formed on the companion with what established its certificate there, and
  every requested entry not formed with the reason; and it is
  marked evaluation-only: it has no run of its own, no consumer admitting
  it as the declaration of a run, and what a consumer decides or computes
  of it beside the declaration it accompanies is the consumer's
  (\xref{meat-core/req:companion-admission}). The artefact
  itself, its target, normalizer, emission record
  (Requirement~\ref{req:emission-record}) and entries, is emitted as it is
  without a companion.
\item \emph{No companion.} Where no block is handed and the normalizer
  moves no class, or no claim of this module establishes (H-REF) for the
  instance, or a reciprocal node of the rule of (I1) receives no enclosure,
  so that Corollary~\ref{cor:reference-companion}\,(iii) gives none, no
  companion is emitted, and the record says which of the three holds.
\item \emph{Factor enclosures.} Where the handed reference carries
  factor enclosures
  (\xref{foundations/def:factor-enclosure}) on some of its blocks
  $G^{(b)}_j$, $b$ a handed block, as the receiver's does
  (\xref{receiver/detector-imperfections/req:dark-composed-reference}), and a
  claim of this module establishes that the rule of (I1) of the instance is a
  non-negative combination (Corollary~\ref{cor:companion-factor-enclosure},
  part (iii) for the identity instance and
  Corollary~\ref{cor:flat-companion} for the flat instance), the engine
  forms, for every element of every data sector of the companion's datum
  every combining block of which carries one, the factor enclosure of
  Corollary~\ref{cor:companion-factor-enclosure}\,(i): the boxes of the
  combining blocks side by side in a fixed order, their centres and radii
  unaltered, and their cores multiplied by the coefficients in exact rational
  arithmetic. It emits each beside the entrywise enclosure of the
  datum of (i), which it does not alter, as the factor enclosure of the target's block
  on that data sector (part (ii) of that corollary). It forms none on any other
  element: none where a combining block carries none, and none for an
  instance no claim of this module shows to have such a rule, in particular
  none through a rule that maps the blocks of a data sector by a conjugation
  before it combines them, whose factor enclosure this module does not form.
  The record of (v) carries, per element of the datum, whether a factor
  enclosure is carried, with its number of columns, or why none is, and it
  names the claim that shows the rule of (I1) to be a non-negative
  combination, Corollary~\ref{cor:companion-factor-enclosure}\,(iii) or
  Corollary~\ref{cor:flat-companion}.
\end{enumerate}
\end{digestrequirement}

% claim reductions/req:party-label
\digestlabel{req:party-label}{17.5}\pdfbookmark[2]{Requirement 17.5 Party partition emitted with the sector map}{digest-reductions/req:party-label}
\renewcommand*{\RestatedName}{\hyperref{contract.pdf}{}{claim-reductions/req:party-label}{Requirement~17.5}}%
\begin{digestrequirement}[Party partition emitted with the sector map]
The engine emits the party partition of each reduction instance with
its sector map (\iface{target_descriptor}): the
default party partition where the instance declares no party label,
and otherwise the declared label (\iface{party_label}) on every factor of
every $\Hs{\Qsys_s}$. Where several parties measure, each factor's artefact
carries that factor's own label with its own sector map
(Theorem~\ref{thm:emitted-product}\,(i)).
\end{digestrequirement}

% claim reductions/req:factorised-emission
\digestlabel{req:factorised-emission}{17.6}\pdfbookmark[2]{Requirement 17.6 Emission of several measuring parties as factors}{digest-reductions/req:factorised-emission}
\renewcommand*{\RestatedName}{\hyperref{contract.pdf}{}{claim-reductions/req:factorised-emission}{Requirement~17.6}}%
\begin{digestrequirement}[Emission of several measuring parties as factors]
Where several measuring parties measure, the engine emits in
\iface{factor_emission} the ordered list of the factors' artefacts, each with
its party label, in the declared order of the factors, and the combination
rule \eqref{eq:producttarget}; it emits, declares and rounds no entry of any
product element $F_v$ (Theorem~\ref{thm:emitted-product}).
\end{digestrequirement}

\section*{18\quad Requirements: the assumption ledger}\pdfbookmark[1]{18 Requirements: the assumption ledger}{digest-group-reductions/req:assumption-ledger}

% claim reductions/req:assumption-ledger
\digestlabel{req:assumption-ledger}{18.1}\pdfbookmark[2]{Requirement 18.1 Rows of the emitted assumption ledger}{digest-reductions/req:assumption-ledger}
\renewcommand*{\RestatedName}{\hyperref{contract.pdf}{}{claim-reductions/req:assumption-ledger}{Requirement~18.1}}%
\begin{digestrequirement}[Rows of the emitted assumption ledger]
The engine emits with every artefact the assumption ledger
\iface{assumption_ledger}: one row, carrying its machine key, for each
assumption of this contract and for each assumption the interface lists as
taken from another module, each under its declared presence rule. The row
of A-SRC (\xref{transmitter/ass:src}) is present only where the emitted chain's
source segment is non-empty, and there under the presence rule the
interface declares for it, the row of Assumption~\ref{ass:constdecl} only
where the emitted constraint list is non-empty, the row of
Assumption~\ref{ass:probe} only where the declared setting set $\Tset$ has
at least two values, the row of Assumption~\ref{ass:tsm} only where $\Tset$
has at least two values or a public choice is applied at the measurement the
emitted family describes (and also where the emission is not told whether
one is), a submodule's own assumption
where that submodule's instance is emitted and any narrower presence rule the
submodule declares for it holds, the row of Assumption~\ref{ass:recv} on
every emission save where Requirement~\ref{req:recv-row-provenance} omits
it, and every other row on every emission.
These rows include every key
of the core's own assumptions, and every key the interface lists as taken
from another module, that the closure of the claims producing the artefact
contains, save the row of Assumption~\ref{ass:constdecl} where the emitted
constraint list is empty, a condition on the constants of emitted entries
holding vacuously of an empty list, the row of Assumption~\ref{ass:probe}
where $\Tset=\{\ast\}$, whose clauses hold of the one setting drawn with
weight one, and the row of Assumption~\ref{ass:tsm} where in addition no
public choice is applied at that measurement, whose clauses then concern no
choice, and the row of Assumption~\ref{ass:recv} where
Requirement~\ref{req:recv-row-provenance} omits it, whose clauses that
requirement shows to hold there (the supplier's keys ride the row of
Assumption~\ref{ass:recv}
instead, under the supplier's own presence rules,
Requirement~\ref{req:supplier-ledger-pass-through}); a row may be present that
no such closure contains.
\end{digestrequirement}

% claim reductions/req:never-verified-assumptions
\digestlabel{req:never-verified-assumptions}{18.2}\pdfbookmark[2]{Requirement 18.2 Never-verified rows and validation of declared data}{digest-reductions/req:never-verified-assumptions}
\renewcommand*{\RestatedName}{\hyperref{contract.pdf}{}{claim-reductions/req:never-verified-assumptions}{Requirement~18.2}}%
\begin{digestrequirement}[Never-verified rows and validation of declared data]
The engine marks as never verifiable (\iface{never_verified}) the row, in the
assumption ledger, of every assumption that the engine never verifies, as
that assumption's verification field says, whether this contract, one of its
submodules or another module states it; no code path decides what any of
them says beyond its clauses on the declared data. The engine decides by
validation the conditions on the declared data that the assumptions of this
contract are read with, such as the declared dimensions of
Definition~\ref{def:system-setting} and the declared class weights of
Definition~\ref{def:joint-family}.
\end{digestrequirement}

% claim reductions/req:supplier-ledger-pass-through
\digestlabel{req:supplier-ledger-pass-through}{18.3}\pdfbookmark[2]{Requirement 18.3 Pass-through of the supplier's assumption ledger}{digest-reductions/req:supplier-ledger-pass-through}
\renewcommand*{\RestatedName}{\hyperref{contract.pdf}{}{claim-reductions/req:supplier-ledger-pass-through}{Requirement~18.3}}%
\begin{digestrequirement}[Pass-through of the supplier's assumption ledger]
On the row of Assumption~\ref{ass:recv} in the assumption ledger
\iface{assumption_ledger}, the engine carries, as machine keys and beside the
artefact's module version \iface{module_version}, every key the supplied
ledger carries, as supplied and in the supplier's order; a key beyond those
the supplier's presence rules require is carried rather than refused. The
engine holds no list of the supplier's keys and adds no presence rule of its
own: what it reads from the supplying module, as it is installed where the
check runs, is that module's machine keys and the set of keys its rules make
conditional, and it evaluates the condition itself on the artefact's own
data.
\end{digestrequirement}

% claim reductions/req:supplier-ledger-refusals
\digestlabel{req:supplier-ledger-refusals}{18.4}\pdfbookmark[2]{Requirement 18.4 Refusal of a supplier ledger with a repeated or missing key}{digest-reductions/req:supplier-ledger-refusals}
\renewcommand*{\RestatedName}{\hyperref{contract.pdf}{}{claim-reductions/req:supplier-ledger-refusals}{Requirement~18.4}}%
\begin{digestrequirement}[Refusal of a supplier ledger with a repeated or missing key]
The engine refuses a supplied ledger that carries one machine key more than
once, and it refuses an artefact whose ledger lacks a key the supplier's own
presence rules make present, rather than passing it on as a silently
shortened list; both refusals carry the reason
\iface{receiver_spec_not_accepted}.
\end{digestrequirement}

% claim reductions/req:supplier-version-accepted
\digestlabel{req:supplier-version-accepted}{18.5}\pdfbookmark[2]{Requirement 18.5 Refusal of a supplier artefact of an unaccepted version}{digest-reductions/req:supplier-version-accepted}
\renewcommand*{\RestatedName}{\hyperref{contract.pdf}{}{claim-reductions/req:supplier-version-accepted}{Requirement~18.5}}%
\begin{digestrequirement}[Refusal of a supplier artefact of an unaccepted version]
The engine refuses a supplier artefact whose declared module version
\iface{module_version} is not in the caller's declared set of accepted
versions \iface{accepted_receiver_versions}, and one that does not carry the
supplying module's validation results; both refusals carry the reason
\iface{receiver_spec_not_accepted}, the reason the ledger refusals of
Requirement~\ref{req:supplier-ledger-refusals} carry as well.
\end{digestrequirement}

% claim reductions/req:supplier-presence-rules
\digestlabel{req:supplier-presence-rules}{18.6}\pdfbookmark[2]{Requirement 18.6 Supplier presence rules that omit no imported premise}{digest-reductions/req:supplier-presence-rules}
\renewcommand*{\RestatedName}{\hyperref{contract.pdf}{}{claim-reductions/req:supplier-presence-rules}{Requirement~18.6}}%
\begin{digestrequirement}[Supplier presence rules that omit no imported premise]
The supplier's presence rules that the engine reads omit a key of the closure
of the facts a reduction instance imports from that supplier only where
the facts imported through it hold of the declaration without it; a change
of those rules, which takes effect in the engine's check without an edit of
this module, must itself meet this requirement.
\end{digestrequirement}

% claim reductions/req:recv-row-provenance
\digestlabel{req:recv-row-provenance}{18.7}\pdfbookmark[2]{Requirement 18.7 Omission of the supplier row for a receiver artefact the run built}{digest-reductions/req:recv-row-provenance}
\renewcommand*{\RestatedName}{\hyperref{contract.pdf}{}{claim-reductions/req:recv-row-provenance}{Requirement~18.7}}%
\begin{digestrequirement}[Omission of the supplier row for a receiver artefact the run built]
The engine omits the row of Assumption~\ref{ass:recv} from the assumption
ledger \iface{assumption_ledger} of an emission only where the
reduction instance read at least one supplier artefact and, for every
supplier artefact it read, all of the following hold:
(i) the supplying module is the receiver, and the artefact is one that the
part of the engine that ran this module obtained in the same run by running
the receiver module, as installed where the run runs, on the receiver
declaration the run holds, and handed to this module as it was emitted;
(ii) the module version \iface{module_version} the artefact declares is the
one that installed module states for the artefacts it emits, and the
artefact passes Requirements~\ref{req:supplier-version-accepted}
and~\ref{req:supplier-ledger-refusals};
(iii) the record of the run carries that artefact's own assumption ledger;
(iv) the artefact records that its declaration has a model admissible
exactly (\xref{receiver/req:exact-route-decision}\,(iv)), and it carries the
evaluation label \iface{rigorous} (\xref{receiver/req:evaluation-label}),
every per-setting weight having been formed by the rigorous evaluation of
\xref{receiver/req:rigorous-weight}.
There clause~(i) of Assumption~\ref{ass:recv} holds by construction, the
acceptance of the version being the effective status of the receiver's
claims in the run's closure, which the record of the run reports, and
clause~(ii) holds of the facts the instance imports, each a theorem of the
receiver about the receiver instrument the declaration is, holding of the
numbers emitted under the ledger~(iii) carries and resting on no premise
about arithmetic. That (i) and (iii) hold is established by the part of the
engine that ran the receiver module and docked its artefact into the
instance, and is passed to this module beside the artefact
(\iface{receiver_provenance}); the engine never reads it from the artefact.
The engine decides (ii) and (iv) on the artefact. Where it omits the row, it
records the omission with the artefact's version
(\iface{receiver_row_omitted_on_provenance}), and the supplier's keys, which
ride that row where it is emitted
(Requirements~\ref{req:supplier-ledger-pass-through}
and~\ref{C-req:supplier-eig-record}), reach the record of the run from the
artefact's own ledger by~(iii). Where \iface{receiver_provenance} is not
passed, or a condition fails for any artefact read, the row is emitted as
Requirement~\ref{req:assumption-ledger} says.
\end{digestrequirement}

% claim reductions/req:output-state-statistics
\digestlabel{req:output-state-statistics}{18.8}\pdfbookmark[2]{Requirement 18.8 Output-state model of the test statistics}{digest-reductions/req:output-state-statistics}
\renewcommand*{\RestatedName}{\hyperref{contract.pdf}{}{claim-reductions/req:output-state-statistics}{Requirement~18.8}}%
\begin{digestrequirement}[Output-state model of the test statistics]
The engine emits no state-independent test fraction, and it declares the
test-statistics model \iface{test_statistics_model} to be
\iface{output_state}.
\end{digestrequirement}

\part*{Flag-state and weight-preserving squashers}\pdfbookmark[0]{Flag-state and weight-preserving squashers}{digest-part-3}

\section*{19\quad Flag-state targets and their chain}\pdfbookmark[1]{19 Flag-state targets and their chain}{digest-group-reductions/wpfss/def:targets}

% claim reductions/wpfss/def:targets
\digestlabel{def:targets}{19.1}\pdfbookmark[2]{Definition 19.1 Certified weights and the flag-state and weight-preserving targets}{digest-reductions/wpfss/def:targets}
\renewcommand*{\RestatedName}{\hyperref{contract.pdf}{}{claim-reductions/wpfss/def:targets}{Definition~19.1}}%
\begin{digestclaim}[Certified weights and the flag-state and weight-preserving targets]
Throughout this definition and the claims of this submodule that use it,
$\Hs{\Bsys}=\Hs{\Bsys_{\leqN}}\dsum\Hs{\Bsys_{\gtN}}$
is an orthogonal two-block decomposition with projections $\Pleq$ and $\Pgt$
--- a photon-number cutoff in the optical instance, where
$\Hs{\Bsys_{\leqN}}=\bigoplus_{m\le N}\Hs{\Chan_m}$ is finite-dimensional ---
and every $M_v$ is block diagonal for it,
$M_v=\blk{M}{\leqN}_v\dsum\blk{M}{\gtN}_v$. $\Hs{\Flag}=\mathbb{C}^{|\Vset|}$
carries the orthonormal basis $\{\ket{v}\}_{v\in\Vset}$, and
$\Hs{\Qsys}=\Hs{\Bsys_{\leqN}}\dsum\Hs{\Flag}$ with block projections
$\Pleq$ (the same symbol on either space) and $\Pflag$.

A family $(p_v)_{v\in\Vset}$ of
\emph{certified weights}
for $M$ at the cutoff consists of non-negative numbers $p_v$ --- on the
emission path, integer multiples of $2^{-\kappa}$, the $p_{(y,t)}$ of
\eqref{eq:weightroute} --- with
\begin{equation*}\tag{\getrefnumber{C-eq:certweights}}\digesteqlabel{eq:certweights}
  \blk{M}{\gtN}_v\;\succeq\;p_v\,\idon{\Bsys_{\gtN}}
  \qquad\text{for every }v\in\Vset ,
\end{equation*}
and we write $p\coloneqq\sum_vp_v$ for their
\emph{total weight}.
Define
\begin{align*}
  \Ffss_v&\coloneqq\blk{M}{\leqN}_v\dsum\dyad{v}_{\Flag},
  \tag{\getrefnumber{C-eq:FSS}}\digesteqlabel{eq:FSS}\\
  \Fwp_v&\coloneqq\blk{M}{\leqN}_v\dsum
     \big(p_v\,\idon{\Flag}+(1-p)\,\dyad{v}_{\Flag}\big).
  \tag{\getrefnumber{C-eq:WPFSS}}\digesteqlabel{eq:WPFSS}
\end{align*}
$\Ffss$ is the
\emph{flag-state target}
and $\Fwp$ the
\emph{weight-preserving target},
the multi-bound weight-preserving flag-state-squasher target, with one
certified weight per outcome. The
\emph{single-bound case}
is the one in which only one $p_v$ is non-zero, and $\Ffss$ is the case
$p_v\equiv0$.
\end{digestclaim}

% claim reductions/wpfss/prop:chain
\digestlabel{prop:chain}{19.4}\pdfbookmark[2]{Proposition 19.4 Flag-state squashers: existence, chain and encoded weights}{digest-reductions/wpfss/prop:chain}
\begin{mainclaim}
\renewcommand*{\RestatedName}{\hyperref{contract.pdf}{}{claim-reductions/wpfss/prop:chain}{Proposition~19.4}}%
\begin{digestclaim}[Flag-state squashers: existence, chain and encoded weights]
Let $(p_v)$ be certified weights with $p<1$ and define the channels
\begin{align*}
  \Lwp(\rho)
   &\coloneqq\Pleq\rho\,\Pleq\ \dsum\
      \sum_{v}\tr\!\big[\Mwp_v\,\Pgt\rho\,\Pgt\big]\,\dyad{v}_{\Flag},
   \qquad
   \Mwp_v\coloneqq\frac{\blk{M}{\gtN}_v-p_v\,\idon{\Bsys_{\gtN}}}{1-p},
   \tag{\getrefnumber{C-eq:Lwp}}\digesteqlabel{eq:Lwp}\\
  \Lfss(\rho)
   &\coloneqq\Pleq\rho\,\Pleq\ \dsum\
      \sum_{v}\tr\!\big[\blk{M}{\gtN}_v\,\Pgt\rho\,\Pgt\big]\,\dyad{v}_{\Flag},
   \tag{\getrefnumber{C-eq:Lfss}}\digesteqlabel{eq:Lfss}\\
  \Lnoise(\sigma)
   &\coloneqq(1-p)\,\sigma+p\,\Pleq\sigma\,\Pleq
      +\tr\!\big[\Pflag\sigma\big]\,\mu,
   \qquad \mu\coloneqq\sum_vp_v\dyad{v}_{\Flag}.
   \tag{\getrefnumber{C-eq:Lnoise}}\digesteqlabel{eq:Lnoise}
\end{align*}
Then:
\begin{enumerate}[label=(\roman*),leftmargin=3.0em]
\item \emph{(existence, with the cutoff projectors transported)}
  \begin{equation*}\tag{\getrefnumber{C-eq:transportproj}}\digesteqlabel{eq:transportproj}
    \big(M;\Pleq,\Pgt\big)\dsq_{\Lwp}\big(\Fwp;\Pleq,\Pflag\big),
    \qquad
    \big(M;\Pleq,\Pgt\big)\dsq_{\Lfss}\big(\Ffss;\Pleq,\Pflag\big);
  \end{equation*}
\item \emph{(the chain)} $\Fwp\dsq_{\Lnoise}\Ffss$, and the channels compose
  exactly, $\Lfss=\Lnoise\circ\Lwp$ on all of $\Tcl{\Bsys}$, so that
  \begin{equation*}\tag{\getrefnumber{C-eq:ladder}}\digesteqlabel{eq:ladder}
    M\ \dsq_{\Lwp}\ \Fwp\ \dsq_{\Lnoise}\ \Ffss\ \dsq_{\Lmp}\ \Ppvm ;
  \end{equation*}
\item \emph{(the encoded constraint)} $\Fwp_v\succeq p_v\,\Pflag$ as operators
  on $\Qsys$, so $\tr[\Fwp_v\sigma]\ge p_v\tr[\Pflag\sigma]$ holds on
  \emph{every} squashed state, reachable or not. Under $\Ffss$ the same
  inequality is only a pull-back constraint: it holds on every reachable state,
  since
  $\Lfss^{\dagger}\big(\Ffss_v-p_v\Pflag\big)
  =\blk{M}{\leqN}_v\dsum\big(\blk{M}{\gtN}_v-p_v\idon{\Bsys_{\gtN}}\big)
  \succeq0$, and it \emph{fails} at $\sigma=\dyad{v'}_{\Flag}$ for any $v'\ne v$
  with $p_v>0$;
\item \emph{(no reverse relation in general)} $\Ffss\dsq\Fwp$ can fail: at
  $\Hs{\Bsys_{\leqN}}=\Hs{\Bsys_{\gtN}}=\mathbb{C}$, $\Vset=\{1,2\}$,
  $M_1=\tfrac12\dsum1$, $M_2=\tfrac12\dsum0$, $p_1=\tfrac12$, $p_2=0$, one has
  $\|\Ffss_2\|=1>\tfrac12=\|\Fwp_2\|$, which
  Lemma~\ref{lem:obstructions}\,(i) forbids.
\end{enumerate}
\end{digestclaim}
\end{mainclaim}

% claim reductions/wpfss/def:routes
\digestlabel{def:routes}{19.5}\pdfbookmark[2]{Definition 19.5 The encoded and the licensed route}{digest-reductions/wpfss/def:routes}
\renewcommand*{\RestatedName}{\hyperref{contract.pdf}{}{claim-reductions/wpfss/def:routes}{Definition~19.5}}%
\begin{digestclaim}[The encoded and the licensed route]
The \emph{encoded} route uses the weight-preserving target $\Fwp$ and
emits $\CList=\varnothing$; the \emph{licensed} route uses the
flag-state target $\Ffss$ and emits the entry of
Theorem~\ref{thm:licence}.
\end{digestclaim}

% claim reductions/wpfss/prop:weight-monotone
\digestlabel{prop:weight-monotone}{19.7}\pdfbookmark[2]{Proposition 19.7 Monotonicity of the weight-preserving target in its weights}{digest-reductions/wpfss/prop:weight-monotone}
\renewcommand*{\RestatedName}{\hyperref{contract.pdf}{}{claim-reductions/wpfss/prop:weight-monotone}{Proposition~19.7}}%
\begin{digestclaim}[Monotonicity of the weight-preserving target in its weights]
Let $(p_v)_{v\in\Vset}$ be certified weights for $M$ with
total weight $p$, let $p'_v\in[0,p_v]$ for every $v\in\Vset$, and
write $p'\coloneqq\sum_vp'_v$. Write $\Fwp[p]$ and $\Fwp[p']$ for the
weight-preserving targets \eqref{eq:WPFSS}
built from $(p_v)$ and from $(p'_v)$.
\begin{enumerate}[label=(\roman*),leftmargin=3.0em]
\item $(p'_v)$ are certified weights for $M$, and $p'\le p$; in particular,
  where $p<1$ the construction of Proposition~\ref{prop:chain} runs
  unchanged with $(p'_v)$.
\item If $p<1$, put
  \begin{equation*}\tag{\getrefnumber{C-eq:weightmono}}\digesteqlabel{eq:weightmono}
    \Psi(\sigma)\coloneqq\lambda\,\sigma+(1-\lambda)\,\Pleq\sigma\Pleq
     +\tr\!\big[\Pflag\sigma\big]\,\nu,
    \qquad
    \lambda\coloneqq\frac{1-p}{1-p'},
    \qquad
    \nu\coloneqq\sum_v\frac{p_v-p'_v}{1-p'}\dyad{v}_{\Flag}.
  \end{equation*}
  Then $\Psi$ is a channel and $\Fwp[p]\dsq_{\Psi}\Fwp[p']$, so by
  Corollary~\ref{cor:monotone} the smaller-weight target is the cleaner one.
\end{enumerate}
\end{digestclaim}

\section*{20\quad The weight route}\pdfbookmark[1]{20 The weight route}{digest-group-reductions/wpfss/lem:descent}

% claim reductions/wpfss/lem:descent
\digestlabel{lem:descent}{20.1}\pdfbookmark[2]{Lemma 20.1 Descent of a dark-count matrix through a coarse-graining}{digest-reductions/wpfss/lem:descent}
\renewcommand*{\RestatedName}{\hyperref{contract.pdf}{}{claim-reductions/wpfss/lem:descent}{Lemma~20.1}}%
\begin{digestclaim}[Descent of a dark-count matrix through a coarse-graining]
Let $R$ be column-stochastic on $2^{[K]}$ and $\cg\colon2^{[K]}\to\Yset$
surjective. Say that \emph{$R$ descends through $\cg$} if the quantity
$P(y\,|\,S')\coloneqq\sum_{S\in\cg^{-1}(y)}R_{S|S'}$ depends on $S'$ only
through $\cg(S')$; the resulting matrix
$R^{\cg}_{y|y'}\coloneqq P(y\,|\,S')$ for any $S'\in\cg^{-1}(y')$ is then well
defined and column-stochastic. In that case, for every family of operators
$\{X_{S}\}$,
\begin{equation*}\tag{\getrefnumber{C-eq:descend}}\digesteqlabel{eq:descend}
  \sum_{S\in\cg^{-1}(y)}\ \sum_{S'}R_{S|S'}X_{S'}
  \;=\;\sum_{y'\in\Yset}R^{\cg}_{y|y'}\ \sum_{S'\in\cg^{-1}(y')}X_{S'},
\end{equation*}
so applying $R$ and then coarse-graining agrees with coarse-graining and then
applying $R^{\cg}$. A sufficient condition, which the canonical
threshold-detector declaration satisfies, is that $R$ be \emph{monotone}
($R_{S|S'}>0\Rightarrow S'\subseteq S$) and that, with
$\MCK\coloneqq\{S:|S|\ge2\}$,
\begin{equation*}\tag{\getrefnumber{C-eq:canonicalcg}}\digesteqlabel{eq:canonicalcg}
  \cg\ \text{is injective on }2^{[K]}\setminus\MCK,
  \qquad
  \cg\ \text{is \emph{constant} on }\MCK,\ \text{say }\cg(\MCK)=\{y_{\mc}\},
  \qquad
  \cg^{-1}(y_{\mc})=\MCK ;
\end{equation*}
the case $\cg=\mathrm{id}$ is trivial.

On declared data both the hypothesis and the conclusion are read exactly. The
declared $R$ is exactly column-stochastic, and the descent
condition, that $P(y|S')=P(y|S'')$ for every $y$ and every pair $S',S''$ with
$\cg(S')=\cg(S'')$, is an identity between the dyadic rationals the
declaration carries and is decided as one
(Definition~\ref{def:substrate}).
The representative is declared, not chosen:
$R^{\cg}_{y|y'}\coloneqq P(y|S'_{y'})$ with $S'_{y'}$ the first element of
$\cg^{-1}(y')$ in the declared order of $2^{[K]}$, so the descended
matrix is a determinate object and not one of several. Column-stochasticity of
$R^{\cg}$ needs no second decision: the identity
$\sum_yR^{\cg}_{y|y'}=\sum_SR_{S|S'}$ of the proof regroups one finite sum,
and in the exact substrate a regrouped sum \emph{is} the sum it regroups. So \eqref{eq:descend} holds of the declared
numbers as an exact identity.
\end{digestclaim}

% claim reductions/wpfss/def:weight-route
\digestlabel{def:weight-route}{20.3}\pdfbookmark[2]{Definition 20.3 The weight route}{digest-reductions/wpfss/def:weight-route}
\renewcommand*{\RestatedName}{\hyperref{contract.pdf}{}{claim-reductions/wpfss/def:weight-route}{Definition~20.3}}%
\begin{digestclaim}[The weight route]
Let a receiver artifact supply, for each setting $t\in\Tset$, preserved
photon-number blocks, $m\le N$, in one of the two input forms below, for a
conditional POVM $\{M_{y|t}\}_{y\in\Yset}$ on $\Chan$ block diagonal at the
cutoff; a
dark-count matrix $R$ on $2^{[K]}$ that descends through the
declared coarse-graining $\cg$ with descended matrix $R^{\cg}$
(Lemma~\ref{lem:descent}); class weights $\{q_t\}_{t\in\Tset}$ as in
Definition~\ref{def:joint-family}; and a slot
$\{(\mathcal{O}_i,(p^{(t)}_i)_{t\in\Tset})\}_{i=1}^{r}$ of certified
high-photon weights, one per entry and per setting of the same $\Tset$;
every declared entry and scalar lies on the declared lattice
\eqref{eq:latinv}, and every one of the supplier's facts about them is held
under Assumption~\ref{ass:recv} (Proposition~\ref{prop:route}\,(iv)). The
\emph{weight route}
is the following three steps, in this order, each with the claim that
licenses it. Steps~1--2 are classical post-processing of the physical outcome
and precede the first squashing link.
\begin{enumerate}[leftmargin=2.2em]
\item \emph{Dark counts, on the emitted alphabet.} Replace the supplied
  conditional POVM by
  \begin{equation*}\tag{\getrefnumber{C-eq:darkstep}}\digesteqlabel{eq:darkstep}
    \widetilde M_{y|t}\;\coloneqq\;\sum_{y'\in\Yset}R^{\cg}_{y|y'}\,M_{y'|t},
    \qquad y\in\Yset,\ t\in\Tset .
  \end{equation*}
  This is post-processing by a column-stochastic matrix
  (Lemma~\ref{lem:linear}\,(i)), so the result is a conditional POVM at each
  $t$ with the same photon-number block structure, and
  Lemma~\ref{lem:descent} shows it equals the coarse-graining of the raw-level
  composition.
\item \emph{Joint alphabet.} Form the
  joint family of Definition~\ref{def:joint-family}\,(b),
  $M_{(y,t)}\coloneqq q_t\,\widetilde M_{y|t}$ on the joint alphabet
  $\Vset=\Yset\times\Tset$;
  this is a POVM by \eqref{eq:jointfamily} (Lemma~\ref{lem:linear}\,(ii);
  the weights are the measuring party's own, Assumption~\ref{ass:probe}), block diagonal at
  the cutoff, and the class-marginal identity \eqref{eq:classmarginal}
  holds of it block by block. Its preserved blocks $\blk{M}{=m}_{(y,t)}$,
  $m\le N$, are the
  \emph{post-processed datum}.
\item \emph{The instance.} The post-processed joint family, with the routed
  weights below, is the datum, the hypotheses and the target data of the
  \emph{flat instance} (Proposition~\ref{prop:flat-instance}), and its
  emission is that instance's.
\end{enumerate}

\emph{The two input forms.} The supplied preserved blocks are declared in one
of two forms. In the
\emph{bare form}
they are the blocks $\blk{M}{=m}_{y|t}$, $m\le N$, of the conditional POVM
above, and step~1 acts on them. In the
\emph{dark-composed form}
they are a declared family $\{\blk{\widetilde M}{=m}_{y|t}\}_{y\in\Yset}$,
$m\le N$, per setting $t$, each block Hermitian, positive semidefinite and
complete on its block, which the model takes below the cutoff after the
dark counts: in this form the family $\widetilde M_{y|t}$ of step~1 is
defined as these blocks below the cutoff and by \eqref{eq:darkstep} above
it, step~1 is not applied to the supplied blocks, and the post-processed
datum is formed from them by step~2 alone. Where no snap intervened between
composing and supplying, this form is the bare form composed in advance
(Proposition~\ref{prop:precomposed-blocks}\,(ii)); where the declaring party
snapped the composed blocks, they are its snapped family, and this
definition asserts of them no relation to a bare family. In either form $R^{\cg}$ is
formed and its conditions are decided as below, \eqref{eq:darkstep} defines
the family above the cutoff, $\blk{\widetilde M}{\gtN}_{y|t}
=\sum_{y'}R^{\cg}_{y|y'}\blk{M}{\gtN}_{y'|t}$, and the routed weights below
read $R^{\cg}$ as they stand. In either form $\{\widetilde M_{y|t}\}_y$ is a
POVM at each $t$, block diagonal at the cutoff: below it by step~1 in the
bare form and by the declared properties of the supplied blocks in the
dark-composed form, and above it by step~1.

\emph{Exactness of steps~1 and~2.} The two steps are written on the declared
numbers, and their POVM conclusions hold \emph{exactly} of them. The declared
$R^{\cg}$ is exactly column-stochastic with non-negative entries, decided on
entry, and $\{q_t\}$ is
exactly a probability distribution, as Definition~\ref{def:joint-family}
requires of class weights; and the combination is formed in the exact
substrate of Lemma~\ref{lem:exactpsd} (Definition~\ref{def:substrate}), so
$\{\widetilde M_{y|t}\}_y$ and $\{M_{(y,t)}\}$ are POVMs and every claim of
this submodule reads them as such. What the two steps can nevertheless do is
leave the lattice: a product of two dyadic rationals adds significand widths,
so the post-processed blocks $\blk{M}{=m}_{(y,t)}$ need not be binary64
numbers even though every factor is. That is the one place in this submodule
where exact arithmetic outruns the representation; it is where the emission
normaliser of Definition~\ref{def:emitnorm} acts, and it acts on numbers this
module computed and is about to declare, never on numbers it was handed.

\emph{The routed weights.} They follow the same route, one bounded letter at a
time. Fix an entry $i$ whose bounded set is exactly a fibre,
$\mathcal{O}_i=\cg^{-1}(y_i)$, and whose letter is absorbing,
$R^{\cg}_{y_i|y_i}=1$; both are declared input validity conditions, and both
are decided exactly. Write
\begin{equation*}\tag{\getrefnumber{C-eq:weightroute}}\digesteqlabel{eq:weightroute}
\begin{gathered}
  \hat p_{(y_i,t)}\ \coloneqq\
  q_t\Big(\big(1-r_{\min,i}\big)p^{(t)}_i+r_{\min,i}\Big),
  \qquad
  p_{(y_i,t)}\ \coloneqq\ \lat{\hat p_{(y_i,t)}},\\
  r_{\min,i}\coloneqq\min_{y'\ne y_i}R^{\cg}_{y_i|y'}
  \ \ \big(\coloneqq0\ \text{when that minimum is over the empty set}\big),
\end{gathered}
\end{equation*}
with $p_{(y,t)}\coloneqq0$ for every other $(y,t)$; the family
$(p_v)_{v\in\Vset}$ is the
\emph{routed weights},
and $p\coloneqq\sum_{v\in\Vset}p_v$. Here
\begin{equation*}\tag{\getrefnumber{C-eq:latdown}}\digesteqlabel{eq:latdown}
  \lat{x}\;\coloneqq\;2^{-\kappa}\big\lfloor2^{\kappa}x\big\rfloor
\end{equation*}
is the largest integer multiple of $2^{-\kappa}$ that is at most $x$: the
\emph{downward} rounding onto the declared lattice, computed exactly on the
rational $\hat p_{(y_i,t)}$ and equal to it whenever it is already a lattice
number. The routed weights are \emph{certified} when the element
$M_{(y_i,t)}$ of step~2 satisfies, on every channel block above the cutoff,
\begin{equation*}\tag{\getrefnumber{C-eq:certroute}}\digesteqlabel{eq:certroute}
  \blk{M}{\gtN}_{(y_i,t)}\ \succeq\ p_{(y_i,t)}\,\idon{}
  \qquad\text{for every }i\text{ and every }t\in\Tset ,
\end{equation*}
which Proposition~\ref{prop:route} proves from the declared input validity
conditions together with the absorbing hypothesis those conditions impose,
and from nothing else.

\emph{The direction of the rounding.} The emitted weight is $p_{(y_i,t)}$ and
not $\hat p_{(y_i,t)}$, and the direction of that rounding is the whole of
what makes it safe: it is the one emitted number of this submodule that a
downstream inequality consumes as a \emph{bound}, the certified-weight
inequality of Definition~\ref{def:targets} is preserved when a certified
weight is lowered and, when one is raised, is no longer proved and can be
false, and $\lat{\cdot}$ only lowers (Proposition~\ref{prop:route};
Proposition~\ref{prop:weight-monotone}\,(i)). Four features of the route are relied on, and none of them may
be dropped.
\emph{(a)} $r_{\min}$ is indexed by the bounded letter: it is a property of row
$y_i$ of $R^{\cg}$ read across the columns $y'\ne y_i$, and one entry's value is
not another's.
\emph{(b)} The absorbing condition is a hypothesis, not a consequence:
without it the bound \eqref{eq:certroute} for the weights of
\eqref{eq:weightroute} is no longer proved.
\emph{(c)} A non-zero $r_{\min,i}$ can occur only at $r=1$
(Corollary~\ref{cor:oneabsorbing}), so for $r\ge2$ \emph{every} entry
contributes the un-improved weight $\lat{q_tp^{(t)}_i}$, which is sound,
merely not an improvement.
\emph{(d)} $p_{(y_i,t)}\ge0$ needs no clamp: $\hat p_{(y_i,t)}$ is $q_t>0$
times $(1-r_{\min,i})p^{(t)}_i+r_{\min,i}$, a convex combination of
$p^{(t)}_i\ge0$ and $1$, since $r_{\min,i}\in[0,1]$ as an entry of the
column-stochastic $R^{\cg}$ or as $0$; and the largest lattice number at most
a non-negative rational is itself non-negative.
\end{digestclaim}

% claim reductions/wpfss/prop:route
\digestlabel{prop:route}{20.4}\pdfbookmark[2]{Proposition 20.4 Certified weights routed from the receiver artifact}{digest-reductions/wpfss/prop:route}
\renewcommand*{\RestatedName}{\hyperref{contract.pdf}{}{claim-reductions/wpfss/prop:route}{Proposition~20.4}}%
\begin{digestclaim}[Certified weights routed from the receiver artifact]
Let $\{M_{y|t}\}_{y\in\Yset}$ be, for each $t\in\Tset$, a POVM on $\Chan$ block
diagonal at the cutoff, and let $R$ descend through $\cg$ with descended matrix
$R^{\cg}$ (Lemma~\ref{lem:descent}). Let $\{q_t\}$ be a probability
distribution, and let $\{(\mathcal{O}_i,(p^{(t)}_i)_{t\in\Tset})\}_{i=1}^{r}$
be the declared slot, with $\mathcal{O}_i=\cg^{-1}(y_i)$ for pairwise distinct
letters $y_1,\dots,y_r$, with $0\le p^{(t)}_i\le1$, and with the per-setting
guarantee $\blk{M}{\gtN}_{y_i|t}\succeq p^{(t)}_i\,\idon{}$ for every
$t\in\Tset$ and every $i$. Suppose each bounded letter is \emph{absorbing},
\begin{equation*}\digesteqlabel{eq:absorbing}
  R^{\cg}_{y_i|y_i}=1
  \qquad\text{for every }i=1,\dots,r,
  \tag{absorbing}
\end{equation*}
and put $r_{\min,i}\coloneqq\min_{y'\ne y_i}R^{\cg}_{y_i|y'}$, with the
convention $r_{\min,i}\coloneqq0$ when the minimum is over the empty set (that
is, when $\Yset=\{y_i\}$), the conservative value, which reduces
\eqref{eq:routebound} to the un-improved bound. Then:
\begin{enumerate}[label=(\roman*),leftmargin=2.4em]
\item The joint-alphabet POVM
  $M_{(y,t)}\coloneqq q_t\sum_{y'}R^{\cg}_{y|y'}M_{y'|t}$ of
  Definitions~\ref{def:joint-family} and~\ref{def:weight-route} satisfies, for
  every $i$ and every $t\in\Tset$,
  \begin{equation*}\tag{\getrefnumber{C-eq:routebound}}\digesteqlabel{eq:routebound}
    \blk{M}{\gtN}_{(y_i,t)}\ \succeq\
    q_t\Big(\big(1-r_{\min,i}\big)\,p^{(t)}_i+r_{\min,i}\Big)\,\idon{}
    \;=\;\hat p_{(y_i,t)}\,\idon{} .
  \end{equation*}
  A slot whose weights do not depend on the setting, $p^{(t)}_i=p_i$ for
  every $t$ (the $t$-uniform form, which every one-setting slot has), is the
  special case, and there \eqref{eq:routebound} is the bound in $p_i$.
\item The \emph{emitted} weight of \eqref{eq:weightroute} is certified, because
  the rounding runs downwards. Every number on the right of
  \eqref{eq:routebound} lies on the declared lattice: $q_t$ and $p^{(t)}_i$ by
  declaration, and $r_{\min,i}$ as an entry of $R^{\cg}$, which is a sum of
  declared entries and lies in $[0,1]$. So $\hat p_{(y_i,t)}$ is an exact
  element of $\Fld$ formed from them; and
  $p_{(y_i,t)}=\lat{\hat p_{(y_i,t)}}\le\hat p_{(y_i,t)}$ by
  \eqref{eq:latdown}, whence
  \begin{equation*}
    \blk{M}{\gtN}_{(y_i,t)}\ \succeq\ p_{(y_i,t)}\,\idon{} ,
  \end{equation*}
  which is \eqref{eq:certroute}; so the emitted $(p_v)$ are
  certified weights in the sense of Definition~\ref{def:targets} and they lie on the declared
  lattice.
\item Under the sufficient condition \eqref{eq:canonicalcg} of
  Lemma~\ref{lem:descent} and with the single entry $\mathcal{O}_1=\MCK$,
  \eqref{eq:absorbing} holds at $y_1=y_{\mc}$, and its per-setting hypothesis
  is the supplier's slot guarantee
  $\sum_{S\in\MCK}\blk{M}{\gtN}_{S|t}\succeq p^{(t)}_1\,\idon{}$ at each $t$,
  read through $\cg^{-1}(y_{\mc})=\MCK$.
\item On the emission path the hypotheses about the supplier's objects hold
  under Assumption~\ref{ass:recv}: the conditional families, and their
  block diagonality in the photon number, block by block and so in
  particular at the cutoff, are at every setting those of the coarse-grained
  elements of the setting family on which the supplier states its
  weights, by item~(a) of \xref{receiver/def:setting-family} and
  \xref{receiver/thm:coarse-graining}; for a receiver without probe inputs
  that family is the one-setting family of
  \xref{receiver/cor:unconditioned-family}, and the conditional families and
  their block diagonality are given by the receiver's
  \xref{receiver/thm:click-povm}\,(iii) and
  \xref{receiver/thm:coarse-graining}; the
  dark-count matrix $R$, and the block diagonality in the photon number of
  the elements with the dark counts composed, by
  \xref{receiver/detector-imperfections/def:darkadm} and
  \xref{receiver/detector-imperfections/thm:dark-counts}; so every element of
  the joint family of (i), and of the post-processed joint family that
  Definition~\ref{def:weight-route} forms in either of its forms, is block
  diagonal in the photon number in every block, below the cutoff as above
  it; and the slot, its
  per-setting guarantee (the first conclusion of
  \xref{receiver/thm:emitted-bound}) and $0\le p^{(t)}_i\le1$,
  $\sum_ip^{(t)}_i\le1$ at every $t$ by \xref{receiver/def:slot},
  \xref{receiver/thm:emitted-bound} and \xref{receiver/prop:slotvalid}. The
  other hypotheses are not supplier facts: that $R$ is column-stochastic and descends through $\cg$, that
  $\mathcal{O}_i=\cg^{-1}(y_i)$ and \eqref{eq:absorbing} are decided exactly
  on the declared data, a declaration failing one being
  refused, and $\{q_t\}$ are the declared class weights
  (Definition~\ref{def:joint-family}).
\end{enumerate}
\end{digestclaim}

% claim reductions/wpfss/prop:precomposed-blocks
\digestlabel{prop:precomposed-blocks}{20.5}\pdfbookmark[2]{Proposition 20.5 Preserved blocks supplied with the dark counts already composed}{digest-reductions/wpfss/prop:precomposed-blocks}
\renewcommand*{\RestatedName}{\hyperref{contract.pdf}{}{claim-reductions/wpfss/prop:precomposed-blocks}{Proposition~20.5}}%
\begin{digestclaim}[Preserved blocks supplied with the dark counts already composed]
Let the data of Definition~\ref{def:weight-route} be given in the
bare form, with supplied preserved blocks $B^{(m)}_{y|t}$, $m\le N$,
and put
\begin{equation*}\tag{\getrefnumber{C-eq:precomposed}}\digesteqlabel{eq:precomposed}
  \widetilde B^{(m)}_{y|t}\;\coloneqq\;\sum_{y'\in\Yset}R^{\cg}_{y|y'}\,B^{(m)}_{y'|t}.
\end{equation*}
\begin{enumerate}[label=(\roman*),leftmargin=3.0em]
\item For every $t$ and $m$, $\{\widetilde B^{(m)}_{y|t}\}_{y\in\Yset}$ is a
  family of Hermitian positive semidefinite matrices complete on its block.
\item The flat instance of Proposition~\ref{prop:flat-instance} formed from
  the blocks $\widetilde B$ in the dark-composed form, and the one
  formed from the blocks $B$ in the bare form, have the same
  handed datum, $A^{(m)}_{(y,t)}=q_t\widetilde B^{(m)}_{y|t}$, the same
  hypotheses (H), model map, routed weights,
  scalar values, channel, transported set and
  entries; so they are one reduction instance, and the
  emission normalizer gives them the same declared datum,
  declared model and record.
\item Let $\{G^{(m)}_{y|t}\}_{y\in\Yset}$ be a reference of the bare
  blocks of setting $t$ with constant $1$, read through the
  entrywise enclosures
  $\big(C^{(m)}_{y|t},r^{(m)}_{y|t}\big)$
  (\xref{foundations/def:exact-reference}), and put
  $\widetilde G^{(m)}_{y|t}\coloneqq\sum_{y'}R^{\cg}_{y|y'}G^{(m)}_{y'|t}$.
  Then $\widetilde G$ is a reference with constant $1$, read through the
  enclosures
  $\big(\sum_{y'}R^{\cg}_{y|y'}C^{(m)}_{y'|t},\,\sum_{y'}R^{\cg}_{y|y'}r^{(m)}_{y'|t}\big)$,
  and it is an exact reference where $G$ is. If $B$ is the
  snapped family of a snap of $G$ and $B'$ that of a snap of
  $\widetilde G$ (\xref{foundations/def:snap}), then the handed data
  $A^{\mathrm{b}}_{(y,t)}=q_t\sum_{y'}R^{\cg}_{y|y'}B_{y'|t}$ of the bare
  form and $A^{\mathrm{d}}_{(y,t)}=q_tB'_{y|t}$ of the dark-composed form
  differ in general. The two flat instances agree in everything but their
  data, so their emissions differ only in their
  data sectors.
\end{enumerate}
\end{digestclaim}

% claim reductions/wpfss/cor:oneabsorbing
\digestlabel{cor:oneabsorbing}{20.6}\pdfbookmark[2]{Corollary 20.6 Dark-count sharpening at a single bounded letter only}{digest-reductions/wpfss/cor:oneabsorbing}
\renewcommand*{\RestatedName}{\hyperref{contract.pdf}{}{claim-reductions/wpfss/cor:oneabsorbing}{Corollary~20.6}}%
\begin{digestclaim}[Dark-count sharpening at a single bounded letter only]
Under \eqref{eq:absorbing}: if $r\ge2$ then $r_{\min,i}=0$ for \emph{every}
$i$, so \eqref{eq:routebound} reads
$\blk{M}{\gtN}_{(y_i,t)}\succeq q_tp^{(t)}_i\,\idon{}$ throughout and the
emitted weight is $\lat{q_tp^{(t)}_i}$. Consequently a positive
$r_{\min,i}$, the only case in which the dark counts sharpen a weight, can
occur for at most one bounded letter, and only when $r=1$.
\end{digestclaim}

\section*{21\quad The flat instance and its emission}\pdfbookmark[1]{21 The flat instance and its emission}{digest-group-reductions/wpfss/cor:fssinstance}

% claim reductions/wpfss/cor:fssinstance
\digestlabel{cor:fssinstance}{21.1}\pdfbookmark[2]{Corollary 21.1 Licensed flag-sector entry of a certified weight}{digest-reductions/wpfss/cor:fssinstance}
\renewcommand*{\RestatedName}{\hyperref{contract.pdf}{}{claim-reductions/wpfss/cor:fssinstance}{Corollary~21.1}}%
\begin{digestclaim}[Licensed flag-sector entry of a certified weight]
In the situation of Proposition~\ref{prop:chain} with the
flag-state target $\Ffss$, fix an outcome $o\in\Vset$ with certified
weight $p_o$ and put
\begin{equation*}\tag{\getrefnumber{C-eq:fssentry}}\digesteqlabel{eq:fssentry}
  \hat O\coloneqq\Ffss_o-p_o\,\Pflag,
  \qquad \bowtie\;=\;\succeq,\qquad c=0 .
\end{equation*}
Then \eqref{eq:certificate} holds, so Theorem~\ref{thm:licence} applies and the
single-entry list $\CList=\{(\hat O,\succeq,0)\}$ is licensed with the
admissible set
\begin{equation*}\tag{\getrefnumber{C-eq:fssattset}}\digesteqlabel{eq:fssattset}
  \AttSet_{\CList}(E',\Qsys E'')
  =\Big\{\Aev:\ \Aev^{\dagger}\big(\Ffss_o\otimes\idon{E''}\big)
  \;\succeq\;p_o\,\Aev^{\dagger}\big(\Pflag\otimes\idon{E''}\big)\Big\}.
\end{equation*}
\end{digestclaim}

% claim reductions/wpfss/prop:flat-instance
\digestlabel{prop:flat-instance}{21.2}\pdfbookmark[2]{Proposition 21.2 The flat instance of the post-processed receiver}{digest-reductions/wpfss/prop:flat-instance}
\renewcommand*{\RestatedName}{\hyperref{contract.pdf}{}{claim-reductions/wpfss/prop:flat-instance}{Proposition~21.2}}%
\begin{digestclaim}[The flat instance of the post-processed receiver]
Let the data of Definition~\ref{def:weight-route} pass every refusal rule of
this module's interface; let
$\{M_v\}_{v\in\Vset}$ be the post-processed joint family of its step~2 on
$\Chan$, with $\Hs{\Chan}=\bigoplus_{m\ge0}\Hs{\Chan_m}$, $\Hs{\Chan_{\leqN}}$
finite-dimensional, every $M_v$ block diagonal at the cutoff $N$; let $(p_v)$
be its routed weights with $p=\sum_vp_v<1$; and let a route, encoded or
licensed, be declared (Definition~\ref{def:routes}). Then the \emph{flat instance}
$\mathcal{J}_{\mathrm{flat}}$, given by
\begin{enumerate}[label=(I\arabic*),leftmargin=3.0em]
\item $\Bsys=\Chan$; the data sectors $(m)$, $0\le m\le N$, in that order,
  with $d_m=\dim\Hs{\Chan_m}$ and the default supports $\Vset^{(m)}=\Vset$;
  the handed datum $A^{(m)}_v\coloneqq\blk{M}{=m}_v$, formed from the declared
  blocks, $R^{\cg}$ and $\{q_t\}$ by steps~1--2 of
  Definition~\ref{def:weight-route} in the bare form, and from the
  declared blocks and $\{q_t\}$ by step~2 alone in the
  dark-composed form; the hypotheses (H): \emph{(H1)} for
  each $t\in\Tset$ the above-cutoff blocks $\{\blk{M}{\gtN}_{y|t}\}_{y\in\Yset}$
  of the supplied conditional POVM form a POVM on $\Chan_{\gtN}$
  (the supplier's fact, Assumption~\ref{ass:recv},
  Proposition~\ref{prop:route}\,(iv)); \emph{(H2)} the routed weights are
  certified, \eqref{eq:certroute} (Proposition~\ref{prop:route});
  \emph{(H3)} $p<1$, decided on entry; and the model map
  \begin{equation*}\tag{\getrefnumber{C-eq:flat-model}}\digesteqlabel{eq:flat-model}
    M_v[A']\;\coloneqq\;\bigoplus_{m\le N}A'^{(m)}_v\ \dsum\ \blk{M}{\gtN}_v,
    \qquad
    \blk{M}{\gtN}_{(y,t)}=q_t\sum_{y'\in\Yset}R^{\cg}_{y|y'}\blk{M}{\gtN}_{y'|t} ;
  \end{equation*}
\item the scalar sectors $(f)$, $f\in\Vset$, in the order of $\Vset$, each of
  dimension $1$, with
  \begin{equation*}\tag{\getrefnumber{C-eq:flat-scalars}}\digesteqlabel{eq:flat-scalars}
    \phi_f(v)\;\coloneqq\;p_v+(1-p)\,\delta_{vf}\quad\text{(encoded route)},
    \qquad
    \phi_f(v)\;\coloneqq\;\delta_{vf}\quad\text{(licensed route)},
  \end{equation*}
  so that $\Hs{\Qsys}=\Hs{\Chan_{\leqN}}\dsum\Hs{\Flag}$ and the target
  \eqref{eq:instance-target} is $F_v[A']=\Fwp_v$ resp.\ $\Ffss_v$ of
  Definition~\ref{def:targets} for the model $M[A']$;
\item $\Lambda=\PinS\circ\Lwp$ resp.\ $\PinS\circ\Lfss$ of
  Proposition~\ref{prop:chain}, independent of the datum, with the
  transported set $\{\Pleq,\Pflag\}\cup\{\Pi_{(m)}\}_{m\le N}\cup
  \{\Pi_{(f)}\}_{f\in\Vset}$ and pull-backs $\Pleq$, $\Pgt$, the
  photon-number projectors $\Pnumon{m}{\Chan}$, and $\Mwp_f$ resp.\
  $\blk{M}{\gtN}_f$ read as operators on $\Chan$ supported in $\Chan_{\gtN}$;
\item an empty entry list on the encoded route, and on the
  licensed route the flag-sector entries
  $\big(\Ffss_o-p_o\Pflag,\succeq,0\big)$ of Corollary~\ref{cor:fssinstance}
  for every bounded letter $o=(y_i,t)$; the class-marginal identity
  discharged,
\end{enumerate}
is a reduction instance over the joint alphabet of
Definition~\ref{def:weight-route}. Its handed model $M[A]$ is the
post-processed joint family $\{M_v\}$ itself where every $M_v$ is also block
diagonal for the decomposition
$\Hs{\Chan_{\leqN}}=\bigoplus_{m\le N}\Hs{\Chan_m}$, the photon-number blocks
the datum reads, and otherwise the family's pinching onto those blocks below
the cutoff; on the emission path the first case holds, the receiver's
elements being block diagonal in every photon number
(Proposition~\ref{prop:route}\,(iv)), so that there the handed model is the
post-processed joint family itself. Its declared model, at the
declared datum $A'$ the normalizer of
Definition~\ref{def:emitnorm} produces from $A$, is
\begin{equation*}\tag{\getrefnumber{C-eq:flat-declared}}\digesteqlabel{eq:flat-declared}
  M'_v\;\coloneqq\;M_v[A']\;=\;\bigoplus_{m\le N}A'^{(m)}_v\ \dsum\ \blk{M}{\gtN}_v ,
\end{equation*}
a POVM on $\Chan$ block diagonal at the cutoff, agreeing with the handed one
above it and satisfying every certified-weight inequality \eqref{eq:certroute}
the handed one does.
\end{digestclaim}

% claim reductions/wpfss/cor:emitted-flat
\digestlabel{cor:emitted-flat}{21.3}\pdfbookmark[2]{Corollary 21.3 Exact emission of the flat instance}{digest-reductions/wpfss/cor:emitted-flat}
\begin{mainclaim}
\renewcommand*{\RestatedName}{\hyperref{contract.pdf}{}{claim-reductions/wpfss/cor:emitted-flat}{Corollary~21.3}}%
\begin{digestclaim}[Exact emission of the flat instance]
Let the declared data pass every refusal rule of this module's
interface, let
$\mathcal{J}_{\mathrm{flat}}$ be the flat instance of
Proposition~\ref{prop:flat-instance} at the declared route with handed datum
$A$, and let $A'$ be the datum the normaliser of
Definition~\ref{def:emitnorm} declares. Then Theorem~\ref{thm:emit} holds of
$\mathcal{J}_{\mathrm{flat}}$, and in particular:
\begin{enumerate}[label=(\roman*),leftmargin=3.0em]
\item every number on the list of Theorem~\ref{thm:emit}\,(i) and, in
  addition, the descended matrix $R^{\cg}$, the constants $r_{\min,i}$, the
  exact weights $\hat p_v$ of \eqref{eq:weightroute}, the total $p$ and its
  complement $1-p$ are formed from the declared entries and scalars by
  finitely many operations in $\Fldi$, floors and comparisons; the
  emission displacement $\Delta_m$ per preserved block is the one
  emitted number that is not;
\item the emitted family, $\Fwp$ resp.\ $\Ffss$ of the declared model
  \eqref{eq:flat-declared}, satisfies conditions (a)--(e) of
  Definition~\ref{def:exactadm} together with \eqref{eq:latinv} of every
  entry and is decided so before it is emitted;
\item for every $v\in\Vset$,
  \begin{equation*}\tag{\getrefnumber{C-eq:emitted-flat}}\digesteqlabel{eq:emitted-flat}
    \Lambda^{\dagger}(F_v)\;=\;M'_v\;=\;\bigoplus_{m\le N}A'^{(m)}_v\ \dsum\ \blk{M}{\gtN}_v
    \qquad\text{exactly, on }\Chan ,
  \end{equation*}
  both sides operators on $\Chan$ itself: the physical system of this
  instance is the channel system and the identity is on it, not on an
  ampliation. $M'$ is a POVM on $\Chan$ block diagonal at the cutoff for which
  $(p_v)$ are certified weights in the sense of
  Definition~\ref{def:targets}, and Proposition~\ref{prop:chain},
  Proposition~\ref{prop:route} and Corollary~\ref{cor:fssinstance} hold
  verbatim of it;
\item the per-outcome weights are $p_v=\lat{\hat p_v}$, rounded \emph{down},
  and \eqref{eq:certroute} holds for $M'$ at those values
  (Proposition~\ref{prop:route}); the complement $1-p$ is \emph{recomputed}
  from the rounded weights and is never rounded on its own; a caller's
  constant off the lattice is rounded toward the enlarging side
  (Corollary~\ref{cor:entries}\,(iii)); and there are no other rounding
  directions;
\item where some $\Delta_m>0$ no bound relates a key length computed from
  $M'$ alone to one computed from the pre-snap post-processed family, or from
  the receiver's own output, and nothing else relates them
  (Theorem~\ref{thm:emit}\,(v)).
\end{enumerate}
The emitted constraint list is \emph{empty on the encoded route}; on the
licensed route it carries the flag-sector entries
$(\Ffss_o-p_o\Pflag,\succeq,0)$ of Corollary~\ref{cor:fssinstance} at $c=0$,
unshifted, and, on request, the class-marginal entries
$\big(\sum_{v\in\Vset_t}F_v,=,q_t\big)$, which the flat instance discharges,
a default carried operator as an entry
(Corollary~\ref{cor:entries}\,(i)),
and any further requested entry that Corollary~\ref{cor:entries}\,(iii)
admits, which for this instance is the two-sector entry
$(\Ffss_e-a\,\Pleq-b\,\Pflag,\succeq,0)$ where its decided certificate of
Corollary~\ref{cor:two-sector-decided} holds, all under the licence of that
route (Corollary~\ref{cor:entries}). The instance's default
carried operators,
the sector projectors and the two projectors $\Pleq$ and $\Pflag$ it
transports besides, \emph{are emitted with the transported set at $c=0$ with
$\succeq$ on either route} (Corollary~\ref{cor:entries}\,(i),
Proposition~\ref{prop:flat-instance}\,(I3)). These statements are about the
instance's link and its emitted family.
\end{digestclaim}
\end{mainclaim}

% claim reductions/wpfss/cor:two-sector-decided
\digestlabel{cor:two-sector-decided}{21.4}\pdfbookmark[2]{Corollary 21.4 The requested two-sector entry, decided on the emitted numbers}{digest-reductions/wpfss/cor:two-sector-decided}
\renewcommand*{\RestatedName}{\hyperref{contract.pdf}{}{claim-reductions/wpfss/cor:two-sector-decided}{Corollary~21.4}}%
\begin{digestclaim}[The requested two-sector entry, decided on the emitted numbers]
Let $\mathcal{J}_{\mathrm{flat}}$ be the flat instance of
Proposition~\ref{prop:flat-instance} on the licensed route, with channel
$\Lambda$ and routed weights $(p_v)$, and let $A''$ be an admissible datum of
it over any substrate of Definition~\ref{def:substrate} over which it is a
reduction instance, $\Fld=\mathbb{R}$ included. Let $e\in\Vset$, let $a$ and $b$ be declared real constants, and put
\begin{equation*}\tag{\getrefnumber{C-eq:two-sector-at}}\digesteqlabel{eq:two-sector-at}
  \hat O[A'']\;\coloneqq\;\Ffss_e[A'']-a\,\Pleq-b\,\Pflag ,
\end{equation*}
written in the form \eqref{eq:instance-entry} with coefficient $1$ for $e$
on every sector, every other coefficient $0$, and $Z=-a\,\Pleq-b\,\Pflag$.
Then:
\begin{enumerate}[label=(\roman*),leftmargin=3.0em]
\item its pull-back is
  \begin{equation*}\tag{\getrefnumber{C-eq:two-sector-pullback}}\digesteqlabel{eq:two-sector-pullback}
    \Lambda^{\dagger}\big(\hat O[A'']\big)
    \;=\;\bigoplus_{m\le N}\big(A''^{(m)}_e-a\,\idon{d_m}\big)
    \ \dsum\ \big(\blk{M}{\gtN}_e-b\,\idon{}\big);
  \end{equation*}
\item if
  \begin{equation*}\tag{\getrefnumber{C-eq:two-sector-condition}}\digesteqlabel{eq:two-sector-condition}
    A''^{(m)}_e-a\,\idon{d_m}\;\succeq\;0\quad(m\le N)
    \qquad\text{and}\qquad b\;\le\;p_e ,
  \end{equation*}
  then $\Lambda^{\dagger}(\hat O[A''])\succeq0$, which is the certificate
  \eqref{eq:certificate} of the entry $(\hat O[A''],\succeq,0)$;
\item the first condition of \eqref{eq:two-sector-condition} is necessary for
  that certificate. At the declared datum $A'$ of Theorem~\ref{thm:emit},
  whose blocks and routed weights are emitted lattice values, and for
  constants $a$ and $b$ that are lattice values, both conditions are decided
  by finitely many field operations and sign comparisons on these numbers,
  and \eqref{eq:two-sector-condition} is a decided certificate for the entry
  $(\hat O[A'],\succeq,0)$ in the sense of
  Corollary~\ref{cor:entries}\,(iii).
\end{enumerate}
\end{digestclaim}

% claim reductions/wpfss/cor:flat-companion
\digestlabel{cor:flat-companion}{21.5}\pdfbookmark[2]{Corollary 21.5 The flat instance has a reference companion}{digest-reductions/wpfss/cor:flat-companion}
\renewcommand*{\RestatedName}{\hyperref{contract.pdf}{}{claim-reductions/wpfss/cor:flat-companion}{Corollary~21.5}}%
\begin{digestclaim}[The flat instance has a reference companion]
Let the flat instance $\mathcal{J}_{\mathrm{flat}}$ of
Proposition~\ref{prop:flat-instance} be given, at either route and in either
form, with $K$ the declared low-photon blocks from which its handed datum is
formed by Definition~\ref{def:weight-route} (the conditional POVM blocks
$\blk{M}{=m}_{y'|t}$, $m\le N$, of each setting $t$ in the bare form, the
supplied dark-composed blocks $\blk{\widetilde M}{=m}_{y|t}$ in the
dark-composed form). Then $\mathcal{J}_{\mathrm{flat}}$ satisfies (H-REF) of
Corollary~\ref{cor:reference-companion}: for every family $G$ over the same
blocks, settings, outcomes and dimensions whose blocks are Hermitian and
positive semidefinite, with $\sum_{y'}G^{(m)}_{y'|t}=\idon{\Chan_m}$ exactly
for every $t$ and $m\le N$, the datum the same rule forms from $G$ is
admissible, and $\mathcal{J}_{\mathrm{flat}}$ is a reduction instance over
the substrate $\Fld=\mathbb{R}$. Its rule of (I1) is a non-negative
combination in the sense of Corollary~\ref{cor:companion-factor-enclosure}:
the element $(y,t)$ of the datum it forms from a family $K'$ with the
properties just listed for $G$ is, on the data sector $(m)$, $q_t\sum_{y'}R^{\cg}_{y|y'}K'^{(m)}_{y'|t}$ in the bare
form and $q_tK'^{(m)}_{y|t}$ in the dark-composed form. At the datum $A^{\mathrm{ref}}$ so formed,
its target is $\Fwp[A^{\mathrm{ref}}]$ on the encoded route and
$\Ffss[A^{\mathrm{ref}}]$ on the licensed route, and on the licensed route its
entries are $\big(\Ffss_o[A^{\mathrm{ref}}]-p_o\Pflag,\succeq,0\big)$ for
every bounded letter $o$, with the routed weights $p_o$ of the handed
instance. On the licensed route, let the request carry besides an entry
$(\hat O,\succeq,0)$ with $\hat O=\Ffss_e-a\,\Pleq-b\,\Pflag$, at an outcome
$e\in\Vset$ and declared real constants $a$ and $b$, written in the form
\eqref{eq:instance-entry} with coefficient $1$ for $e$ on every sector, every
other coefficient $0$, and $Z=-a\,\Pleq-b\,\Pflag$, so that its companion form
(Corollary~\ref{cor:reference-companion}\,(vi)) is
$\hat O^{\mathrm{ref}}=\Ffss_e[A^{\mathrm{ref}}]-a\,\Pleq-b\,\Pflag$. Let
$(C,R)$ be an entrywise enclosure of $A^{\mathrm{ref}}$ whose centres are
Hermitian, and put $\rho_m\coloneqq\max_i\sum_lR^{(m)}_{e,il}$ for $m\le N$.
If
\begin{equation*}\tag{\getrefnumber{C-eq:companion-two-sector}}\digesteqlabel{eq:companion-two-sector}
  C^{(m)}_e-(a+\rho_m)\,\idon{d_m}\;\succeq\;0\quad(m\le N)
  \qquad\text{and}\qquad b\;\le\;p_e ,
\end{equation*}
then $\Lambda^{\dagger}(\hat O^{\mathrm{ref}})\succeq0$, with $\Lambda$ the
instance's channel, which is the same at every datum: the companion form
carries the certificate of the entry it is formed from.
\end{digestclaim}

% claim reductions/wpfss/cor:normalized-flat
\digestlabel{cor:normalized-flat}{21.6}\pdfbookmark[2]{Corollary 21.6 The flat instance of a normalized station}{digest-reductions/wpfss/cor:normalized-flat}
\renewcommand*{\RestatedName}{\hyperref{contract.pdf}{}{claim-reductions/wpfss/cor:normalized-flat}{Corollary~21.6}}%
\begin{digestclaim}[The flat instance of a normalized station]
Let $\Recv$ be a receiver instrument without declared source modes, with
efficiencies $\eta_k$ and $\eta_{\max}>0$, let $t\in[\eta_{\max},1]$, and let
$\Recv'$ be a normalized station of $\Recv$ at the
factored loss $t$
(\xref{receiver/detector-imperfections/def:normalized-station}), both with
the channel system $\Chan$. Let the data of
Definition~\ref{def:weight-route} be supplied for $\Recv'$, with
$\Tset=\{\off\}$ and $q_{\off}=1$, pass every refusal rule of this module's
interface, and let $\mathcal{J}'_{\mathrm{flat}}$ be the flat instance of
Proposition~\ref{prop:flat-instance} at the declared route. Put
$\Phi\coloneqq\Lambda^{\mathrm{loss}}_t$, the pure-loss channel of
transmittance $t$ on $\Chan$
(\xref{receiver/detector-imperfections/def:pure-loss}), and
$\mathcal{J}_{\mathrm{flat}}\coloneqq\mathcal{J}'_{\mathrm{flat}}\circ\Phi$.
\begin{enumerate}[label=(\roman*),leftmargin=3.0em]
\item \emph{(the instance)} $\mathcal{J}_{\mathrm{flat}}$ is a reduction
  instance on $\Chan$ with the datum, sectors, scalar values, target,
  transported set, entries, constants and class-marginal flag of
  $\mathcal{J}'_{\mathrm{flat}}$; its channel is that of
  $\mathcal{J}'_{\mathrm{flat}}$ precomposed with $\Phi$, the
  pull-backs of its transported set are $\Phi^{\dagger}$ of
  those of $\mathcal{J}'_{\mathrm{flat}}$, and its reachable states
  (Definition~\ref{def:sound}) are those of that composite channel. Where the supplied blocks are those of
  $\Recv'$'s coarse-grained click-pattern POVM, in the bare form, or of its
  dark-composed family with the declared dark-count matrix, in the
  dark-composed form, its handed model is the post-processed joint family of
  steps~1--2 of Definition~\ref{def:weight-route} formed from $\Recv$'s own
  measurement with the same dark-count matrix.
\item \emph{(the emission)} Theorem~\ref{thm:emit} holds of
  $\mathcal{J}_{\mathrm{flat}}$: its declared datum, emitted family, routed
  weights, entries and constants are those Corollary~\ref{cor:emitted-flat}
  gives for $\mathcal{J}'_{\mathrm{flat}}$, its channel is
  $\Lambda\circ\Phi$ with $\Lambda$ that of $\mathcal{J}'_{\mathrm{flat}}$,
  and for every $v\in\Vset$
  \begin{equation*}\tag{\getrefnumber{C-eq:normalized-flat}}\digesteqlabel{eq:normalized-flat}
    (\Lambda\circ\Phi)^{\dagger}(F_v)\;=\;\Phi^{\dagger}(M'_v)
    \qquad\text{exactly, on }\Chan,
  \end{equation*}
  with $M'$ the declared model \eqref{eq:flat-declared} of
  $\mathcal{J}'_{\mathrm{flat}}$. Where moreover the normalizer kept every
  class, the right-hand side is the family of (i) formed from $\Recv$'s
  measurement.
\item \emph{(several measuring factors)} Where each factor $f$ of
  Theorem~\ref{thm:emitted-product} is such an instance
  $\mathcal{J}^{f}_{\mathrm{flat}}=\mathcal{J}'^{f}_{\mathrm{flat}}\circ
  \Lambda^{\mathrm{loss}}_{t_f}$, each at its own factored loss $t_f$, that
  theorem applies, with channel
  $\bigotimes_f\big(\Lambda_f\circ\Lambda^{\mathrm{loss}}_{t_f}\big)$ and every
  entry at the constant of its factor.
\item \emph{(the reference companion)} $\mathcal{J}_{\mathrm{flat}}$
  satisfies (H-REF) of Corollary~\ref{cor:reference-companion}, and its
  reference companion at a family $G$ is that of
  $\mathcal{J}'_{\mathrm{flat}}$ at $G$ precomposed with $\Phi$: the same
  target and entries, the companion's channel precomposed with $\Phi$, and
  $\Phi^{\dagger}$ of the companion's pull-backs.
\end{enumerate}
\end{digestclaim}

\section*{22\quad Requirements}\pdfbookmark[1]{22 Requirements}{digest-group-reductions/wpfss/req:one-route-per-artifact}

% claim reductions/wpfss/req:one-route-per-artifact
\digestlabel{req:one-route-per-artifact}{22.1}\pdfbookmark[2]{Requirement 22.1 One route per artefact}{digest-reductions/wpfss/req:one-route-per-artifact}
\renewcommand*{\RestatedName}{\hyperref{contract.pdf}{}{claim-reductions/wpfss/req:one-route-per-artifact}{Requirement~22.1}}%
\begin{digestrequirement}[One route per artefact]
The engine emits every artefact on exactly one route: the request's
\iface{route} is \iface{encoded} or \iface{licensed}, the emitted payload is
single-valued, and a request asking for both is refused with
\iface{route_not_single_valued}. A caller wanting both issues two requests and
receives two artefacts.
\end{digestrequirement}

% claim reductions/wpfss/req:total-weight-refusal
\digestlabel{req:total-weight-refusal}{22.2}\pdfbookmark[2]{Requirement 22.2 Refusal of a total weight of one or more}{digest-reductions/wpfss/req:total-weight-refusal}
\renewcommand*{\RestatedName}{\hyperref{contract.pdf}{}{claim-reductions/wpfss/req:total-weight-refusal}{Requirement~22.2}}%
\begin{digestrequirement}[Refusal of a total weight of one or more]
The engine refuses certified weights whose total weight
$p=\sum_vp_v$ is not below one, with \iface{total_weight_not_below_one}. It
decides this for every declaration, whether or not $\Hs{\Bsys_{\gtN}}$ is
trivial, and it constructs no squasher at $p=1$, where the construction of
Proposition~\ref{prop:chain} would divide by $1-p=0$.
\end{digestrequirement}

% claim reductions/wpfss/req:largest-certified-weight
\digestlabel{req:largest-certified-weight}{22.3}\pdfbookmark[2]{Requirement 22.3 Use and record of the largest certified weight}{digest-reductions/wpfss/req:largest-certified-weight}
\renewcommand*{\RestatedName}{\hyperref{contract.pdf}{}{claim-reductions/wpfss/req:largest-certified-weight}{Requirement~22.3}}%
\begin{digestrequirement}[Use and record of the largest certified weight]
The engine uses, for each outcome, the largest weight the artefact certifies
and records the weights it used in \iface{wpfss_data}, and it never silently
substitutes a larger one.
\end{digestrequirement}

% claim reductions/wpfss/req:descent-decided-exactly
\digestlabel{req:descent-decided-exactly}{22.4}\pdfbookmark[2]{Requirement 22.4 Exact decision of monotonicity and descent}{digest-reductions/wpfss/req:descent-decided-exactly}
\renewcommand*{\RestatedName}{\hyperref{contract.pdf}{}{claim-reductions/wpfss/req:descent-decided-exactly}{Requirement~22.4}}%
\begin{digestrequirement}[Exact decision of monotonicity and descent]
The engine decides exactly, on the declared dark-count matrix $R$, that $R$
is monotone, $R_{S|S'}>0\Rightarrow S'\subseteq S$, and that it descends
through the declared coarse-graining $\cg$ in the sense of
Lemma~\ref{lem:descent}, $P(y|S')=P(y|S'')$ for every $y$ and every pair
$S',S''$ with $\cg(S')=\cg(S'')$, each an identity or a comparison between the
dyadic rationals the declaration carries (Definition~\ref{def:substrate}). It
tests the descent property itself, and treats the sufficient condition
\eqref{eq:canonicalcg} only as a shortcut that must itself be checked in
full. A declaration that fails either condition is refused with the reason
\iface{dark_counts_do_not_descend}.
\end{digestrequirement}

% claim reductions/wpfss/req:absorbing-bounded-letters
\digestlabel{req:absorbing-bounded-letters}{22.5}\pdfbookmark[2]{Requirement 22.5 Absorbing bounded letters as an input validity condition}{digest-reductions/wpfss/req:absorbing-bounded-letters}
\renewcommand*{\RestatedName}{\hyperref{contract.pdf}{}{claim-reductions/wpfss/req:absorbing-bounded-letters}{Requirement~22.5}}%
\begin{digestrequirement}[Absorbing bounded letters as an input validity condition]
The engine takes \eqref{eq:absorbing} as a declared input validity condition:
it decides $R^{\cg}_{y_i|y_i}=1$ exactly at every bounded letter $y_i$ of the
slot, and refuses a declaration in which some bounded letter fails it, with
the reason \iface{bounded_letter_not_absorbing}. The refusal is of the whole
artifact: the failing entry's weight is not set to zero and the rest kept.
\end{digestrequirement}

% claim reductions/wpfss/req:fibre-equality
\digestlabel{req:fibre-equality}{22.6}\pdfbookmark[2]{Requirement 22.6 Certified sets equal to fibres of the coarse-graining}{digest-reductions/wpfss/req:fibre-equality}
\renewcommand*{\RestatedName}{\hyperref{contract.pdf}{}{claim-reductions/wpfss/req:fibre-equality}{Requirement~22.6}}%
\begin{digestrequirement}[Certified sets equal to fibres of the coarse-graining]
The engine requires every bounded set of the slot to be exactly a fibre of
the declared coarse-graining, $\mathcal{O}_i=\cg^{-1}(y_i)$, validates
the equality at input, and refuses a slot that fails it with the reason
\iface{bound_set_not_a_fibre}.
\end{digestrequirement}

% claim reductions/wpfss/req:column-stochastic-refusal
\digestlabel{req:column-stochastic-refusal}{22.7}\pdfbookmark[2]{Requirement 22.7 Exact column-stochasticity of the declared dark-count matrix}{digest-reductions/wpfss/req:column-stochastic-refusal}
\renewcommand*{\RestatedName}{\hyperref{contract.pdf}{}{claim-reductions/wpfss/req:column-stochastic-refusal}{Requirement~22.7}}%
\begin{digestrequirement}[Exact column-stochasticity of the declared dark-count matrix]
The engine decides \emph{exactly} that the declared dark-count matrix $R$ is
column-stochastic, $R_{S|S'}\ge0$ for every entry and $\sum_SR_{S|S'}=1$ for
every column, and refuses a declaration that is not, with the reason
\iface{dark_counts_not_column_stochastic}. No second rule decides the
property of the descended matrix $R^{\cg}$ of Lemma~\ref{lem:descent}.
\end{digestrequirement}

% claim reductions/wpfss/req:supplier-version
\digestlabel{req:supplier-version}{22.8}\pdfbookmark[2]{Requirement 22.8 Receiver versions from which per-setting weights are read}{digest-reductions/wpfss/req:supplier-version}
\renewcommand*{\RestatedName}{\hyperref{contract.pdf}{}{claim-reductions/wpfss/req:supplier-version}{Requirement~22.8}}%
\begin{digestrequirement}[Receiver versions from which per-setting weights are read]
The engine reads the per-setting weights
$p^{(t)}_i$ of a receiver artifact (\iface{high_photon_bounds}) only where
the artifact's module version resolves to an accepted receiver version that
states them: one whose \xref{receiver/thm:emitted-bound} carries the
per-setting inequality as its first conclusion and whose
\xref{receiver/def:slot} carries a weight per entry and setting. An artifact
whose version does not state them, a probe-free one included, is refused with
\iface{high_photon_slot_ill_formed}: the per-setting values it reports beside
the minimum are not read from it, and its minimum is not read at every
setting in their place.
\end{digestrequirement}

% claim reductions/wpfss/req:normalized-link-record
\digestlabel{req:normalized-link-record}{22.10}\pdfbookmark[2]{Requirement 22.10 The link of a normalized station records the factored loss}{digest-reductions/wpfss/req:normalized-link-record}
\renewcommand*{\RestatedName}{\hyperref{contract.pdf}{}{claim-reductions/wpfss/req:normalized-link-record}{Requirement~22.10}}%
\begin{digestrequirement}[The link of a normalized station records the factored loss]
Where the receiver artefact's \iface{efficiency_record} carries a
factored loss $t<1$, the engine emits the artefact of the precomposed
instance $\mathcal{J}'_{\mathrm{flat}}\circ\Lambda^{\mathrm{loss}}_t$ of
Corollary~\ref{cor:normalized-flat}, and its link records, in place of the
normalized station's: (i) the channel witness as the flat instance's witness
preceded by the pure-loss channel of transmittance $t$ on the
channel modes, with $t$ recorded under \iface{factored_loss}; (ii) the
source family as the receiver's measurement; (iii) each pull-back of the
carried tuple as the pull-back through that loss of the flat instance's
pull-back, not as the flat instance's pull-back itself; and (iv) the
reachable-set statement as the one for that composite channel
(Definition~\ref{def:sound}). The emitted target, routed weights,
transported operators, entries, constants, certificate cases and
licence record are those emitted for the normalized station, unchanged. The reference
companion's link is recorded the same way: its channel and pull-backs carry
the same loss in front, and its target and entries are unchanged.
\end{digestrequirement}

% claim reductions/wpfss/req:normalized-branch-refusal
\digestlabel{req:normalized-branch-refusal}{22.11}\pdfbookmark[2]{Requirement 22.11 A normalized receiver is not read as a branch of a branching instance}{digest-reductions/wpfss/req:normalized-branch-refusal}
\renewcommand*{\RestatedName}{\hyperref{contract.pdf}{}{claim-reductions/wpfss/req:normalized-branch-refusal}{Requirement~22.11}}%
\begin{digestrequirement}[A normalized receiver is not read as a branch of a branching instance]
An artefact that embeds the flat instance as a branch of a branching
instance, the register instance included, reads only a receiver artefact whose \iface{efficiency_record} carries the
factored loss $t=1$, and refuses one with $t<1$ with
\iface{normalized_receiver_in_branch}.
\end{digestrequirement}

\part*{Domination links}\pdfbookmark[0]{Domination links}{digest-part-4}

\section*{23\quad The domination link and its noise channel}\pdfbookmark[1]{23 The domination link and its noise channel}{digest-group-reductions/domination/def:domination-link}

% claim reductions/domination/def:domination-link
\digestlabel{def:domination-link}{23.1}\pdfbookmark[2]{Definition 23.1 The domination link}{digest-reductions/domination/def:domination-link}
\renewcommand*{\RestatedName}{\hyperref{contract.pdf}{}{claim-reductions/domination/def:domination-link}{Definition~23.1}}%
\begin{digestclaim}[The domination link]
Let $\SecSet$ be a finite ordered set of sectors with dimensions $d_s\ge1$,
$\Hs{\Qsys}=\bigoplus_{s\in\SecSet}\mathbb{C}^{d_s}$ with sector projectors
$\Pi_s$, and let $G=\{G_v\}_{v\in\Vset}$ and $F=\{F_v\}_{v\in\Vset}$ be POVMs
on $\Qsys$ over one finite alphabet $\Vset$, both block diagonal for
$\{\Pi_s\}$, $G_v=\bigoplus_sG^{(s)}_v$ and $F_v=\bigoplus_sF^{(s)}_v$. Let
$\vartheta=(\vartheta_s)_{s\in\SecSet}$ with $\vartheta_s\in[0,1)$ be weights at which
$G$ dominates $F$, every block complete at $1$
(\xref{foundations/def:domination}):
\begin{equation*}\tag{\getrefnumber{C-eq:dom-link}}\digesteqlabel{eq:dom-link}
  G^{(s)}_v-(1-\vartheta_s)\,F^{(s)}_v\;\succeq\;0
  \qquad\text{for every }v\in\Vset\text{ and }s\in\SecSet ;
\end{equation*}
they are the
\emph{domination weights}
of the pair. Let $R_v\coloneqq\bigoplus_sR^{(s)}_v$ with $R^{(s)}$ the
residual family of \xref{foundations/def:domination} on sector $s$.
A family $F$ on $\Qsys$ has
\emph{indicator sectors}
when there are pairwise distinct sectors $(f)\in\SecSet$, one for each
$f\in\Vset$, with $d_{(f)}=1$ and $F^{((f))}_v=\delta_{vf}$ for every $v$;
write $\Pi_{(f)}$ for the projector onto the sector $(f)$, a rank-one
projector.
\begin{enumerate}[label=(\alph*),leftmargin=2.4em]
\item \emph{(added flag sectors)} Let $\Qsys^{+}$ be the system with
  $\Hs{\Qsys^{+}}=\Hs{\Qsys}\dsum\Hs{\Flag^{+}}$, $\Hs{\Flag^{+}}=\mathbb{C}^{|\Vset|}$
  with orthonormal basis $\{\ket{v}\}_{v\in\Vset}$, whose sector set is
  $\SecSet$ followed by $|\Vset|$ one-dimensional sectors $(f)^{+}$, $f\in\Vset$,
  in the order of $\Vset$, the sector $(f)^{+}$ spanned by $\ket{f}$. Put
  $F^{+}_v\coloneqq F_v\dsum\dyad{v}_{\Flag^{+}}$, so that the added sectors are
  indicator sectors of $F^{+}$, and
  \begin{equation*}\tag{\getrefnumber{C-eq:noise-fresh}}\digesteqlabel{eq:noise-fresh}
    \Phi^{+}_{\vartheta}(\sigma)\;\coloneqq\;
    \sum_{s\in\SecSet}(1-\vartheta_s)\,\Pi_s\sigma\Pi_s\ \dsum\
    \sum_{v\in\Vset}\Big(\sum_{s\in\SecSet}\vartheta_s
      \tr\big[R_v\,\Pi_s\sigma\Pi_s\big]\Big)\dyad{v}_{\Flag^{+}}
  \end{equation*}
  for $\sigma\in\Tcl{\Qsys}$.
\item \emph{(shared flag sectors)} Suppose $F$ has indicator sectors $(f)$ and
  $\vartheta_{(f)}=0$ for every $f\in\Vset$. Put
  \begin{equation*}\tag{\getrefnumber{C-eq:noise-shared}}\digesteqlabel{eq:noise-shared}
    \Phi_{\vartheta}(\sigma)\;\coloneqq\;
    \sum_{s\in\SecSet}(1-\vartheta_s)\,\Pi_s\sigma\Pi_s\ +\
    \sum_{f\in\Vset}\Big(\sum_{s\in\SecSet}\vartheta_s
      \tr\big[R_f\,\Pi_s\sigma\Pi_s\big]\Big)\Pi_{(f)} .
  \end{equation*}
\end{enumerate}
A
\emph{domination link}
is the tuple \eqref{eq:link} of a receiver-stage typed link
(Definition~\ref{def:link}) with $\textsf{stage}=\textsf{receiver}$, kind
$\textsf{family-dominating}$ and source family $(\Qsys,G)$; with
target $(\Qsys^{+},F^{+})$ and channel witness
$\Phi^{+}_{\vartheta}$ in case (a), or target $(\Qsys,F)$ and channel witness
$\Phi_{\vartheta}$ in case (b); with carried operators the projectors
$\Pi_s\dsum0$, $s\in\SecSet$, and the flag projector
$\Pflag^{+}\coloneqq0\dsum\idon{\Flag^{+}}$ in case (a), and the projectors
$\Pi_s$ of the sectors that are not indicator sectors and the flag projector
$\Pflag\coloneqq\sum_{f}\Pi_{(f)}$ in case (b), each with its pull-back under
the channel witness; with the reachable-set statement of
Definition~\ref{def:sound} for that channel; with the constraint entries and
licence record of the chain it ends; and with
$\textsf{status}=\textsf{exact}$.
\end{digestclaim}

% claim reductions/domination/thm:noise-channel
\digestlabel{thm:noise-channel}{23.2}\pdfbookmark[2]{Theorem 23.2 The noise channel of a domination}{digest-reductions/domination/thm:noise-channel}
\begin{mainclaim}
\renewcommand*{\RestatedName}{\hyperref{contract.pdf}{}{claim-reductions/domination/thm:noise-channel}{Theorem~23.2}}%
\begin{digestclaim}[The noise channel of a domination]
In the setting of Definition~\ref{def:domination-link}, let $G$ dominate $F$
at the domination weights $\vartheta$ and let $R$ be the residual family.
\begin{enumerate}[label=(\roman*),leftmargin=3.0em]
\item $R=\{R_v\}_{v\in\Vset}$ is a POVM on $\Qsys$, block diagonal for
  $\{\Pi_s\}$, and
  $G^{(s)}_v=(1-\vartheta_s)F^{(s)}_v+\vartheta_sR^{(s)}_v$ for every $v$ and $s$.
\item \emph{(added flag sectors)} $\Phi^{+}_{\vartheta}$ of \eqref{eq:noise-fresh} is a
  channel $\Tcl{\Qsys}\to\Tcl{\Qsys^{+}}$, and
  \begin{equation*}\tag{\getrefnumber{C-eq:noise-fresh-relation}}\digesteqlabel{eq:noise-fresh-relation}
    \big(G;\,(1-\vartheta_s)\Pi_s\,(s\in\SecSet),\,\textstyle\sum_s\vartheta_s\Pi_s\big)
    \ \dsq_{\Phi^{+}_{\vartheta}}\
    \big(F^{+};\,\Pi_s\dsum0\,(s\in\SecSet),\,0\dsum\idon{\Flag^{+}}\big).
  \end{equation*}
\item \emph{(shared flag sectors)} If $F$ has indicator sectors $(f)$ with
  $\vartheta_{(f)}=0$ for every $f$, then
  $\Phi_{\vartheta}$ of \eqref{eq:noise-shared} is a channel
  $\Tcl{\Qsys}\to\Tcl{\Qsys}$ with $G\dsq_{\Phi_{\vartheta}}F$, and its adjoint
  takes the sector projectors and the flag projector to
  \begin{equation*}\tag{\getrefnumber{C-eq:noise-shared-pullbacks}}\digesteqlabel{eq:noise-shared-pullbacks}
  \begin{gathered}
    \Phi_{\vartheta}^{\dagger}(\Pi_s)=(1-\vartheta_s)\Pi_s\quad
    (s\text{ not an indicator sector}),\\
    \Phi_{\vartheta}^{\dagger}(\Pi_{(f)})=\Pi_{(f)}+\sum_{s}\vartheta_s\,
      \Pi_sR_f\Pi_s,\qquad
    \Phi_{\vartheta}^{\dagger}(\Pflag)=\Pflag+\sum_{s}\vartheta_s\Pi_s .
  \end{gathered}
  \end{equation*}
\item \emph{(output)} The outputs of both channels are block diagonal for the
  sectors of their target, so either channel is unchanged when the sector
  pinching of its target is applied after it.
\item \emph{(weight zero)} If $\vartheta_s=0$ for every $s$, then $F=G$, and
  $\Phi_{\vartheta}$, where defined, is the sector pinching
  $\sigma\mapsto\sum_s\Pi_s\sigma\Pi_s$.
\end{enumerate}
In particular the domination link of Definition~\ref{def:domination-link}, in
either case, is a receiver-stage link of Definition~\ref{def:link}: its
channel witness relates its source family and carried operators to its
target and carried operators exactly, as \eqref{eq:link} requires with
$\textsf{status}=\textsf{exact}$.
\end{digestclaim}
\end{mainclaim}

\section*{24\quad Domination in a reduction chain}\pdfbookmark[1]{24 Domination in a reduction chain}{digest-group-reductions/domination/cor:dom-chain}

% claim reductions/domination/cor:dom-chain
\digestlabel{cor:dom-chain}{24.1}\pdfbookmark[2]{Corollary 24.1 A squashing link followed by a domination link}{digest-reductions/domination/cor:dom-chain}
\renewcommand*{\RestatedName}{\hyperref{contract.pdf}{}{claim-reductions/domination/cor:dom-chain}{Corollary~24.1}}%
\begin{digestclaim}[A squashing link followed by a domination link]
Let $\mathcal{J}$ be a reduction instance with sector set
$\SecSet=\SecSet_{\mathrm{dat}}\sqcup\SecSet_{\mathrm{sc}}$, target family
$F[\cdot]$ and model map $M[\cdot]$
(Definition~\ref{def:reduction-instance}), and let $A^{\circ}$ and $A'$ be
data for its data sectors such that
\begin{enumerate}[label=(\alph*),leftmargin=2.4em]
\item $A'\in\mathcal{D}_{\mathrm{adm}}$; the blocks $A^{\circ(s)}_v$ are
  Hermitian, positive semidefinite and class complete, so that the
  target family $F[A^{\circ}]$ of \eqref{eq:instance-target} is a POVM
  on $\Qsys$; and there are a POVM $M^{\circ}$ on $\Bsys$, a channel
  $\Lambda^{\circ}\colon\Tcl{\Bsys}\to\Tcl{\Qsys}$ and a finite set
  $\hat O^{\circ}$ of operators on $\Qsys$ containing every sector projector,
  with pull-backs $O^{\circ}$, such that
  $\big(M^{\circ};O^{\circ}\big)\dsq_{\Lambda^{\circ}}
  \big(F[A^{\circ}];\hat O^{\circ}\big)$;
\item $A^{\circ}$ dominates $A'$ at weights
  $\vartheta_s\in[0,1)$, $s\in\SecSet_{\mathrm{dat}}$:
  $A^{\circ(s)}_v-(1-\vartheta_s)A'^{(s)}_v\succeq0$ for every $v\in\Vset$
  (\xref{foundations/def:domination}).
\end{enumerate}
Put $\vartheta_s\coloneqq0$ on every scalar sector. Then:
\begin{enumerate}[label=(\roman*),leftmargin=3.0em]
\item $F[A^{\circ}]$ dominates $F[A']$ at the domination weights
  $\vartheta=(\vartheta_s)_{s\in\SecSet}$, as POVMs on $\Qsys$ block diagonal for
  the sectors.
\item Let $\Lnk_{\mathrm{dom}}$ be the domination link
  (Definition~\ref{def:domination-link}) with source family
  $(\Qsys,F[A^{\circ}])$ and target family $F[A']$, with added flag sectors
  (case (a) of that definition), or with shared flag sectors (case (b)) where
  $F[A']$ has indicator sectors among its scalar sectors; write $F^{\mathrm{dom}}$ and $\Phi$ for its target and channel
  witness. Let $\Lnk_1$ be the receiver-stage link with source family
  $(\Bsys,M^{\circ})$, target $(\Qsys,F[A^{\circ}])$, channel witness
  $\Lambda^{\circ}$ and target-side carried operators the pull-backs
  under $\Phi$ of the carried operators of $\Lnk_{\mathrm{dom}}$, with
  $\textsf{status}=\textsf{exact}$. Then $(\Lnk_1,\Lnk_{\mathrm{dom}})$ is a
  composable and tuple-chained receiver segment, and with an empty source
  segment a reduction chain (Definition~\ref{def:chain}).
\item Its receiver composite $\Phi\circ\Lambda^{\circ}$ satisfies
  $M^{\circ}\dsq_{\Phi\circ\Lambda^{\circ}}F^{\mathrm{dom}}$ exactly,
  its composite status is $\textsf{exact}$, and
  Theorem~\ref{thm:factorization} and Theorem~\ref{thm:transfer} apply to
  it: the protocol whose measurement is $M^{\circ}$ and the squashed
  protocol with target $F^{\mathrm{dom}}$ have the same output state, and
  security of the squashed protocol over analysed inputs that contain the
  composite's image transfers to the protocol with $M^{\circ}$ with the
  same $\varepsilon$.
\item An entry $(\hat O,\bowtie,c)$ declared on $F^{\mathrm{dom}}$ is sound for
  the composite, and Theorem~\ref{thm:licence} applies to the list of such
  entries with the composite channel, when
  $(\Lambda^{\circ})^{\dagger}\big(\Phi^{\dagger}(\hat O)\big)\bowtie
  c\,\idon{\Bsys}$. Every projector onto a sector of the target has a
  positive semidefinite composite pull-back.
\item Where $\sum_{v\in\Vset_t}M^{\circ}_v=q_t\idon{\Bsys}$ for every
  $t\in\Tset$, the class-marginal identity
  \eqref{eq:instance-classmarginal} of $M^{\circ}$, the composite pulls
  $\sum_{v\in\Vset_t}F^{\mathrm{dom}}_v$ back to $q_t\idon{\Bsys}$ for every
  $t\in\Tset$, so the class-marginal entry of
  Definition~\ref{def:reduction-instance}\,(I4) is certified for the
  composite.
\end{enumerate}
\end{digestclaim}

\bibliographystyle{plain}
\bibliography{../../../refs/references}

\end{document}
